% engine_manual.tex —— MIPSolvers 求解引擎（engine）模块技术手册（专著级）
% 规范：docs/modules/README.md 与 docs/modules/THEORY_WRITING_GUIDE.md。
% 编译（本目录下，XeLaTeX 两遍）：
%   xelatex -interaction=nonstopmode engine_manual.tex
%   xelatex -interaction=nonstopmode engine_manual.tex
\documentclass{ctexart}
\usepackage{../../_manual_common/mipsolvers_manual}

\title{MIPSolvers 求解引擎模块技术手册}
\author{MIPSolvers 文档组}
\date{2026 年 9 月 13 日}

\begin{document}
\maketitle

\begin{abstract}
本手册是 MIPSolvers 求解引擎（engine）模块的\textbf{专著级}技术文档。模块公共头文件
位于 \srcpath{include/mipsolvers/engine/}，实现位于 \srcpath{src/engine/}；它对外提供统一
门面 \texttt{SolverEngine}（LE/NLE/LP/QP/NLP/MILP/MINLP/锥规划八类问题），对内包含
原生对偶单纯形 LP 内核、Mehrotra 预测-校正 LP 内点、filter 驱动 NLP 内点、锥内点
（含 chordal 分解）、KKT 装配与稀疏线性代数后端、原生 MILP Branch-and-Cut 树搜索、
LP 预求解，以及 HiGHS/Ipopt/SCIP/Gurobi/CPLEX 外部适配器层。

手册分两层：第 3--6 章为通用理论层（专著级：完整推导与关键定理的严格证明，非代码
转写）；第 7--10 章为源码等价转写层（公式只转写当前源码实际执行的赋值、残差、阈值
与分支）；第 11 章给出数值验证证据，第 12 章列出接口与参数，第 13 章声明验证状态与
边界局限，第 14 章为跨模块共享的工业级评价基线。所有 \srcpath{文件:函数名} 锚点均为
仓库相对路径、禁止行号，由 \texttt{tools/doc\_anchor\_check.py} 回归校验。
\end{abstract}

\tableofcontents
\newpage

\section{阅读指南与约定}

\subsection{数据来源}

本手册实现层（第 7--10 章）的唯一事实来源是下列头/实现文件对，转写前已逐函数通读；
凡代码未消费的输入字段、近似路径、回退与硬编码阈值均在相邻正文明确声明：

\begin{center}
\begin{footnotesize}
\begin{tabular}{@{}ll@{}}
\toprule
\textbf{头文件（声明）} & \textbf{实现文件}\\
\midrule
\srcpath{include/mipsolvers/engine/api/solver.hpp} & \srcpath{src/engine/api/solver.cpp}（门面）\\
\srcpath{include/mipsolvers/engine/api/session.hpp} & \srcpath{src/engine/api/session.cpp}（增量会话）\\
\srcpath{include/mipsolvers/engine/api/options.hpp}、\srcpath{api/problem.hpp}、\srcpath{api/result.hpp} & 仅头文件；经门面与适配器消费\\
\srcpath{include/mipsolvers/engine/solve\_context.hpp} & 资源契约（时限/线程/取消令牌），由各适配器消费\\
\srcpath{include/mipsolvers/engine/solver/solver\_adapter.hpp}、\srcpath{solver/adapter\_registry.hpp} & \srcpath{src/engine/solver/solver\_adapter.cpp}、\srcpath{src/engine/solver/adapter\_registry.cpp}\\
\srcpath{include/mipsolvers/engine/strategy/dispatcher.hpp} & \srcpath{src/engine/strategy/dispatcher.cpp} 及 \srcpath{src/engine/strategy/} 辅助件\\
\srcpath{include/mipsolvers/engine/kernel/lp\_kernel/dual\_simplex.hpp}、\srcpath{lp\_kernel/backend.hpp} & \srcpath{src/engine/kernel/lp\_kernel/dual\_simplex.cpp}、\srcpath{dual\_simplex\_api.cpp}、\srcpath{native\_dual/}\\
\srcpath{include/mipsolvers/engine/kernel/ipm/ipm\_lp\_solver.hpp} & \srcpath{src/engine/kernel/ipm/ipm\_lp\_solver.cpp}、\srcpath{ipm\_lp\_solver\_cached.cpp}\\
\srcpath{include/mipsolvers/engine/kernel/ipm/ipm\_solver.hpp} & \srcpath{src/engine/kernel/ipm/ipm\_solver.cpp} 等 NLP 内点族\\
\srcpath{include/mipsolvers/engine/kernel/ipm/conic\_ipm\_solver.hpp}、\srcpath{ipm/cones.hpp} & \srcpath{src/engine/kernel/ipm/conic\_ipm\_solver.cpp}、\srcpath{cones.cpp}、\srcpath{chordal\_decomposition.cpp}\\
\srcpath{include/mipsolvers/engine/kernel/kkt/kkt\_system.hpp} & \srcpath{src/engine/kernel/kkt/kkt\_system.cpp}\\
\srcpath{include/mipsolvers/engine/kernel/linear\_algebra/} 各头 & \srcpath{src/engine/kernel/linear\_algebra/}（CHOLMOD、HFactor、稀疏 LU）\\
\srcpath{include/mipsolvers/engine/presolve/lp\_presolve.hpp} & \srcpath{src/engine/presolve/lp\_presolve.cpp}\\
\srcpath{include/mipsolvers/engine/bc/api.hpp}、\srcpath{bc/options.hpp}、\srcpath{bc/enums.hpp} & \srcpath{src/engine/solver/native/milp/bc/api.cpp} 与 \srcpath{bc/legacy/} 树搜索族\\
\srcpath{include/mipsolvers/engine/solver/native/native\_adapters.hpp} & \srcpath{src/engine/solver/native/}（LE/NLE/LP 选择器、PDLP、锥适配器）\\
\srcpath{include/mipsolvers/engine/solver/external/adapters.hpp} & \srcpath{src/engine/solver/external/adapters.cpp}（HiGHS/Ipopt/SCIP/Gurobi/CPLEX）\\
— & \srcpath{src/engine/util/problem\_validation.cpp}（输入校验，模块内部）\\
\bottomrule
\end{tabular}
\end{footnotesize}
\end{center}

\subsection{索引空间与单位制}

\begin{itemize}
  \item \textbf{索引}：全模块一律为 \textbf{0 起始}整数索引（C++ 语义）。变量、约束、
        基下标、锥块成员、树节点编号均从 0 开始；割平面暖启动池
        \texttt{BCRootCuts} 的列下标明文规定为预求解前原始 LP 列空间
        （\srcpath{include/mipsolvers/engine/bc/options.hpp}）。
  \item \textbf{单位}：engine 为通用数值求解模块，模型系数本身无量纲，单位制由
        调用方（如 scuc 模块的 MW/MWh 口径）自带。引擎内部的物理量纲仅出现在
        资源契约上：时间一律为墙钟秒（\texttt{double}，字段名以
        \texttt{\_sec} 结尾），内存预算为字节（\fld{memory\_limit\_bytes}），
        线程数为个。容差类字段为绝对或相对无量纲阈值，逐字段口径见第 12 章。
  \item \textbf{矩阵存储}：原生 LP 内核的标准形矩阵 \texttt{StandardColumnMatrix}
        为列主序 CSC，索引宽度自适应（32/64 位）；行视图 \texttt{StandardRowMatrix}
        与之共享值数组（\srcpath{include/mipsolvers/engine/kernel/lp\_kernel/dual\_simplex.hpp}）。
\end{itemize}

\subsection{诚实口径声明}

\begin{itemize}
  \item 本手册为\textbf{专著级}：理论章（第 3--6 章）含完整推导与关键定理的严格
        证明，章首以 theorynote 声明``通用理论，非代码转写''，章末附``与实现的
        对应''表，超出实现的内容就地以 gapnote 标记；
  \item 实现转写章（第 7--10 章）与第 12 章参数表以 \textbf{2026-09-13 工作树}为准，
        公式与字段经逐一核实；代码演进后必须重新核验并以锚点回归兜底；
  \item 数值结论只引用第 11 章所列真实测试目标、基线文件与可复现命令；未执行的
        交叉验证不写成已完成；
  \item 与用户手册 \texttt{docs/manual/05-solvers-engines.md} 行文存在差异之处，
        以代码为准并在 \ref{sec:impl-api-adapters} 章集中声明。
\end{itemize}

\subsection{排版宏}

参数表使用 \texttt{paramtable} 环境（字段名|符号|单位|典型范围|默认值|说明）；字段名以
\verb|\fld{}| 排版；源码锚点以 \verb|\srcpath{}| 排版；结构体元信息以
\verb|\compmeta{名称}{头文件}| 排版；理论章三要素（theorynote、与实现的对应表、
gapnote）按 \texttt{docs/modules/THEORY\_WRITING\_GUIDE.md} 执行。

\section{功能定位与架构}

\subsection{功能定位}

engine 是 MIPSolvers 的\textbf{求解内核与调度层}：对上以统一门面
\texttt{SolverEngine} 接收八类问题的结构化模型，按策略在原生内核与外部求解器之间
选择、回退并归一化结果；对下包含自研的 LP/MILP/NLP/锥求解内核与稀疏线性代数栈。
它不承担建模职责（代数建模属 aml 模块），也不绑定任何应用语义（SCUC 等应用层模块
经门面调用本模块）。

\subsection{对外入口一览}

\begin{footnotesize}
\begin{longtable}{@{}L{3.6cm}L{5.6cm}L{5.0cm}@{}}
\toprule
\textbf{入口} & \textbf{锚点} & \textbf{说明}\\
\midrule
\endhead
统一求解 \texttt{solve} & \srcpath{src/engine/api/solver.cpp:SolverEngine::solve\_normalized} &
\texttt{ProblemVariant} 总入口：规范化、校验、分派、结果映射\\
\texttt{solve\_lp} & \srcpath{src/engine/api/solver.cpp:SolverEngine::solve\_lp} & 线性规划\\
\texttt{solve\_qp} & \srcpath{src/engine/api/solver.cpp:SolverEngine::solve\_qp} & 二次规划\\
\texttt{solve\_nlp} & \srcpath{src/engine/api/solver.cpp:SolverEngine::solve\_nlp} & 非线性规划\\
\texttt{solve\_milp} & \srcpath{src/engine/api/solver.cpp:SolverEngine::solve\_milp} & 混合整数线性规划\\
\texttt{solve\_minlp} & \srcpath{src/engine/api/solver.cpp:SolverEngine::solve\_minlp} & 混合整数非线性规划\\
\texttt{solve\_conic} & \srcpath{src/engine/api/solver.cpp:SolverEngine::solve\_conic} & 锥规划（SOCP/SDP 等）\\
\texttt{solve\_le} & \srcpath{src/engine/api/solver.cpp:SolverEngine::solve\_le} & 稀疏线性方程组\\
\texttt{solve\_nle} & \srcpath{src/engine/api/solver.cpp:SolverEngine::solve\_nle} & 非线性方程组\\
B\&C 直接入口 & \srcpath{src/engine/solver/native/milp/bc/api.cpp:solve\_milp\_bc}、
\srcpath{src/engine/solver/native/milp/bc/api.cpp:solve\_minlp\_bc} &
原生 Branch-and-Cut / 空间分支定界，支持暖启动与 ML 回调\\
增量会话 & \srcpath{src/engine/api/session.cpp:LPModelSession::update\_objective} 等 &
\texttt{LPModelSession}：目标/界/右端项的增量修改与重解\\
\bottomrule
\end{longtable}
\end{footnotesize}

原生内核层（经适配器挂入门面，亦可直接使用）：
\begin{itemize}
  \item 原生对偶单纯形 LP 内核：
        \srcpath{src/engine/kernel/lp\_kernel/dual\_simplex.cpp:solve\_lp\_from\_sf}
        （标准形装配见 \srcpath{src/engine/kernel/lp\_kernel/dual\_simplex\_api.cpp:build\_standard\_form\_lp}，
        内核主循环在 \srcpath{src/engine/kernel/lp\_kernel/native\_dual/solver.cpp}）；
  \item LP 内点（Mehrotra + Gondzio）：
        \srcpath{src/engine/kernel/ipm/ipm\_lp\_solver.cpp:NativeIPMLPAdapter::solve\_lp}；
  \item NLP 内点（filter 全局化）：
        \srcpath{src/engine/kernel/ipm/ipm\_solver.cpp:NativeIPMAdapter::solve\_nlp}；
  \item 锥内点：\srcpath{src/engine/kernel/ipm/conic\_ipm\_solver.cpp:ConicIPMSolver::solve}；
  \item PDLP（一阶方法）：\srcpath{src/engine/solver/native/lp/pdlp\_solver.cpp:NativePDLPAdapter::solve\_lp}；
  \item LP 原生选择器/portfolio：
        \srcpath{src/engine/solver/native/native\_lp\_selector.cpp:NativeDualSimplexLPAdapter::solve\_lp}。
\end{itemize}

外部适配器层由 \srcpath{src/engine/solver/external/adapters.cpp} 实现
（HiGHS/Ipopt/SCIP/Gurobi/CPLEX），与原生适配器一同注册进
\srcpath{src/engine/solver/adapter\_registry.cpp:AdapterRegistry::register\_adapter}，
由 \srcpath{src/engine/strategy/dispatcher.cpp} 的 \texttt{StrategyDispatcher} 按
候选排序与回退策略分派。

\subsection{关键结构体与类}

\compmeta{SolverEngine}{include/mipsolvers/engine/api/solver.hpp}
统一求解门面：持有 \texttt{AdapterRegistry} 与 \texttt{StrategyDispatcher}，提供
上表全部 \texttt{solve\_*} 入口；构造函数可批量注册默认适配器。

\compmeta{SolveOptions}{include/mipsolvers/engine/api/options.hpp}
调用级选项：首选求解器、回退开关、策略枚举、时限/线程/内存预算、随机种子、
组合模式与取消令牌（字段明细见第 12 章）。

\compmeta{ProblemVariant / Result}{include/mipsolvers/engine/api/problem.hpp}
问题变体（八类模型的 \texttt{std::variant}）与统一结果结构
（\srcpath{include/mipsolvers/engine/api/result.hpp}）：状态、目标值、原始/对偶
向量、残差与后端名。

\compmeta{SolveContext}{include/mipsolvers/engine/solve\_context.hpp}
求解资源契约：把时限、线程与取消令牌下发给适配器。

\compmeta{LPModelSession}{include/mipsolvers/engine/api/session.hpp}
LP 增量会话：在保留缓存状态的前提下更新目标、变量界与右端项并重解。

\compmeta{StandardFormLP / SimplexOptions / SimplexBasis}{include/mipsolvers/engine/kernel/lp\_kernel/dual\_simplex.hpp}
原生对偶单纯形核心的公共契约：标准形（含自适应宽度 CSC 存储与 Ruiz 缩放）、
单纯形选项、可跨节点携带的基快照（含 DSE 权重与基操作缓存）。

\compmeta{NativeIPMLPAdapter / IPMLPOptions}{include/mipsolvers/engine/kernel/ipm/ipm\_lp\_solver.hpp}
LP 内点适配器与选项：Mehrotra 预测-校正、Gondzio 校正、Ruiz 均衡、可选 crossover
与节点 LP 缓存求解。

\compmeta{IPMOptions}{include/mipsolvers/engine/kernel/ipm/ipm\_solver.hpp}
NLP 内点选项：容差族、可接受水平、步长上限与外部回退开关（字段明细见第 12 章）。

\compmeta{ConicIPMSolver / ConicIPMOptions}{include/mipsolvers/engine/kernel/ipm/conic\_ipm\_solver.hpp}
锥内点求解器与选项：cvxopt 标准形锥模型、NT 方向、KKT 精化与 chordal 分解开关。

\compmeta{solve\_milp\_bc / BCOptions}{include/mipsolvers/engine/bc/api.hpp}
原生 B\&C 入口与选项束（\srcpath{include/mipsolvers/engine/bc/options.hpp}）：
节点/迭代/时限预算、容差、割平面族、分支与节点选择、启发式、并行与回退策略。

\compmeta{SolverAdapter / AdapterRegistry}{include/mipsolvers/engine/solver/adapter\_registry.hpp}
适配器基类（\srcpath{include/mipsolvers/engine/solver/solver\_adapter.hpp}）与按
问题类索引的注册表。

\compmeta{StrategyDispatcher}{include/mipsolvers/engine/strategy/dispatcher.hpp}
候选排序与回退分派器：按策略枚举生成候选序列，失败时顺延并把剩余时限传给下一个
候选。

\subsection{架构分层}

\begin{itemize}
  \item \textbf{API 层}：\texttt{api/}（\texttt{SolverEngine}、\texttt{LPModelSession}、
        \texttt{SolveOptions}/\texttt{ProblemVariant}/\texttt{Result}）——只做规范化、
        校验与结果映射，不含算法；
  \item \textbf{调度层}：\texttt{solver/} 与 \texttt{strategy/}——适配器接口、注册表、
        候选排序与回退、presolve/postsolve/warmstart 管理器；
  \item \textbf{原生内核层}：\texttt{kernel/}（对偶单纯形、LP/NLP/锥内点、KKT 装配、
        稀疏线性代数）与 \texttt{solver/native/}（把内核包装为适配器，含 LP 选择器、
        PDLP、B\&C）；
  \item \textbf{外部适配器层}：\texttt{solver/external/}——第三方求解器的模型映射与
        结果解析；
  \item \textbf{支撑层}：\texttt{presolve/}（原生 LP 预求解）与 \texttt{util/}（输入
        校验）。
\end{itemize}

% theory_linear_programming.tex —— engine 通用理论：线性规划
% （标准形与变换、对偶定理、最优性条件、对偶单纯形法、LP 预求解理论）
% 本文件为理论层章节，遵循 docs/modules/THEORY_WRITING_GUIDE.md。

\section{线性规划：对偶理论、单纯形方法与预求解}\label{sec:thlp}

\begin{theorynote}
本章为通用数学理论，按线性规划领域的标准模型与文献体系撰写，覆盖：标准形
与等价变换、弱/强对偶定理及其证明、最优性条件、对偶单纯形法（对偶可行基
维持、比率测试、退化与循环、Bland 规则、steepest edge / Devex / BFRT 定价
原理与参考值框架更新推导、数值稳定性）以及 Andersen \& Andersen 意义下的
LP 预求解理论。本章公式不要求逐式对应代码；与平台实现
（\srcpath{src/engine/kernel/lp\_kernel/native\_dual/}、
\srcpath{src/engine/presolve/lp\_presolve.cpp}）的对应关系见本章末
``与实现的对应''表，超出实现的内容以``未实现扩展说明''框就地标记。
\end{theorynote}

\subsection{记号与约定}\label{subsec:thlp-notation}

本章固定如下记号；与手册其余部分冲突时以本章声明为准。

\begin{itemize}
  \item 矩阵 $A\in\mathbb{R}^{m\times n}$（约束系数矩阵），$b\in\mathbb{R}^m$
        （右端项），$c\in\mathbb{R}^n$（目标系数）；$A$ 的第 $j$ 列记作 $a_j$，
        第 $i$ 行记作 $a^i$（行向量）。不作特殊说明时假设 $\mathrm{rank}(A)=m\le n$。
  \item 下标 $j\in\mathbb{R}$ 的变量界：下界 $l_j\in\mathbb{R}\cup\{-\infty\}$，
        上界 $u_j\in\mathbb{R}\cup\{+\infty\}$，$l_j\le u_j$；
        $l\le x\le u$ 按分量理解。$x_j$ 称为 free 变量当且仅当
        $l_j=-\infty$ 且 $u_j=+\infty$；fixed 变量当且仅当 $l_j=u_j$。
  \item 基（basis）：列指标集 $B\subseteq\{1,\dots,n\}$，$|B|=m$ 且基矩阵
        $A_B$（亦简记 $B$，上下文可区分）非奇异；非基指标集
        $N=\{1,\dots,n\}\setminus B$。
  \item 对基 $B$ 的典则量：原始基本解 $x_B=B^{-1}b$（记 $\bar b$），
        对偶解 $y=B^{-T}c_B$，既约费用（reduced cost）
        $\bar c=c-A^{T}y$； tableau 行/列
        $\bar a_j=B^{-1}a_j$、$w^i=e_i^{T}B^{-1}$（$e_i$ 为单位向量）。
  \item $\|\cdot\|$ 未注明时下标指明范数；$\kappa(M)=\|M\|\,\|M^{-1}\|$ 为条件数；
        $\varepsilon_{\mathrm{mach}}$ 为双精度机器精度（$\approx 2.2\times10^{-16}$）。
  \item 定理环境按节编号；``字典（dictionary）''指单纯形迭代中把基变量与目标
        显式表示为非基变量仿射函数的方程组。
\end{itemize}

\subsection{标准形与等价变换}\label{subsec:thlp-standard}

\begin{definition}[一般形式与标准形]\label{def:thlp-std}
线性规划（LP）的一般形式为
\begin{equation}\label{eq:thlp-general}
  \min_{x}\; c^{T}x
  \quad\text{s.t.}\quad
  b^{L}\le Ax\le b^{U},\quad l\le x\le u,
\end{equation}
其中行侧 $b^{L}_i\in\mathbb{R}\cup\{-\infty\}$、$b^{U}_i\in\mathbb{R}\cup\{+\infty\}$。
本章用于对偶与单纯形理论推导的\textbf{标准形}为
\begin{equation}\label{eq:thlp-std}
  \min_{x}\; c^{T}x \quad\text{s.t.}\quad Ax=b,\quad x\ge 0,
\end{equation}
其实现侧使用的\textbf{箱式有界标准形}为
\begin{equation}\label{eq:thlp-boxed}
  \min_{x}\; c^{T}x \quad\text{s.t.}\quad Ax=b,\quad l\le x\le u .
\end{equation}
\end{definition}

\begin{proposition}[等价变换]\label{prop:thlp-transform}
一般形式 \eqref{eq:thlp-general} 可在多项式（线性）规模的变换下化为标准形
\eqref{eq:thlp-std}，且两问题的可行/不可行状态、最优值与最优解集通过显式
仿射映射一一对应。具体地：
\begin{enumerate}
  \item \textbf{最大化变最小化}：$\max c^{T}x=-\min(-c)^{T}x$。
  \item \textbf{不等式行变等式}：行 $a^{i}x\le b^{U}_i$ 引入松弛变量
        $s_i\ge0$：$a^{i}x+s_i=b^{U}_i$；双侧行
        $b^{L}_i\le a^{i}x\le b^{U}_i$ 引入 $0\le s_i\le b^{U}_i-b^{L}_i$
        （ranged row 的松弛变量自身带箱式界）。
  \item \textbf{自由变量拆分}：$x_j$ free 时令 $x_j=x_j^{+}-x_j^{-}$，
        $x_j^{+},x_j^{-}\ge0$。
  \item \textbf{一般下界平移}：$x_j\ge l_j$（$l_j$ 有限）令
        $x_j'=x_j-l_j\ge0$，右端项吸收常数 $b\leftarrow b-l_j a_j$；
        有限上界由第 2 步作为不等式处理。
\end{enumerate}
\end{proposition}

\begin{proof}
每步都是变量与约束的显式仿射替换，替换映射在可行点集上双射且保持目标值
（第 1 步差一个符号与常数）。以第 2 步为例：若 $(x,s)$ 满足
$a^{i}x+s_i=b^{U}_i$、$s_i\ge0$，则 $a^{i}x=b^{U}_i-s_i\le b^{U}_i$；
反之若 $a^{i}x\le b^{U}_i$ 则取 $s_i=b^{U}_i-a^{i}x\ge0$。其余各步同理。
\end{proof}

\begin{remark}\label{rem:thlp-transform-cost}
变换 3 把 $n$ 个变量至多扩为 $2n$ 个，变换 2 每行至多加一个变量，故标准形
规模与原问题线性相关。工业实现通常不真的做拆分/松弛展开，而是直接在箱式
有界形 \eqref{eq:thlp-boxed} 上运行（界 handled implicitly），这正是
\srcpath{src/engine/kernel/lp\_kernel/native\_dual/} 的
\srcpath{include/mipsolvers/engine/kernel/lp\_kernel/dual\_simplex.hpp:StandardFormLP} 采用的表示：等式约束 + 箱式变量界，
非基变量只能处于下界、上界或（free 变量时）零点三种状态之一。
\end{remark}

\subsection{基、顶点与字典}\label{subsec:thlp-basis}

\begin{definition}[基与基本解]\label{def:thlp-bfs}
对标准形 \eqref{eq:thlp-std}，若列指标集 $B$ 使 $A_B$ 非奇异，则称 $B$ 为
基；$x_B=A_B^{-1}b$、$x_N=0$ 称为对应的基本解；若 $x_B\ge0$ 则称
\textbf{基本可行解}（BFS），$B$ 为可行基。若 $x_B$ 的某分量为 $0$，称该
BFS（原始）退化。
\end{definition}

\begin{lemma}[锥的 Carath\'eodory 定理]\label{lem:thlp-caratheodory}
设 $b=\sum_{j=1}^{k}\lambda_j a_j$，$\lambda_j>0$。则存在指标子集
$S\subseteq\{j:\lambda_j>0\}$，$\{a_j\}_{j\in S}$ 线性无关，且
$b=\sum_{j\in S}\mu_j a_j$，$\mu_j>0$。特别地 $|S|\le m$。
\end{lemma}

\begin{proof}
取正系数个数最少的表示。若对应列 $\{a_j:\lambda_j>0\}$ 线性相关，则存在
不全为零的 $\nu_j$（仅在正系数指标上非零）使 $\sum_j\nu_j a_j=0$。取
$\theta=\min\{\lambda_j/|\nu_j|:\nu_j\ne0\}$ 并选择符号使
$\lambda_j'=\lambda_j-\theta\nu_j\ge0$ 且至少一个 $\lambda_j'$ 变为 $0$，
得到正系数更少的表示，矛盾。故最少表示的列线性无关，从而个数不超过
$\mathrm{rank}(A)\le m$。
\end{proof}

\begin{proposition}[基本可行解的存在性]\label{prop:thlp-bfs-exist}
若标准形 \eqref{eq:thlp-std} 可行且 $\mathrm{rank}(A)=m$，则存在 BFS。
\end{proposition}

\begin{proof}
任取可行 $x$，$b=Ax=\sum_{j:x_j>0}x_j a_j$。由引理
\ref{lem:thlp-caratheodory}，可用线性无关列子集 $S$（$|S|\le m$）正系数
表示 $b$。因 $\mathrm{rank}(A)=m$，可把 $S$ 扩充为含 $m$ 列的基 $B$
（Steinitz 换入），对应基本解满足 $x_B\ge0$（$S$ 上分量仍为正，其余为 $0$）。
\end{proof}

\begin{definition}[字典]\label{def:thlp-dict}
对基 $B$，等式系统 $Ax=b$ 等价于
\begin{equation}\label{eq:thlp-dict}
  x_B=\bar b-\bar N x_N,\qquad
  z=v+\bar c_N^{T}x_N,
\end{equation}
其中 $\bar N=A_B^{-1}A_N$，$\bar b=A_B^{-1}b$，$v=c_B^{T}\bar b$，
$\bar c_N=c_N-\bar N^{T}c_B$。称 \eqref{eq:thlp-dict} 为基 $B$ 的字典。
字典是把 $z=c^{T}x$ 与 $x_B$ 在仿射子空间 $\{x:Ax=b\}$ 上用坐标
$x_N$ 表达的仿射表示；同一可行点用任何字典计算目标值都得到同一实数，
这一事实在 Bland 终止性证明中起关键作用（定理 \ref{thm:thlp-bland}）。
\end{definition}

\begin{remark}\label{rem:thlp-boxed-nonbasic}
箱式有界形 \eqref{eq:thlp-boxed} 中非基变量 $x_j$（$j\in N$）不必为 $0$：
它被钉在其下界、上界或（free 时）零点上。此时
$\bar b=B^{-1}(b-\sum_{j\in N}a_j x_j)$，既约费用符号条件按非基状态分类
（见命题 \ref{prop:thlp-boxed-kkt}）。单纯形迭代允许非基变量在两个界之间
直接翻转（bound flip）而不换基，这是箱式实现相对教科书形式的主要结构差异。
\end{remark}

\subsection{对偶理论}\label{subsec:thlp-duality}

对标准形原问题（P）：$\min\{c^{T}x:Ax=b,\ x\ge0\}$，其对偶（D）为
$\max\{b^{T}y:A^{T}y\le c\}$。对偶变量 $y$ 对应等式约束的乘子，自由取值。

\begin{theorem}[弱对偶]\label{thm:thlp-weak}
若 $x$ 对（P）可行、$y$ 对（D）可行，则 $b^{T}y\le c^{T}x$。
\end{theorem}

\begin{proof}
$b^{T}y=y^{T}b=y^{T}Ax=(A^{T}y)^{T}x\le c^{T}x$，其中不等式由
$A^{T}y\le c$ 与 $x\ge0$ 分量相乘保持方向得到。
\end{proof}

\begin{corollary}\label{cor:thlp-weak-conseq}
（i）若（P）无界则（D）不可行；（ii）若（D）无界则（P）不可行；
（iii）若可行对 $(x,y)$ 满足 $c^{T}x=b^{T}y$，则二者各自最优。
\end{corollary}

强对偶的证明依赖于 Farkas 引理与单纯形方法的存在性结论。我们先给出
Farkas 引理的完整证明（投影论证），再从单纯形终止性推出强对偶。

\begin{lemma}[Farkas]\label{lem:thlp-farkas}
恰好下列二者之一成立：
（i）存在 $x\ge0$ 使 $Ax=b$；
（ii）存在 $y$ 使 $A^{T}y\ge0$ 且 $b^{T}y<0$。
\end{lemma}

\begin{proof}
\textbf{不能同时成立}：若 $x\ge0$、$Ax=b$ 与 $A^{T}y\ge0$ 同时存在，则
$0\le y^{T}Ax=y^{T}b<0$，矛盾。

\textbf{至少其一成立}：设锥 $K=\{Ax:x\ge0\}$。由引理
\ref{lem:thlp-caratheodory}，$K$ 是有限个闭锥
$K_S=\{A_S\mu:\mu\ge0\}$（$S$ 取遍线性无关列子集）的并，故 $K$ 为闭凸锥。
设 $b\notin K$。取 $z_0\in K$（如 $z_0=0$），令
$R=\|z_0-b\|$；紧集 $K\cap\{z:\|z-b\|\le R\}$ 上连续函数 $\|z-b\|$ 取得
最小值点 $p\in K$，且 $p$ 是 $K$ 中距 $b$ 最近的点。令 $y=p-b\ne0$。
对任意 $z\in K$ 与 $t\in[0,1]$，凸性给出 $p+t(z-p)\in K$，由 $p$ 的最近性，
\begin{equation*}
  \|p+t(z-p)-b\|^{2}\ge\|p-b\|^{2}
  \;\Longrightarrow\;
  2t\,y^{T}(z-p)+t^{2}\|z-p\|^{2}\ge0 .
\end{equation*}
除以 $t>0$ 并令 $t\downarrow0$ 得 $y^{T}(z-p)\ge0$，即
$y^{T}z\ge y^{T}p$（$\forall z\in K$）。取 $z=0$ 得 $y^{T}p\le0$；
取 $z=2p$ 得 $y^{T}p\ge0$；故 $y^{T}p=0$。取 $z=a_j=Ae_j\in K$ 得
$y^{T}a_j\ge0$，即 $A^{T}y\ge0$。最后
$b^{T}y=(p-y)^{T}y=-\|y\|^{2}<0$。于是（ii）成立。
\end{proof}

\begin{remark}\label{rem:thlp-farkas-geo}
（ii）中的 $y$ 是把 $b$ 与锥 $K$ 严格分离的超平面法向：$K$ 在其非负侧而
$b$ 在其严格负侧。Farkas 引理是 LP 不可行性的代数证书（certificate）形式：
求解器报 ``infeasible'' 时应能输出这样的乘子，见
\srcpath{certificate.cpp:phase\_one\_farkas\_certificate}。
\end{remark}

\begin{theorem}[强对偶]\label{thm:thlp-strong}
若（P）存在最优解，则（D）存在最优解，且二者最优值相等。
\end{theorem}

\begin{proof}
设（P）有最优解 $x^{*}$，最优值 $z^{*}$ 有限。由命题
\ref{prop:thlp-bfs-exist}，（P）存在 BFS。从任一可行基出发执行带 Bland
规则的原始单纯形法；由 Bland 有限终止定理（定理 \ref{thm:thlp-bland}，
其证明不依赖对偶理论），迭代必在有限步终止于某个基 $B$，使既约费用
$\bar c=c-A^{T}y\ge0$，其中 $y=B^{-T}c_B$（若中途检测到无界，与
$z^{*}$ 有限矛盾，故不可能）。于是 $A^{T}y\le c$，即 $y$ 对偶可行；且
\begin{equation*}
  b^{T}y=b^{T}B^{-T}c_B=(B^{-1}b)^{T}c_B=c_B^{T}x_B^{*}=c^{T}x^{*},
\end{equation*}
最后一步因 $x^{*}_N=0$。由弱对偶（定理 \ref{thm:thlp-weak}），$y$ 对偶
最优，且两最优值相等。
\end{proof}

\begin{corollary}[互补松弛]\label{cor:thlp-cs}
设 $x$、$y$ 分别为（P）、（D）可行。则二者皆最优当且仅当
\begin{equation}\label{eq:thlp-cs}
  x_j\bigl(c-A^{T}y\bigr)_j=0,\qquad j=1,\dots,n .
\end{equation}
\end{corollary}

\begin{proof}
弱对偶证明中的不等式链给出
$c^{T}x-b^{T}y=(c-A^{T}y)^{T}x=\sum_j x_j(c-A^{T}y)_j\ge0$，
求和各项均非负。由强对偶，最优当且仅当对偶间隙为零，当且仅当每项为零。
\end{proof}

\begin{proposition}[箱式有界形的最优性条件]\label{prop:thlp-boxed-kkt}
对箱式有界形 \eqref{eq:thlp-boxed}，$x$ 最优当且仅当存在
$y\in\mathbb{R}^m$ 使 $\bar c=c-A^{T}y$ 满足
\begin{equation}\label{eq:thlp-boxed-kkt}
  \begin{cases}
    \bar c_j\ge0, & x_j=l_j\ (\text{仅贴下界}),\\
    \bar c_j\le0, & x_j=u_j\ (\text{仅贴上界}),\\
    \bar c_j=0, & l_j<x_j<u_j .
  \end{cases}
\end{equation}
\end{proposition}

\begin{proof}
经命题 \ref{prop:thlp-transform} 的等价变换把 \eqref{eq:thlp-boxed} 化为
标准形：下界平移 $\tilde x_j=x_j-l_j\ge0$，上界 $x_j\le u_j$ 引入松弛
$s_j=u_j-x_j\ge0$（作为额外行 $x_j+s_j=u_j$）。对变换后问题应用推论
\ref{cor:thlp-cs}：原变量与松弛的互补条件 $x_j^{+}\bar c_j^{+}=0$、
$s_j\bar c_j^{-}=0$ 合并即得 \eqref{eq:thlp-boxed-kkt} 的三分支。具体地，
等式行乘子合并为 $y$，$x_j$ 贴下界时对应 $\tilde x_j=0$ 的互补条件
$\bar c_j\ge0$；贴上界时对应 $s_j=0$ 的互补条件 $\bar c_j\le0$；
严格内部时两者同时施加，得 $\bar c_j=0$。
\end{proof}

\begin{remark}\label{rem:thlp-duality-cases}
（P）/（D）的可行与有界组合共四种可能：皆有最优（值相等）、（P）不可行且
（D）无界、（D）不可行且（P）无界、两者皆不可行。``（P）不可行而（D）有
最优''不可能发生（弱对偶），等等。工业求解器返回的状态枚举
（optimal / infeasible / unbounded / interrupted）正是该分类的计算对应。
\end{remark}

\subsection{对偶单纯形法}\label{subsec:thlp-dual-simplex}

\subsubsection{对偶可行基与迭代结构}

\begin{definition}[对偶可行基]\label{def:thlp-dualfeas}
对标准形 \eqref{eq:thlp-std}，基 $B$ 称为\textbf{对偶可行}的，若既约费用
$\bar c=c-A^{T}y\ge0$（$y=B^{-T}c_B$）。对箱式有界形
\eqref{eq:thlp-boxed}，对偶可行指 $\bar c$ 满足 \eqref{eq:thlp-boxed-kkt}
中与非基状态对应的符号条件。
\end{definition}

对偶单纯形法维持对偶可行，迭代消除原始不可行性（$x_B$ 的界违反），直到
原始也可行——此时由推论 \ref{cor:thlp-cs} 已达最优；或检测到原始不可行
（Farkas 证书）。为推导清晰，本小节先在标准形 $x\ge0$ 上演示比率测试，
再回到箱式有界形给出实现所用的 BFRT。

\begin{proposition}[对偶单纯形比率测试的正确性]\label{prop:thlp-bfrt-std}
设基 $B$ 对偶可行（$\bar c\ge0$），$x_B=\bar b$ 的第 $p$ 个分量
$\bar b_p<0$。记第 $p$ 行 tableau 行向量
$\rho^{T}=e_p^{T}B^{-1}A$，其非基分量 $\bar a_{pj}=\rho_j$（$j\in N$）。
\begin{enumerate}
  \item 若对所有 $j\in N$ 都有 $\bar a_{pj}\ge0$，则（P）不可行。
  \item 否则，取
  \begin{equation}\label{eq:thlp-bfrt-std}
    q\in\operatorname*{arg\,min}_{j\in N:\,\bar a_{pj}<0}
    \frac{\bar c_j}{|\bar a_{pj}|},
  \end{equation}
  以 $(p,q)$ 为主元换基（$x_q$ 入基、$x_{B_p}$ 出基于 $0$），则新基仍对偶
  可行，且对偶目标值 $b^{T}y$ 不减。
\end{enumerate}
\end{proposition}

\begin{proof}
（1）取 $y'=B^{-T}e_p$。则 $(y')^{T}b=e_p^{T}B^{-1}b=\bar b_p<0$；
对任意 $x\ge0$ 若 $Ax=b$，则第 $p$ 行等式给出
$\bar b_p=x_{B_p}+\sum_{j\in N}\bar a_{pj}x_j\ge0$（因 $x\ge0$、
$\bar a_{pj}\ge0$），与 $\bar b_p<0$ 矛盾。等价地，$A^{T}y'\ge0$、
$b^{T}y'<0$ 即 Farkas 证书（引理 \ref{lem:thlp-farkas} 的（ii））。

（2）换基后新既约费用
\begin{equation}\label{eq:thlp-rc-update}
  \bar c_j'=\bar c_j-\frac{\bar a_{pj}}{\bar a_{pq}}\,\bar c_q ,
  \qquad j\ne q,
\end{equation}
而 $\bar c_q'=0$。当 $\bar a_{pj}\ge0$ 时，
$-\bar a_{pj}\bar c_q/\bar a_{pq}\ge0$（因 $\bar a_{pq}<0$、
$\bar c_q\ge0$），故 $\bar c_j'\ge\bar c_j\ge0$。当 $\bar a_{pj}<0$ 时，
\eqref{eq:thlp-bfrt-std} 的选择保证
$\bar c_q/\bar a_{pq}\ge\bar c_j/\bar a_{pj}$（两边同除以负数
$|\bar a_{pj}||\bar a_{pq}|$ 即见），从而 $\bar c_j'\ge0$。
新对偶解 $y'=y-t\,B^{-T}e_p$，$t=\bar c_q/\bar a_{pq}\le0$，
对偶目标变化 $b^{T}(y'-y)=-t\,\bar b_p\ge0$（$t\le0$、$\bar b_p<0$）。
\end{proof}

\begin{remark}\label{rem:thlp-dual-progress}
命题 \ref{prop:thlp-bfrt-std}（2）中，若 $\bar c_q>0$ 且 $\bar b_p<0$
严格，则对偶目标严格上升（$t<0$）；若 $\bar c_q=0$（对偶退化）则
$t=0$，迭代只换基而目标不动，称为\textbf{停滞}（stalling）。对偶退化是
对偶单纯形循环的根源，见 \ref{subsubsec:thlp-cycling} 小节。
\end{remark}

\subsubsection{箱式有界形上的 BFRT}

实现侧在箱式有界形上工作，比率测试推广为\textbf{有界变量比率测试}
（bounded-variable ratio test，实现中代号 BFRT）：出基行 $p$ 的基变量违反
界（$\beta=x_{B_p}\notin[l_{B_p},u_{B_p}]$），以移动方向
$\sigma=+1$（从下侧违反推向 $u$）或 $\sigma=-1$ 刻画；对每个非基列
$j\in N$ 计算断点（breakpoint）
\begin{equation}\label{eq:thlp-bfrt-boxed}
  t_j=\frac{\bar c_j}{\sigma\,\bar a_{pj}}
  \quad\text{（仅当 }\sigma\,\bar a_{pj}\text{ 的符号使 }t_j>0
  \text{ 且移动缩小 }|\bar c_j|\text{ 朝 }0\text{ 时入候选集）},
\end{equation}
比率步长 $t^{*}=\min_j t_j$，候选集随 $t$ 增长依次达到符号边界；首个达到
断点的列入基（或发生 bound flip）。该规则是命题 \ref{prop:thlp-bfrt-std}
在非基变量可取上下界情形下的直接推广：既约费用沿 $t$ 的更新仍是
\eqref{eq:thlp-rc-update} 的形式，符号条件逐分支验证即可，此处从略。

\begin{remark}[Harris 两段比率测试]\label{rem:thlp-harris}
精确的最小比率 $t^{*}$ 在数值上是脆弱的：两个断点相差
$O(\varepsilon_{\mathrm{mach}})$ 量级时，取哪个主元应由数值质量而非算术
次序决定。Harris 的\textbf{两段过程}（two-pass ratio test）\cite{lpharris1973}：
第一段在\textbf{放宽的}容差 $\delta$（如 $10^{-6}$ 倍尺度）下求最大可接受
步长 $\hat t$；第二段在所有断点不超过 $\hat t$ 的候选中，选主元元素
$|\bar a_{pj}|$ 最大者。第一段保证对偶可行性维持在容差内，第二段把选择
自由度用于数值稳定性（大主元减小换基后的误差放大）。该思想同时适用于
原始与对偶比率测试；实现中原始侧为
\srcpath{primal.cpp:choose\_leaving\_harris}，对偶侧 BFRT 的两段结构内嵌于
\srcpath{pricing.cpp:choose\_entering\_bfrt}。
\end{remark}

\subsubsection{退化、停滞与循环}\label{subsubsec:thlp-cycling}

\begin{definition}[对偶退化与循环]\label{def:thlp-cycling}
对偶单纯形迭代中，若某个非基既约费用 $\bar c_j=0$，称当前基
\textbf{对偶退化}。若比率测试的最小比率 $t^{*}=0$，称该次迭代
\textbf{停滞}。若基序列 $B_0,B_1,\dots$ 出现 $B_{k+T}=B_k$（$T\ge1$），
称发生\textbf{循环}（cycling）。
\end{definition}

由于同一基唯一确定字典（定义 \ref{def:thlp-dict}），循环意味着算法在有限
个字典上无限打转。Beale 的经典例子 \cite{lpbeale1955} 表明：使用
``最小比率 + 固定次序打破平局'' 的朴素主元规则，对偶（与原始）单纯形确实
可以循环。实践中最小索引平局打破即足以在多数问题上避免循环，但理论上需要
可证的反循环规则。

\begin{theorem}[Bland 规则的有限终止性]\label{thm:thlp-bland}
对标准形原始单纯形（$\max c^{T}x$，$Ax=b$，$x\ge0$；字典
$z=v+\sum_{j\in N}\bar c_j x_j$），采用 Bland 规则
\cite{lpbland1977}：
\begin{enumerate}
  \item 入基：取满足 $\bar c_j>0$ 的\textbf{最小}下标 $j$；
  \item 出基：在比率测试取得最小值的行中，取基变量下标\textbf{最小}者。
\end{enumerate}
则单纯形法在有限步内终止（最优或检测到无界）。
\end{theorem}

\begin{proof}
反设发生循环。循环中目标值 $v$ 不变（否则严格递增的 $v$ 不可能重复出现
同一基），故循环内所有主元步的比率均为 $0$，即所有参与换基的变量在循环
内取值恒为 $0$（退化主元）；称这些变量为\textbf{循环变量}，其余变量在
循环中保持基/非基身份与取值不变。令
$t$ 为循环中入基过的变量下标的最大值；$t$ 必也在循环中某次出基（基只有
$m$ 个位置，循环回到起点）。取循环中两个字典：
\begin{itemize}
  \item $D_1$：$t$ 在该步\textbf{入基}。记其既约费用为 $\bar c$，
        tableau 元为 $\bar a_{ij}$。由入基规则，
        \begin{equation}\label{eq:thlp-bland-d1}
          \bar c_t>0,\qquad \bar c_j\le0\quad(j\in N_1,\ j<t).
        \end{equation}
  \item $D_2$：$t$ 在该步\textbf{出基}，入基变量为 $s$，主元行 $r$ 为
        $t$ 所在行。由入基规则
        \begin{equation}\label{eq:thlp-bland-d2}
          \bar d_s>0,\qquad \bar d_j\le0\quad(j\in N_2,\ j<s),
        \end{equation}
        其中 $\bar d$、$\hat a_{ij}$ 为 $D_2$ 的既约费用与 tableau 元。
        $s$ 在循环中入基，故 $s\le t$；又 $s\ne t$（$t$ 此时在基内），
        故 $s<t$。
\end{itemize}
\textbf{关键观察}：对 $D_2$ 中任一循环基变量 $x_j$（$j\ne t$，其所在行
记 $i$），若 $\hat a_{is}>0$，则行 $i$ 与行 $r$ 同为比率 $0$ 的平局行
（循环行右端项为 $0$），出基规则应取下标最小者，而已选 $t$，故 $j\ge t$；
但 $j$ 是循环变量，$j\le t$，于是 $j=t$，矛盾。因此
\begin{equation}\label{eq:thlp-bland-rowsign}
  j\ \text{为}\ D_2\ \text{的循环基变量且}\ j\ne t
  \;\Longrightarrow\; \hat a_{js}\le0 ,
\end{equation}
其中 $\hat a_{js}$ 简记 $j$ 所在行的 $s$ 列元。

现在沿 $D_2$ 构造参数化点：令 $x_s=\varepsilon$（$\varepsilon\in\mathbb{R}$
任意），其余 $D_2$ 非基变量取 $0$，基变量由等式确定：
$x_j^{*}=\hat b_j-\hat a_{js}\varepsilon$（循环行 $\hat b_j=0$）。
该点对所有 $\varepsilon$ 满足 $Ax=b$（不要求非负），其目标值用 $D_2$
计算为 $z(\varepsilon)=v+\bar d_s\varepsilon$；用 $D_1$ 计算为
\begin{equation*}
  z(\varepsilon)=v+\bar c_s\,\mathbf{1}[s\in N_1]\,\varepsilon
  \;-\!\!\sum_{\substack{j\in N_1\cap B_2\\ j\ \text{循环}}}\!\!
  \bar c_j\,\hat a_{js}\,\varepsilon ,
\end{equation*}
其中 $\mathbf{1}[\cdot]$ 为示性函数（条件成立取 $1$，否则取 $0$）；
$N_1\cap B_2$ 的循环基变量恰是那些需代入的项（$j$ 在 $D_1$ 中非基
而在 $D_2$ 中为基，故必在循环中换过身份）。两字典表示同一仿射函数
（定义 \ref{def:thlp-dict} 的说明），比较 $\varepsilon$ 的系数得
\begin{equation}\label{eq:thlp-bland-id}
  \bar d_s
  =\bar c_s\,\mathbf{1}[s\in N_1]
  \;-\!\!\sum_{\substack{j\in N_1\cap B_2\\ j\ \text{循环}}}
  \bar c_j\,\hat a_{js}.
\end{equation}
逐项定号：
\begin{itemize}
  \item 若 $s\in N_1$：由 $s<t$ 与 \eqref{eq:thlp-bland-d1}，$\bar c_s\le0$。
  \item 求和项中 $j\ne t$：$j<t$（$j$ 循环且 $j\ne t$），由
        \eqref{eq:thlp-bland-d1} 得 $\bar c_j\le0$，由
        \eqref{eq:thlp-bland-rowsign} 得 $\hat a_{js}\le0$，故
        $-\bar c_j\hat a_{js}\le0$。
  \item 求和项中 $j=t$（注意 $t\in N_1\cap B_2$ 且 $t$ 循环）：
        $\bar c_t>0$（\eqref{eq:thlp-bland-d1}），
        $\hat a_{ts}=\hat a_{rs}>0$（$D_2$ 主元元为正），故
        $-\bar c_t\hat a_{ts}<0$。
\end{itemize}
于是 \eqref{eq:thlp-bland-id} 右端严格小于 $0$，而左端 $\bar d_s>0$
（\eqref{eq:thlp-bland-d2}），矛盾。故循环不可能发生，迭代在有限步内终止。
\end{proof}

\begin{remark}\label{rem:thlp-bland-min-form}
定理 \ref{thm:thlp-bland} 按最大化陈述；最小化问题等价于最大化 $-c^{T}x$，
入基条件变为 $\bar c_j<0$ 中的最小下标，证明逐字适用。对偶单纯形的反循环
通过对偶问题应用同一规则获得：对对偶 LP 而言，``入基/出基''角色互换，
即\textbf{对偶 Bland 规则}要求出基行与入基列都按最小下标平局打破；有限
终止性由定理 \ref{thm:thlp-bland} 作用于对偶 LP 保证。Bland 规则的理论
意义在于\textbf{可证终止}，工程上它速度远慢于 steepest-edge 类规则，因此
实现把它作为检测到循环迹象后的兜底（fallback）而非常态规则。
\end{remark}

\begin{remark}[摄动与移位]\label{rem:thlp-perturbation}
另一类反循环机制是对右端项（原始侧）或目标系数（对偶侧）施加小摄动
$\varepsilon$ 字典序（lexicographic）规则：把平局打破替换为对
$(b+\varepsilon^1,\dots,\varepsilon^m)$ 的字典序比较，可证与原问题在
$\varepsilon\to0$ 极限下一致且终止 \cite{lpchvatal1983}。实现侧的等价物是
代价移位（cost shift）：给目标加确定性的小扰动，使对偶退化被结构性打破；
详见 \srcpath{state.cpp:initialize\_stabilized\_cost} 与
\srcpath{state.cpp:reperturb\_stabilized\_cost}。此外实现还维护基于基签名哈希
的循环检测与 taboo 表（
\srcpath{state.cpp:initialize\_cycle\_guard}、
\srcpath{state.cpp:record\_cycle\_arrival}、
\srcpath{state.cpp:is\_taboo\_row}），属于对 Bland 保证的工程化补充。
\end{remark}

\subsubsection{定价：Dantzig、steepest edge、Devex 与 DSE}\label{subsubsec:thlp-pricing}

定价（pricing）指在每次迭代中选择入基列（原始）或出基行（对偶）的规则。
Dantzig 原始规则选 $\bar c_j$ 最负者；其单位工作量小但迭代数多。下列规则
以``沿边的目标改进率''为准则，显著减少迭代数。

\begin{definition}[steepest edge]\label{def:thlp-se}
对原始单纯形，非基列 $j$ 的\textbf{边方向}为
$\eta_j=\bigl(-B^{-1}a_j;\ e_j\bigr)\in\mathbb{R}^{n}$（沿 $x_j$ 增大而保持
等式的位移方向）。steepest edge 规则选取使
$|\bar c_j|/\|\eta_j\|_2$ 最大的入基列 \cite{lpgoldfarb1977}。权重
$w_j=\|\eta_j\|_2^2$ 称为 edge weight。对偶单纯形的对偶版本（DSE）为出基行
$i$ 定义权重 $w_i=\|e_i^{T}B^{-1}\|_2^2$（结构列部分），选取使
原始违反量与 $\sqrt{w_i}$ 之比最大的出基行
\cite{lpforrest1992}。
\end{definition}

朴素地每次迭代重算全部权重代价过高。steepest edge 的实用性来自下列
精确更新递推。

\begin{theorem}[steepest edge 权重的换基更新]\label{thm:thlp-se-update}
设换基把列 $q$ 换入位置 $p$（基矩阵 $B'=B+(a_q-Be_p)e_p^{T}$）。记
$u=B^{-1}a_q$（主元列，FTRAN 结果），$u_p=e_p^{T}u\ne0$，
$t=B^{-T}e_p$（出基行的 BTRAN）。则对 $i\ne p$，
\begin{equation}\label{eq:thlp-se-rowupdate}
  e_i^{T}(B')^{-1}
  =e_i^{T}B^{-1}-\frac{u_i}{u_p}\,e_p^{T}B^{-1},
  \qquad
  e_p^{T}(B')^{-1}=\frac{1}{u_p}\,e_p^{T}B^{-1},
\end{equation}
从而对偶侧行权重满足精确递推
\begin{equation}\label{eq:thlp-se-wupdate}
  w_i'=w_i-2\,\frac{u_i}{u_p}\,\tau_i+\Bigl(\frac{u_i}{u_p}\Bigr)^{2}w_p,
  \qquad
  w_p'=\frac{w_p}{u_p^{2}},
\end{equation}
其中 $\tau_i=\bigl(e_i^{T}B^{-1}\bigr)t$ 可由一次 BTRAN（$t$）加逐行内积
获得。
\end{theorem}

\begin{proof}
 Sherman--Morrison 公式：$B'=B\bigl(I+(u-e_p)e_p^{T}\bigr)$（因
$B(u-e_p)=a_q-Be_p$），而
$\bigl(I+ve_p^{T}\bigr)^{-1}=I-ve_p^{T}/(1+e_p^{T}v)$（$v=u-e_p$，
$1+e_p^{T}v=u_p$），故
$(B')^{-1}=\bigl(I-\frac{(u-e_p)e_p^{T}}{u_p}\bigr)B^{-1}$。
左乘 $e_i^{T}$（$i\ne p$ 时 $e_i^{T}(u-e_p)=u_i$；$i=p$ 时为 $u_p-1$，
代入化简）即得 \eqref{eq:thlp-se-rowupdate}。对第一式取模方，
\begin{equation*}
  w_i'=\|e_i^{T}(B')^{-1}\|_2^2
  =w_i-2\frac{u_i}{u_p}\bigl(e_i^{T}B^{-1}\bigr)\bigl(B^{-T}e_p\bigr)
   +\Bigl(\frac{u_i}{u_p}\Bigr)^{2}w_p ,
\end{equation*}
中间项正是 $\tau_i$；$w_p'$ 由第二式取模方得到。
\end{proof}

\begin{remark}[计算代价与参考值框架]\label{rem:thlp-se-cost}
定理 \ref{thm:thlp-se-update} 把每次迭代的权重维护降为一个 FTRAN
（$u$，换基本就需要）、一个额外线性求解（$t$ 或等价量）与逐候选内积。
对原始 steepest edge，列权重的对应递推（Goldfarb--Reid
\cite{lpgoldfarb1977}）为
\begin{equation*}
  w_j'=w_j-2\,\frac{\bar a_{pj}}{\bar a_{pq}}\,\bigl(a_j^{T}t'\bigr)
  +\Bigl(\frac{\bar a_{pj}}{\bar a_{pq}}\Bigr)^{2}w_q,
  \qquad t'=B^{-T}u/u_p\ \text{的对应量},
\end{equation*}
结构与 \eqref{eq:thlp-se-wupdate} 相同：``标量比率项 + 与固定向量的内积 +
出基权重二次项''。Harris 的 Devex 规则 \cite{lpharris1973} 放弃精确内积，
只在\textbf{参考值框架}（reference framework，一个固定的行/列子集）上
维护近似权重，换基时用
$w_j'\approx\max\bigl\{w_j,\ (\bar a_{pj}/\bar a_{pq})^{2}w_q\bigr\}$
之类的保守更新，并在框架漂移过大时整体重置。DSE（dual steepest edge）
是定义 \ref{def:thlp-se} 的对偶权重在
\eqref{eq:thlp-se-wupdate} 下的工程实现，常以 Devex 式参考框架初始化并按
需切换为精确计算。实现对应：权重模式枚举
\srcpath{model.hpp:EdgeWeightMode}（SteepestEdge / Devex），初始化
\srcpath{state.cpp:initialize\_exact\_edge\_weights}、
\srcpath{state.cpp:initialize\_devex\_framework}，DSE 权重计算
\srcpath{pricing.cpp:compute\_dse\_weights}，精确权重
\srcpath{factor.cpp:BasisFactor::compute\_exact\_edge\_weights}。
\end{remark}

\subsubsection{数值稳定性}\label{subsubsec:thlp-stability}

单纯形迭代的每步核心线性代数是四次基矩阵求解：FTRAN（$Bu=a_q$，
主元列；及比率/更新所需的更多右端）与 BTRAN（$B^{T}t=e_p$，定价行）。
实现侧的设计原则如下（推导细节见内部文档
\srcpath{docs/archive/numerical\_methods.md}，此处只给结论与条件）：

\begin{enumerate}
  \item \textbf{稀疏 LU + 乘积式更新}：基改变后不重分解，而把换基表示为
        eta 矩阵序列（product-form update）；随 eta 链增长，求解误差与耗时
        均上升，故定期\textbf{重分解}（refactorization，重构
        \srcpath{factor.cpp:BasisFactor::rebuild}，增量更新
        \srcpath{factor.cpp:BasisFactor::update}）。
  \item \textbf{阈值主元}：LU 分解采用 Markowitz 计数择主 + 阈值容差
        （主元不小于列最大元的 $u$ 倍，典型 $u\in[0.01,0.1]$），在稀疏性
        与增长因子之间折中。
  \item \textbf{残差校验与迭代精化}：对线性解 $\hat x$，计算残差
        $r=b-B\hat x$；向后稳定意义下
        \begin{equation}\label{eq:thlp-refinement}
          \frac{\|\hat x-x\|}{\|x\|}
          \ \lesssim\ \kappa(B)\,\frac{\|r\|}{\|b\|},
        \end{equation}
        故以相对残差作求解质量门（
        \srcpath{factor.cpp:BasisFactor::solve\_checked}、
        \srcpath{factor.cpp:BasisFactor::refine\_checked}、
        \srcpath{factor.cpp:BasisFactor::backward\_error\_acceptable}），
        不达标时做一次精化步 $B\delta=r$、$\hat x\leftarrow\hat x+\delta$，
        仍不达标则触发重分解。
  \item \textbf{容差分层}：可行性容差（primal/dual feasibility
        tolerance）与最优性容差分离；比率测试在放宽容差下进行（Harris
        两段，注记 \ref{rem:thlp-harris}），终态审计用严格容差复核
        （\srcpath{state.cpp:audit}）。
\end{enumerate}

\begin{gapnote}
下列对偶单纯形主题为文献中的成熟技术，但在当前代码中尚未实现或仅部分
实现：
（i）\textbf{子优化 / 多重定价}（suboptimization，多入基列同时换基）；
（ii）\textbf{参数化与灵敏度分析}（rhs/cost ranging 的后最优分析）；
（iii）\textbf{字典序摄动的全程维持}（实现以代价移位 + taboo + Bland
兜底组合替代，见注记 \ref{rem:thlp-perturbation}）。
对应表中相应条目标注为 ``未实现''。
\end{gapnote}

\subsection{LP 预求解理论}\label{subsec:thlp-presolve}

预求解（presolve）在调用核心算法前，对 \eqref{eq:thlp-general} 做一系列
保等价性的约简，产出更小的 LP 与一个\textbf{后求解映射}（postsolve），
把约简问题的解映回原问题解。本节按 Andersen \& Andersen
\cite{lpandersen1995} 的分类陈述主要约简规则及其正确性依据。

\subsubsection{活动界与行约简}

\begin{definition}[行活动界]\label{def:thlp-activity}
对行 $i$，其在变量界上的\textbf{最小/最大活动}为区间算术
\begin{equation}\label{eq:thlp-activity}
  a_i^{-}=\sum_{j:a_{ij}>0}a_{ij}l_j+\sum_{j:a_{ij}<0}a_{ij}u_j,
  \qquad
  a_i^{+}=\sum_{j:a_{ij}>0}a_{ij}u_j+\sum_{j:a_{ij}<0}a_{ij}l_j ,
\end{equation}
遇无穷界时相应项使活动界取 $\pm\infty$（约定 $0\cdot\infty=0$，因为
$a_{ij}=0$ 的项不参与求和）。
\end{definition}

\begin{proposition}[行约简规则的正确性]\label{prop:thlp-rowrules}
对行 $b_i^{L}\le a^{i}x\le b_i^{U}$：
\begin{enumerate}
  \item \textbf{空行}：无非零系数。若 $0\notin[b_i^{L},b_i^{U}]$ 则问题
        不可行，否则该行冗余可删。
  \item \textbf{冗余行}：若 $[a_i^{-},a_i^{+}]\subseteq[b_i^{L},b_i^{U}]$，
        行对所有可行 $x$ 自动成立，可删。
  \item \textbf{forcing 行}：若 $a_i^{-}=b_i^{U}$（或 $a_i^{+}=b_i^{L}$），
        则可行迫使每个正系数变量贴上界、负系数变量贴下界（对称情形反之），
        行可删且相应变量固定。
  \item \textbf{singleton 行}（行只有一个非零元 $a_{ij}$）：行约束等价于
        对 $x_j$ 的一个界，并入变量界后删行。
\end{enumerate}
\end{proposition}

\begin{proof}
（1）（2）由定义 \ref{def:thlp-activity} 直接成立：$a_i^{-}\le a^{i}x\le a_i^{+}$
对所有界内 $x$ 成立。（3）以 $a_i^{-}=b_i^{U}$ 为例：可行性要求
$a^{i}x\le b_i^{U}=a_i^{-}$，而 $a^{i}x\ge a_i^{-}$，故 $a^{i}x=a_i^{-}$；
求和 \eqref{eq:thlp-activity} 中每一项 $a_{ij}x_j$ 都已取其区间端点，
任何 $x_j$ 偏离端点都使对应项严格内移、总和严格大于 $a_i^{-}$，矛盾，
故所有变量钉在端点上。（4）$a_{ij}>0$ 时行约束等价于
$b_i^{L}/a_{ij}\le x_j\le b_i^{U}/a_{ij}$（$a_{ij}<0$ 时不等号翻转），
与原有变量界取交即可。
\end{proof}

\subsubsection{列约简}

\begin{proposition}[列约简规则的正确性]\label{prop:thlp-colrules}
\begin{enumerate}
  \item \textbf{固定列}：$l_j=u_j$ 时把 $a_j x_j$ 并入右端项常数，删列。
  \item \textbf{free/implied-free singleton 列}：$x_j$ 只出现在行 $i$。
        若 $x_j$ free（或经活动界分析隐含 free：行 $i$ 除去 $x_j$ 后的剩余
        活动区间允许 $x_j$ 在其整个界内自由移动），则对任意其余变量的取
        值，行 $i$ 总可通过调整 $x_j$ 满足，且最优时 $x_j$ 由行等式唯一
        确定；故可删去行 $i$ 与列 $j$，postsolve 由行方程回代 $x_j$。
  \item \textbf{doubleton 方程代入}：行 $i$ 为等式且只含两列 $j,k$：
        $a_{ij}x_j+a_{ik}x_k=b_i$。解出
        $x_j=(b_i-a_{ik}x_k)/a_{ij}$ 代入目标与其他行，消去 $x_j$ 与行
        $i$，并把 $x_j$ 的界投影为对 $x_k$ 的附加界
        （$b_i-a_{ik}x_k\in a_{ij}[l_j,u_j]$）。
  \item \textbf{singleton 等式列代入}：列 $j$ 唯一出现于等式行 $i$ 时
        同理消元；与（3）的区别在于选择主元 $a_{ij}$ 时以数值稳定性
        （$|a_{ij}|$ 不过小、代入后行侧膨胀受控）为准。
\end{enumerate}
\end{proposition}

\begin{proof}
（1）平凡。（2）设其余变量固定，行 $i$ 对 $x_j$ 的要求为
$a_{ij}x_j\in[b_i^{L}-r_i^{-},b_i^{U}-r_i^{+}]$（$r$ 为剩余活动）。free /
implied-free 条件保证该区间与 $[l_j,u_j]$ 的交非空且含内点区间，故可行性
不依赖 $x_j$ 的具体值；目标中 $c_jx_j$ 一项在最优时把 $x_j$ 推向可行
区间端点或由行方程钉住，两种情况都可由 postsolve 回代恢复。
（3）（4）代入是线性方程组的高斯消元一步：等式约束把 $x_j$ 表达为其余
变量的仿射函数，替换后可行域在该坐标子空间上的投影即约简问题的可行域；
界投影来自 $x_j\in[l_j,u_j]$ 对代入表达式的约束。数值条件（主元大小、
界膨胀上限）不改变正确性，只控制误差传播。
\end{proof}

\begin{remark}[约简的整体正确性原则]\label{rem:thlp-presolve-sound}
每条规则的作用可归纳为三种之一：
（a）\textbf{证书型}：直接给出原问题不可行/无界的证书（如空行）；
（b）\textbf{等价型}：约简问题与原问题可行域之间存在显式仿射双射
（保持目标值），如代入类规则；
（c）\textbf{蕴含型}：删去的约束是其余约束与变量界的逻辑蕴含，如冗余行、
forcing 行。
因此任意规则序列的复合保持：原问题不可行 $\iff$ 约简问题不可行；
否则最优值相等且 postsolve 把约简最优解映回原问题最优解。实现把规则组织
为定点迭代（pass 到无新约简为止），并以 dirty 队列避免全量重扫：
\srcpath{lp\_presolve.cpp:lp\_presolve\_run}（主驱动）、
\srcpath{lp\_presolve.cpp:row\_sweep}（行规则）、
\srcpath{lp\_presolve.cpp:col\_sweep}（列规则）、
\srcpath{lp\_presolve.cpp:implied\_bound\_sweep}（活动界收紧）、
\srcpath{lp\_presolve.cpp:substitution\_sweep} 与
\srcpath{lp\_presolve.cpp:substitute\_singleton\_equality\_column}（代入类）、
\srcpath{lp\_presolve.cpp:postsolve\_primal}（后求解回代）、
\srcpath{lp\_presolve.cpp:rebuild\_activity}（活动界维护）。
\end{remark}

\begin{boundarynote}
预求解的数值语义是\textbf{容差下的逻辑蕴含}：活动界比较、零系数判定都在
缩放后的绝对/相对容差下进行（实现侧
\srcpath{lp\_presolve.cpp:row\_tol} 与
\srcpath{lp\_presolve.cpp:lp\_presolve\_publication\_tol\_scale}）。容差外的
``近冗余'' 不作删除，以保证约简不改变可行状态；这与命题
\ref{prop:thlp-rowrules} 的精确算术表述之间的差异是实现层的审慎放大，
不是理论漏洞。
\end{boundarynote}

\subsection{与实现的对应}\label{subsec:thlp-mapping}

\begin{longtable}{@{}p{0.40\textwidth}p{0.56\textwidth}@{}}
\toprule
理论条目 & 实现状态 \\
\midrule
\endhead
箱式有界标准形 \eqref{eq:thlp-boxed} 与非基状态分类
（注记 \ref{rem:thlp-boxed-nonbasic}）
& \implfull{include/mipsolvers/engine/kernel/lp\_kernel/dual\_simplex.hpp:StandardFormLP} \\
等价变换（命题 \ref{prop:thlp-transform}）
& \implpart{实现不显式拆分自由变量/展开松弛，直接在箱式形上运行；行侧不等式在
前端建模层处理} \\
BFS 存在性与基表示（定义 \ref{def:thlp-bfs}）
& \implfull{src/engine/kernel/lp\_kernel/native\_dual/state.cpp:build\_logical\_basis} \\
Farkas 引理与不可行证书（引理 \ref{lem:thlp-farkas}）
& \implfull{src/engine/kernel/lp\_kernel/native\_dual/certificate.cpp:phase\_one\_farkas\_certificate} \\
原始不可行证书（Farkas，对偶侧）
& \implfull{src/engine/kernel/lp\_kernel/native\_dual/certificate.cpp:primal\_infeasibility\_certificate} \\
强对偶/最优性审计（定理 \ref{thm:thlp-strong}、命题
\ref{prop:thlp-boxed-kkt}）
& \implpart{定理本身不``实现''；最优性由终态审计按 \eqref{eq:thlp-boxed-kkt}
三分支符号条件复核}，锚点
\srcpath{src/engine/kernel/lp\_kernel/native\_dual/state.cpp:audit} \\
对偶单纯形主迭代（命题 \ref{prop:thlp-bfrt-std} 的循环）
& \implfull{src/engine/kernel/lp\_kernel/native\_dual/solver.cpp:minor\_iteration} \\
相位组织（dual Phase-1 / Phase-2 与切换）
& \implfull{src/engine/kernel/lp\_kernel/native\_dual/solver.cpp:run\_phase}；
Phase-1 构造
\srcpath{src/engine/kernel/lp\_kernel/native\_dual/state.cpp:initialize\_dual\_phase\_one}，
切换
\srcpath{src/engine/kernel/lp\_kernel/native\_dual/state.cpp:transition\_dual\_phase\_one\_to\_two} \\
出基行选择（DSE merit + 堆）
& \implfull{src/engine/kernel/lp\_kernel/native\_dual/pricing.cpp:choose\_leaving} \\
箱式 BFRT 比率测试 \eqref{eq:thlp-bfrt-boxed}
& \implfull{src/engine/kernel/lp\_kernel/native\_dual/pricing.cpp:choose\_entering\_bfrt} \\
Harris 两段比率测试（注记 \ref{rem:thlp-harris}）
& \implpart{原始侧独立函数
\srcpath{src/engine/kernel/lp\_kernel/native\_dual/primal.cpp:choose\_leaving\_harris}；
对偶侧两段结构内嵌于 \srcpath{pricing.cpp:choose\_entering\_bfrt}，无独立函数} \\
Bland 规则（定理 \ref{thm:thlp-bland}）
& \implpart{仅作循环兜底（\texttt{bland\_active} 触发），非常态定价规则}，
锚点
\srcpath{src/engine/kernel/lp\_kernel/native\_dual/pricing.cpp:choose\_leaving\_bland} \\
代价移位摄动（注记 \ref{rem:thlp-perturbation}）
& \implfull{src/engine/kernel/lp\_kernel/native\_dual/state.cpp:initialize\_stabilized\_cost}；
重摄动
\srcpath{src/engine/kernel/lp\_kernel/native\_dual/state.cpp:reperturb\_stabilized\_cost}；
移位起点决策
\srcpath{src/engine/kernel/lp\_kernel/native\_dual/state.cpp:decide\_cost\_shifted\_dual\_start} \\
循环检测与 taboo（注记 \ref{rem:thlp-perturbation}）
& \implfull{src/engine/kernel/lp\_kernel/native\_dual/state.cpp:record\_cycle\_arrival}；
\srcpath{src/engine/kernel/lp\_kernel/native\_dual/state.cpp:is\_taboo\_row} \\
steepest edge 权重递推（定理 \ref{thm:thlp-se-update}）
& \implpart{精确权重按需计算并带重建门控，非每步无条件全量更新}，锚点
\srcpath{src/engine/kernel/lp\_kernel/native\_dual/factor.cpp:BasisFactor::compute\_exact\_edge\_weights}、
\srcpath{src/engine/kernel/lp\_kernel/native\_dual/pricing.cpp:compute\_dse\_weights} \\
Devex 参考值框架（注记 \ref{rem:thlp-se-cost}）
& \implfull{src/engine/kernel/lp\_kernel/native\_dual/state.cpp:initialize\_devex\_framework} \\
FTRAN / BTRAN 与基求解
& \implfull{src/engine/kernel/lp\_kernel/native\_dual/factor.cpp:BasisFactor::ftran}；
\srcpath{src/engine/kernel/lp\_kernel/native\_dual/factor.cpp:BasisFactor::btran} \\
LU 重分解与乘积式更新
& \implfull{src/engine/kernel/lp\_kernel/native\_dual/factor.cpp:BasisFactor::rebuild}；
\srcpath{src/engine/kernel/lp\_kernel/native\_dual/factor.cpp:BasisFactor::update} \\
残差校验与迭代精化 \eqref{eq:thlp-refinement}
& \implfull{src/engine/kernel/lp\_kernel/native\_dual/factor.cpp:BasisFactor::solve\_checked}；
\srcpath{src/engine/kernel/lp\_kernel/native\_dual/factor.cpp:BasisFactor::refine\_checked} \\
原始单纯形（作为重启/后备路径）
& \implfull{src/engine/kernel/lp\_kernel/native\_dual/primal.cpp:run\_primal\_phase}；
无界射线证书
\srcpath{src/engine/kernel/lp\_kernel/native\_dual/primal.cpp:certify\_primal\_unbounded\_ray} \\
预求解定点迭代与规则（命题 \ref{prop:thlp-rowrules}、
\ref{prop:thlp-colrules}）
& \implfull{src/engine/presolve/lp\_presolve.cpp:lp\_presolve\_run}；
行规则 \srcpath{src/engine/presolve/lp\_presolve.cpp:row\_sweep}；
列规则 \srcpath{src/engine/presolve/lp\_presolve.cpp:col\_sweep} \\
活动界维护（定义 \ref{def:thlp-activity}）
& \implfull{src/engine/presolve/lp\_presolve.cpp:rebuild\_activity}；
增量更新 \srcpath{src/engine/presolve/lp\_presolve.cpp:update\_column\_activity} \\
代入类约简（命题 \ref{prop:thlp-colrules}（3）（4））
& \implfull{src/engine/presolve/lp\_presolve.cpp:substitution\_sweep}；
singleton 等式列
\srcpath{src/engine/presolve/lp\_presolve.cpp:substitute\_singleton\_equality\_column} \\
postsolve 回代
& \implpart{当前实现恢复原始变量与行乘子主路径}，锚点
\srcpath{src/engine/presolve/lp\_presolve.cpp:postsolve\_primal} \\
子优化/多重定价、灵敏度分析（gapnote）
& \implnone \\
\bottomrule
\end{longtable}

\subsection{参考文献}

\begin{thebibliography}{19}\small
\bibitem{lpdantzig1963}
G. B. Dantzig, \emph{Linear Programming and Extensions}, Princeton University
Press, 1963.
\bibitem{lpchvatal1983}
V. Chv\'atal, \emph{Linear Programming}, W. H. Freeman, 1983.
\bibitem{lpbertsimas1997}
D. Bertsimas, J. N. Tsitsiklis, \emph{Introduction to Linear Optimization},
Athena Scientific, 1997.
\bibitem{lpbland1977}
R. G. Bland, New finite pivoting rules for the simplex method,
\emph{Mathematics of Operations Research} 2(2):103--107, 1977.
\bibitem{lpbeale1955}
E. M. L. Beale, Cycling in the dual simplex algorithm,
\emph{Naval Research Logistics Quarterly} 2(4):269--275, 1955.
\bibitem{lpharris1973}
P. M. J. Harris, Pivot selection methods of the Devex LP code,
\emph{Mathematical Programming} 5(1):1--28, 1973.
\bibitem{lpgoldfarb1977}
D. Goldfarb, J. K. Reid, A practicable steepest-edge simplex algorithm,
\emph{Mathematical Programming} 12(1):361--371, 1977.
\bibitem{lpforrest1992}
J. J. H. Forrest, D. Goldfarb, Steepest-edge simplex algorithms for linear
programming, \emph{Mathematical Programming} 57(1):341--374, 1992.
\bibitem{lpandersen1995}
E. D. Andersen, K. D. Andersen, Presolving in linear programming,
\emph{Mathematical Programming} 71(2):221--245, 1995.
\bibitem{lpmaros2003}
I. Maros, \emph{Computational Techniques of the Simplex Method}, Kluwer
Academic Publishers, 2003.
\end{thebibliography}


% theory_interior_point.tex —— engine 通用理论：原始-对偶内点法
% （障碍问题与中心路径、牛顿方向与预测-校正、KKT 线性代数、
%   惯性校正与正则化、滤波线搜索与可行性恢复、步长与终止）
% 本文件为理论层章节，遵循 docs/modules/THEORY_WRITING_GUIDE.md。

\section{原始-对偶内点法：中心路径、牛顿系统与全局化}\label{sec:thipm}

\begin{theorynote}
本章为通用数学理论，按原始-对偶内点法（primal-dual interior-point method,
IPM）领域的标准模型与文献体系撰写，覆盖：对数障碍问题与中心路径（含 LP
中心路径存在唯一性定理及其证明）、原始-对偶牛顿方向、Mehrotra 预测-校正、
增广 KKT 系统与 normal equations 的推导与权衡、惯性校正与正则化
（W\"achter--Biegler 框架）、滤波线搜索与可行性恢复、步长规则与终止判据，
以及 LP 与 NLP 两种形式的统一视角。本章公式不要求逐式对应代码；与平台实现
（\srcpath{src/engine/kernel/ipm/}、\srcpath{src/engine/kernel/kkt/}）的
对应关系见本章末``与实现的对应''表，超出实现的内容以``未实现扩展说明''框
就地标记。
\end{theorynote}

\subsection{记号与约定}\label{subsec:thipm-notation}

\begin{itemize}
  \item NLP 标准写法：
        \begin{equation}\label{eq:thipm-nlp}
          \min_{x\in\mathbb{R}^{n}}\; f(x)
          \quad\text{s.t.}\quad c_E(x)=0,\quad c_I(x)\ge0,
        \end{equation}
        $f$、$c_E$、$c_I$ 二次连续可微；等式 $m_E$ 个、不等式 $m_I$ 个。
        引入松弛 $s=c_I(x)\ge0$ 后不等式化为 $c_I(x)-s=0$、$s\ge0$。
  \item LP 特例：$f(x)=c^{T}x$，$c_E(x)=Ax-b$，无 $c_I$ 时变量界
        $x\ge0$ 充当不等式；写为
        \begin{equation}\label{eq:thipm-lp}
          \min_{x}\;c^{T}x\quad\text{s.t.}\quad Ax=b,\quad x\ge0 .
        \end{equation}
  \item 大写对角矩阵约定：$X=\mathrm{diag}(x)$，$Z=\mathrm{diag}(z)$，
        $S=\mathrm{diag}(s)$；$e=(1,\dots,1)^{T}$。
  \item 对偶变量：等式乘子 $y$（NLP 中也记 $\lambda$），界/不等式乘子
        $z\ge0$；Lagrange 函数
        $\mathcal{L}(x,y,z)=f(x)-y^{T}c_E(x)-z^{T}c_I(x)$（LP 情形
        $\mathcal{L}=c^{T}x-y^{T}(Ax-b)-z^{T}x$）。
  \item barrier parameter $\mu>0$；对偶间隙度量 $\mu$ 的 LP 形式为
        $x^{T}z/n$，NLP 形式为 $s^{T}z/m_I$（归一化按互补对个数）。
  \item KKT 残差：原始可行性 $r_p$、对偶驻定性 $r_d$、互补性 $r_c$；
        终止判据对三者分别取缩放后的无穷范数。
\end{itemize}

\subsection{障碍问题与中心路径}\label{subsec:thipm-central-path}

\begin{definition}[对数障碍问题]\label{def:thipm-barrier}
对 \eqref{eq:thipm-lp}，barrier parameter $\mu>0$ 的对数障碍问题为
\begin{equation}\label{eq:thipm-barrier}
  \min_{x}\; \varphi_\mu(x)=c^{T}x-\mu\sum_{j=1}^{n}\ln x_j
  \quad\text{s.t.}\quad Ax=b,\quad x>0 .
\end{equation}
对 NLP \eqref{eq:thipm-nlp}（松弛化后）对应为
\begin{equation}\label{eq:thipm-barrier-nlp}
  \min_{x,s}\; f(x)-\mu\sum_{i=1}^{m_I}\ln s_i
  \quad\text{s.t.}\quad c_E(x)=0,\quad c_I(x)-s=0,\quad s>0 .
\end{equation}
障碍项在边界 $x_j\downarrow0$（或 $s_i\downarrow0$）处趋于 $+\infty$，
把不等式约束``软化''进目标。
\end{definition}

\begin{definition}[扰动 KKT 系统与中心路径]\label{def:thipm-perturbed-kkt}
\eqref{eq:thipm-barrier} 的驻点条件（一阶最优性）写成原始-对偶形式：
\begin{equation}\label{eq:thipm-primaldual}
  Ax=b,\quad x>0;\qquad
  A^{T}y+z=c,\quad z>0;\qquad
  XZe=\mu e .
\end{equation}
对每个 $\mu>0$，\eqref{eq:thipm-primaldual} 的解
$(x(\mu),y(\mu),z(\mu))$（若存在唯一）形成的轨迹称为\textbf{中心路径}
（central path）。第三式称为扰动互补条件：$\mu=0$ 时退化为 LP 的互补松弛
条件。
\end{definition}

\begin{lemma}[障碍问题的严格凸性]\label{lem:thipm-strictconvex}
在 $\mathcal{F}^{\circ}=\{x:Ax=b,\ x>0\}$ 上，$\varphi_\mu$ 沿可行方向
严格凸：对任意 $x\in\mathcal{F}^{\circ}$ 与 $d\ne0$、$Ad=0$，
$d^{T}\nabla^{2}\varphi_\mu(x)\,d>0$。
\end{lemma}

\begin{proof}
$\nabla^{2}\varphi_\mu(x)=\mu X^{-2}$（线性项 $c^{T}x$ 无曲率），
$X^{-2}=\mathrm{diag}(x_j^{-2})$ 正定，故
$d^{T}\nabla^{2}\varphi_\mu d=\mu\sum_j d_j^{2}/x_j^{2}>0$（$d\ne0$）。
\end{proof}

\begin{theorem}[LP 中心路径的存在唯一性]\label{thm:thipm-centralpath}
设原始问题 \eqref{eq:thipm-lp} 与其对偶
$\max\{b^{T}y:A^{T}y+z=c,\ z\ge0\}$ 都存在严格可行点（即存在
$x>0$、$Ax=b$ 与 $(y,z)$、$z>0$、$A^{T}y+z=c$）。则对每个 $\mu>0$，
扰动 KKT 系统 \eqref{eq:thipm-primaldual} 存在唯一解
$(x(\mu),y(\mu),z(\mu))$；且对偶间隙
\begin{equation}\label{eq:thipm-gap}
  c^{T}x(\mu)-b^{T}y(\mu)=x(\mu)^{T}z(\mu)=n\mu .
\end{equation}
特别地，沿中心路径 $c^{T}x(\mu)\to z^{*}$（原问题最优值），$\mu\downarrow0$。
\end{theorem}

\begin{proof}
\textbf{存在性}：固定对偶严格可行点 $(\bar y,\bar z)$，$\bar z>0$。对任意
$x\in\mathcal{F}^{\circ}$，
$c^{T}x=\bar y^{T}Ax+\bar z^{T}x=b^{T}\bar y+\bar z^{T}x$，故
\begin{equation*}
  \varphi_\mu(x)
  =b^{T}\bar y+\sum_{j=1}^{n}\bigl(\bar z_j x_j-\mu\ln x_j\bigr).
\end{equation*}
一元函数 $t\mapsto \bar z_j t-\mu\ln t$（$t>0$）在 $t=\mu/\bar z_j$ 处取
最小值 $\mu-\mu\ln(\mu/\bar z_j)$，且当 $t\downarrow0$ 或 $t\to\infty$ 时
趋于 $+\infty$。因此 $\varphi_\mu$ 在 $\mathcal{F}^{\circ}$ 上下方有界，
且任一子水平集 $\{x\in\mathcal{F}^{\circ}:\varphi_\mu(x)\le M\}$ 有界
（每个分量 $x_j$ 都被控制）且与边界保持距离；又
$\mathcal{F}^{\circ}$ 相对闭包中 $\varphi_\mu$ 连续延拓到边界时取
$+\infty$，故 $\varphi_\mu$ 在 $\mathcal{F}^{\circ}$ 上取得最小值
$x(\mu)$。\textbf{唯一性}：由引理 \ref{lem:thipm-strictconvex}，
$\varphi_\mu$ 在凸集 $\mathcal{F}^{\circ}$ 上严格凸，最小值点唯一。
\textbf{对偶变量的存在与唯一}：$x(\mu)$ 满足驻点条件：存在 $y$ 使
$\nabla\varphi_\mu(x)=c-\mu X^{-1}e=A^{T}y$（等式约束的乘子法则，$A$
行满秩时 $y$ 唯一；否则 $y$ 不唯一但 $z$ 唯一）。令
$z=\mu X^{-1}e>0$，则 $A^{T}y+z=c$ 且 $XZe=\mu e$。反之任何
\eqref{eq:thipm-primaldual} 的解给出 $\varphi_\mu$ 的驻点即最小点。
\textbf{间隙}：$x^{T}z=e^{T}XZe=n\mu$，而
$c^{T}x-b^{T}y=(c-A^{T}y)^{T}x=z^{T}x$。令 $\mu\downarrow0$：
由强对偶（对偶可行给出 $b^{T}y(\mu)\le z^{*}\le c^{T}x(\mu)$）夹逼得
$c^{T}x(\mu)\to z^{*}$。
\end{proof}

\begin{remark}\label{rem:thipm-analytic-center}
中心路径当 $\mu\downarrow0$ 时的极限不仅是最优点，而且是最优面的
\textbf{解析中心}（analytic center）：在严格互补意义下，极限点在每个起作用
维度上均衡分布。该结论的证明需要最优面分解（Goldman--Tucker 严格互补定理）
的工具，超出本章范围，见 \cite{ipmwright1997} 第 2 章。定理
\ref{thm:thipm-centralpath} 的存在性证明本质上是 Fiacco--McCormick
障碍函数理论 \cite{ipmfiacco1968} 在 LP 上的显式化。
\end{remark}

\begin{gapnote}
\textbf{齐次自对偶嵌入}（homogeneous self-dual embedding, HSD）把可行性
检测与最优性求解统一在一个恒严格可行的增广 LP 中，使 IPM 无需 Phase-1 且
能给出不可行/无界证书 \cite{ipmye1994,ipmandersen2000}。当前实现
（\srcpath{src/engine/kernel/ipm/ipm\_lp\_solver.cpp}）采用不可行起点 +
Mehrotra 修正的经典框架，\textbf{未}实现 HSD；不可行/无界的判定依赖
最优性审计与残差诊断（
\srcpath{ipm\_lp\_solver.cpp:audit\_ipm\_lp\_optimality}），而非嵌入系统的
证书。
\end{gapnote}

\subsection{原始-对偶牛顿方向}\label{subsec:thipm-newton}

中心路径条件 \eqref{eq:thipm-primaldual} 是关于 $w=(x,y,z)$ 的非线性方程组
$F_\mu(w)=0$。IPM 的每次迭代对当前点（不必可行，只要求 $x,z>0$）做一步
阻尼牛顿法。牛顿方向由线性化系统给出：
\begin{equation}\label{eq:thipm-newtonsystem}
  \begin{pmatrix}
    0 & A^{T} & I\\
    A & 0 & 0\\
    Z & 0 & X
  \end{pmatrix}
  \begin{pmatrix}\Delta x\\ \Delta y\\ \Delta z\end{pmatrix}
  =-
  \begin{pmatrix}
    A^{T}y+z-c\\ Ax-b\\ XZe-\mu e
  \end{pmatrix}.
\end{equation}

\begin{lemma}[牛顿系统的可解性]\label{lem:thipm-nonsingular}
若 $A$ 行满秩且 $x>0$、$z>0$，则 \eqref{eq:thipm-newtonsystem} 的系数矩阵
非奇异。
\end{lemma}

\begin{proof}
设 $(\Delta x,\Delta y,\Delta z)$ 在其零空间中。第三行给出
$\Delta z=-X^{-1}Z\Delta x$；代入第一行：
$A^{T}\Delta y=X^{-1}Z\Delta x$。左乘 $\Delta x^{T}$ 并用第二行
$A\Delta x=0$：
$0=\Delta x^{T}A^{T}\Delta y=\Delta x^{T}X^{-1}Z\Delta x$。
$X^{-1}Z$ 为严格正对角阵，故 $\Delta x=0$；于是 $A^{T}\Delta y=0$，
由 $A$ 行满秩得 $\Delta y=0$，进而 $\Delta z=0$。零空间平凡即非奇异。
\end{proof}

\begin{remark}\label{rem:thipm-nlp-newton}
NLP 情形的牛顿系统结构相同，只是 $(1,1)$ 块不再是零：
线性化 Lagrange 驻定性得块
$W=\nabla_{xx}^{2}\mathcal{L}(x,y,z)$，等式/松弛化不等式的 Jacobian 记
$J$（行满秩是 LICQ 的线性代数形式），互补块为
$S\Delta z+Z\Delta s=\mu e-SZe$。消去 $\Delta s$、$\Delta z$ 后得到
增广系统
\begin{equation}\label{eq:thipm-augmented-nlp}
  \begin{pmatrix}
    W+\Sigma & J^{T}\\
    J & 0
  \end{pmatrix}
  \begin{pmatrix}\Delta x\\ \Delta y\end{pmatrix}
  =
  -\begin{pmatrix} r_d'\\ r_p\end{pmatrix},
  \qquad \Sigma=S^{-1}Z\succ0,
\end{equation}
其中 $r_d'$ 吸收了互补残差的回代。与 LP 的本质差别仅在 $W\ne0$：
$W$ 不定或奇异时，$(1,1)$ 块失去拟定性，这是 NLP-IPM 需要惯性校正
（\ref{subsec:thipm-inertia} 小节）而 LP-IPM 不需要的根本原因。
\end{remark}

\subsection{Mehrotra 预测-校正}\label{subsec:thipm-mehrotra}

纯牛顿步（对固定的 $\mu$ 解 \eqref{eq:thipm-newtonsystem}）有两个缺陷：
（i）$\mu$ 的缩减节奏盲目；（ii）线性化忽略 $XZe$ 的二次项，步长常被边界
截断得很小。Mehrotra 预测-校正 \cite{ipmmehrotra1992} 是实践中的标准解法：

\begin{enumerate}
  \item \textbf{预测（affine-scaling）步}：取 $\mu=0$ 解
        \eqref{eq:thipm-newtonsystem}，得方向
        $(\Delta x^{\mathrm{aff}},\Delta y^{\mathrm{aff}},\Delta z^{\mathrm{aff}})$。
        按 fraction-to-boundary 规则计算最大可行步长
        \begin{equation}\label{eq:thipm-ftb}
          \alpha_p^{\mathrm{aff}}
          =\min\Bigl\{1,\ \min_{j:\,\Delta x_j^{\mathrm{aff}}<0}
          -\frac{x_j}{\Delta x_j^{\mathrm{aff}}}\Bigr\},
          \qquad
          \alpha_d^{\mathrm{aff}}
          =\min\Bigl\{1,\ \min_{j:\,\Delta z_j^{\mathrm{aff}}<0}
          -\frac{z_j}{\Delta z_j^{\mathrm{aff}}}\Bigr\},
        \end{equation}
        并估计仿射步后的间隙
        $\mu_{\mathrm{aff}}=(x+\alpha_p^{\mathrm{aff}}\Delta x^{\mathrm{aff}})^{T}
        (z+\alpha_d^{\mathrm{aff}}\Delta z^{\mathrm{aff}})/n$。
  \item \textbf{中心化参数}：$\sigma=(\mu_{\mathrm{aff}}/\mu)^{3}$。
        若仿射步几乎消除间隙（$\mu_{\mathrm{aff}}\ll\mu$），$\sigma$ 接近
        $0$，下一步偏攻最优性；若仿射步受阻，$\sigma$ 接近 $1$，下一步
        偏回中心路径。指数 $3$ 是 Mehrotra 的经验选择，理论上
        $[0,1]$ 内任一连续递减响应都可行。
  \item \textbf{校正步}：以右端项
        \begin{equation}\label{eq:thipm-corrector}
          r_c=XZe-\sigma\mu e+\Delta X^{\mathrm{aff}}\Delta Z^{\mathrm{aff}}e
        \end{equation}
        重解同一系数矩阵的牛顿系统（只需一次新的回代），得总方向。其中
        二阶项 $\Delta X^{\mathrm{aff}}\Delta Z^{\mathrm{aff}}e$ 补偿了
        线性化丢弃的交叉项，这正是``预测-校正''名称的由来。
  \item \textbf{步长与更新}：对总方向按带回退因子 $\tau$ 的
        fraction-to-boundary 取
        $\alpha_p,\alpha_d$（$\tau$ 的典型取法见
        \ref{subsec:thipm-steps} 小节），更新
        $(x,y,z)\leftarrow(x+\alpha_p\Delta x,\ y+\alpha_d\Delta y,\
        z+\alpha_d\Delta z)$。
\end{enumerate}

\begin{remark}[Gondzio 多重中心校正]\label{rem:thipm-gondzio}
预测-校正之后，只要新的校正方向能显著增大步长且不破坏中心化，就可以用
同一次因式分解反复回代，追加多个校正 \cite{ipmgondzio1996}；每个校正的
收益递减，实现以校正次数上限与步长增益门限控制（
\srcpath{ipm\_lp\_solver.cpp:NativeIPMLPAdapter::solve\_lp\_impl} 内的校正
循环）。
\end{remark}

\subsection{KKT 线性代数：增广系统与法方程}\label{subsec:thipm-kkt}

牛顿系统 \eqref{eq:thipm-newtonsystem} 每步都要以新右端求解一次，其线性
代数占 IPM 总耗时的绝大部分。两种标准约简如下。由第三行
$\Delta z=X^{-1}(r_c-Z\Delta x)$ 代入第一行，得\textbf{增广系统}
（augmented system）
\begin{equation}\label{eq:thipm-augmented}
  \underbrace{\begin{pmatrix}
    -D^{-2} & A^{T}\\
    A & 0
  \end{pmatrix}}_{K}
  \begin{pmatrix}\Delta x\\ \Delta y\end{pmatrix}
  =
  \begin{pmatrix} -\hat r_d\\ \hat r_p\end{pmatrix},
  \qquad D^{2}=XZ^{-1}\succ0 ,
\end{equation}
再消去 $\Delta x$ 得\textbf{法方程}（normal equations）
\begin{equation}\label{eq:thipm-normal}
  AD^{2}A^{T}\,\Delta y=\tilde r,
  \qquad
  \Delta x=D^{2}(A^{T}\Delta y-\hat r_d'),\quad
  \Delta z=X^{-1}(r_c-Z\Delta x).
\end{equation}

\begin{lemma}[增广矩阵的惯性]\label{lem:thipm-inertia}
约定惯性三元组 $\mathrm{In}(M)=(\text{正},\text{负},\text{零})$。设
$D^{2}\succ0$ 对角、$A$ 行满秩。则 \eqref{eq:thipm-augmented} 的矩阵
$K$ 是对称拟定（symmetric quasidefinite）矩阵 \cite{ipmvanderbei1995}，
存在不带换主元崩溃的 $LDL^{T}$ 分解（$L$ 单位下三角、$D$ 对角非奇异），
且
\begin{equation}\label{eq:thipm-inertia-lp}
  \mathrm{In}(K)=(m,\ n,\ 0),
\end{equation}
即 $m$ 个正、$n$ 个负特征值、无零特征值。等价地，把第一块行与第一块列同
乘 $-1$ 得到的 congruence 矩阵
$K'=\bigl[\begin{smallmatrix} D^{-2} & -A^{T}\\ -A & 0\end{smallmatrix}\bigr]$
满足 $\mathrm{In}(K')=(n,\ m,\ 0)$。
\end{lemma}

\begin{proof}
取 $P=\bigl[\begin{smallmatrix} I & 0\\ AD^{2} & I\end{smallmatrix}\bigr]$
（下三角、对角线为 $I$，故 $P$ 非奇异），直接验证 congruence
\begin{equation*}
  P K P^{T}
  =
  \begin{pmatrix} -D^{-2} & 0\\ 0 & AD^{2}A^{T}\end{pmatrix}:
\end{equation*}
$(2,1)$ 块 $=-AD^{2}D^{-2}+A=0$；$(1,2)$ 块
$=-D^{-2}D^{2}A^{T}+A^{T}=0$；$(2,2)$ 块 $=AD^{2}A^{T}$。
$-D^{-2}\prec0$ 贡献 $n$ 个负方向；$A$ 行满秩且 $D^{2}\succ0$ 蕴含 Schur
补 $AD^{2}A^{T}\succ0$，贡献 $m$ 个正方向；右端为块对角非奇异矩阵。由
Sylvester 惯性定律，$\mathrm{In}(K)=\mathrm{In}(PKP^{T})=(m,n,0)$。
$K'$ 与 $K$ 相差第一块行列的符号翻转，是一次换符号的 congruence，故
$\mathrm{In}(K')=(n,m,0)$。零惯性为零保证了 $LDL^{T}$ 分解的存在性
（拟定的标准性质，\cite{ipmvanderbei1995} 定理 2.1）。
\end{proof}

\begin{remark}[法方程 vs 增广系统的权衡]\label{rem:thipm-ne-vs-aug}
法方程 \eqref{eq:thipm-normal} 的矩阵 $AD^{2}A^{T}$ 对称正定，可用
Cholesky；但其稀疏模式是 $AA^{T}$ 的模式：$A$ 中任一``稠密列''（如自由
变量拆分、大量界约束对应的列）会在 $AD^{2}A^{T}$ 中产生稠密块，需按
列密度分裂处理。增广系统保持 $A$ 的原始稀疏性，代价是矩阵不定、需要支持
惯性报告的 $LDL^{T}$（Bunch--Kaufman 型或带惯性统计的稀疏多波前/超节点
分解）。$D^{2}=XZ^{-1}$ 的对角元在迭代后期\textbf{严重分离}（贴边界的
$x_j$ 使比值趋 $0$ 或 $\infty$），这是 IPM 线性系统条件数恶化的结构性
来源；两种形式都受影响，但增广形式的误差集中在可通过迭代精化恢复的通道，
法方程形式则直接平方了条件数，$\kappa(AD^{2}A^{T})$ 可达
$O(\kappa^{2})$ 量级。数值讨论见 \cite{ipmwright1997} 第 8 章；实现侧
增广装配为
\srcpath{kkt\_system.cpp:assemble\_augmented\_kkt}，法方程求解为
\srcpath{kkt\_system.cpp:solve\_kkt\_reduced\_sparse}，因式分解复用为
\srcpath{kkt\_system.cpp:factor\_current\_kkt}。
\end{remark}

\subsection{惯性校正与正则化（NLP 情形）}\label{subsec:thipm-inertia}

对 NLP，增广矩阵 \eqref{eq:thipm-augmented-nlp} 的 $(1,1)$ 块
$W+\Sigma$ 不再保证拟定：非凸问题的 $W$ 有负曲率，分解可能出现错误的
惯性或病态主元。W\"achter--Biegler 框架 \cite{ipmwachter2006} 的处理分两层：

\begin{enumerate}
  \item \textbf{惯性判据}：约束化简后的 KKT 矩阵
        \begin{equation}\label{eq:thipm-kkt-nlp}
          K(\delta_w,\delta_c)=
          \begin{pmatrix}
            W+\Sigma+\delta_w I & J^{T}\\
            J & -\delta_c I
          \end{pmatrix}
        \end{equation}
        的 $LDL^{T}$ 惯性应为 $(n,\,m,\,0)$（$n$ 为原始变量含松弛的维数、
        $m$ 为约束数）。其依据与引理 \ref{lem:thipm-inertia} 相同：块消去
        后惯性等于 $(1,1)$ 块与负 Schur 补的惯性之和；惯性正确时，得到的
        原始方向 $\Delta x$ 是 barrier 目标沿约束流形的下降方向
        \cite{ipmwachter2006}（第 3.1 节）。
  \item \textbf{$\delta_w$ 递增协议}：分解后检查惯性；不符则按
        $\delta_w\leftarrow\kappa_w^{+}\delta_w$（首次失败取
        $\delta_w^{0}$）放大 $(1,1)$ 块正则化并重分解；反复失败达上限时
        触发 $\delta_c>0$ 的约束正则化或转入可行性恢复。该协议是
        \eqref{eq:thipm-kkt-nlp} 保持拟定结构（从而可分解）与方向质量
        之间的显式权衡：$\delta_w$ 过大使方向偏向最速下降、收敛变慢。
\end{enumerate}

\begin{remark}[正则化与 PMM 的边界]\label{rem:thipm-pmm-boundary}
把 $\delta_w$、$\delta_c$ 同时取为随 $\mu$ 缩减的正则项并配合不动点内
循环，即得原始-对偶正则化内点法（PDPM/IP-PMM）的骨架；但仅把 $\delta_w$
当数值兜底（惯性不符才启用）\textbf{不是} PMM——前者是算法结构、后者是
数值防护。当前实现属后者：惯性驱动的 $\delta_w$ 升级为
\srcpath{kkt\_system.cpp:factor\_kkt\_inertia\_corrected\_sparse}、
\srcpath{kkt\_system.cpp:solve\_kkt\_inertia\_corrected\_sparse}，参数解析为
\srcpath{kkt\_system.cpp:resolve\_inertia\_settings}，递增规则为
\srcpath{kkt\_system.cpp:increased\_primal\_regularization}；凝聚形式
（condensed）的正定性证书为
\srcpath{kkt\_system.cpp:certify\_primal\_positive\_definite}。
\end{remark}

\subsection{全局化：滤波线搜索与可行性恢复}\label{subsec:thipm-filter}

NLP-IPM 需要全局化机制协调``减少不可行性''与``减少（barrier）目标''两个
冲突目标。定义约束违反度与 barrier 目标
\begin{equation}\label{eq:thipm-theta-phi}
  \theta(x)=\|c_E(x)\|+\|c_I(x)-s\|\quad(\text{约定范数}),
  \qquad
  \phi_\mu(x,s)=f(x)-\mu\sum_{i}\ln s_i .
\end{equation}

\begin{definition}[滤波与可接受性]\label{def:thipm-filter}
滤波（filter）是 $(\theta,\phi)$ 平面上互不支配的点集
\cite{ipmfletcher2002,ipmwachter2006}：试验点 $(\theta_t,\phi_t)$ 被滤波
接受，当且仅当对每个滤波条目 $(\theta_j,\phi_j)$，
\begin{equation}\label{eq:thipm-filter-acc}
  \theta_t\le(1-\gamma_\theta)\,\theta_j
  \qquad\text{或}\qquad
  \phi_t\le\phi_j-\gamma_\phi\,\theta_j ,
\end{equation}
其中 $\gamma_\theta,\gamma_\phi\in(0,1)$ 为小常数。条目 $(\theta_j,\phi_j)$
支配所有同时劣于它的点，滤波即这些支配区域的并的补。
\end{definition}

\begin{remark}[切换条件与 Armijo 判据]\label{rem:thipm-switching}
仅滤波可能排斥``目标显著改进''的好步，故当步长足够小且方向对 barrier
目标有足够下降时切换为纯目标判据：设
$m_k(\alpha)=\nabla\phi_\mu^{T}d(\alpha)$ 为方向导数（线性模型），若
\begin{equation}\label{eq:thipm-switching}
  m_k(\alpha)<0
  \quad\text{且}\quad
  \alpha\,[-m_k(\alpha)]^{s_\phi}>\delta\,[\theta_k]^{s_\theta}
  \quad(\delta>0,\ s_\theta>1,\ s_\phi\ge1),
\end{equation}
则放弃滤波检验，改用 Armijo 充分下降
\begin{equation}\label{eq:thipm-armijo}
  \phi_\mu(x_k+\alpha d_k)\le\phi_\mu(x_k)+\eta_\mu\,\alpha\,m_k(\alpha),
  \qquad \eta_\mu\in(0,\tfrac12),
\end{equation}
并要求试验点不劣化 $\theta$ 超过滤波上界
$\theta_{\max}$（滤波初始化/扩张机制）。若连 Armijo 也失败，步长回退
（backtracking）；回退到低于阈值仍无进展则进入可行性恢复。
\end{remark}

\begin{theorem}[滤波线搜索 IPM 的全局收敛性（陈述）]\label{thm:thipm-filter-conv}
在 \cite{ipmwachter2006} 的标准假设下（约束函数与目标在迭代轨道所经水平
集上连续可微且梯度 Lipschitz、乘子有界、KKT 矩阵分解满足惯性判据与有界
性条件），上述滤波 + 切换 + Armijo + 可行性恢复的 IPM 满足：要么有限步
终止于近似 KKT 点/近似不可行证书，要么任一极限点要么是 KKT 点，要么是
约束违反度函数 $\theta$ 的驻点（即恢复阶段的收敛点）。
\end{theorem}

\begin{proof}
完整证明（\cite{ipmwachter2006} 第 4--5 节，约二十页的精细分析）超出本章
范围；此处只勾勒其结构：滤波机制保证 $\theta$ 与 $\phi$ 至少其一严格
单调有下界，故 $\{(\theta_k,\phi_k)\}$ 收敛；切换条件
\eqref{eq:thipm-switching} 保证当 $\theta_k\to0$ 后目标序列本身充分下降；
可行性恢复阶段把``无法继续''的迭代转化为 $\theta$ 的一个光滑极小化子
问题，其收敛性由标准无约束/界约束优化理论给出。三部分拼接即得定理。
\end{proof}

\begin{definition}[可行性恢复]\label{def:thipm-restoration}
当线搜索无法找到可接受步长时，暂时放弃原目标，求解可行性恢复问题
\begin{equation}\label{eq:thipm-restoration}
  \min_{x,s}\; \theta(x)\ \Bigl(=\|c_E(x)\|+\|c_I(x)-s\|\Bigr)
  \quad\text{的正则化光滑形式},
\end{equation}
实现上把 $\theta$ 的不可微范数用附加非负变量
（$c_E(x)=p-n$，$p,n\ge0$ 等）重写为光滑 NLP，目标取
$\zeta\|D_R(x-x_R)\|^{2}+\eta\sum(p_i+n_i)$ 一类的带锚点正则项
（$x_R$ 为进入恢复时的参考点），仍用同一 IPM 求解
\cite{ipmwachter2006}（第 3.3 节）。恢复收敛后以其解（或不可行证书）回到
主问题。
\end{definition}

实现对应：滤波数据结构 \srcpath{ipm\_filter.cpp:Filter::is\_acceptable}、
\srcpath{ipm\_filter.cpp:Filter::add\_entry}、
\srcpath{ipm\_filter.cpp:Filter::reset\_with\_theta\_upper\_bound}；
$\theta/\phi$ 与斜率计算
\srcpath{ipm\_solver.cpp:compute\_theta}、
\srcpath{ipm\_solver.cpp:compute\_barrier\_phi}、
\srcpath{ipm\_solver.cpp:barrier\_descent\_slope}；试验点评估
\srcpath{ipm\_solver.cpp:evaluate\_filter\_trial\_values}；主循环
\srcpath{ipm\_solver.cpp:solve\_nlp\_filter\_impl}；恢复 NLP 构造
\srcpath{ipm\_restoration.cpp:build\_restoration\_nlp} 与暖启动恢复
\srcpath{ipm\_restoration.cpp:recover\_restoration\_warm\_start}。

\subsection{步长规则与终止判据}\label{subsec:thipm-steps}

\textbf{Fraction-to-boundary}：方向 $d$ 确定后，步长按
\eqref{eq:thipm-ftb} 的形式取
$\alpha=\min\{1,\ \tau\alpha_{\max}\}$，其中 $\alpha_{\max}$ 是保持
$x+\alpha\Delta x\ge0$（及 $z+\alpha\Delta z\ge0$）的最大步长，回退因子
\begin{equation}\label{eq:thipm-tau}
  \tau=\max\{\tau_{\min},\ 1-\mu\},\qquad \tau_{\min}\in(0.99,1)
\end{equation}
随 $\mu$ 缩减而趋 $1$：迭代后期允许贴近边界，这是超线性收敛行为的必要
条件之一 \cite{ipmwright1997}。NLP 情形还需满足滤波/Armijo 的额外约束，
故 NLP 的步长是 \eqref{eq:thipm-ftb} 与线搜索回溯的复合。

\textbf{终止判据}：以缩放 KKT 残差判定近似最优。记
\begin{equation}\label{eq:thipm-scaled-residual}
  \kappa_d=\frac{\|\nabla f(x)+J^{T}y-z\|_\infty}{s_d},
  \qquad
  \kappa_c=\frac{\|SZe-\mu e\|_\infty}{s_c},
  \qquad
  \kappa_p=\|c_E(x)\|_\infty+\|c_I(x)-s\|_\infty ,
\end{equation}
其中缩放因子（Ipopt 惯例，\cite{ipmwachter2006} 第 5 节）
$s_d=\max\{s_{\max},\ (\|y\|_1+\|z\|_1)/(m+n)\}/s_{\max}$，
$s_c=\max\{s_{\max},\ \|z\|_1/n\}/s_{\max}$（$s_{\max}\ge1$ 为常数）。
当 $\max\{\kappa_d,\kappa_c,\kappa_p\}\le$ tol 时声明收敛；另设
``acceptable'' 层级（容差放宽若干数量级）与发散/不可行/资源中断的
分支判定。缩放的意义：把``残差绝对小''改为``相对当前乘子尺度小''，避免
大乘子问题被绝对容差错误拒绝。实现对应：
\srcpath{ipm\_solver.cpp:evaluate\_ipm\_termination}、
\srcpath{ipm\_solver.cpp:apply\_optimality\_scaling}、
\srcpath{ipm\_solver.cpp:strict\_termination}、
\srcpath{ipm\_solver.cpp:acceptable\_termination}；LP 侧的独立审计为
\srcpath{ipm\_lp\_solver.cpp:audit\_ipm\_lp\_optimality}。

\textbf{初始化}：起点须严格内点（$x,s,z>0$ 于被障碍的坐标）；常见做法是
把起点向界内投影（interiorization）并以最小二乘估计初始乘子，实现对应
\srcpath{ipm\_solver.cpp:interiorize\_initial\_point}、
\srcpath{ipm\_solver.cpp:initialize\_cold\_primal\_dual\_state}。

\subsection{LP 与 NLP 形式的统一视角}\label{subsec:thipm-unify}

LP 是 \eqref{eq:thipm-nlp} 中 $f$ 线性、$c_E$、$c_I$ 仿射的特例。统一
框架下 LP-IPM 的简化逐项对应：

\begin{itemize}
  \item $W=\nabla^{2}\mathcal{L}=0$：增广矩阵恒为拟定（引理
        \ref{lem:thipm-inertia}），惯性校正与 $\delta_w$ 协议退化为可选的
        数值防护，而非算法必需；
  \item 约束是仿射的：牛顿系统的线性化是\textbf{精确的}，不存在模型误差，
        因此无需线搜索保证收敛性——步长由 fraction-to-boundary
        \eqref{eq:thipm-ftb} 与中心化参数 $\sigma$ 完全决定，Mehrotra
        预测-校正即是这一自由度的最优实践；
  \item 收敛理论可达多项式复杂度：对 $\mu$ 缩减受控的短步/长步方法有
        $O(\sqrt{n}\,L)$、$O(nL)$ 迭代界（$L$ 为输入位长），见
        \cite{ipmwright1997}；这些界依赖位长模型，对浮点实现只有定性
        指导意义，实践中 IPM 迭代数对问题规模近似 $O(\log n)$ 级别稳定；
  \item NLP 侧 $W\ne0$ 与模型线性化误差是全部复杂性的来源：惯性校正
        （\ref{subsec:thipm-inertia}）、滤波全局化
        （\ref{subsec:thipm-filter}）、拟牛顿近似（当解析 Hessian 不可得，
        实现对应 \srcpath{ipm\_solver.cpp:update\_quasi\_newton\_state}、
        \srcpath{ipm\_solver.cpp:build\_lagrangian\_hessian}）。
\end{itemize}

\begin{gapnote}
下列主题为文献中的成熟技术，当前实现未覆盖或仅部分覆盖：
（i）\textbf{crossover}：IPM 终点向最优基/顶点的转换（单纯形 push 阶段），
实现仅有 active-set KKT 抛光
（\srcpath{ipm\_solver.cpp:try\_active\_set\_kkt\_polish}，部分实现）；
（ii）\textbf{HSD 嵌入}（见 \ref{subsec:thipm-central-path} 节的
gapnote）；
（iii）\textbf{watchdog 非单调线搜索}：允许偶发目标回升以换取长步，
当前滤波实现未包含该机制。
\end{gapnote}

\subsection{与实现的对应}\label{subsec:thipm-mapping}

\begin{longtable}{@{}p{0.40\textwidth}p{0.56\textwidth}@{}}
\toprule
理论条目 & 实现状态 \\
\midrule
\endhead
LP 障碍问题与中心路径（定理 \ref{thm:thipm-centralpath}）
& \implpart{以箱式界的对数障碍 + 不可行起点实现，非可行起点路径跟踪}，锚点
\srcpath{src/engine/kernel/ipm/ipm\_lp\_solver.cpp:NativeIPMLPAdapter::solve\_lp\_impl} \\
Mehrotra 预测-校正（\ref{subsec:thipm-mehrotra}）
& \implfull{src/engine/kernel/ipm/ipm\_lp\_solver.cpp:NativeIPMLPAdapter::solve\_lp\_impl} \\
Gondzio 多重校正（注记 \ref{rem:thipm-gondzio}）
& \implfull{src/engine/kernel/ipm/ipm\_lp\_solver.cpp:NativeIPMLPAdapter::solve\_lp\_impl} \\
牛顿系统可解性（引理 \ref{lem:thipm-nonsingular}）与 LP 牛顿步
& \implfull{src/engine/kernel/ipm/ipm\_lp\_solver.cpp:NativeIPMLPAdapter::solve\_lp\_impl} \\
增广系统装配 \eqref{eq:thipm-augmented}（NLP 侧）
& \implfull{src/engine/kernel/ipm/ipm\_solver.cpp:assemble\_augmented\_newton}；
KKT 层 \srcpath{src/engine/kernel/kkt/kkt\_system.cpp:assemble\_augmented\_kkt} \\
法方程求解 \eqref{eq:thipm-normal}
& \implfull{src/engine/kernel/kkt/kkt\_system.cpp:solve\_kkt\_reduced\_sparse} \\
因式分解复用与结构缓存
& \implfull{src/engine/kernel/kkt/kkt\_system.cpp:factor\_current\_kkt}；
节点级 LP 缓存
\srcpath{src/engine/kernel/ipm/ipm\_lp\_solver\_cached.cpp:prepare\_for\_node\_solves}、
\srcpath{src/engine/kernel/ipm/ipm\_lp\_solver\_cached.cpp:solve\_cached\_node\_lp} \\
惯性校正与 $\delta_w$ 协议（\ref{subsec:thipm-inertia}）
& \implfull{src/engine/kernel/kkt/kkt\_system.cpp:factor\_kkt\_inertia\_corrected\_sparse}；
求解
\srcpath{src/engine/kernel/kkt/kkt\_system.cpp:solve\_kkt\_inertia\_corrected\_sparse}；
参数
\srcpath{src/engine/kernel/kkt/kkt\_system.cpp:resolve\_inertia\_settings}、
\srcpath{src/engine/kernel/kkt/kkt\_system.cpp:increased\_primal\_regularization} \\
凝聚形式正定性证书
& \implfull{src/engine/kernel/kkt/kkt\_system.cpp:certify\_primal\_positive\_definite} \\
滤波与可接受性（定义 \ref{def:thipm-filter}）
& \implfull{src/engine/kernel/ipm/ipm\_filter.cpp:Filter::is\_acceptable}；
条目维护 \srcpath{src/engine/kernel/ipm/ipm\_filter.cpp:Filter::add\_entry} \\
切换条件与 Armijo（注记 \ref{rem:thipm-switching}）
& \implpart{判据实现在主循环内联，无量独立函数}，锚点
\srcpath{src/engine/kernel/ipm/ipm\_solver.cpp:solve\_nlp\_filter\_impl}；
斜率计算 \srcpath{src/engine/kernel/ipm/ipm\_solver.cpp:barrier\_descent\_slope} \\
$\theta/\phi$ 计算 \eqref{eq:thipm-theta-phi}
& \implfull{src/engine/kernel/ipm/ipm\_solver.cpp:compute\_theta}；
\srcpath{src/engine/kernel/ipm/ipm\_solver.cpp:compute\_barrier\_phi} \\
可行性恢复（定义 \ref{def:thipm-restoration}）
& \implfull{src/engine/kernel/ipm/ipm\_restoration.cpp:build\_restoration\_nlp}；
暖启动
\srcpath{src/engine/kernel/ipm/ipm\_restoration.cpp:recover\_restoration\_warm\_start} \\
缩放后终止判据 \eqref{eq:thipm-scaled-residual}
& \implfull{src/engine/kernel/ipm/ipm\_solver.cpp:evaluate\_ipm\_termination}；
\srcpath{src/engine/kernel/ipm/ipm\_solver.cpp:strict\_termination}、
\srcpath{src/engine/kernel/ipm/ipm\_solver.cpp:acceptable\_termination} \\
LP 最优性审计
& \implfull{src/engine/kernel/ipm/ipm\_lp\_solver.cpp:audit\_ipm\_lp\_optimality} \\
起点 interiorization 与冷启动
& \implfull{src/engine/kernel/ipm/ipm\_solver.cpp:interiorize\_initial\_point}；
\srcpath{src/engine/kernel/ipm/ipm\_solver.cpp:initialize\_cold\_primal\_dual\_state} \\
问题缩放（equilibration）
& \implpart{目标/约束梯度范数缩放，非完整 Ruiz 迭代均衡}，锚点
\srcpath{src/engine/kernel/ipm/ipm\_scaling.cpp:compute\_scaling\_factors}、
\srcpath{src/engine/kernel/ipm/ipm\_scaling.cpp:build\_scaled\_nlp\_model} \\
解析/拟牛顿 Hessian
& \implfull{src/engine/kernel/ipm/ipm\_solver.cpp:build\_lagrangian\_hessian}；
拟牛顿 \srcpath{src/engine/kernel/ipm/ipm\_solver.cpp:update\_quasi\_newton\_state} \\
crossover / HSD / watchdog（gapnote）
& \implpart{crossover 仅有 active-set KKT 抛光
\srcpath{src/engine/kernel/ipm/ipm\_solver.cpp:try\_active\_set\_kkt\_polish}；
HSD 与 watchdog} \implnone \\
\bottomrule
\end{longtable}

\subsection{参考文献}

\begin{thebibliography}{19}\small
\bibitem{ipmwright1997}
S. J. Wright, \emph{Primal-Dual Interior-Point Methods}, SIAM, 1997.
\bibitem{ipmmehrotra1992}
S. Mehrotra, On the implementation of a primal-dual interior point method,
\emph{SIAM Journal on Optimization} 2(4):575--601, 1992.
\bibitem{ipmwachter2006}
A. W\"achter, L. T. Biegler, On the implementation of an interior-point
filter line-search algorithm for large-scale nonlinear programming,
\emph{Mathematical Programming} 106(1):25--57, 2006.
\bibitem{ipmfletcher2002}
R. Fletcher, S. Leyffer, Nonlinear programming without a penalty function,
\emph{Mathematical Programming} 91(2):239--269, 2002.
\bibitem{ipmgondzio1996}
J. Gondzio, Multiple centrality corrections in a primal-dual method for
linear programming, \emph{Computational Optimization and Applications}
6(2):137--156, 1996.
\bibitem{ipmnocedal2006}
J. Nocedal, S. J. Wright, \emph{Numerical Optimization}, 2nd ed., Springer,
2006.
\bibitem{ipmfiacco1968}
A. V. Fiacco, G. P. McCormick, \emph{Nonlinear Programming: Sequential
Unconstrained Minimization Techniques}, Wiley, 1968.
\bibitem{ipmvanderbei1995}
R. J. Vanderbei, Symmetric quasidefinite matrices,
\emph{SIAM Journal on Optimization} 5(1):100--113, 1995.
\bibitem{ipmye1994}
Y. Ye, M. J. Todd, S. Mizuno, An $O(\sqrt{n}L)$-iteration homogeneous and
self-dual linear programming algorithm,
\emph{Mathematics of Operations Research} 19(1):53--67, 1994.
\bibitem{ipmandersen2000}
E. D. Andersen, K. D. Andersen, The MOSEK interior point optimizer for
linear programming: an implementation of the homogeneous algorithm, in
\emph{High Performance Optimization} (H. Frenk et al., eds.), Kluwer
Academic Publishers, 197--232, 2000.
\bibitem{ipmforsgren2002}
A. Forsgren, P. E. Gill, M. H. Wright, Interior methods for nonlinear
optimization, \emph{SIAM Review} 44(4):525--597, 2002.
\end{thebibliography}


% theory_conic.tex —— engine 通用理论：锥规划（conic programming）
% 本文件为理论层章节，遵循 docs/modules/THEORY_WRITING_GUIDE.md。

\section{锥规划通用理论}

\begin{theorynote}
本章为通用数学理论，按锥规划领域的标准模型与文献体系撰写，覆盖锥与对偶锥、
自和谐障碍函数（self-concordant barrier）、Nesterov--Todd（NT）缩放、
SOCP/SDP 的 svec/smat 打包、稀疏 SDP 的弦分解（chordal decomposition）原理，
以及原始/对偶不可行证书。本章公式不要求逐式对应代码；与平台实现
（\texttt{src/engine/kernel/ipm/}、\texttt{src/engine/solver/native/}）的对应关系
见本章末"与实现的对应"表，超出实现的内容以"未实现扩展说明"框就地标记。
\end{theorynote}

\subsection{记号与约定}

\begin{itemize}
  \item $\mathbb{R}^n$：$n$ 维欧氏空间；$\mathbb{R}^l_+ = \{u \in \mathbb{R}^l : u \ge 0\}$ 为非负象限。
  \item $\mathcal{Q}^k = \{u \in \mathbb{R}^k : u_0 \ge \|u_1\|_2\}$：$k$ 维二阶锥
        （second-order cone, SOC；亦称 Lorentz cone 或 ice-cream cone），
        其中 $u = (u_0, u_1)$，$u_0 \in \mathbb{R}$，$u_1 \in \mathbb{R}^{k-1}$。
  \item $\mathcal{S}^p$：$p \times p$ 实对称矩阵空间；
        $\mathcal{S}^p_+ \subset \mathcal{S}^p$ 为半定锥（positive-semidefinite cone）。
  \item $\langle X, Y \rangle_F = \operatorname{tr}(X^\top Y)$：Frobenius 内积；
        对 $X, Y \in \mathcal{S}^p$ 有 $\langle X, Y \rangle_F = \operatorname{tr}(XY)$。
  \item $\mathcal{K}$：本章总指三类锥的笛卡尔积
        $\mathbb{R}^l_+ \times \prod_i \mathcal{Q}^{q_i} \times \prod_j \mathcal{S}^{s_j}_+$。
  \item $u \succeq_{\mathcal{K}} 0$（$u \succ_{\mathcal{K}} 0$）：$u$ 属于 $\mathcal{K}$
        （属于其内部 $\operatorname{int}\mathcal{K}$）。
  \item $\operatorname{int}$、$\operatorname{cl}$：拓扑内部与闭包；$I$：单位阵；$e$：Jordan 代数单位元。
\end{itemize}

\subsection{锥与对偶锥}

\begin{definition}[锥与真锥]
集合 $\mathcal{C} \subseteq \mathbb{E}$（$\mathbb{E}$ 为有限维欧氏空间）称为\textbf{锥}（cone），
若 $u \in \mathcal{C}$ 且 $\alpha \ge 0$ 蕴含 $\alpha u \in \mathcal{C}$。若 $\mathcal{C}$ 还是
闭的、凸的、有非空内部（solid）且不含直线（pointed，即
$\mathcal{C} \cap (-\mathcal{C}) = \{0\}$），则称为\textbf{真锥}（proper cone）。
真锥在 $\mathbb{E}$ 上诱导偏序 $u \succeq_{\mathcal{C}} v \iff u - v \in \mathcal{C}$。
\end{definition}

\begin{definition}[对偶锥]
锥 $\mathcal{C}$ 的\textbf{对偶锥}（dual cone）为
\[
\mathcal{C}^* = \{ z \in \mathbb{E} : \langle z, u \rangle \ge 0 \;\; \forall\, u \in \mathcal{C} \}.
\]
若 $\mathcal{C}^* = \mathcal{C}$，则称 $\mathcal{C}$ \textbf{自对偶}（self-dual）。
\end{definition}

对偶锥总是闭凸锥；当真锥 $\mathcal{C}$ 为闭凸锥时 $\mathcal{C}^{**} = \mathcal{C}$
（双极定理，见 \cite{boyd2004} 第 2.6 节）。笛卡尔积的对偶为各因子对偶的笛卡尔积。

\begin{proposition}[三类锥的自对偶性]
\label{prop:conic-self-dual}
非负象限 $\mathbb{R}^l_+$、二阶锥 $\mathcal{Q}^k$ 与半定锥 $\mathcal{S}^p_+$ 均为自对偶真锥。
\end{proposition}
\begin{proof}
逐锥验证 $\mathcal{C} \subseteq \mathcal{C}^*$ 与 $\mathcal{C}^* \subseteq \mathcal{C}$。

（i）$\mathbb{R}^l_+$：$z \ge 0$ 时显然 $z^\top s \ge 0$ 对一切 $s \ge 0$ 成立；
反之若某分量 $z_i < 0$，取 $s = e_i \ge 0$ 得 $z^\top s = z_i < 0$，故 $z \notin (\mathbb{R}^l_+)^*$。

（ii）$\mathcal{Q}^k$：设 $z, s \in \mathcal{Q}^k$。由 Cauchy--Schwarz，
\[
z^\top s = z_0 s_0 + z_1^\top s_1 \;\ge\; z_0 s_0 - \|z_1\|\,\|s_1\| \;\ge\; 0,
\]
故 $\mathcal{Q}^k \subseteq (\mathcal{Q}^k)^*$。反之，若 $z_0 < \|z_1\|$，
取 $s = (\|z_1\|, -z_1)$，则 $s \in \mathcal{Q}^k$ 而
$z^\top s = z_0 \|z_1\| - \|z_1\|^2 < 0$；若 $z_0 < 0 \le \|z_1\|$ 类似构造，
故 $(\mathcal{Q}^k)^* \subseteq \mathcal{Q}^k$。

（iii）$\mathcal{S}^p_+$：$S, Z \succeq 0$ 时
$\langle S, Z \rangle_F = \operatorname{tr}(Z^{1/2} S Z^{1/2}) \ge 0$，
其中 $Z^{1/2} S Z^{1/2} \succeq 0$。反之若 $S$ 有负特征值
$\lambda(v) = v^\top S v < 0$（$\|v\| = 1$），取 $Z = vv^\top \succeq 0$
得 $\operatorname{tr}(SZ) = v^\top S v < 0$。
\end{proof}

\begin{remark}
自对偶性意味着原始与对偶松弛向量可以落在\textbf{同一个}锥中，
这是统一内点框架同时跟踪 $s$ 与 $z$ 的代数前提。进一步地，这三类锥还是
\textbf{自尺度}（self-scaled，即齐性对称）锥：对任意内点 $u$ 存在自同构把
$\mathcal{K}$ 映到自身并将 $u$ 映到单位元。自对偶加自尺度正是
Nesterov--Todd 缩放（\ref{subsec:nt-scaling} 小节）存在的充要理论背景
\cite{nesterov1997}。
\end{remark}

\subsection{锥规划标准型与对偶性}

本章采用与 CVXOPT \cite{vandenberghe2010} 一致的原始-对偶标准型：
\begin{equation}
\label{eq:conic-primal-dual}
\begin{aligned}
\text{(P)}\quad & \min_{x \in \mathbb{R}^n}\; c^\top x
  && \text{(D)}\quad & \max_{y, z}\; -h^\top z - b^\top y \\
\text{s.t.}\quad & Gx + s = h,\; Ax = b,\; s \in \mathcal{K}
  && \text{s.t.}\quad & G^\top z + A^\top y + c = 0,\; z \in \mathcal{K},
\end{aligned}
\end{equation}
其中 $G \in \mathbb{R}^{m \times n}$、$A \in \mathbb{R}^{m_{eq} \times n}$，
$\mathcal{K}$ 为三类锥的笛卡尔积。注意 (D) 中 $z$ 与原始松弛 $s$ 同锥——
这依赖命题 \ref{prop:conic-self-dual} 的自对偶性。

\begin{theorem}[弱对偶]
\label{thm:conic-weak-duality}
对 \eqref{eq:conic-primal-dual} 的任意原始可行 $(x, s)$ 与对偶可行 $(y, z)$，
\[
c^\top x - \big(-h^\top z - b^\top y\big) = s^\top z \ge 0.
\]
等号（最优性）当且仅当互补松弛 $s^\top z = 0$ 成立。
\end{theorem}
\begin{proof}
利用 $Gx + s = h$ 与 $Ax = b$ 作恒等变形：
\[
c^\top x = -\big(G^\top z + A^\top y\big)^\top x
 = -z^\top (h - s) - y^\top b = -h^\top z - b^\top y + s^\top z.
\]
$s, z \in \mathcal{K}$ 且 $\mathcal{K}$ 自对偶，故 $s^\top z \ge 0$。
\end{proof}

量 $s^\top z$ 称为\textbf{对偶间隙}（duality gap），它是内点法实际跟踪的收敛度量。

\begin{theorem}[强对偶，条件性结论]
\label{thm:conic-strong-duality}
若 (P) 与 (D) 均\textbf{严格可行}（Slater 条件：存在原始可行点使
$s \in \operatorname{int}\mathcal{K}$，且存在对偶可行点使 $z \in \operatorname{int}\mathcal{K}$），
则两者最优值相等且均可达到，存在最优对 $(x^*, s^*)$、$(y^*, z^*)$ 满足
$(s^*)^\top z^* = 0$。
\end{theorem}
\begin{proof}
证明略；见 \cite{nesterov1994} 或 \cite{boyd2004} 第 5 章。
条件缺一不可：仅单侧严格可行时可能只保证单侧可达，两侧都不严格可行时
间隙可以为正（锥规划的固有困难，LP 中不出现的病态）。
\end{proof}

\subsection{自和谐障碍函数与中心路径}

\begin{definition}[自和谐障碍]
设 $\mathcal{C}$ 为真锥。凸函数 $f : \operatorname{int}\mathcal{C} \to \mathbb{R}$
称为 $\mathcal{C}$ 上参数为 $\nu_f$ 的\textbf{自和谐障碍}
（self-concordant barrier），若：
\begin{enumerate}
  \item $f$ 三阶连续可微，且当 $u$ 趋近边界时 $f(u) \to +\infty$；
  \item 自和谐性：对一切 $u \in \operatorname{int}\mathcal{C}$ 与方向 $d$，
        \[
        \big| D^3 f(u)[d,d,d] \big| \le 2 \big( D^2 f(u)[d,d] \big)^{3/2};
        \]
  \item 对数齐次性（logarithmic homogeneity）：对一切 $u \in \operatorname{int}\mathcal{C}$、$t > 0$，
        \[
        f(tu) = f(u) - \nu_f \ln t.
        \]
\end{enumerate}
常数 $\nu_f$ 称为\textbf{障碍参数}（barrier parameter），也称锥度。
\end{definition}

对数齐次性蕴含两个在算法中反复使用的恒等式（对 $t$ 求导并取 $t = 1$）：
\[
\langle \nabla f(u), u \rangle = -\nu_f, \qquad
\nabla^2 f(u)\, u = -\nabla f(u).
\]

三类锥的标准障碍与障碍参数为
\begin{equation}
\label{eq:conic-barriers}
\begin{array}{lll}
\mathbb{R}^l_+: & f(u) = -\sum_{i=1}^{l} \ln u_i, & \nu_f = l;\\[2pt]
\mathcal{Q}^k: & f(u) = -\ln\big(u_0^2 - \|u_1\|_2^2\big), & \nu_f = 2;\\[2pt]
\mathcal{S}^p_+: & f(U) = -\ln\det U, & \nu_f = p.
\end{array}
\end{equation}
其自和谐性证明见 \cite{nesterov1994}（第 2 章、第 5 章）。
笛卡尔积的障碍为各块障碍之和，障碍参数相加。注意：SOC 障碍的参数为 $2$，
但与 Jordan 代数框架一致的算法约定常按\textbf{每锥计 1} 归一
（见 \ref{subsec:nt-scaling} 小节末），二者只差全局缩放约定，不影响理论。

\begin{definition}[中心路径]
以 $\mu > 0$ 为障碍参数，对 \eqref{eq:conic-primal-dual} 定义扰动 KKT 系统
\[
Gx + s = h, \qquad Ax = b, \qquad G^\top z + A^\top y + c = 0, \qquad
s \circ z = \mu e, \qquad s, z \succ_{\mathcal{K}} 0,
\]
其中 $\circ$ 为各锥对应的 Euclidean Jordan 代数乘积
（象限为逐分量乘积；SOC 见 \ref{subsec:nt-scaling} 小节；
半定锥为对称积 $S \circ Z = \tfrac12(SZ + ZS)$），$e$ 为单位元。
其解曲线 $\{(x(\mu), y(\mu), z(\mu)) : \mu > 0\}$ 称为\textbf{中心路径}
（central path）。
\end{definition}

\begin{theorem}[中心路径的存在性与收敛性，条件性结论]
在定理 \ref{thm:conic-strong-duality} 的严格可行条件下，中心路径对每个
$\mu > 0$ 唯一存在，且 $\mu \to 0$ 时收敛到一对严格互补的原始-对偶最优解。
以障碍参数 $\nu$ 归一的短步原始-对偶路径跟踪算法具有迭代复杂度
$O\!\big(\sqrt{\nu}\, \ln(1/\varepsilon)\big)$。
\end{theorem}
\begin{proof}
证明略；存在唯一性见 \cite{nesterov1994} 第 3 章，多项式复杂度见
\cite{nesterov1994} 与 \cite{nesterov1998}。核心机制：自和谐性给出
Newton 步的二次收敛邻域，对数齐次性把每次 $\mu$ 缩减所需的 Newton 步数
控制在 $O(\sqrt{\nu}\, \delta)$（$\delta$ 为缩减因子）。
\end{proof}

\subsection{Nesterov--Todd 缩放方向}
\label{subsec:nt-scaling}

原始-对偶内点法每次迭代对扰动 KKT 系统取 Newton 方向。对自尺度锥，
Nesterov 与 Todd \cite{nesterov1997,nesterov1998} 证明：对任意严格内部对
$(s, z)$，存在\textbf{唯一}的缩放元素（scaling point）$w \succ_{\mathcal{K}} 0$，
其二次表示（quadratic representation）$P(w)$ 满足
\begin{equation}
\label{eq:nt-scaling}
P(w)\, z = \lambda, \qquad P(w)^{-1} s = \lambda,
\end{equation}
即原始点与对偶点被映到\textbf{同一个}缩放点 $\lambda$。在缩放纵标中，
互补条件线性化后的 Newton 方程对 $ds$、$dz$ 对称，方向由组合算子
$H = P^{-1} D$（$P$ 原始缩放、$D$ 对偶缩放）决定。以下逐锥给出显式构造。

\paragraph{非负象限。}逐分量取
\[
d_i = \sqrt{s_i / z_i}, \qquad \lambda_i = \sqrt{s_i z_i}, \qquad W = \operatorname{diag}(d),
\]
则 $\lambda = W z = W^{-1} s$，组合算子为对角阵
$H = W^\top W = \operatorname{diag}(s_i / z_i)$。

\paragraph{二阶锥。}取 Minkowski 符号阵 $J = \operatorname{diag}(1, -1, \ldots, -1)$，
定义 Jordan 代数
\[
\det(u) = u_0^2 - \|u_1\|^2, \qquad
u \circ v = \big(u^\top v,\; u_0 v_1 + v_0 u_1\big), \qquad
u^{-1} = \frac{Ju}{\det(u)}, \qquad e = (1, 0).
\]
$u \succ_{\mathcal{Q}} 0 \iff u_0 > 0 \land \det(u) > 0$。记
$\rho(u) = \sqrt{\det(u)}$ 并归一化 $\bar s = s/\rho(s)$、$\bar z = z/\rho(z)$，
则归一化 NT 点为
\begin{equation}
\label{eq:nt-soc}
\gamma = \sqrt{\frac{1 + \bar s^\top \bar z}{2}}, \qquad
\bar w = \frac{\bar s + J \bar z}{2\gamma}, \qquad
v = \sqrt[\circ]{\bar w}, \qquad
\beta = \sqrt{\frac{\rho(s)}{\rho(z)}},
\end{equation}
其中 $\sqrt[\circ]{\cdot}$ 为 Jordan 平方根（
$v \circ v = u$ 当 $v_0 = \sqrt{(u_0 + \rho(u))/2}$、$v_1 = u_1/(2v_0)$）。
缩放算子为 $v$ 的二次表示
$Wx = \beta\big(2v(v^\top x) - Jx\big)$，满足定义恒等式
$\lambda = Wz = W^{-1}s$。组合算子有闭式
\begin{equation}
\label{eq:nt-soc-h}
H = \beta^2\big(2\bar w \bar w^\top - J\big), \qquad
H^{-1} = \beta^{-2}\big(2(J\bar w)(J\bar w)^\top - J\big),
\end{equation}
两者都是"$J$ 的秩-1 修正"，作用一次只需 $O(k)$。
推导与恒等式验证见 \cite{vandenberghe2010} 与内部推导文档
\srcpath{docs/archive/conic\_sdp.md}（第 4.2 节；已对照当前代码核验一致）。

\paragraph{半定锥。}设 $S = L_s L_s^\top$、$Z = L_z L_z^\top$ 为 Cholesky 分解，
对 $L_z^\top L_s$ 做奇异值分解 $L_z^\top L_s = U\Sigma V^\top$（$\sigma_i > 0$），定义
\begin{equation}
\label{eq:nt-sdp-r}
R = L_s V \Sigma^{-1/2}, \qquad R^{-T} = L_s^{-T} V \Sigma^{1/2}.
\end{equation}

\begin{proposition}[NT 合同恒等式]
\label{prop:nt-sdp}
式 \eqref{eq:nt-sdp-r} 定义的合同变换同时对角化 $S$ 与 $Z$：
$R^\top Z R = R^{-1} S R^{-T} = \Sigma$。
\end{proposition}
\begin{proof}
由 $L_s^\top L_z = V\Sigma U^\top$（SVD 转置），
\[
R^\top Z R = \Sigma^{-1/2} V^\top (L_s^\top L_z)(L_z^\top L_s) V \Sigma^{-1/2}
= \Sigma^{-1/2} V^\top (V\Sigma U^\top)(U\Sigma V^\top) V \Sigma^{-1/2} = \Sigma.
\]
同理
$R^{-1} S R^{-T} = \Sigma^{1/2} V^\top L_s^{-1} L_s L_s^\top L_s^{-T} V \Sigma^{1/2}
= \Sigma$。
\end{proof}

于是缩放点为对角阵 $\lambda = \operatorname{svec}(\operatorname{diag}(\Sigma))$；
组合算子作用为两次合同乘，$H : X \mapsto (RR^\top) X (RR^\top)$，
$H^{-1} : X \mapsto MXM$ 且 $M = (RR^\top)^{-1}$，故从不显式成形
$\mathbb{R}^{p(p+1)/2}$ 上的稠密矩阵。

\paragraph{锥度与步长参数。}中心路径参数按
$\nu = l + N_q + \sum_j p_j$（$N_q$ 为 SOC 块数）归一，并令
$\mu = s^\top z / \nu$。每个 SOC 块不论尺寸均计 1——这是 Jordan 代数约定
（$\mathcal{Q}^k$ 的秩为 2 但障碍参数约定按块计），与式
\eqref{eq:conic-barriers} 的教科书参数相差常数倍，不影响中心路径本身。

\subsection{svec/smat 打包与 Frobenius 内积}

半定块以 \textbf{svec}（symmetric vectorization）打包为向量以纳入统一框架：
$p \times p$ 对称矩阵按列主序下三角展开为长度 $p(p+1)/2$ 的向量，
且\textbf{非对角元乘以 $\sqrt 2$}：
\begin{equation}
\label{eq:svec}
\operatorname{svec}(X) = \big(X_{11}, \sqrt2 X_{21}, \ldots, \sqrt2 X_{p1},
X_{22}, \sqrt2 X_{32}, \ldots, X_{pp}\big)^\top .
\end{equation}
逆变换 $\operatorname{smat}$ 把打包向量还原为对称矩阵（非对角除以 $\sqrt 2$）。

\begin{proposition}[svec 的等距性]
\label{prop:svec-isometry}
对任意 $S, Z \in \mathcal{S}^p$，
\[
\operatorname{svec}(S)^\top \operatorname{svec}(Z) = \operatorname{tr}(SZ)
= \langle S, Z \rangle_F .
\]
\end{proposition}
\begin{proof}
展开左边：
\[
\operatorname{svec}(S)^\top \operatorname{svec}(Z)
= \sum_{i} S_{ii} Z_{ii} + \sum_{i > j} (\sqrt2 S_{ij})(\sqrt2 Z_{ij})
= \sum_i S_{ii}Z_{ii} + 2\sum_{i>j} S_{ij}Z_{ij}
= \operatorname{tr}(SZ),
\]
末步利用 $S, Z$ 对称把 Frobenius 内积的对角项与成对非对角项分开计数。
\end{proof}

\begin{remark}
等距性的实用后果：锥成员性（$S \succeq 0$ 等价于打包向量属于打包锥）、
障碍梯度/Hessian、NT 合同算子都可以直接在打包向量上表达；
Schur 补元素因此可写成
$\operatorname{svec}(X_i)^\top H^{-1} \operatorname{svec}(X_j)
= \operatorname{tr}(X_i M X_j M)$，完全在矩阵空间计算而无需成形
$H^{-1} \in \mathbb{R}^{p(p+1)/2 \times p(p+1)/2}$。这正是实现
\srcpath{src/engine/kernel/ipm/cones.cpp:svec} 与
\srcpath{src/engine/kernel/ipm/cones.cpp:smat} 所固化的约定。
\end{remark}

\subsection{Newton 方向与 Schur 补结构}

在当前 NT 缩放下线性化互补条件 $\lambda \circ \lambda = \mu e$，得每次迭代
要求解的 $3 \times 3$ 系统
\begin{equation}
\label{eq:conic-newton}
\begin{bmatrix} 0 & A^\top & G^\top \\ A & 0 & 0 \\ G & 0 & -H \end{bmatrix}
\begin{bmatrix} dx \\ dy \\ dz \end{bmatrix}
=
\begin{bmatrix} -r_c \\ -r_x \\ -r_s - W r_\lambda \end{bmatrix},
\qquad
\begin{aligned}
r_c &= G^\top z + A^\top y + c,\\
r_x &= Ax - b,\\
r_s &= Gx + s - h,
\end{aligned}
\end{equation}
其中 $r_\lambda$ 为缩放空间右端（仿射步 $r_\lambda = -\lambda$；Mehrotra
校正步含二阶项 $-L(\lambda)^{-1}(d\lambda_s \circ d\lambda_z) + \sigma\mu\,\lambda^{-1}$，
见 \cite{vandenberghe2010}）。消去 $dz$ 得对称拟定（quasidefinite）系统
\begin{equation}
\label{eq:conic-schur}
\begin{bmatrix} G^\top H^{-1} G + \delta I & A^\top \\ A & -\delta I \end{bmatrix}
\begin{bmatrix} dx \\ dy \end{bmatrix}
=
\begin{bmatrix} -r_c + G^\top H^{-1} t \\ -r_x \end{bmatrix},
\qquad
dz = H^{-1}(G\,dx - t),\quad ds = -r_s - G\,dx ,
\end{equation}
其中 $t = -r_s - W r_\lambda$，$\delta > 0$ 为正则化参数。

\begin{proposition}[拟定正则化的可分解性，条件性结论]
\label{prop:quasidefinite}
设 $H^{-1} \succ 0$（三类锥的 NT 组合算子均满足）且 $\delta > 0$。
则无等式时 \eqref{eq:conic-schur} 的系数阵对称正定，可做 $LL^\top$ 分解；
有等式时系数阵为拟定鞍点阵：对任意对称排列都存在无主元交换的
$LDL^\top$ 分解，且负主元恰好出现在等式块。
\end{proposition}
\begin{proof}
正定性：$G^\top H^{-1}G \succeq 0$（$H^{-1} \succ 0$ 的合同保持），加
$\delta I$ 得严格正定。拟定阵的无主元分解存在性是 Vanderbei 定理，
见 \cite{vanderbei1995}。
\end{proof}

按锥块分解 $B = G^\top H^{-1} G$：象限行贡献对角加权的稀疏外积
$\sum_r d_r^{-2} g_r g_r^\top$；由式 \eqref{eq:nt-soc-h}，SOC 块贡献
"稀疏部分 $-\beta^{-2} G_q^\top J G_q$ 加秩-1 修正 $2\beta^{-2} uu^\top$
（$u = G_q^\top J\bar w$）"；由命题 \ref{prop:svec-isometry}，半定块贡献
稠密团 $(B_s)_{ij} = \operatorname{tr}(X_i M X_j M)$。同一半定块触及的所有
变量在 $B$ 中构成稠密团，这是 SDP 每迭代装配与分解成本的结构来源，
也是下节弦分解的动机。

\subsection{稀疏 SDP 的 Chordal 分解原理}

\begin{definition}[聚合稀疏图]
设 SDP 松弛矩阵 $S(x) = S_0 + \sum_j x_j S_j \in \mathcal{S}^p$。
其\textbf{聚合稀疏图}（aggregate sparsity graph）为无向图
$G_{agg} = (V, E)$，$V = \{1,\ldots,p\}$，
$\{i,j\} \in E$ 当且仅当某个 $S_j$ 或 $S_0$ 在 $(i,j)$ 处非零。
图 $G$ 为\textbf{弦图}（chordal graph），若每个长度 $\ge 4$ 的环都有一条弦。
\end{definition}

\begin{theorem}[Grone--Johnson--S\'a--Wolkowicz，条件性结论]
\label{thm:psd-completion}
设 $G = (V, E)$ 为弦图。定义于 $E$ 与对角上的部分对称矩阵 $X$
可取为半定完备（PSD completable）当且仅当 $X$ 在 $G$ 的每个\textbf{最大团}
（maximal clique）$C_k$ 上的主子阵半定：$X[C_k, C_k] \succeq 0$。
\end{theorem}
\begin{proof}
证明略；见 \cite{grone1984}。非弦图结论不成立——这是必须先做弦扩展
（chordal extension，添加填充边使图成为弦图）的理论原因。最小度等
消元序启发式生成小填充的弦扩展与对应的完美消元序列（perfect
elimination ordering）。
\end{proof}

\begin{theorem}[Agler 分解，条件性结论]
\label{thm:agler}
对弦扩展 $E^+$ 的最大团族 $\{C_k\}$，$S \in \mathcal{S}^p$ 属于
"在 $E^+$ 上可半定完备的锥"当且仅当存在 $S_k \succeq 0$（$S_k$ 为
$|C_k|$ 阶矩阵）使
\begin{equation}
\label{eq:agler}
S = \sum_k E_{C_k}^\top S_k E_{C_k},
\end{equation}
其中 $E_{C_k}$ 为选取团 $C_k$ 坐标行的 0-1 提取矩阵。
\end{theorem}
\begin{proof}
证明略；见 \cite{agler1988} 与 \cite{fukuda2001}。方向"$\Leftarrow$"：
半定矩阵之和半定。方向"$\Rightarrow$"依赖定理 \ref{thm:psd-completion}
沿完美消元序列的递归完备构造。
\end{proof}

\paragraph{精确转换与近似松弛的区别。}把 $S(x)$ 替换为式 \eqref{eq:agler}
右端并补上\textbf{一致性等式}（consistency equalities：和式在每条弦边
——含填充边——上等于原仿射矩阵的对应项），得到与原问题\textbf{精确等价}
的锥规划，其 SDP 块为若干小团块。若仅要求各团主子阵半定而不加一致性等式，
得到的是定理 \ref{thm:psd-completion} 的\textbf{完备化松弛}：填充边上的值
自由，松弛一般严格。实现采用前者（见"与实现的对应"表）；求解后沿完美消元
序列对偶完备（dual completion）恢复原模型的对偶矩阵
\cite{fukuda2001}。弦分解把单个 $p$ 阶稠密块换成若干 $|C_k|$ 阶块，
Schur 装配与分解成本由最大团阶而非 $p$ 决定，代价是新增的团变量与一致性
等式；因此对"最大团接近原阶"的近稠密 SDP 无收益，需结构性门控。

\subsection{原始/对偶不可行检测与证书}

锥规划可能原始不可行或对偶不可行（原始无界）。内点法通过对偶间隙的
发散方向构造渐近 Farkas 型证书。

\begin{theorem}[原始不可行证书]
\label{thm:conic-primal-infeas}
若存在 $(y, z)$ 满足
\[
G^\top z + A^\top y = 0, \qquad z \in \mathcal{K}, \qquad h^\top z + b^\top y < 0,
\]
则 \eqref{eq:conic-primal-dual}(P) 不可行。
\end{theorem}
\begin{proof}
反设存在原始可行 $(x, s)$：$Gx + s = h$、$Ax = b$、$s \in \mathcal{K}$。则
\[
0 < -\big(h^\top z + b^\top y\big)
= -(Gx + s)^\top z - (Ax)^\top y
= -x^\top(G^\top z + A^\top y) - s^\top z = -s^\top z \le 0,
\]
末步用 $s, z \in \mathcal{K}$ 的自对偶非负性，矛盾。
\end{proof}

\begin{theorem}[对偶不可行（原始无界）证书]
\label{thm:conic-dual-infeas}
若存在 $(x, s)$ 满足
\[
Ax = 0, \qquad Gx + s = 0, \qquad s \in \mathcal{K}, \qquad c^\top x < 0,
\]
则 (P) 若可行即无界，从而 (D) 不可行。
\end{theorem}
\begin{proof}
$(x, s)$ 是 (P) 的下降射线：任取原始可行 $(x_0, s_0)$，则对一切
$t \ge 0$，$(x_0 + tx, s_0 + ts)$ 仍可行且目标值
$c^\top x_0 + t\, c^\top x \to -\infty$。故 (P) 无界；由定理
\ref{thm:conic-weak-duality}，(D) 的任何可行点都会给出下界，故 (D) 不可行。
\end{proof}

实践中内点迭代只渐近满足齐次等式，因此实现用\textbf{归一化残差}判据：
当 $h^\top z + b^\top y < 0$ 且
$\|G^\top z + A^\top y\|_\infty / \big(-(h^\top z + b^\top y)\big) \le \varepsilon_{\mathrm{feas}}$
时认定近似原始不可行证书（对偶情形对称，分母为 $-c^\top x$）。
归一化使判据尺度不变：证书射线乘以任意正常数仍满足同一不等式。

\begin{gapnote}
本章证书定理针对精确算术。实现
（\srcpath{src/engine/kernel/ipm/conic\_ipm\_solver.cpp:ConicIPMSolver::solve}）
采用上述 $\infty$-范数归一化近似判据，其"证书"是数值意义下的近似证书，
不构成严格算术下的不可行性证明；需要严格证明时应配合有理算术或
事后验证流程，当前代码未提供该层。
\end{gapnote}

\subsection{与实现的对应}

\begin{longtable}{p{0.42\textwidth} p{0.50\textwidth}}
\toprule
理论条目 & 实现状态与锚点 \\
\midrule
\endhead
三类锥的笛卡尔积标准型 \eqref{eq:conic-primal-dual} &
\implfull{src/engine/kernel/ipm/conic\_ipm\_solver.cpp:ConicIPMSolver::solve}；
问题描述 \srcpath{include/mipsolvers/engine/problem\_types.hpp:ConicModel} \\
锥块布局与锥度 $\nu = l + N_q + \sum p_j$ &
\implfull{src/engine/kernel/ipm/cones.cpp:ConeLayout::build} \\
自对偶性（命题 \ref{prop:conic-self-dual}）&
\implpart{自对偶是算法的代数前提，不作为运行期检查；代码直接假定三锥自对偶} \\
svec/smat 等距打包（命题 \ref{prop:svec-isometry}）&
\implfull{src/engine/kernel/ipm/cones.cpp:svec}、
\srcpath{src/engine/kernel/ipm/cones.cpp:smat} \\
象限 NT 缩放 $H = \operatorname{diag}(s_i/z_i)$ &
\implfull{src/engine/kernel/ipm/cones.cpp:ConeNtScaling::compute}；
尺寸-1 SOC 折叠为象限行 \\
SOC Jordan 代数与 NT 点 \eqref{eq:nt-soc}、$H/H^{-1}$ 闭式 \eqref{eq:nt-soc-h} &
\implfull{src/engine/kernel/ipm/cones.cpp:q\_jordan\_sqrt}、
\srcpath{src/engine/kernel/ipm/cones.cpp:q\_jordan\_solve}、
\srcpath{src/engine/kernel/ipm/cones.cpp:ConeNtScaling::apply\_h}、
\srcpath{src/engine/kernel/ipm/cones.cpp:ConeNtScaling::apply\_h\_inv} \\
半定 NT 合同 \eqref{eq:nt-sdp-r}（命题 \ref{prop:nt-sdp}）&
\implfull{src/engine/kernel/ipm/cones.cpp:ConeNtScaling::compute}；
算子经 \srcpath{src/engine/kernel/ipm/cones.cpp:ConeNtScaling::apply\_w} 等
六个合同成员实现，从不显式成形 $H$ \\
Newton $3\times3$ 系统 \eqref{eq:conic-newton} 与拟定消元
\eqref{eq:conic-schur} & \implfull{src/engine/kernel/ipm/conic\_ipm\_solver.cpp:ConicIPMSolver::solve}
内部的 \texttt{solve\_direction} lambda 与
\srcpath{src/engine/kernel/ipm/conic\_ipm\_solver.cpp:KktAssembler::assemble} \\
拟定正则化与 SPD/LDL$^\top$ 路径（命题 \ref{prop:quasidefinite}）&
\implfull{src/engine/kernel/ipm/conic\_ipm\_solver.cpp:KktAssembler::apply\_delta}、
\srcpath{src/engine/kernel/ipm/conic\_ipm\_solver.cpp:KktBackend}
（CHOLMOD 超节点 $LL^\top$ / simplicial $LDL^\top$ / MUMPS 回退） \\
Mehrotra 预测-校正、步长与最大步公式 &
\implfull{src/engine/kernel/ipm/conic\_ipm\_solver.cpp:ConicIPMSolver::solve}；
逐锥最大步 \srcpath{src/engine/kernel/ipm/cones.cpp:l\_max\_step}、
\srcpath{src/engine/kernel/ipm/cones.cpp:q\_max\_step}、
\srcpath{src/engine/kernel/ipm/cones.cpp:s\_max\_step}；
校正右端 \srcpath{src/engine/kernel/ipm/conic\_ipm\_solver.cpp:combined\_lambda\_rhs} \\
弦分解：定理 \ref{thm:psd-completion}、\ref{thm:agler} 与一致性等式精确转换 &
\implfull{src/engine/kernel/ipm/chordal\_decomposition.cpp:chordal\_decompose}、
\srcpath{src/engine/kernel/ipm/chordal\_decomposition.cpp:chordal\_recover}
（含对偶完备 \srcpath{src/engine/kernel/ipm/chordal\_decomposition.cpp:complete\_dual}）；
未采用完备化松弛路径 \\
不可行证书（定理 \ref{thm:conic-primal-infeas}、\ref{thm:conic-dual-infeas}）&
\implpart{实现为 $\infty$-范数归一化的数值近似判据，见本章 gapnote；
锚点 \srcpath{src/engine/kernel/ipm/conic\_ipm\_solver.cpp:ConicIPMSolver::solve}} \\
指数锥、幂锥等其他自尺度锥 & \implnone \\
适配器层对偶布局 $[z \mid y]$ &
\implfull{src/engine/solver/native/conic\_ipm\_adapter.cpp:NativeConicIPMAdapter::solve\_conic} \\
\bottomrule
\end{longtable}

\subsection{参考文献}

\begin{thebibliography}{9}\small
\bibitem{nesterov1994} Yu. Nesterov, A. Nemirovskii,
\emph{Interior-Point Polynomial Algorithms in Convex Programming}, SIAM, 1994.
\bibitem{nesterov1997} Yu. Nesterov, M. J. Todd,
Self-scaled barriers and interior-point methods for convex programming,
\emph{Mathematics of Operations Research} 22(1):1--42, 1997.
\bibitem{nesterov1998} Yu. Nesterov, M. J. Todd,
Primal-dual interior-point methods for self-scaled cones,
\emph{SIAM Journal on Optimization} 8(2):324--364, 1998.
\bibitem{vandenberghe2010} L. Vandenberghe,
\emph{The CVXOPT linear and quadratic cone program solvers},
技术报告, 2010.
\bibitem{boyd2004} S. Boyd, L. Vandenberghe,
\emph{Convex Optimization}, Cambridge University Press, 2004.
\bibitem{vanderbei1995} R. J. Vanderbei,
Symmetric quasidefinite matrices,
\emph{SIAM Journal on Optimization} 5(1):100--113, 1995.
\bibitem{grone1984} R. Grone, C. R. Johnson, E. M. S\'a, H. Wolkowicz,
Positive definite completions of partial Hermitian matrices,
\emph{Linear Algebra and its Applications} 58:109--124, 1984.
\bibitem{agler1988} J. Agler, J. W. Helton, S. McCullough, L. Rodman,
Positive semidefinite matrices with a given sparsity pattern,
\emph{Linear Algebra and its Applications} 107:101--149, 1988.
\bibitem{fukuda2001} M. Fukuda, M. Kojima, K. Murota, K. Nakata,
Exploiting sparsity in semidefinite programming via matrix completion I:
General framework, \emph{SIAM Journal on Optimization} 11(3):647--674, 2001.
\end{thebibliography}


% theory_milp.tex —— engine 通用理论：混合整数线性规划（MILP）
% 本文件为理论层章节，遵循 docs/modules/THEORY_WRITING_GUIDE.md。

\section{混合整数线性规划通用理论}

\begin{theorynote}
本章为通用数学理论，按 MILP 领域的标准模型与文献体系（专著级）撰写，覆盖
分支定界框架与界定理、LP 松弛与对偶界、强分支与伪成本、可靠性分支
（reliability branching）推导、割平面理论（Chv\'atal--Gomory 割、Gomory
分数割、MIR、CGLP）、domain propagation、clique/cover 不等式、原始启发式
思想与 presolve 约简。本章公式不要求逐式对应代码；与平台实现
（\texttt{src/engine/solver/native/milp/bc/}）的对应关系见本章末
"与实现的对应"表，超出实现的内容以"未实现扩展说明"框就地标记。
平台侧的推导/实验记录见 \srcpath{docs/archive/} 下
\texttt{native\_milp\_*\_2026-08-13.md} 系列，本章引用处均已对照当前代码核验。
\end{theorynote}

\subsection{记号与约定}

\begin{itemize}
  \item MILP 标准型：$\min\{c^\top x : Ax \le b,\; l \le x \le u,\;
        x_j \in \mathbb{Z}\ \forall j \in I\}$，其中
        $A \in \mathbb{R}^{m \times n}$，$I \subseteq \{1,\ldots,n\}$ 为整数
        变量下标集；$B \subseteq I$ 为 0-1 变量下标集。
  \item $X = \{x \in \mathbb{R}^n : Ax \le b,\ l \le x \le u\}$：LP 松弛可行域；
        $X_I = \{x \in X : x_j \in \mathbb{Z}\ \forall j \in I\}$：MILP 可行集；
        $z^* = \min\{c^\top x : x \in X_I\}$（不可行时 $z^* = +\infty$）。
  \item 对节点 $k$ 施加的局部界记 $[l^k, u^k]$，其 LP 松弛最优值记
        $z_{LP}^k$；$\lfloor t \rfloor$、$\lceil t \rceil$ 为下/上取整，
        $\operatorname{frac}(t) = t - \lfloor t \rfloor$ 为小数部。
  \item 行 $i$ 的活动度（activity）：$a_i x$ 在 $[l,u]$ 上的最小/最大可达值
        记 $\alpha_i^{\min}, \alpha_i^{\max}$（\ref{subsec:domain-propagation} 小节定义）。
  \item 对任意实数 $r$ 与整数标量情形，$(r)^+ = \max\{r, 0\}$。
\end{itemize}

\subsection{MILP 模型、LP 松弛与对偶界}

\begin{definition}[LP 松弛]
去掉整数约束 $x_j \in \mathbb{Z}$（$j \in I$）得到的线性规划称为 MILP 的
\textbf{LP 松弛}（LP relaxation）。
\end{definition}

\begin{proposition}[松弛界]
\label{prop:milp-relaxation-bound}
若 $X_I \neq \emptyset$ 且 LP 松弛有有限最优值 $z_{LP}$，则
$z_{LP} \le z^*$；且若 LP 松弛最优解 $\hat x$ 恰满足全部整数约束，
则 $\hat x$ 即 MILP 最优解。
\end{proposition}
\begin{proof}
$X_I \subseteq X$，在更大集合上取最小只会更小：$z_{LP} = \min_{X} c^\top x
\le \min_{X_I} c^\top x = z^*$。若 $\hat x \in X_I$ 则
$c^\top \hat x \ge z^*$，与 $c^\top \hat x = z_{LP} \le z^*$ 合并得等号。
\end{proof}

更一般地，任何对偶可行解都给出合法下界：

\begin{proposition}[对偶界]
\label{prop:milp-dual-bound}
设 $\pi \ge 0$ 满足 $\pi^\top A \le c^\top$ 的分量条件
（对 $Ax \le b$ 型约束的 LP 对偶可行），则
$z^* \ge \pi^\top b + \sum_j \min_{l_j \le x_j \le u_j}
(c_j - \pi^\top A_{\cdot j}) x_j$。特别地 LP 弱对偶给出
$z_{LP} = \max_\pi \{\ldots\} \ge$ 任何对偶可行点的值，故任何对偶可行点的
目标值都是 $z^*$ 的有效下界。
\end{proposition}
\begin{proof}
对任意 $x \in X_I \subseteq X$，
$c^\top x \ge \pi^\top A x + \sum_j (c_j - \pi^\top A_{\cdot j}) x_j
\ge \pi^\top b + \sum_j \min_{[l_j,u_j]}(c_j - \pi^\top A_{\cdot j}) x_j$，
第一步用 $\pi \ge 0$、$Ax \le b$ 与对偶可行性逐项放缩，
第二步把每个分量放缩到其界上的最小值。对 $x$ 取最小即得结论。
\end{proof}

\begin{remark}
命题 \ref{prop:milp-relaxation-bound} 与 \ref{prop:milp-dual-bound} 是
全部"界"的来源：分支定界维护的全局下界、割平面增强的松弛界、以及节点
剪枝（pruning by bound）都归结到"某松弛/对偶可行对象给出的下界不超过
当前 incumbent 值则可丢弃子树"。实现锚点
\srcpath{src/engine/solver/native/milp/bc/legacy/bc\_relaxation.cpp:solve\_lp\_relaxation}。
\end{remark}

\subsection{分支定界框架与界定理}
\label{subsec:bnb}

分支定界（branch and bound, Land--Doig \cite{landdoig1960}）把 MILP 递归分解为
子问题构成的搜索树：每个节点对应附加了局部变量界 $[l^k, u^k]$ 的子问题。

\begin{definition}[节点处理与剪枝]
对节点 $k$ 求解其 LP 松弛。记 incumbent 值（当前最好可行解目标）为
$\bar z$。节点被\textbf{剪枝}（fathomed），若出现以下三者之一：
\begin{enumerate}
  \item LP 松弛不可行（infeasibility）；
  \item $z_{LP}^k \ge \bar z$（by bound：该子树不可能产生更优解）；
  \item LP 松弛最优解满足全部整数约束（by integrality，并可能更新
        $\bar z$）。
\end{enumerate}
否则选取某个分数整数变量 $x_j$（$j \in I$，
$\operatorname{frac}(\hat x_j) \in (0,1)$）\textbf{分支}：生成两个子节点，
分别附加 $x_j \le \lfloor \hat x_j \rfloor$ 与
$x_j \ge \lceil \hat x_j \rceil$。
\end{definition}

\begin{lemma}[分支完备性]
\label{lem:branching-completeness}
设节点 $k$ 的可行集为 $X_I^k$。对分支变量 $x_j$ 生成的两个子节点可行集
$X_I^{k,-}$、$X_I^{k,+}$ 满足
$X_I^k = X_I^{k,-} \cup X_I^{k,+}$（对 $x_j \in \mathbb{Z}$ 的 $x$，
$x_j \le \lfloor\hat x_j\rfloor$ 与 $x_j \ge \lceil\hat x_j\rceil$ 穷尽一切
整数值，且当前 LP 最优解 $\hat x$ 同被两侧排除）。
\end{lemma}
\begin{proof}
$x_j$ 取整数值时必满足 $x_j \le \lfloor\hat x_j\rfloor$ 或
$x_j \ge \lfloor\hat x_j\rfloor + 1 = \lceil\hat x_j\rceil$ 之一
（$\hat x_j$ 分数，故 $\lfloor\hat x_j\rfloor < \hat x_j < \lceil\hat x_j\rfloor$），
故并集等于 $X_I^k$；$\hat x$ 两侧皆不满足故被排除。
\end{proof}

\begin{theorem}[界定理]
\label{thm:bnb-bound}
设 $\mathcal{L}$ 为某时刻全部未剪枝且未展开的开放节点集，
$\bar z$ 为当前 incumbent 值（无 incumbent 时 $\bar z = +\infty$）。则
\[
\underline z := \min_{k \in \mathcal{L}} z_{LP}^k \;\le\; z^* \;\le\; \bar z,
\]
即\textbf{全局对偶界}等于开放节点松弛界的最小值。
\end{theorem}
\begin{proof}
上界：任何 incumbent 都是可行解，故 $z^* \le \bar z$。下界：由引理
\ref{lem:branching-completeness} 的归纳，任何未被搜索树排除的可行解必落在
某个开放节点（或已剪枝节点）的子树内；被 by bound 剪枝的节点子树中
可行解目标值 $\ge z_{LP} \ge \bar z \ge z^*$，故若其中存在最优解则
$\bar z = z^*$ 已成立，结论平凡。否则每个含优于 $\bar z$ 之可行解的子树
必经过某个开放节点 $k$，且 $z^* \ge z_{LP}^k \ge \min_{k' \in \mathcal{L}}
z_{LP}^{k'} = \underline z$，末步因该 $k$ 在 $\mathcal{L}$ 中
（命题 \ref{prop:milp-relaxation-bound} 应用于节点子问题给出第一步）。
\end{proof}

\begin{theorem}[分支定界的有限终止性，条件性结论]
\label{thm:bnb-termination}
若所有整数变量有有限界（$l_j, u_j$ 有限，$j \in I$），且每个节点的 LP
松弛被求到最优，则分支定界在有限步内终止，返回最优解或证明不可行。
\end{theorem}
\begin{proof}
每个节点的深度由整数变量界的收缩严格限制：沿从根到叶的任何路径，
每次分支把分支变量的可行整数区间至少缩小一个整数值（由
$\lfloor\hat x_j\rfloor < \hat x_j < \lceil\hat x_j\rceil$），
故路径长度不超过 $\sum_{j \in I}(u_j - l_j) + |I| < \infty$。
在有限深度处变量全部被固定，LP 松弛要么不可行要么给出整点，
节点必被剪枝。搜索树为有限分支度的有限深度树，故节点总数有限；
每次迭代至少展开一个节点，故算法有限步终止。终止时 $\mathcal{L} = \emptyset$，
由定理 \ref{thm:bnb-bound}，$\bar z = z^*$（或无可行解）。
\end{proof}

\begin{remark}
有限终止是\textbf{存在性}结论，不含任何多项式保证——MILP 为 NP-hard，
树的规模可对输入指数增长。工程效率几乎全部来自：更强的松弛界（割）、
更少的节点（好的分支决策与传播）与更早的 incumbent（启发式）。
定理 \ref{thm:bnb-bound} 同时解释了 best-bound 节点选择的流行原因：
取 $\arg\min_{k \in \mathcal{L}} z_{LP}^k$ 展开，是贪心提升全局下界
$\underline z$ 的策略。实现锚点：主循环
\srcpath{src/engine/solver/native/milp/bc/legacy/branch\_and\_cut.cpp:BCSolveState::run}，
入口 \srcpath{src/engine/solver/native/milp/bc/api.cpp:solve\_milp\_bc}。
\end{remark}

\subsection{强分支、伪成本与可靠性分支}
\label{subsec:branching}

分支变量的选取对树大小有决定性影响（计算证据见
\cite{linderoth1999,achterberg2005}）。以下记父节点 LP 值 $z_{LP}$，
候选分数变量 $j$ 的向下/向上子节点 LP 值为 $z_j^-, z_j^+$，
增益 $\Delta_j^- = z_j^- - z_{LP}$、$\Delta_j^+ = z_j^+ - z_{LP}$
（剪枝的子节点按约定记为充分大的增益）。

\begin{definition}[强分支]
\label{def:strong-branching}
\textbf{强分支}（strong branching）：对候选集 $\mathcal{C}$ 中每个 $j$
实际求解两个子 LP，按组合评分选变量，常用乘积评分
\begin{equation}
\label{eq:product-score}
\operatorname{score}(j) =
\big(\max\{\Delta_j^-, \varepsilon\}\big)^{1-\mu}
\big(\max\{\Delta_j^+, \varepsilon\}\big)^{\mu},
\end{equation}
$\mu \in [0,1]$（常取 $\mu = 1/6$ 的加权几何形式，$\varepsilon > 0$ 为
防零下限）。强分支信息准确但每个候选两次 LP 求解，代价高。
\end{definition}

\begin{definition}[伪成本]
\label{def:pseudocost}
对每个整数变量 $j$ 维护\textbf{伪成本}（pseudocost）历史：
\begin{equation}
\label{eq:pseudocost}
\psi_j^- = \frac{1}{|O_j^-|}\sum_{o \in O_j^-}
\frac{\Delta_o^-}{f_o}, \qquad
\psi_j^+ = \frac{1}{|O_j^+|}\sum_{o \in O_j^+}
\frac{\Delta_o^+}{1 - f_o},
\end{equation}
其中 $O_j^-, O_j^+$ 为历史上对 $j$ 向下/向上分支的观测集，
$f_o = \operatorname{frac}(\hat x_j^{(o)})$ 为当时的分数部。
\textbf{估计增益} $\widehat\Delta_j^- = \psi_j^- f_j$、
$\widehat\Delta_j^+ = \psi_j^+ (1 - f_j)$ 代入式 \eqref{eq:product-score}
得伪成本分支。
\end{definition}

伪成本的统计依据是如下简化模型：对固定的 $j$，"单位分数距离的对偶界增益"
近似为与节点无关的常数，故历史均值是可用估计量。该假设不严格成立
（增益依赖于节点的局部界与 LP 基），这是可靠性分支的动机。

\begin{definition}[可靠性分支]
\label{def:reliability-branching}
给定可靠性阈值 $\eta_{rel} \ge 1$。称变量 $j$ 的方向
$d \in \{-, +\}$ \textbf{可靠}（reliable），若该方向已积累至少
$\eta_{rel}$ 次观测（$|O_j^d| \ge \eta_{rel}$）。\textbf{可靠性分支}
（reliability branching，Achterberg--Koch--Martin \cite{achterberg2005}）：
\begin{enumerate}
  \item 取候选集 $\mathcal{C}$（如按伪成本评分排序的前若干名）；
  \item 对存在不可靠方向的候选 $j$，用强分支\textbf{实测}其不可靠方向的
        增益，并把观测并入伪成本历史；
  \item 全部方向可靠后用伪成本估计评分，选评分最高者。
\end{enumerate}
\end{definition}

\begin{proposition}[可靠性分支的估计结构]
\label{prop:reliability}
设候选 $j$ 两方向均可靠，则其评分由历史伪成本给出
$\widehat{\operatorname{score}}(j)$；若方向 $d$ 不可靠，实测值
$\Delta_j^d$ 既是本次评分依据，又使 $|O_j^d|$ 加一。于是对每个变量，
强分支求解次数以 $2\eta_{rel}$ 为界（每方向至多 $\eta_{rel}$ 次实测），
之后退化为 $O(1)$ 的查表评分。
\end{proposition}
\begin{proof}
由定义 \ref{def:reliability-branching} 的步骤 2--3 直接成立：方向可靠后
不再触发该方向的实测；每方向计数单调增加到阈值后停止。
故每个变量的总实测次数 $\le 2\eta_{rel}$，后续评分只查表。
\end{proof}

\begin{remark}[内部实验结论：局部增益不等于全局界进展]
\label{rem:reliability-mismatch}
可靠性分支的理论保证是\textbf{局部}的（命题 \ref{prop:reliability}：
有限实测后评分成本退化），它不含"固定节点数下全局下界更快上升"的断言。
平台实验记录
\srcpath{docs/archive/native\_milp\_reliability\_branching\_2026-08-13.md}
确认了这一区分：在 proof-tail（best-bound 队列）模式下移除对可靠性分支的
绕行后，探测 LP 求解数按成本模型预测出现（实现忠实、成本模型准确），
但 500 节点处的相对间隙反而略微变差——因为 best-bound 队列的全局下界
由\textbf{整个前沿的最小值}（定理 \ref{thm:bnb-bound}）决定，局部更强的
分裂改变了后代的分布却不必提高固定预算下的前沿最小值。
该失配归因于假设违背而非实现缺陷，已写回上述记录。
实现锚点
\srcpath{src/engine/solver/native/milp/bc/legacy/bc\_branching.cpp:choose\_branch\_var\_reliability\_impl}
（含深度上限与探测预算按深度折减的工程化策略）。
\end{remark}

对应实现还包括乘积评分
\srcpath{src/engine/solver/native/milp/bc/legacy/bc\_branching.cpp:compute\_branch\_product\_score}、
伪成本更新
\srcpath{src/engine/solver/native/milp/bc/legacy/bc\_branching.cpp:update\_pseudocost\_from\_branch\_evidence}
与 most-infeasible 基线
\srcpath{src/engine/solver/native/milp/bc/legacy/bc\_branching.cpp:most\_infeasible\_score}。

\subsection{割平面理论}
\label{subsec:cuts}

\begin{definition}[有效不等式]
不等式 $\alpha^\top x \le \beta$ 称为对 $X_I$ \textbf{有效}（valid），
若 $X_I \subseteq \{x : \alpha^\top x \le \beta\}$。若它切除当前 LP 松弛
最优解 $\hat x$（$\alpha^\top \hat x > \beta$），则称之为\textbf{割}
（cut / cutting plane）；把割加入松弛得更强对偶界
（命题 \ref{prop:milp-relaxation-bound} 应用于收紧后的 $X$）。
\end{definition}

\subsubsection{Chv\'atal--Gomory 割}

\begin{lemma}[CG 舍入]
\label{lem:cg-rounding}
设 $x \ge 0$ 全部分量为整数变量，$\sum_j a_j x_j \le b$ 有效，则对任意
$\lambda > 0$ 有\textbf{Chv\'atal--Gomory 割}
\begin{equation}
\label{eq:cg-cut}
\sum_j \big\lfloor \lambda a_j \big\rfloor\, x_j \;\le\; \lfloor \lambda b \rfloor.
\end{equation}
\end{lemma}
\begin{proof}
由 $\sum_j a_j x_j \le b$ 与 $\lambda > 0$ 得
$\sum_j \lambda a_j x_j \le \lambda b$。因 $x_j \ge 0$，
$\sum_j \lfloor\lambda a_j\rfloor x_j \le \sum_j \lambda a_j x_j \le \lambda b$。
左端为整数（整数变量之整系数组合），整数 $\le \lambda b$ 蕴含
$\le \lfloor\lambda b\rfloor$。
\end{proof}

取 $\lambda$ 遍历对偶乘子向量并对多行取非负组合，即得一般 CG 割
$u^\top A$ 系数的舍入形式。对有理数据的多面体 $X$，反复施加 CG 闭包在
有限秩内收敛到 $\operatorname{conv}(X_I)$（Chv\'atal 定理，条件性结论；
证明见 \cite{cornuejols2008} 或 \cite{nemhauser1988}）。

\subsubsection{Gomory 分数割}

\begin{theorem}[Gomory 分数割的推导与正确性]
\label{thm:gfc}
设 LP 松弛（标准等式型 $Ax = b$，$x \ge 0$，全部变量整数）的某最优基
对应单纯形表行
\begin{equation}
\label{eq:tableau-row}
x_{B_i} + \sum_{j \in N} \bar a_{ij} x_j = \bar b_i,
\end{equation}
其中 $N$ 为非基变量集，当前基解 $\hat x$ 中
$\hat x_{B_i} = \bar b_i \notin \mathbb{Z}$、$\hat x_j = 0$（$j \in N$）。记
$f_{ij} = \operatorname{frac}(\bar a_{ij})$、$f_{i0} = \operatorname{frac}(\bar b_i)$。
则 \textbf{Gomory 分数割}（Gomory fractional cut, GFC）
\begin{equation}
\label{eq:gfc}
\sum_{j \in N} f_{ij}\, x_j \;\ge\; f_{i0}
\end{equation}
对 $X_I$ 有效，且被 $\hat x$ 违反（左端 $0 < f_{i0}$）。
\end{theorem}
\begin{proof}
违反性直接：$\hat x_j = 0$（$j \in N$），左端为 0，而
$f_{i0} > 0$（$\bar b_i$ 非整数）。

有效性：由式 \eqref{eq:tableau-row}，
$x_{B_i} = \bar b_i - \sum_{j \in N}\bar a_{ij} x_j$。把系数拆分为整数与
分数部 $\bar a_{ij} = \lfloor\bar a_{ij}\rfloor + f_{ij}$、
$\bar b_i = \lfloor\bar b_i\rfloor + f_{i0}$，整理得
\begin{equation}
\label{eq:gfc-key}
x_{B_i} + \sum_{j \in N}\lfloor\bar a_{ij}\rfloor x_j - \lfloor\bar b_i\rfloor
= f_{i0} - \sum_{j \in N} f_{ij} x_j .
\end{equation}
对任何 $x \in X_I$（全整数、$x \ge 0$ 且满足式 \eqref{eq:tableau-row}），
式 \eqref{eq:gfc-key} 左端为整数，故右端亦为整数。而右端
$\le f_{i0} < 1$（因 $f_{ij} \ge 0$、$x_j \ge 0$），整数且 $< 1$ 故
$\le 0$，即 $f_{i0} - \sum_j f_{ij} x_j \le 0$，移项即式 \eqref{eq:gfc}。
\end{proof}

对混合整数情形（连续变量存在、基变量整数），同样的表行导出更强的
\textbf{Gomory 混合整数割}（GMI）：对表行 \eqref{eq:tableau-row} 按
系数与 $f_{i0}$ 的相对大小分配合适的斜率，得到对每个非基变量分段线性
的割系数。GMI 的标准推导见 \cite{cornuejols2008} 第 2 节；其形态是
下节 MIR 函数的特例。

\subsubsection{MIR 不等式}

\begin{definition}[MIR 函数]
对标量 $b \notin \mathbb{Z}$，记 $f_0 = \operatorname{frac}(b)$。定义
\textbf{MIR 函数}（mixed-integer rounding function）
\[
F_{f_0}(r) = \lfloor r \rfloor + \frac{\big(\operatorname{frac}(r) - f_0\big)^+}{1 - f_0}.
\]
\end{definition}

\begin{theorem}[两行 MIR 不等式]
\label{thm:mir}
考虑混合整数集
$Q = \{(x, y, s) : x + y \le b + s,\ x \in \mathbb{Z},\ y \ge 0,\ s \ge 0\}$
的弱化形式
$\{(v, y) : v \le b + y,\ v \in \mathbb{Z},\ y \ge 0\}$。则不等式
\begin{equation}
\label{eq:mir-basic}
v \;\le\; \lfloor b \rfloor + \frac{y}{1 - f_0}
\end{equation}
对该集合有效。
\end{theorem}
\begin{proof}
分两种情形。若 $v \le \lfloor b \rfloor$，因 $y/(1-f_0) \ge 0$，式
\eqref{eq:mir-basic} 成立。若 $v \ge \lfloor b \rfloor + 1 = \lceil b\rceil$，
则 $y \ge v - b \ge \lceil b \rceil - b = 1 - f_0$，故
$\lfloor b\rfloor + y/(1-f_0) \ge \lfloor b\rfloor + 1 \ge v$。
两情形穷尽（$v$ 整数），证毕。
\end{proof}

对一般的单行（或聚集行）$\sum_j a_j x_j \le b$（$x_j \ge 0$，
部分变量整数），通过\textbf{补余}（complementation：把 $x_j$ 换为
$u_j - x_j$ 以选择有利的符号）与系数缩放把行变换为定理 \ref{thm:mir}
的两行形式，回代即得\textbf{补余 MIR 割}（complemented MIR,
c-MIR \cite{marchand2001}）。c-MIR 严格强于朴素 CG 割，且可从单纯形表行
系统化生成；实现锚点
\srcpath{src/engine/solver/native/milp/bc/legacy/bc\_cuts.cpp:add\_row\_mir\_cuts}、
\srcpath{src/engine/solver/native/milp/bc/legacy/bc\_cuts.cpp:build\_sparse\_complemented\_mir}、
\srcpath{src/engine/solver/native/milp/bc/legacy/bc\_cuts.cpp:add\_basis\_mir\_cuts}
（GMI 锚点
\srcpath{src/engine/solver/native/milp/bc/legacy/bc\_cuts.cpp:add\_basis\_gomory\_cuts}）。

\subsubsection{CGLP：析取割与 lift-and-project}

\begin{definition}[析取与 CGLP]
对 0-1 变量 $x_j$，析取
$\big(X \cap \{x_j = 0\}\big) \cup \big(X \cap \{x_j = 1\}\big)$ 的凸包
是对 $X_I$ 的有效收紧。给定当前点 $\hat x$（$0 < \hat x_j < 1$），
\textbf{割生成 LP}（cut generation LP, CGLP；Balas--Ceria--Cornu\'ejols
\cite{balas1993}）在割系数空间求解：变量为 $(\alpha, \beta)$ 与两分支
的 Farkas 乘子 $(u^0, v^0)$、$(u^1, v^1)$，约束
\[
\alpha \le u^{0\top} A^{0},\quad \alpha \le u^{1\top} A^{1},\quad
\beta \ge u^{0\top} b^{0},\quad \beta \ge u^{1\top} b^{1},
\]
（$A^{0} x \le b^{0}$、$A^{1} x \le b^{1}$ 为两分支系统）加正规化约束，
目标为最大化违反度 $\alpha^\top \hat x - \beta$。
\end{definition}

\begin{proposition}[CGLP 割的有效性]
\label{prop:cglp}
CGLP 的任何可行点 $(\alpha, \beta)$ 给出的不等式
$\alpha^\top x \le \beta$ 对
$\operatorname{conv}\big((X\cap\{x_j=0\}) \cup (X\cap\{x_j=1\})\big)$ 有效；
最优 CGLP 解给出在最大化违反意义下的最强析取割。
\end{proposition}
\begin{proof}
有效性是 Farkas 引理的直接推论：$\alpha \le u^{d\top} A^{d}$、
$\beta \ge u^{d\top} b^{d}$（$u^d \ge 0$）表明
$\alpha^\top x \le \beta$ 是分支 $d$ 系统 $A^{d} x \le b^{d}$ 的
非负组合的推论，故在两分支上都有效，从而对其并的凸包有效。
最优性：CGLP 的可行域经正规化后等价于析取多面体极割的系数锥截面，
最大化违反即选出分离 $\hat x$ 的最深割；严格陈述与证明见
\cite{balas1993,cornuejols2008}。
\end{proof}

\begin{remark}
CGLP 的求解成本为一次规模约 $2m + n$ 的 LP，通常只在根节点对少数
候选变量运行。实现
\srcpath{src/engine/solver/native/milp/bc/legacy/bc\_cglp\_model.cpp:build\_cglp\_lp}、
\srcpath{src/engine/solver/native/milp/bc/legacy/bc\_cglp.cpp:generate\_cglp\_cut}、
\srcpath{src/engine/solver/native/milp/bc/legacy/bc\_cglp.cpp:add\_cglp\_cuts\_root}
包含候选选择、效力过滤与\textbf{逐分支有效性校验}的保守准入
（\srcpath{src/engine/solver/native/milp/bc/legacy/bc\_cglp.cpp:validate\_cut\_on\_branch}），
以及插值回退
（\srcpath{src/engine/solver/native/milp/bc/legacy/bc\_cglp.cpp:fallback\_interpolation\_cut}）。
\end{remark}

\subsubsection{割平面循环的不动点语义}

\begin{remark}[退化 LP 上的割平面不动点]
在退化（degenerate）LP 上，一张被采纳的割可以不改变目标值却切去当前
最优面的一部分，使后续分离面对新的最优面。因此"一批被采纳的割未立即
提升目标值"\textbf{不是}分离循环的不动点判据；正确的不动点是分离器/割池
生命周期在配置的不动预算内\textbf{不再产生任何进展}。该语义结论的平台化
实验记录见
\srcpath{docs/archive/native\_milp\_degenerate\_cutpool\_fixed\_point\_2026-08-13.md}：
保留零提升割确实给出预测的更强根界（对偶侧断言成立），但改变了退化根
点/基，使下游原始启发式消耗了整个预算——端到端契约失败，该改动按
失配协议回滚。理论教训：对偶界增强与端到端运行时间由同一退化面耦合，
不能分别建模。相关的割归属语义（生成 $\to$ 割池候选 $\to$ LP 采纳且
单调重优化成功后才成为全局有效证据行）记录于
\srcpath{docs/archive/native\_milp\_cutpool\_proof\_ownership\_2026-08-13.md}。
\end{remark}

\subsection{Domain propagation}
\label{subsec:domain-propagation}

\begin{definition}[行活动度与蕴含界]
对行 $i$：$a_i^\top x \le b_i$（$\le$ 型不失一般），其在 $[l,u]$ 上的
\textbf{活动度}为区间算术
\[
\alpha_i^{\min} = \sum_{j: a_{ij} > 0} a_{ij} l_j + \sum_{j: a_{ij} < 0} a_{ij} u_j,
\qquad
\alpha_i^{\max} = \sum_{j: a_{ij} > 0} a_{ij} u_j + \sum_{j: a_{ij} < 0} a_{ij} l_j .
\]
\textbf{界收紧}（bound tightening）：对每个 $j$，由
$a_{ij} x_j \le b_i - \big(\alpha_i^{\min} - \min\{0, a_{ij}l_j, a_{ij}u_j\}\big)$
解出 $x_j$ 的\textbf{蕴含界}（implied bound），若强于当前界则收紧。
\end{definition}

\begin{proposition}[蕴含界的有效性]
\label{prop:implied-bounds}
按上述规则收紧后的界不切除 $X_I$ 的任何点。
\end{proposition}
\begin{proof}
设 $a_{ij} > 0$。对任何满足行 $i$ 与界的 $x$，
$a_{ij} x_j = a_i^\top x - \sum_{k \ne j} a_{ik} x_k
\le b_i - \alpha_{i,\neg j}^{\min}$，其中
$\alpha_{i,\neg j}^{\min} = \alpha_i^{\min} - a_{ij} l_j$
是除 $j$ 外各项之和的下界，故
$x_j \le \big(b_i - \alpha_{i,\neg j}^{\min}\big)/a_{ij}$，
即该上界为行与界的逻辑推论，对所有可行点成立。$a_{ij} < 0$ 与下界情形
对称。整数变量再把上界下取整、下界上取整，仍有效（整点满足更强条件）。
\end{proof}

\textbf{探测}（probing，Savelsbergh \cite{savelsbergh1994}）：对 0-1 变量
$x_j$ 分别暂设 $x_j = 0$ 与 $x_j = 1$ 运行界收紧。若某侧不可行则
$x_j$ 可固定；若两侧都可行，则两侧收紧结果的并给出其他变量的隐含界
（如 $x_j = 0 \Rightarrow x_k \ge \alpha$ 且
$x_j = 1 \Rightarrow x_k \ge \beta$ 蕴含 $x_k \ge \min\{\alpha, \beta\}$），
并可抽取蕴含关系 $x_j = 1 \Rightarrow x_k \ge \alpha$ 供 clique 表与
割生成使用。传播以不动点迭代执行：界收紧触发活动度更新，再触发新的收紧，
直至无进展或不可行。实现锚点
\srcpath{src/engine/solver/native/milp/bc/legacy/bc\_domain\_probe.cpp:BCDomainProbeWorkspace::propagate}、
\srcpath{src/engine/solver/native/milp/bc/legacy/bc\_domain\_probe.cpp:BCDomainProbeWorkspace::probe}、
\srcpath{src/engine/solver/native/milp/bc/legacy/bc\_implied\_bounds.cpp:compute\_implied\_column\_bounds\_from\_rows}。

\subsection{Clique 与 cover 不等式}
\label{subsec:clique-cover}

\begin{definition}[冲突图与 clique 不等式]
对 0-1 变量，若任意可行点至多使文字（literal，$x_j$ 或 $1 - x_j$）集合
$C$ 中的一个取真，则
\begin{equation}
\label{eq:clique}
\sum_{j \in C^+} x_j + \sum_{j \in C^-} (1 - x_j) \;\le\; 1
\end{equation}
为 \textbf{clique 不等式}。文字间的两两冲突（"至多一个为真"）构成
\textbf{冲突图}（conflict graph）的边，clique 不等式对应其团。
\end{definition}

\begin{proposition}[clique 不等式的有效性]
式 \eqref{eq:clique} 对 $X_I$ 有效；且当 $C$ 在冲突图中为极大团时，
该不等式不被同形式的更强不等式支配。
\end{proposition}
\begin{proof}
有效性即定义：任意 $x \in X_I$ 至多为 $C$ 中一个文字赋值真，
左端计数取真文字数，故 $\le 1$。极大性：若 $C$ 非极大，存在文字
$\ell \notin C$ 与 $C$ 全部冲突，则 $C \cup \{\ell\}$ 的不等式逐项
$\ge$ 原不等式且严格多一项，支配原式；反之极大团无此支配者。
\end{proof}

\begin{definition}[knapsack cover]
对单行 knapsack 约束 $\sum_j a_j x_j \le b$（$a_j > 0$，$x_j \in \{0,1\}$），
集合 $C$ 称为\textbf{覆盖}（cover），若 $\sum_{j \in C} a_j > b$；
称\textbf{最小覆盖}（minimal cover），若去掉任一元素即非覆盖。
\textbf{覆盖不等式}
\begin{equation}
\label{eq:cover}
\sum_{j \in C} x_j \;\le\; |C| - 1
\end{equation}
对 knapsack 可行集有效。
\end{definition}

\begin{proposition}[覆盖不等式的有效性]
\label{prop:cover-validity}
式 \eqref{eq:cover} 对 $\{x \in \{0,1\}^n : \sum_j a_j x_j \le b\}$ 有效。
\end{proposition}
\begin{proof}
若式 \eqref{eq:cover} 被违反，则 $C$ 中全部 $|C|$ 个变量取 1，
左端 $\sum_{j \in C} a_j x_j = \sum_{j \in C} a_j > b$，与约束矛盾。
\end{proof}

\begin{remark}[lift 与强度]
最小覆盖不等式一般不是面（facet）；通过\textbf{提升}（lifting）——把
$C$ 外变量按 knapsack 子问题算出的系数序贯加入——可得（在条件下）面的
加强形式。序贯提升系数由求解 $|C|$ 个 knapsack 子问题确定；标准处理见
\cite{cornuejols2008} 第 5 节与 \cite{crowder1983}。实现的 cover/clique
分离锚点
\srcpath{src/engine/solver/native/milp/bc/legacy/bc\_cuts.cpp:add\_cover\_cuts}、
\srcpath{src/engine/solver/native/milp/bc/legacy/bc\_cuts.cpp:add\_clique\_cuts}，
冲突图维护与传播
\srcpath{src/engine/solver/native/milp/bc/legacy/bc\_clique\_table.cpp:CliqueTable::build}、
\srcpath{src/engine/solver/native/milp/bc/legacy/bc\_clique\_table.cpp:CliqueTable::propagate}、
\srcpath{src/engine/solver/native/milp/bc/legacy/bc\_clique\_table.cpp:CliqueTable::find\_violated\_cliques}。
\end{remark}

\subsection{原始启发式思想}
\label{subsec:primal-heuristics}

对偶侧（界）只能证明最优性，\textbf{ incumbent} 须由原始侧提供。
主要思想类别：

\begin{enumerate}
  \item \textbf{舍入启发式}（rounding）：把 LP 松弛解的分数整数变量逐个
        舍入并沿行传播修复；成本低、成功率依问题而定。实现锚点
        \srcpath{src/engine/solver/native/milp/bc/legacy/bc\_relaxation.cpp:try\_rounding\_heuristic}。
  \item \textbf{潜水}（diving）：从当前 LP 解出发，反复把一个分数变量
        固定到舍入值并重优化 LP（热启动），直到整可行或不可行回溯。
  \item \textbf{可行性泵}（feasibility pump，Fischetti--Glover--Lodi
        \cite{fischetti2005}）：交替执行"把当前 LP 解舍入到最近整点"与
        "重优化 LP 使目标为到该整点的 $\ell_1$ 距离"，直至收敛或
        循环检测触发扰动/重启。
  \item \textbf{邻域搜索}（LNS / RINS / local branching）：给定 incumbent
        $\bar x$，RINS \cite{danna2005} 固定 LP 解与 $\bar x$ 一致的整数
        变量，对剩余子 MILP 限时求解；local branching 以
        $\sum_{j: \bar x_j = 0} x_j + \sum_{j: \bar x_j = 1}(1-x_j) \le k$
        定义 Hamming 邻域并在其中搜索改进。
  \item \textbf{交叉}（crossover）与目标驱动潜水等组合变体。
\end{enumerate}

这些启发式都不改变任何界或证书的合法性：采纳的 incumbent 仍须通过
原始模型的独立可行性审计。实现的根节点启发式管线内联于
\srcpath{src/engine/solver/native/milp/bc/legacy/branch\_and\_cut.cpp:BCSolveState::run}
（可行性泵、潜水、LNS、探测等按选项与进展门控）。

\subsection{Presolve 对 MILP 的约简}
\label{subsec:milp-presolve}

Presolve 在求解前对模型做\textbf{保持最优解集等价}的约简；每条规则都附带
postsolve 逆映射以恢复原空间解。标准规则及其有效性：

\begin{enumerate}
  \item \textbf{固定变量与空行/空列删除}：$l_j = u_j$ 的变量代入消去；
        恒可行（$\alpha_i^{\max} \le b_i$）或恒紧的行删除。
  \item \textbf{单子行}（singleton row）：行 $i$ 只含 $x_j$ 时直接给出
        $x_j$ 的最强蕴含界（命题 \ref{prop:implied-bounds} 的退化情形），
        行随后删除。
  \item \textbf{单子列}（singleton column）：$x_j$ 只出现于行 $i$ 时，
        依目标系数符号把 $x_j$ 压到使行最松的界，行删除、变量由
        postsolve 重构。
  \item \textbf{强制行}（forcing row）：$\alpha_i^{\max} = b_i$（$\le$ 型）
        时，行等号迫使每个变量取其贡献达到最大的界；全部固定后行删除。
  \item \textbf{支配列}（dominated column）：若 $x_j$ 的各约束系数与目标
        被 $x_k$ 逐项支配，则存在不含 $x_j$（在其界上）的最优解，可固定。
  \item \textbf{双元等式}（doubleton equation）：$a x_j + b x_k = c$ 用于
        聚合消去一个变量（保持整数性需系数整除性检查）。
  \item \textbf{系数加强与整数右端取整}：全整数行右端下取整
        （引理 \ref{lem:cg-rounding} 的特例）；由活动度反推过大的系数。
  \item \textbf{探测}：见 \ref{subsec:domain-propagation} 小节。
\end{enumerate}

\begin{proposition}[单子行收紧的完备性]
对只含 $x_j$ 的行 $a_{ij} x_j \le b_i$：若 $a_{ij} > 0$ 则
$x_j \le b_i / a_{ij}$ 与当前上界取小为该行蕴含的\textbf{最强}上界；
$a_{ij} < 0$ 对称。收紧后删除行不损失任何可行点。
\end{proposition}
\begin{proof}
行约束在 $x_j$ 上单独起作用（无其他项），故 $x_j \le b_i/a_{ij}$
（$a_{ij} > 0$ 时）是该行的重述，凡满足行的点必满足新界；反之在
$[l_j, u_j]$ 收紧后的任何点代回仍满足原行。故可行集不变。
\end{proof}

\begin{remark}[presolved 坐标中的仿射搬运]
presolve 建立原空间与约简空间之间的逐变量仿射映射
$x = s \odot y + c$（scale/shift）。任何原空间的证据结构——例如
variable-bound 蕴含 $x_t \le a x_z + b$——要用于约简空间，必须经同一映射
\textbf{精确搬运}：
\[
y_t \;\le\; \frac{a\, s_z}{s_t}\, y_z \;+\; \frac{a\, c_z + b - c_t}{s_t},
\qquad s_t < 0 \text{ 时不等号反向},
\]
任一端点被消去、映射缺失或 $s_t = 0$ 时拒绝该条记录。直接以原列号充当
约简列号是无效坐标回放，会产生错误蕴含并抬高松弛界越过真最优值——
该缺陷的完整推导、代价模型与修正验证记录于
\srcpath{docs/archive/native\_milp\_presolve\_varbound\_coordinates\_2026-08-13.md}
（结论：仿射搬运是\textbf{正确性}前提，但不单独复现已知求解器的全部
根切割强度）。有效性证明：把 $x_t = s_t y_t + c_t$、$x_z = s_z y_z + c_z$
代入原不等式并除以 $s_t$（负号反向）即得，逐点保持可行集。
实现锚点
\srcpath{src/engine/solver/native/milp/bc/legacy/bc\_implied\_bounds.cpp:build\_variable\_bound\_table\_from\_rows}，
postsolve 映射
\srcpath{src/engine/solver/native/milp/bc/milp\_presolve.cpp:MILPPresolve::postsolve}、
\srcpath{src/engine/solver/native/milp/bc/milp\_presolve.cpp:MILPPresolve::forward\_map}。
\end{remark}

实现锚点汇总：入口
\srcpath{src/engine/solver/native/milp/bc/milp\_presolve.cpp:presolve\_mip} 与
\srcpath{src/engine/solver/native/milp/bc/milp\_presolve.cpp:MILPPresolve::run}；
各规则对应
\srcpath{src/engine/solver/native/milp/bc/milp\_presolve.cpp:MILPPresolve::process\_singleton\_rows}、
\srcpath{src/engine/solver/native/milp/bc/milp\_presolve.cpp:MILPPresolve::process\_singleton\_columns}、
\srcpath{src/engine/solver/native/milp/bc/milp\_presolve.cpp:MILPPresolve::process\_doubleton\_equations}、
\srcpath{src/engine/solver/native/milp/bc/milp\_presolve.cpp:MILPPresolve::detect\_forcing\_rows}、
\srcpath{src/engine/solver/native/milp/bc/milp\_presolve.cpp:MILPPresolve::detect\_dominated\_columns}、
\srcpath{src/engine/solver/native/milp/bc/milp\_presolve.cpp:MILPPresolve::tighten\_bounds}、
\srcpath{src/engine/solver/native/milp/bc/milp\_presolve.cpp:MILPPresolve::run\_probing}。

\subsection{与实现的对应}

\begin{longtable}{p{0.42\textwidth} p{0.50\textwidth}}
\toprule
理论条目 & 实现状态与锚点 \\
\midrule
\endhead
分支定界主框架（\ref{subsec:bnb} 小节，定理
\ref{thm:bnb-bound}/\ref{thm:bnb-termination}）&
\implfull{src/engine/solver/native/milp/bc/api.cpp:solve\_milp\_bc}、
\srcpath{src/engine/solver/native/milp/bc/legacy/branch\_and\_cut.cpp:BCSolveState::run} \\
LP 松弛与对偶界（命题 \ref{prop:milp-relaxation-bound}、
\ref{prop:milp-dual-bound}）&
\implfull{src/engine/solver/native/milp/bc/legacy/bc\_relaxation.cpp:solve\_lp\_relaxation}
（vendored HiGHS / 原生内核双后端） \\
伪成本（定义 \ref{def:pseudocost}）与乘积评分
\eqref{eq:product-score} &
\implfull{src/engine/solver/native/milp/bc/legacy/bc\_branching.cpp:compute\_branch\_product\_score}、
\srcpath{src/engine/solver/native/milp/bc/legacy/bc\_branching.cpp:update\_pseudocost\_from\_branch\_evidence} \\
可靠性分支（定义 \ref{def:reliability-branching}、命题
\ref{prop:reliability}）&
\implpart{实现了可靠性阈值、双向实测与历史伪成本切换，但带工程化限制：
深度上限 6、探测预算按深度折减、持久 LP 热启动复用；锚点
\srcpath{src/engine/solver/native/milp/bc/legacy/bc\_branching.cpp:choose\_branch\_var\_reliability\_impl}} \\
most-infeasible 分支 & \implfull{src/engine/solver/native/milp/bc/legacy/bc\_branching.cpp:most\_infeasible\_score} \\
Gomory 分数/混合整数割（定理 \ref{thm:gfc} 及 GMI）&
\implfull{src/engine/solver/native/milp/bc/legacy/bc\_cuts.cpp:add\_basis\_gomory\_cuts}；
另有 zero-half 变体
\srcpath{src/engine/solver/native/milp/bc/legacy/bc\_cuts.cpp:add\_zero\_half\_cuts} \\
补余 MIR 割（定理 \ref{thm:mir} 及 c-MIR）&
\implfull{src/engine/solver/native/milp/bc/legacy/bc\_cuts.cpp:add\_row\_mir\_cuts}、
\srcpath{src/engine/solver/native/milp/bc/legacy/bc\_cuts.cpp:build\_sparse\_complemented\_mir}、
\srcpath{src/engine/solver/native/milp/bc/legacy/bc\_cuts.cpp:add\_basis\_mir\_cuts}、
\srcpath{src/engine/solver/native/milp/bc/legacy/bc\_cuts.cpp:add\_mir\_like\_cuts} \\
CGLP 析取割（命题 \ref{prop:cglp}）&
\implpart{根节点 lift-and-project 主 LP + 候选选择 + 效力过滤 +
逐分支保守有效性校验 + 插值回退；锚点
\srcpath{src/engine/solver/native/milp/bc/legacy/bc\_cglp\_model.cpp:build\_cglp\_lp}、
\srcpath{src/engine/solver/native/milp/bc/legacy/bc\_cglp.cpp:generate\_cglp\_cut}、
\srcpath{src/engine/solver/native/milp/bc/legacy/bc\_cglp.cpp:add\_cglp\_cuts\_root}} \\
domain propagation 与探测（命题 \ref{prop:implied-bounds}）&
\implfull{src/engine/solver/native/milp/bc/legacy/bc\_domain\_probe.cpp:BCDomainProbeWorkspace::propagate}、
\srcpath{src/engine/solver/native/milp/bc/legacy/bc\_domain\_probe.cpp:BCDomainProbeWorkspace::probe}、
\srcpath{src/engine/solver/native/milp/bc/legacy/bc\_implied\_bounds.cpp:compute\_implied\_column\_bounds\_from\_rows} \\
clique 不等式与冲突图 \eqref{eq:clique} &
\implfull{src/engine/solver/native/milp/bc/legacy/bc\_clique\_table.cpp:CliqueTable::build}、
\srcpath{src/engine/solver/native/milp/bc/legacy/bc\_clique\_table.cpp:CliqueTable::propagate}、
\srcpath{src/engine/solver/native/milp/bc/legacy/bc\_cuts.cpp:add\_clique\_cuts} \\
cover 不等式 \eqref{eq:cover}（命题 \ref{prop:cover-validity}）&
\implpart{实现了 cover 分离与投影容量变体，未实现系统的序贯提升；锚点
\srcpath{src/engine/solver/native/milp/bc/legacy/bc\_cuts.cpp:add\_cover\_cuts}、
\srcpath{src/engine/solver/native/milp/bc/legacy/bc\_cuts.cpp:add\_projected\_capacity\_cover\_cuts}} \\
舍入启发式 & \implfull{src/engine/solver/native/milp/bc/legacy/bc\_relaxation.cpp:try\_rounding\_heuristic} \\
可行性泵 / 潜水 / LNS / RINS 等根与树内启发式 &
\implpart{以 \texttt{BCSolveState::run} 内联段实现并按选项门控
（可行性泵、LP 潜水、LNS、RINS 频率触发等）；锚点
\srcpath{src/engine/solver/native/milp/bc/legacy/branch\_and\_cut.cpp:BCSolveState::run}} \\
MILP presolve 全套规则（\ref{subsec:milp-presolve} 小节）&
\implfull{src/engine/solver/native/milp/bc/milp\_presolve.cpp:presolve\_mip}、
\srcpath{src/engine/solver/native/milp/bc/milp\_presolve.cpp:MILPPresolve::run} \\
variable-bound 蕴含的坐标仿射搬运 &
\implfull{src/engine/solver/native/milp/bc/legacy/bc\_implied\_bounds.cpp:build\_variable\_bound\_table\_from\_rows}；
postsolve 映射
\srcpath{src/engine/solver/native/milp/bc/milp\_presolve.cpp:MILPPresolve::postsolve} \\
割池两阶段归属与退化不动点语义（\ref{subsec:cuts} 小节末 remark）&
\implpart{割生成、选择、采纳与生命周期在
\srcpath{src/engine/solver/native/milp/bc/legacy/bc\_cuts.cpp:add\_cuts}
统一调度；两阶段归属语义已实现，退化不动点保留策略按记录回滚，见对应
archive 记录} \\
约简空间割生成（transformed cuts）&
\implfull{src/engine/solver/native/milp/bc/legacy/bc\_cuts\_transformed.cpp:add\_transformed\_tableau\_cuts} \\
树搜索的并行化 & \implpart{并行路径存在于
\srcpath{src/engine/solver/native/milp/bc/legacy/bc\_parallel.cpp}，
本章理论为串行框架；并行不改变界定理 \ref{thm:bnb-bound} 的语义} \\
\bottomrule
\end{longtable}

\subsection{参考文献}

\begin{thebibliography}{9}\small
\bibitem{landdoig1960} A. H. Land, A. G. Doig,
An automatic method of solving discrete programming problems,
\emph{Econometrica} 28(3):497--520, 1960.
\bibitem{achterberg2007} T. Achterberg,
\emph{Constraint Integer Programming}, 博士论文, TU Berlin, 2007.
\bibitem{achterberg2005} T. Achterberg, T. Koch, A. Martin,
Branching rules revisited,
\emph{Operations Research Letters} 33(1):42--54, 2005.
\bibitem{linderoth1999} J. T. Linderoth, M. W. P. Savelsbergh,
A computational study of search strategies for mixed integer programming,
\emph{INFORMS Journal on Computing} 11(2):173--187, 1999.
\bibitem{marchand2001} H. Marchand, L. A. Wolsey,
Aggregation and mixed integer rounding to solve MIPs,
\emph{Operations Research} 49(3):363--371, 2001.
\bibitem{cornuejols2008} G. Cornu\'ejols,
Valid inequalities for mixed integer linear programs,
\emph{Mathematical Programming} 112(1):3--44, 2008.
\bibitem{nemhauser1988} G. L. Nemhauser, L. A. Wolsey,
\emph{Integer and Combinatorial Optimization}, Wiley, 1988.
\bibitem{balas1993} E. Balas, S. Ceria, G. Cornu\'ejols,
A lift-and-project cutting plane algorithm for mixed 0-1 programs,
\emph{Mathematical Programming} 58(1--3):295--324, 1993.
\bibitem{crowder1983} H. Crowder, E. L. Johnson, M. Padberg,
Solving large-scale zero-one linear programming problems,
\emph{Operations Research} 31(5):803--834, 1983.
\bibitem{savelsbergh1994} M. W. P. Savelsbergh,
Preprocessing and probing techniques for mixed integer programming problems,
\emph{ORSA Journal on Computing} 6(4):445--454, 1994.
\bibitem{fischetti2005} M. Fischetti, F. Glover, A. Lodi,
The feasibility pump,
\emph{Mathematical Programming} 104(1):91--104, 2005.
\bibitem{danna2005} E. Danna, E. Rothberg, C. Le Pape,
Exploring relaxation induced neighborhoods to improve MIP solutions,
\emph{Mathematical Programming} 102(1):71--90, 2005.
\end{thebibliography}


% source_equivalent_lp_kernel.tex —— engine 模块实现转写章：原生对偶单纯形 LP 内核
% 本文件为实现转写章（非理论章），遵循 docs/modules/README.md 数学模型规范第 1/2/4 条。

\section{原生对偶单纯形 LP 内核：源码等价转写}
\label{sec:engine-source-lp-kernel}

\subsection{转写范围与口径}
\label{subsec:engine-lp-kernel-scope}

本章转写以下工作树文件在 2026-09-13 版本实际执行的计算：

\begin{itemize}
  \item 内核主循环与状态机：\srcpath{src/engine/kernel/lp\_kernel/native\_dual/solver.cpp}、
        \srcpath{src/engine/kernel/lp\_kernel/native\_dual/state.cpp}、
        \srcpath{src/engine/kernel/lp\_kernel/native\_dual/primal.cpp}、
        \srcpath{src/engine/kernel/lp\_kernel/native\_dual/certificate.cpp}；
  \item 定价与比率测试：\srcpath{src/engine/kernel/lp\_kernel/native\_dual/pricing.cpp}、
        \srcpath{src/engine/kernel/lp\_kernel/native\_dual/primal\_pricing.hpp}；
  \item 基分解封装与数据结构：\srcpath{src/engine/kernel/lp\_kernel/native\_dual/factor.cpp}、
        \srcpath{src/engine/kernel/lp\_kernel/native\_dual/factor.hpp}、
        \srcpath{src/engine/kernel/lp\_kernel/native\_dual/model.hpp}、
        \srcpath{src/engine/kernel/lp\_kernel/native\_dual/partitioned\_row\_matrix.hpp}、
        \srcpath{src/engine/kernel/lp\_kernel/native\_dual\_core.hpp}；
  \item legacy 包装层：\srcpath{src/engine/kernel/lp\_kernel/dual\_simplex.cpp}、
        \srcpath{src/engine/kernel/lp\_kernel/dual\_simplex\_api.cpp}；
  \item portfolio 选择器：\srcpath{src/engine/solver/native/native\_lp\_selector.cpp}。
\end{itemize}

每组公式后以``对应源码为 \texttt{文件:函数名}''给出锚点；匿名命名空间内的函数按名称引用，
lambda 按 \texttt{lambda 名（lambda）} 引用。本章只转写代码执行的赋值、残差、阈值与分支；
算法动机（Harris 1973、Bland 1977、Maros 2003、Goldfarb--Reid 1977 等）以代码内引用为准，
不作独立推导。代码中由环境变量开启的替代路径（\texttt{MIPSOLVERS\_DS\_*} 系列）在首次出现处
就地声明。

\begin{boundarynote}
本章不覆盖 vendored HiGHS 单纯形桥（\texttt{solve\_standard\_form\_with\_vendored\_highs}，
仅在选择器中作为后端分支提及）、锥/非线性内核（另章转写），也不覆盖
\texttt{MIPSOLVERS\_DS\_PROFILE} 纯测量通道——该通道只累积耗时与计数，不改变任何数值结果。
\end{boundarynote}

% ─────────────────────────────────────────────────────────────────────────────
\subsection{标准形与包装层}
\label{subsec:engine-lp-stdform}

\paragraph{标准形构造.}
\compmeta{StandardFormLP}{include/mipsolvers/engine/kernel/lp\_kernel/dual\_simplex.hpp}
内核消费的标准形为
\begin{equation}
  \max_{x}\; \bar{c}^{\top} x \quad \text{s.t.}\quad A x = b,\quad
  0 \le x \le u,
  \label{eq:lpkernel:stdform}
\end{equation}
其中 $u$ 允许取 $+\infty$。构造规则（原问题 \texttt{LPModel} 含不等式块 $A x \le b$、
等式块 $A_{eq} x = b_{eq}$ 与变量界 $[\ell_j, u_j]$）：

\begin{enumerate}
  \item 下界平移：$\sigma_j = \ell_j$（$\ell_j$ 有限），否则 $\sigma_j = u_j$
        （$u_j$ 有限），否则 $\sigma_j = 0$；RHS 平移为
        $\hat b = b - A\sigma$、$\hat b_{eq} = b_{eq} - A_{eq}\sigma$；
  \item 行定向：rhs 为 $+\infty$ 而 lhs 有限的 G 行保留 lhs 为右端（surplus 形式）；
        rhs 为负的 L/E 行整体乘 $-1$ 并翻号 sense；双边行
        $\mathrm{lhs} \le a^{\top}x \le \mathrm{rhs}$ 保留为一行加有界松弛，
        松弛上界 $\max(0, \mathrm{rhs}-\mathrm{lhs})$，若平移后 rhs 为负则从下端重写；
  \item 列布局：$[x_{\text{orig}}\;|\; s_{\text{slack}}\;|\; s_{\text{surplus}}\;|\;
        s_{\text{artificial}}\;|\; q_{\text{negative}}]$；G 行取 surplus（系数 $-1$）
        加 artificial，E 行取 artificial，L 行取 slack（系数 $+1$）；
        无有限下界的原变量分裂 $x_j = x_j^{+} - q_j$，$q$ 列系数取负；
  \item 矩阵装配按 CSC 终序直接写入，系数绝对值 $\le 10^{-15}$ 的项被丢弃
        （硬编码阈值 \texttt{1e-15}）；
  \item 目标：最大化形 $\bar c$ 满足 $\bar c_j = -c^{\min}_j$（原变量）、
        $\bar c_{q_j} = c^{\min}_j$，其余为 $0$；常数项
        $\mathrm{obj}_{\text{const}} = (c^{\min})^{\top}\sigma$；上界数组
        $u_j = u_j^{\text{orig}} - \sigma_j$（有限时），ranged 松弛取行宽，
        其余辅助列 $+\infty$。
\end{enumerate}
对应源码为 \srcpath{src/engine/kernel/lp\_kernel/dual\_simplex\_api.cpp:build\_standard\_form\_lp}。
界哨兵约定：公开模型层的 $\ge 10^{19}$ 有限哨兵（如 $10^{20}$）在边界处统一归一为
数学 $+\infty$：\texttt{canonical\_upper\_bound} 与 \texttt{has\_finite\_upper\_bound}
均以 \texttt{kModelBoundSentinel}${}=10^{19}$ / \texttt{kUpperBoundSentinel}${}=10^{19}$ 判定。
对应源码为 \srcpath{src/engine/kernel/lp\_kernel/native\_dual/state.cpp:canonical\_upper\_bound}、
\srcpath{src/engine/kernel/lp\_kernel/dual\_simplex\_api.cpp:has\_finite\_upper\_bound}。

\paragraph{Ruiz 均衡.}
\texttt{solve\_lp\_with\_basis\_impl} 在构造标准形后固定执行 10 轮 Ruiz 均衡：
每轮先取当前矩阵各行的行内最大模 $r_i$ 与各列的列内最大模 $c_j$，置
\begin{equation}
  D_r^{(t)} = \operatorname{diag}\bigl(1/\sqrt{r_i}\bigr),\qquad
  D_c^{(t)} = \operatorname{diag}\bigl(1/\sqrt{c_j}\bigr),
\end{equation}
（模 $\le 10^{-15}$ 的行列因子取 $1$），然后 $A \leftarrow D_r^{(t)} A D_c^{(t)}$、
$b \leftarrow D_r^{(t)} b$、$\bar c \leftarrow D_c^{(t)} \bar c$、
$u_j \leftarrow u_j / D^{(t)}_{c,j}$（有限时），并把因子累乘进
\texttt{row\_scale}/\texttt{col\_scale}。早停：第 $t \ge 2$ 轮起，若
$\max\{|r_i - 1|, |c_j - 1|\} < 10^{-2}$ 则终止。
对应源码为 \srcpath{src/engine/kernel/lp\_kernel/dual\_simplex\_api.cpp:ruiz\_scale\_standard\_form}。
解还原 $x^{\text{orig}} = \sigma + D_c\, x^{\text{std}}$（分裂列再减去 $q$ 分量）。
对应源码为 \srcpath{src/engine/kernel/lp\_kernel/dual\_simplex.cpp:extract\_solution}。

\paragraph{节点 LP 的增量界更新.}
B\&C 节点只改少量下界时，\texttt{update\_standard\_form\_bounds\_incremental}
对下界变化 $d = \sigma_j^{\text{new}} - \sigma_j^{\text{old}}$ 执行
\begin{equation}
  b_i \leftarrow b_i - r_i \cdot s_i \cdot A_{ij} \cdot d
  \quad (i \in \operatorname{supp} A_{\cdot j}),\qquad
  \mathrm{obj}_{\text{const}} \mathrel{+}= \pm c_j d,\qquad
  u_j \leftarrow (u_j^{\text{old,orig}} - \sigma_j^{\text{new}})/D_{c,j},
\end{equation}
其中 $r_i$ 为行缩放、$s_i$ 为行符号；$|d| < 10^{-15}$ 时跳过数值更新但仍更新界。
先校验整条变更序列（旧值与缓存状态须在
$10^{-12}\max(1,|\cdot|)$ 内一致）再修改，任何一个下界失去有限性即返回失败并
要求整体重建。\texttt{StandardFormBoundTransaction} 在应用前快照受影响列的
$(\sigma_j, u_j)$、受影响行的 $b_i$ 与目标常数，\texttt{rollback} 逐项恢复。
对应源码为 \srcpath{src/engine/kernel/lp\_kernel/dual\_simplex\_api.cpp:update\_standard\_form\_bounds\_incremental}、
\srcpath{src/engine/kernel/lp\_kernel/dual\_simplex\_api.cpp:StandardFormBoundTransaction::apply}、
\srcpath{src/engine/kernel/lp\_kernel/dual\_simplex\_api.cpp:StandardFormBoundTransaction::rollback}。

\paragraph{原空间发布审计.}
每个自称成功的求解都要通过原空间残差审计
\begin{equation}
  \mathrm{viol} \le \tau \cdot \mathrm{scale},\qquad
  \mathrm{scale} = \max\bigl(1, \max_i |b_i|, \max_i |b^{eq}_i|\bigr),
\end{equation}
其中 viol 聚合行两侧违反、等式残差与变量界违反；$|\mathrm{side}| \ge 10^{19}$
（\texttt{kSideSentinel}）的行端点既不参与 viol 也不参与 scale
（防 $10^{20}$ 哨兵放大接受域）。审计失败则把结果改写为
``Residual audit rejected''。对应源码为
\srcpath{src/engine/kernel/lp\_kernel/dual\_simplex.cpp:lp\_solution\_residual\_acceptable}、
\srcpath{src/engine/kernel/lp\_kernel/dual\_simplex.cpp:solve\_lp\_with\_basis}；
标准形坐标版本为 \srcpath{src/engine/kernel/lp\_kernel/dual\_simplex.cpp:sf\_solution\_residual\_acceptable}
（行/列量分别除以 $D_r$、乘以 $D_c$ 还原原单位，人工列残余以同一门限检查）。

% ─────────────────────────────────────────────────────────────────────────────
\subsection{内核状态与不变量}
\label{subsec:engine-lp-state}

\compmeta{detail::State}{src/engine/kernel/lp\_kernel/native\_dual/model.hpp}
内核在 \eqref{eq:lpkernel:stdform} 上维护：基位置序 \texttt{basis}（第 $p$ 个基位置
存放列号）、成员标记 \texttt{basic}、非基侧标 \texttt{move}（$\mathrm{Up}=+1$ 取下界、
$\mathrm{Down}=-1$ 取上界、$\mathrm{Fixed}=0$）、基变量值 $x_B$、既约成本
$d$、每行 DSE 权重 $w$、工作成本 $\tilde c$ 及其相对原成本的扰动/平移日志
（\texttt{SparseCostJournal}：稀疏表项超过约 $n/8$ 时稠密化，
见 \srcpath{src/engine/kernel/lp\_kernel/native\_dual/model.hpp:SparseCostJournal}）。

每次 minor 迭代保持的不变量由 \texttt{audit} 定义并在各终止点强制重检：
\begin{align}
  \|A x - b\|_{\infty} &\le \varepsilon_{\text{feas}}
      \max\bigl(1, \|b\|_{\infty}\bigr),
  \label{eq:lpkernel:audit-eq}\\
  \max_{p} |d_{\texttt{basis}[p]}| &\le \varepsilon_{\text{opt}},\\
  \ell_{\texttt{basis}[p]} - \varepsilon_{\text{feas}} \le x_{B,p}
      &\le u_{\texttt{basis}[p]} + \varepsilon_{\text{feas}}
      &&(\text{要求原始可行时}),\\
  \operatorname{sign}(m_j)\, d_j &\le \varepsilon_{\text{opt}}
      &&(\text{要求对偶可行时}),
\end{align}
其中 $x$ 由 \texttt{full\_primal} 组装（非基取侧标指定的界，基变量取 $x_B$）。
\texttt{exact\_residual=true} 时 \eqref{eq:lpkernel:audit-eq} 的残差用
\texttt{long double} 累加。对应源码为
\srcpath{src/engine/kernel/lp\_kernel/native\_dual/state.cpp:audit}、
\srcpath{src/engine/kernel/lp\_kernel/native\_dual/state.cpp:equation\_residual\_vector}、
\srcpath{src/engine/kernel/lp\_kernel/native\_dual/state.cpp:full\_primal}。
单点查询：\texttt{primal\_infeasibility} 返回
$\max(\ell - x_B, 0)$（side $=-1$）或 $\max(x_B - u, 0)$（side $=+1$），
\texttt{dual\_infeasibility} 返回 $\max(0, \operatorname{sign}(m_j) d_j)$。
对应源码为 \srcpath{src/engine/kernel/lp\_kernel/native\_dual/state.cpp:primal\_infeasibility}、
\srcpath{src/engine/kernel/lp\_kernel/native\_dual/state.cpp:dual\_infeasibility}。

\paragraph{行分区矩阵.}
\texttt{PartitionedRowMatrix} 持有标准形矩阵的 CSR 视图，每行的存储槽被划分为
非基前缀 $[\texttt{row\_start}, \texttt{nonbasic\_end})$ 与基后缀；每次进/出基以
$O(\mathrm{nnz}(j))$ 的槽交换增量维护（\texttt{set\_basic}/\texttt{set\_nonbasic}），
使 dual PRICE 只扫非基列。\texttt{build} 从 CSC 转置一次并记录每个槽的 CSC 源位置
（\texttt{csc\_pos\_}）以支持 Ruiz 重缩放后的数值重绑定
（\texttt{rebind\_values}）。\texttt{verify} 为调试期全量一致性审计
（\texttt{MIPSOLVERS\_DS\_VERIFY\_PARTITION} 开启时才在提交路径调用）。
对应源码为 \srcpath{src/engine/kernel/lp\_kernel/native\_dual/partitioned\_row\_matrix.hpp:PartitionedRowMatrix}、
\srcpath{src/engine/kernel/lp\_kernel/native\_dual/state.cpp:resync\_partition\_row}、
\srcpath{src/engine/kernel/lp\_kernel/native\_dual/state.cpp:apply\_partition\_row\_swap}。

% ─────────────────────────────────────────────────────────────────────────────
\subsection{因子化契约：BasisFactor}
\label{subsec:engine-lp-factor}

\compmeta{detail::BasisFactor}{src/engine/kernel/lp\_kernel/native\_dual/factor.hpp}
\texttt{BasisFactor} 包装 vendored HFactor 后端（详见
第~\ref{sec:engine-source-ipm-kkt}~章线性代数节），对外提供：

\paragraph{重建与秩修复.}
\texttt{rebuild} 先校验基集合（基数 $=m$、列号在界内、无重复），然后调用
\texttt{factorize\_with\_logicals} 做秩揭示分解：先把每行逻辑列映射为隐式单位列
参与分解，若报告秩亏则把报告奇异位置的基列替换为该行的逻辑列，再在真实列空间
重建一次。秩亏数累加进 \texttt{rank\_repairs}。若秩亏为零但基被改变则视为失败。
对应源码为 \srcpath{src/engine/kernel/lp\_kernel/native\_dual/factor.cpp:BasisFactor::rebuild}；
后端的隐式逻辑列修复为
\srcpath{src/engine/kernel/linear\_algebra/hfactor\_backend.cpp:HFactorBackend::factorize\_with\_logicals\_impl}。

\paragraph{检查型求解与后向误差门.}
\texttt{checked\_ftran}/\texttt{checked\_btran} 的接受准则是分量无关的后向误差：
\begin{equation}
  \|r\|_{\infty} \;\le\; 256\,\varepsilon \cdot
  \max\bigl(1,\; \|B\|_{\infty}\|x\|_{\infty} + \|b\|_{\infty}\bigr),
  \qquad r = b - Bx\ \text{或}\ b - B^{\top} y,
  \label{eq:lpkernel:backward}
\end{equation}
其中 $\|B\|_{\infty}$（转置情形用 $\|B^{\top}\|_{\infty}$）由当前基列绝对值和
维护，进/出基时增量更新且只增不减（保守上界，INVERT 时精确重算）。
对应源码为 \srcpath{src/engine/kernel/lp\_kernel/native\_dual/factor.cpp:BasisFactor::backward\_error\_acceptable}、
\srcpath{src/engine/kernel/lp\_kernel/native\_dual/factor.cpp:BasisFactor::update\_norms\_after\_exchange}；
常数 \texttt{kBackwardErrorMultiplier}${}=256$ 见
\srcpath{src/engine/kernel/lp\_kernel/native\_dual/factor.cpp:BasisFactor::solve\_checked} 同文件匿名命名空间。
不满足 \eqref{eq:lpkernel:backward} 时执行一轮迭代改进：残差用
\texttt{long double} 重组装，以同一因子解校正量并加到解上，再复检
\eqref{eq:lpkernel:backward}；改进路径置 \texttt{needs\_rebuild=true}。
对应源码为 \srcpath{src/engine/kernel/lp\_kernel/native\_dual/factor.cpp:BasisFactor::refine\_checked}。

\paragraph{基更新.}
生产路径 \texttt{update\_indexed} 不接收显式方向向量：它要求此前恰好各执行过一次
\texttt{ftran\_for\_update} 与 \texttt{btran\_for\_update}（capture 语义），
由后端在其捕获的 AQ/EP 工作区上完成 Forrest--Tomlin 更新；捕获序号与因子代数
不匹配则拒绝（响亮失败）。对应源码为
\srcpath{src/engine/kernel/lp\_kernel/native\_dual/factor.cpp:BasisFactor::update\_indexed}、
\srcpath{src/engine/kernel/linear\_algebra/hfactor\_backend.cpp:HFactorBackend::update\_captured}。
\texttt{needs\_rebuild} 透传后端的 \texttt{refactor\_hint\_}（HFactor 在
Update 阶段报告的数值不稳定提示）。
对应源码为 \srcpath{include/mipsolvers/engine/kernel/linear\_algebra/hfactor\_backend.hpp:HFactorBackend}。

\paragraph{跨调用复用.}
仅当矩阵对象地址不变、基位置序逐元素相同、且持久化的 eta 计数等于当前更新计数时，
\texttt{reusable\_for} 允许下一次求解复用该因子（revised-simplex 重优化契约）。
对应源码为 \srcpath{src/engine/kernel/lp\_kernel/native\_dual/factor.cpp:BasisFactor::reusable\_for}。

% ─────────────────────────────────────────────────────────────────────────────
\subsection{主循环：步进序列与不变量}
\label{subsec:engine-lp-mainloop}

\texttt{solve} $\to$ \texttt{solve\_impl} $\to$ \texttt{run\_phase}。
\texttt{run\_phase} 是 major/minor 双层循环：外层每次先做
\texttt{major\_rebuild}（INVERT + 状态重构，见 \ref{subsec:engine-lp-rebuild} 节），
内层反复执行 \texttt{minor\_iteration}，直到出基/入基失败、达到限值或触发重建。

\paragraph{minor 迭代：提案—认证—提交.}
一个 dual minor 迭代按严格三段执行（S3 管线）：
\begin{enumerate}
  \item \textbf{提案}（只读状态 + scratch）：\texttt{choose\_leaving} 选出基行；
        以打包 PRICE 计算主元行 $\hat a^{\top} = \rho^{\top} A$
        （$\rho = B^{-\top} e_p$ 为 pivotal BTRAN 结果）；
        \texttt{choose\_entering\_bfrt} 完成 BFRT 比率测试并产出事务
        \texttt{PivotTransaction}（入基列、翻转集、工作成本平移集、BFRT RHS）；
        对入基列做 FTRAN 得方向 $d = B^{-1} A_q$；
  \item \textbf{认证}（不改状态）：主元行有限性；
        \texttt{certify\_bfrt\_dual\_feasibility} 解析验证更新后既约成本
        \begin{equation}
          \tilde d_j = d_j + \Delta_{\text{dual}}\, \hat a_j + s_j,
          \qquad \Delta_{\text{dual}} = \mathrm{side}\cdot\theta,
          \label{eq:lpkernel:rc-update}
        \end{equation}
        对每个非基、非入基列满足 $\operatorname{sign}(m_j') \tilde d_j \le
        \varepsilon_{\text{opt}}$（$s_j$ 为本事务的工作成本平移，翻转列的
        $m_j'$ 已反号）；
        \texttt{certify\_bfrt\_leaving\_agreement} 验证 BFRT 翻转的 FTRAN 像
        $\delta = B^{-1} r_{\text{flip}}$ 与出基行变化一致：剩余步
        \begin{equation}
          \Delta_{\text{rem}} = x_{B,p} - \delta_p - \ell_{\text{leave}}
          \quad(\text{或上界}),\qquad
          \mathrm{side}\cdot\Delta_{\text{rem}} \ge
          -\, 1024\,\varepsilon \max\bigl(1, |x_{B,p}|, |\delta_p|,
          |\mathrm{bound}|\bigr),
          \label{eq:lpkernel:leaving-agreement}
        \end{equation}
        容差内的错号噪声截断为 $0$；完全覆盖（\texttt{covered\_violation}
        $=$ \texttt{violation}）时 $\Delta_{\text{rem}} = 0$；
  \item \textbf{提交}（原子）：\texttt{update\_indexed} 是唯一可失败的提交步骤；
        成功后依次发布：翻转列 \texttt{move} 反号、出基列 \texttt{move} 按其
        可入基侧重置、入基列 \texttt{move=Fixed}、\texttt{basis}$[p]\leftarrow q$、
        循环签名增量更新、DSE 权重事务写入、\texttt{basic} 位翻转、行分区交换、
        原始值事务 $x_B \leftarrow x_B - \delta - d\cdot\Delta_{\text{primal}}$
        （出基位置取 $\mathrm{bound} + \Delta_{\text{primal}}$，
        $\Delta_{\text{primal}} = \Delta_{\text{rem}} / d_p$）、
        既约成本按 \eqref{eq:lpkernel:rc-update} 应用（基列、入基列置 $0$）、
        工作成本平移累加进 \texttt{cost\_shift} 日志并置
        \texttt{costs\_shifted}。
\end{enumerate}
对应源码为 \srcpath{src/engine/kernel/lp\_kernel/native\_dual/solver.cpp:minor\_iteration}、
\srcpath{src/engine/kernel/lp\_kernel/native\_dual/solver.cpp:certify\_bfrt\_dual\_feasibility}、
\srcpath{src/engine/kernel/lp\_kernel/native\_dual/solver.cpp:certify\_bfrt\_leaving\_agreement}、
\srcpath{src/engine/kernel/lp\_kernel/native\_dual/solver.cpp:build\_primal\_transaction}。

\begin{boundarynote}
提交段的任何认证失败都返回 \texttt{NumericalTrouble} 并触发外层重建；只有
\texttt{update\_indexed} 失败会留下``因子已更新、状态未发布''的半提交，但此时
因子仍与旧基一致以外的唯一变化是拒绝交换——HFactor 端拒绝即未修改，因此状态
保持自洽。该顺序是刻意的：所有可失败步骤要么在提交前（认证），要么是唯一提交点。
\end{boundarynote}

\paragraph{入基列主元读取.}
$d_p$ 优先从后端捕获的 AQ 工作区 $O(1)$ 读取（\texttt{captured\_aq\_value}），
捕获失效时退回打包 FTRAN 像的坐标查询；$d_p$ 非有限或为 $0$ 均为数值失败。
对应源码为 \srcpath{src/engine/kernel/lp\_kernel/native\_dual/solver.cpp:minor\_iteration}、
\srcpath{src/engine/kernel/lp\_kernel/native\_dual/factor.cpp:BasisFactor::captured\_aq\_value}。

\paragraph{限值与中切断界.}
内层循环首检 \texttt{max\_iter} 与墙钟/\texttt{cancel\_flag}；中断时调用
\texttt{certify\_interrupted\_dual\_bound}：仅 Phase II 下，丢弃全部稳定化增量
（\texttt{restore\_original\_cost}）、以检查型求解重构状态、按原既约成本重分类
非基侧标、并以同一 \texttt{audit} 认证对偶可行。成功时
\texttt{dual\_bound\_certified=true}，\texttt{objective\_const - max\_objective}
是原 LP 最优值的严格下界（弱对偶）。Phase II 且提供
\texttt{incumbent\_bound} 时每 \texttt{incumbent\_check\_interval} 次迭代检查
目标截断 $\mathrm{obj} \ge \mathrm{incumbent} - \varepsilon_{\text{opt}}$。
对应源码为 \srcpath{src/engine/kernel/lp\_kernel/native\_dual/solver.cpp:certify\_interrupted\_dual\_bound}、
\srcpath{src/engine/kernel/lp\_kernel/native\_dual/solver.cpp:run\_phase}。

% ─────────────────────────────────────────────────────────────────────────────
\subsection{出基选择（CHUZR）}
\label{subsec:engine-lp-chuzr}

\paragraph{DSE 启发式与惰性堆.}
默认出基按 DSE 价值量
\begin{equation}
  \mathrm{merit}(p) = \frac{v_p^{2}}{w_p},
  \qquad v_p = \texttt{primal\_infeasibility}(p),
  \label{eq:lpkernel:chuzr-merit}
\end{equation}
在惰性最大堆中取最大者（同价值量取较小行号）；堆项带版本号，过期项在弹堆时丢弃。
\texttt{primal\_changes} 事务行增量刷新堆（\texttt{refresh\_leaving\_heap}），
堆项数超过 $4m$ 或任何刷新失败则整堆失效、下次全量重建。
弹出候选若是 taboo 行则记录跳过并继续；所有候选均 taboo 时退回最早到期的
taboo 行并调用 \texttt{release\_taboo\_row}（aspiration：把定价时钟拨过其任期）。
对应源码为 \srcpath{src/engine/kernel/lp\_kernel/native\_dual/pricing.cpp:choose\_leaving}、
\srcpath{src/engine/kernel/lp\_kernel/native\_dual/pricing.cpp:push\_leaving\_row}、
\srcpath{src/engine/kernel/lp\_kernel/native\_dual/pricing.cpp:refresh\_leaving\_heap}、
\srcpath{src/engine/kernel/lp\_kernel/native\_dual/state.cpp:release\_taboo\_row}。

\paragraph{权重核验.}
选出出基行后执行 pivotal BTRAN 得 $\rho$；SteepestEdge 模式下比较缓存权重与
精确权重 $\|\rho\|_2^2$，若
\begin{equation}
  \bigl|\|\rho\|_2^2 - w_p\bigr| > 2048\,\varepsilon \max\bigl(1, \|\rho\|_2^2\bigr)
\end{equation}
则以精确值覆写缓存并重推堆项。该权重只是排序启发，不是可行性证书——主元有效性
由行/列一致性与重构独立保证。对应源码为
\srcpath{src/engine/kernel/lp\_kernel/native\_dual/pricing.cpp:choose\_leaving}。

\paragraph{认证精确 CHUZR（可选模式）.}
\texttt{DualEdgeWeightInitialization::CertifiedExact} 下改用
\texttt{choose\_leaving\_certified\_exact}：先用 Cauchy--Schwarz 与检查型 BTRAN
接受域构造每行的权重下界
\begin{equation}
  \underline{w}_p = \Bigl(\frac{1 - 256\varepsilon}
      {\|A_{\texttt{basis}[p]}\|_2 + (256\varepsilon + \gamma_m)
      \|B^{\top}\|_{\infty}}\Bigr)^{2},
  \qquad \gamma_m = \frac{m\varepsilon}{1 - m\varepsilon},
\end{equation}
得价值量上界 $v_p^2/\underline{w}_p$，按降序只对可能夺冠的行做检查型
\texttt{basis\_inverse\_row\_sparse\_entries}，以精确 $\|\rho\|_2^2$ 重算价值量并
淘汰（\texttt{certified\_dse\_rejections} 计数）。对应源码为
\srcpath{src/engine/kernel/lp\_kernel/native\_dual/pricing.cpp:choose\_leaving\_certified\_exact}、
\srcpath{src/engine/kernel/lp\_kernel/native\_dual/factor.cpp:BasisFactor::basis\_inverse\_row\_sparse\_entries}。

\paragraph{Bland 回退.}
\texttt{bland\_active} 时 \texttt{choose\_leaving\_bland} 直接取最小列号的
原始不可行基行，绕过堆与 taboo 表，并使堆失效。对应源码为
\srcpath{src/engine/kernel/lp\_kernel/native\_dual/pricing.cpp:choose\_leaving\_bland}。

% ─────────────────────────────────────────────────────────────────────────────
\subsection{入基选择：BFRT 比率测试}
\label{subsec:engine-lp-bfrt}

\texttt{choose\_entering\_bfrt} 是有界变量翻转比率测试（bound-flipping ratio test）
的精确断点实现。记出基侧标 $s_r \in \{-1, +1\}$、非基列 $j$ 的方向
$d^{m}_j = \operatorname{sign}(m_j)$、主元行分量 $\hat a_j$。

\paragraph{候选分类.}
对每个活跃列（非基且 $m_j \ne 0$）：
\begin{equation}
  \alpha_j = s_r \, d^{m}_j \, \hat a_j,\qquad
  e^{\text{cheap}}_j = C_j \cdot \max_k |\rho_k|,
  \label{eq:lpkernel:bfrt-alpha}
\end{equation}
其中 $C_j$ 是 \texttt{initialize} 预计算的每列主导系数
\begin{equation}
  C_j = 2\bigl(\gamma_{k_j} + 256\varepsilon\bigr) \|A_j\|_1,\qquad
  \gamma_{k} = \frac{k\varepsilon}{1 - k\varepsilon},
  \label{eq:lpkernel:bfrt-coef}
\end{equation}
保证 $C_j \max|\rho| \ge$ 精确点积误差界
$\texttt{dot\_error\_bound}(j) = \gamma_{k_j} \sum_i |\rho_i A_{ij}| +
256\varepsilon \sum_i |\rho_i A_{ij}|$。
对应源码为 \srcpath{src/engine/kernel/lp\_kernel/native\_dual/pricing.cpp:dot\_error\_bound}、
\srcpath{src/engine/kernel/lp\_kernel/native\_dual/state.cpp:initialize}。

稳定性预滤使用随因子龄期升档的主元门限
\begin{equation}
  \tau_{\text{stab}} =
  \begin{cases}
    10^{-9}, & \texttt{updates\_since\_rebuild} < 10,\\
    3\times 10^{-8}, & 10 \le \texttt{updates\_since\_rebuild} < 20,\\
    10^{-6}, & \texttt{updates\_since\_rebuild} \ge 20.
  \end{cases}
  \label{eq:lpkernel:stable-tol}
\end{equation}
分类规则（Phase II 快速路径）：$\alpha_j \le 0$ 且
$\alpha_j + e^{\text{cheap}}_j \le 0$ 的列直接排除；$0 < \alpha_j \le
\tau_{\text{stab}}$ 的列计入 \texttt{unstable\_pivot\_rejections} 并置
\texttt{stability\_blocked}；$\alpha_j > \tau_{\text{stab}}$ 且
$\alpha_j > e^{\text{cheap}}_j$ 的列获廉价认证（默认策略 M）；仅边界情形
（\texttt{MIPSOLVERS\_DS\_EXACT\_RATIO} 开启的策略 E，或 lean 模式下
$\alpha_j \le e^{\text{cheap}}_j$ 者）才消费精确 \texttt{dot\_error\_bound}。
认证列的 margin 与断点为
\begin{equation}
  \mathrm{margin}_j = \max\bigl(0, -d^{m}_j d_j\bigr),\qquad
  \theta_j = \mathrm{margin}_j / \alpha_j
  \quad(\text{long double}),
  \label{eq:lpkernel:breakpoint}
\end{equation}
$\theta_j$ 非有限即失败。对应源码为
\srcpath{src/engine/kernel/lp\_kernel/native\_dual/pricing.cpp:choose\_entering\_bfrt}。

\begin{boundarynote}
Phase I（\texttt{Phase::DualOne}）或未持误差系数时走通用的
\texttt{BfrtScanEvaluation} 逐列求值路径（同一组判定式，支持
\texttt{kernel\_threads}$>1$ 且候选数 $\ge 8192$ 时线程池并行；合并按原 PRICE
顺序进行，所有容量和、计数与插入序与线程数逐位无关）。默认 Phase II 走
\texttt{classify\_candidate} 分支自由路径；\texttt{MIPSOLVERS\_DS\_SPLIT\_BFRT\_SCAN}
另启 aarch64 NEON 预滤变体（同一谓词的向量化实现）。
对应源码为 \srcpath{src/engine/kernel/lp\_kernel/native\_dual/pricing.cpp:prefilter\_active\_phase\_two}。
\end{boundarynote}

\paragraph{容量行走与翻转组选择.}
候选按 $\theta$ 升序分组（组内按列号升序；前 8 组用选择算法，之后全排序）；
从 $\mathrm{covered} = 10^{-12}$（\texttt{kInitialTotalChange}，EXPAND 容差）
起逐组累加组容量 $\sum \alpha_j \mathrm{range}_j$，首个使
\begin{equation}
  \mathrm{covered} + \mathrm{group\_change} \ge v_p
  \quad\text{或组内含无限 range}
  \label{eq:lpkernel:capacity-walk}
\end{equation}
的组为选中组；此前各组全部有限 range 候选进入翻转集。容量耗尽
（无选中组）时令事务为空，由调用方走不可行证书分支；Phase I 下还会先做
投影残差局部证明
\begin{equation}
  v_p - \mathrm{covered} \le |\rho^{\top}(b - Ax)| +
  \Bigl|\sum_p r^{B}_p (x_{B,p} - x^{0}_{\texttt{basis}[p]})\Bigr|
  + (k + m + 256)\varepsilon\max(\cdot) + E_{\text{stab}},
\end{equation}
成立才接受最终恰好覆盖组。对应源码为
\srcpath{src/engine/kernel/lp\_kernel/native\_dual/pricing.cpp:choose\_entering\_bfrt}。

\paragraph{Harris 两段.}
第一段在选中组起的候选后缀上放宽比率界：
\begin{equation}
  \bar\theta = \min_{j \ge \text{选中组}}\;
  \frac{\varepsilon_{\text{opt}} - d^{m}_j d_j}{\alpha_j}
  \quad(\text{long double}).
  \label{eq:lpkernel:harris1}
\end{equation}
第二段在 $\theta_j \le \bar\theta$ 的候选中取最大 $\alpha_j$（同值取小列号）；
taboo 候选被跳过并计入 \texttt{taboo\_rejections}，全部 taboo 时置
\texttt{all\_harris\_candidates\_taboo} 由调用方把出基行加入 taboo 表。
胜出者给出
\begin{equation}
  q = \mathrm{col},\quad
  \texttt{entering.pivot} = s_r\, d^{m}_q\, \alpha_q,\quad
  \theta = \theta_q .
  \label{eq:lpkernel:harris2}
\end{equation}
对应源码为 \srcpath{src/engine/kernel/lp\_kernel/native\_dual/pricing.cpp:choose\_entering\_bfrt}；
Bland 模式下改为取最小断点（同值取小列号）、无翻转、忽略 taboo，见同函数内
\texttt{bland\_active} 分支。

\paragraph{翻转/平移与对偶步区间.}
终扫把每个``更新后带号既约成本仍超差''的非基列
$d^{m}_j d_j + \alpha_j \theta > \varepsilon_{\text{opt}}$ 分两路处理：
有限正 range 者追加进翻转集，单边列（range 无限）追加工作成本平移
$\Delta c_j = -\bigl(d_j + s_r \theta \hat a_j\bigr)$（更新后既约成本精确归零，
由终止清理移除）。同一遍扫描累积可接受对偶步区间
\begin{equation}
  \theta \in \Bigl[\max_{j\ \text{翻转}}
  \frac{-\varepsilon_{\text{opt}} - d^{m}_j d_j}{\alpha_j},\;
  \min_{j\ \text{未翻转}}
  \frac{\varepsilon_{\text{opt}} - d^{m}_j d_j}{\alpha_j}\Bigr],
  \qquad \alpha_j > 0,
  \label{eq:lpkernel:step-interval}
\end{equation}
$\theta$ 越界即数值失败（事务自相矛盾）。对应源码为
\srcpath{src/engine/kernel/lp\_kernel/native\_dual/pricing.cpp:choose\_entering\_bfrt}。

\paragraph{BFRT 右端与覆盖截断.}
翻转集物化右端
\begin{equation}
  r_{\text{flip}} = \sum_{j \in \text{flips}}
  \operatorname{sign}(m_j^{\text{old}})\, \mathrm{range}_j\, A_{\cdot j},
  \label{eq:lpkernel:bfrt-rhs}
\end{equation}
按行号排序后同行累加。覆盖量
$\mathrm{cov} = s_r \cdot \rho^{\top} r_{\text{flip}}$ 与存储的 $v_p$ 求和次序不同，
故接受域带舍入包络
\begin{equation}
  -S \le \mathrm{cov} \le v_p + S,\qquad
  S = (\mathrm{nnz}(r_{\text{flip}}) + 256)\,\varepsilon
      \max\bigl(1, |\rho^{\top} r_{\text{flip}}|_{\text{abs}}, v_p\bigr),
  \label{eq:lpkernel:coverage}
\end{equation}
然后 \texttt{covered\_violation} $= \min(\max(\mathrm{cov}, 0), v_p)$。
对应源码为 \srcpath{src/engine/kernel/lp\_kernel/native\_dual/pricing.cpp:choose\_entering\_bfrt}。

% ─────────────────────────────────────────────────────────────────────────────
\subsection{DSE / Devex 权重更新}
\label{subsec:engine-lp-dse}

\paragraph{Goldfarb--Reid 递推.}
SteepestEdge 模式下 \texttt{compute\_dse\_weights} 先以捕获模式的辅助 FTRAN 计算
$\tau = B^{-1}\rho$（仅定义在入基方向 $d$ 的支撑上，并要求两者支撑等长），
然后对 $d$ 支撑上每个非出基行
\begin{equation}
  w_i' = w_i - 2\,\frac{d_i}{d_p}\,\tau_i
         + \Bigl(\frac{d_i}{d_p}\Bigr)^{2} w_p,\qquad
  w_i \leftarrow \max\bigl(w_i',\; R_i\bigr),
  \label{eq:lpkernel:dse-recurrence}
\end{equation}
舍入地板 $R_i = 1024\,\varepsilon \max(1, |w_i|, |2 d_i \tau_i / d_p|,
|(d_i/d_p)^2 w_p|)$；出基行取 $w_p' = w_p / d_p^{2}$。
默认在 \texttt{sizeof(long double)==sizeof(double)} 的平台（MSVC、arm64）上
以双精度按同一次序求值；\texttt{MIPSOLVERS\_DS\_LEAN\_DSE} 启用单term地板
$1024\varepsilon\max(1, |w_i|)$；\texttt{MIPSOLVERS\_DS\_REPEATED\_DSE\_TERMS}
强制四term长双精度参考实现。对应源码为
\srcpath{src/engine/kernel/lp\_kernel/native\_dual/pricing.cpp:compute\_dse\_weights}。

\paragraph{Devex 框架.}
Devex 模式下精确权重只在参考框架内求和：出基行的
$\hat w = \max(1, \sum_{j \in \text{ref}} \hat a_j^2)$，若缓存值与其比值超过
$9$ 倍则重启框架（\texttt{initialize\_devex\_framework}：权重全 $1$、参考集重置为
当前基、\texttt{devex\_restarts} 计数）；非出基行取
$w_i \leftarrow \max\bigl(w_i, \hat w_p / d_p^2 \cdot d_i^2\bigr)$，出基行取
$\max(1, \hat w_p / d_p^2)$。对应源码为
\srcpath{src/engine/kernel/lp\_kernel/native\_dual/pricing.cpp:compute\_dse\_weights}、
\srcpath{src/engine/kernel/lp\_kernel/native\_dual/state.cpp:initialize\_devex\_framework}。

\paragraph{初始化.}
\texttt{initialize\_uncached\_edge\_weights} 按
\texttt{DualEdgeWeightInitialization} 选项分派：\texttt{FullExact}/
\texttt{CertifiedExact} 对每行做检查型 BTRAN 求 $\|B^{-\top} e_i\|_2^2$
（\texttt{compute\_exact\_edge\_weights}，批量，累计精化次数）；
\texttt{StructuralExact} 先尝试结构精确初始化——当每个基列都是单非零列时
$w_i = 1/A_{\texttt{basis}[i], i}^{2}$，否则退回 Devex 框架；\texttt{Production}
默认即 StructuralExact（除非 \texttt{exact\_dse\_initialization}）。
对应源码为 \srcpath{src/engine/kernel/lp\_kernel/native\_dual/state.cpp:initialize\_uncached\_edge\_weights}、
\srcpath{src/engine/kernel/lp\_kernel/native\_dual/state.cpp:initialize\_structural\_exact\_edge\_weights}、
\srcpath{src/engine/kernel/lp\_kernel/native\_dual/factor.cpp:BasisFactor::compute\_exact\_edge\_weights}。

% ─────────────────────────────────────────────────────────────────────────────
\subsection{major 重建与 INVERT 频率控制}
\label{subsec:engine-lp-rebuild}

\paragraph{频率.}
minor 循环在三种情形下跳出到 major 重建：
(i) 因子报告 \texttt{needs\_rebuild}（HFactor 更新提示，记
\texttt{RebuildReason::SyntheticWork}）；
(ii) 更新计数达到
\begin{equation}
  L_{\text{update}} = \max\bigl(50,\; \min\bigl(200,\; \max(1, m/4)\bigr)\bigr),
  \label{eq:lpkernel:update-limit}
\end{equation}
（\texttt{RebuildReason::UpdateLimit}）；
(iii) 出基/入基失败或数值故障，且当前不是全新重建（先重建一次再下结论）。
\texttt{reinvert\_after\_pivot}（上支 \texttt{update\_indexed} 时因子自报）也
强制 NumericalTrouble 重建。对应源码为
\srcpath{src/engine/kernel/lp\_kernel/native\_dual/solver.cpp:run\_phase}。

\paragraph{major\_rebuild 序列.}
\begin{enumerate}
  \item 需要时 \texttt{factor->rebuild}（INVERT，含秩修复），\texttt{reconstruct}
        全量重构（见下）；
  \item 记录漂移统计：$\|x_B^{\text{old}} - x_B\|_{\infty}$、
        $\|d^{\text{old}} - d\|_{\infty}$、目标差；
  \item 非基侧标全量重分类：$d_j > \varepsilon_{\text{opt}}$ 且上界有限
        $\Rightarrow$ Down，否则 Up（boxed 列借此恢复对偶可行侧，不改工作成本）；
        随后以精确残差再做一次原始重构（\texttt{exact\_residual=true}）；
  \item 剩余单边非基的对偶不可行以工作成本平移 $\Delta c_j = -d_j$ 精确归零
        （记入 \texttt{cost\_shift}，boxed 列到此仍不可行则失败）；
        有平移时再做一次对偶重构；
  \item \texttt{audit} 门检；重置更新计数；全量重算循环签名。
\end{enumerate}
对应源码为 \srcpath{src/engine/kernel/lp\_kernel/native\_dual/state.cpp:major\_rebuild}。

\paragraph{reconstruct 与缺陷校正.}
原始重构：$\hat b = b - A_N x_N(s)$（dual Phase I 锚定坐标下改从
$z = x - x^0$ 组装零右端），$x_B = B^{-1}\hat b$ 走检查型 FTRAN；
随后按周期（首两次重建或每 \texttt{intermediate\_audit\_interval} 代，或调用方
要求精确时）做 canonical 缺陷校正：当
$\|b - Ax\|_{\infty} > \varepsilon_{\text{feas}} \max(1, \|b\|_{\infty})$ 时反复
$x_B \mathrel{+}= B^{-1} r$，每轮要求无穷范数严格下降，最多
\texttt{digits}（双精度 53）轮，耗尽即失败。
对偶重构：$y = B^{-\top} c_B$（检查型 BTRAN），$d = \tilde c - A^{\top} y$，
基列既约成本置 $0$；目标按
$\tilde c^{\top}\ell + \sum_{j:\,\mathrm{Down}} \tilde c_j (u_j - \ell_j)
+ \sum_p \tilde c_{\texttt{basis}[p]} (x_{B,p} - \ell_{\texttt{basis}[p]})$ 重组。
对应源码为 \srcpath{src/engine/kernel/lp\_kernel/native\_dual/state.cpp:reconstruct}、
\srcpath{src/engine/kernel/lp\_kernel/native\_dual/state.cpp:correct\_canonical\_primal\_residual}、
\srcpath{src/engine/kernel/lp\_kernel/native\_dual/state.cpp:reconstruct\_objective}
（同文件静态函数 \texttt{reconstruct\_reduced\_costs}）。

% ─────────────────────────────────────────────────────────────────────────────
\subsection{退化处理：扰动、taboo 与 Bland}
\label{subsec:engine-lp-degen}

\paragraph{启动稳定化扰动.}
\texttt{initialize\_stabilized\_cost} 仅在 Phase II 且初始最大原始不可行
$> \varepsilon_{\text{feas}}$ 时扰动工作成本。逐列伪随机分数
$\phi_j \in [0,1)$ 由 splitmix64 型混合（列号加黄金比例常数，三次 xorshift-乘，
取高 53 位乘 $2^{-53}$）确定：
\begin{equation}
  \delta_j =
  \begin{cases}
    \operatorname{sign}(c_j)\,(1+\phi_j)(|c_j|+1)\, b_s, & j < n_{\text{orig}},\
      u_j\ \text{有限},\\[2pt]
    -(1+\phi_j)(|c_j|+1)\, b_s, & j < n_{\text{orig}},\ u_j = +\infty,\\[2pt]
    (0.5 - \phi_j)\, 10^{-12}, & j\ \text{为逻辑列},
  \end{cases}
  \label{eq:lpkernel:perturb}
\end{equation}
$b_s = 5\times 10^{-7} \max(1, \max_j |c_j|)$，且 $\max|c_j| > 100$ 时先取其
四次方根压缩尺度。扰动符号约定保持对偶可行（向有限侧推）。
对应源码为 \srcpath{src/engine/kernel/lp\_kernel/native\_dual/state.cpp:initialize\_stabilized\_cost}、
\srcpath{src/engine/kernel/lp\_kernel/native\_dual/state.cpp:deterministic\_fraction}（匿名命名空间）。

\paragraph{再扰动.}
Bland 计数器饱和（$\ge 2\max(50, m/2)$）而循环仍在时，
\texttt{reperturb\_stabilized\_cost} 换种子重播 \eqref{eq:lpkernel:perturb}，
幅度按 $2^{\min(6, k)}$ 几何升级（$k$ 为第几次再扰动；上限
$64 \times 5\times10^{-7} = 3.2\times10^{-5}$），次数上限
\texttt{kMaxReperturbations}${}=12$。对应源码为
\srcpath{src/engine/kernel/lp\_kernel/native\_dual/state.cpp:reperturb\_stabilized\_cost}、
\srcpath{src/engine/kernel/lp\_kernel/native\_dual/solver.cpp:run\_phase}。

\paragraph{循环签名与 taboo 表.}
循环签名是（基指派， 非基侧标）集合的 Zobrist 式双 64 位 XOR 散列：
基元素令牌 $(p, j)$、侧标令牌 $(j, \operatorname{sign}(m_j))$ 分别与盐
\texttt{kCycleBasisSaltA/B}、\texttt{kCycleMoveSaltA/B} 混合后 XOR 进两个累加器；
提交时以 $O(1)$ 令牌增量维护（一次基交换 + 每个翻转列两次 toggle），
major 重建后全量重同步。\texttt{record\_cycle\_arrival} 检出签名重复即计
\texttt{cycles\_detected} 并把该次离开的（出基列， 入基列）写入 taboo 变更表；
taboo 任期 $\min(2m+1, 4096)$，表容量 $65536$（FIFO + 票序压实）。
对应源码为 \srcpath{src/engine/kernel/lp\_kernel/native\_dual/state.cpp:record\_cycle\_arrival}、
\srcpath{src/engine/kernel/lp\_kernel/native\_dual/state.cpp:compute\_cycle\_signature}（匿名命名空间）、
\srcpath{src/engine/kernel/lp\_kernel/native\_dual/state.cpp:add\_taboo\_row}、
\srcpath{src/engine/kernel/lp\_kernel/native\_dual/model.hpp:State}。

\paragraph{Bland 触发与耗尽.}
\texttt{record\_cycle\_arrival} 报告精确基重复（或一次 CycleBlocked 空转迭代）
使停滞计数器加一（上限 $2T$，$T = \max(50, m/2)$），新鲜基使其减一；
计数 $\ge T$ 时 \texttt{bland\_active=true}，回到 $0$ 时退出。
Bland 激活、再扰动配额耗尽、计数器封顶三者同时成立时，一次 CycleBlocked
直接以 \texttt{IterationLimit} 终止并附 \texttt{certify\_interrupted\_dual\_bound}
的中断界——这是防 livelock 的耗尽守卫（代码注释记录了 blend2 上
$9.8\times10^{6}$ 次空转对 21 次主元的实测动机）。退化步统计：
$|\theta| < 10^{-9}$ 计 \texttt{degenerate\_dual\_steps}，
$|\Delta_{\text{primal}}| < 10^{-9}$ 计 \texttt{degenerate\_primal\_steps}。
对应源码为 \srcpath{src/engine/kernel/lp\_kernel/native\_dual/solver.cpp:run\_phase}、
\srcpath{src/engine/kernel/lp\_kernel/native\_dual/solver.cpp:minor\_iteration}。

% ─────────────────────────────────────────────────────────────────────────────
\subsection{证书生成}
\label{subsec:engine-lp-cert}

\paragraph{区间 Farkas 证书.}
给定乘子 $\pi \in \mathbb{R}^m$，\texttt{certificate\_from\_multiplier} 计算
\begin{equation}
  \rho = \pi^{\top} b,\qquad
  \kappa_j = \pi^{\top} A_{\cdot j},\qquad
  [\underline{A}, \overline{A}] =
  \sum_{j:\,\kappa_j \ge 0} \kappa_j [\ell_j, u_j] +
  \sum_{j:\,\kappa_j < 0} \kappa_j [u_j, \ell_j],
  \label{eq:lpkernel:farkas-interval}
\end{equation}
（无界端按 $\pm\infty$ 传播），误差界
$E = 1024\,\varepsilon \max\bigl(1, |\rho| + \sum_{j:\,u_j\,\text{有限}}
|\kappa_j u_j|\bigr)$；分离裕量取
$\max(\underline{A} - \rho,\; \rho - \overline{A})$，较大者在异侧时整体反号。
证书有效当且仅当乘子与 $\rho$ 有限且 $\mathrm{margin} > E$。
对应源码为 \srcpath{src/engine/kernel/lp\_kernel/native\_dual/certificate.cpp:certificate\_from\_multiplier}
（匿名命名空间）。

\paragraph{两个来源.}
(i) dual Phase II 终端（CHUZC 无入基列）：乘子取已物化的 pivotal 行
$\rho = B^{-\top} e_p$（factor-resident 行在此稀有出口被打包）；
(ii) 原始 Phase I 人工目标最优且人工变量和 $> \varepsilon_{\text{feas}}$：
乘子取 $B^{-\top} c_B$（人工成本向量）的检查型 BTRAN。
对应源码为 \srcpath{src/engine/kernel/lp\_kernel/native\_dual/certificate.cpp:primal\_infeasibility\_certificate}、
\srcpath{src/engine/kernel/lp\_kernel/native\_dual/certificate.cpp:phase\_one\_farkas\_certificate}；
组装到结果为 \srcpath{src/engine/kernel/lp\_kernel/native\_dual/solver.cpp:attach\_farkas\_certificate}
（匿名命名空间）。包装层把乘子按行符号还原并按 不等式/等式 拆分为
\texttt{farkas\_ray}/\texttt{farkas\_ray\_eq}。
对应源码为 \srcpath{src/engine/kernel/lp\_kernel/dual\_simplex.cpp:solve\_lp\_from\_sf\_impl}。

\paragraph{无界射线证书.}
原始单纯形比率测试找不到有限限制界时，
\texttt{certify\_primal\_unbounded\_ray} 组装
\begin{equation}
  r_q = m_q,\qquad r_{\texttt{basis}[p]} = -m_q\, d_p
  \ \ (p \in \operatorname{supp} d),
\end{equation}
按无穷范数归一，然后三重校验：每个有限界方向的衰退锥分量不越
$\varepsilon_{\text{feas}}$；$\|A r\|_{\infty} \le \varepsilon_{\text{feas}}
\max(1, \max_i \sum_j |A_{ij} r_j|)$（绝对活动度用 \texttt{long double} 累加）；
目标增益 $c^{\top} r > \varepsilon_{\text{opt}} \max(1,
\|c\|_{\infty} \|r\|_{\infty})$。仅在原成本 Phase II 下出具。
对应源码为 \srcpath{src/engine/kernel/lp\_kernel/native\_dual/primal.cpp:certify\_primal\_unbounded\_ray}
（匿名命名空间）。

% ─────────────────────────────────────────────────────────────────────────────
\subsection{求解编排：冷/热启动与清理}
\label{subsec:engine-lp-driver}

\paragraph{冷启动.}
无基提示时：\texttt{build\_logical\_basis} 取每行 slack/artificial 逻辑列；
\texttt{apply\_certified\_singleton\_crash} 把``列单非零且
$b_i/a_{ij} \in [0, u_j]$''的原变量列换入基（同行取最小列号），换入后重建因子、
重构并初始化权重。然后依次尝试：
\begin{enumerate}
  \item \texttt{normalize\_nonbasic\_moves}：按 $d_j > \varepsilon_{\text{opt}}$
        且上界有限分派 Down，否则 Up；单边正既约成本列无法分派时失败；
        成功则审计并直接进入 dual Phase II；
  \item \texttt{decide\_cost\_shifted\_dual\_start} 自适应门：需要的平移数
        （单边非基、$d_j > \varepsilon_{\text{opt}}$、上界无限的列数）
        低于 $\max(\texttt{dual\_shift\_start\_max\_count},
        \lceil f \cdot n \rceil)$（$f$ =
        \texttt{dual\_shift\_start\_max\_fraction}）时，
        \texttt{initialize\_cost\_shifted\_dual\_start} 把这些列的工作成本
        平移 $\Delta c_j = -d_j$、既约成本归零并置 Up，直接进 Phase II；
        环境变量 \texttt{MIPSOLVERS\_DUAL\_SHIFT\_START}=on/off 可强制；
  \item 否则进入真 dual Phase I：锚点 $x^0$ 由逻辑列独子结构
        $x^0_{\texttt{logical}[i]} = b_i / a_{i,\texttt{logical}[i]}$ 构造并校验
        $\|A x^0 - b\|_{\infty} \le \varepsilon_{\text{feas}} \max(1,
        \|b\|_{\infty})$；在齐次坐标 $z = x - x^0$ 上，无有限上界的可入基列
        取区间 $[x^0_j, x^0_j + 1]$，boxed/fixed 列固定于 $x^0_j$；Phase I 目标为
        $\sum_{j\ \text{非基}} (v_j - x^0_j)\, d_j$（\texttt{long double} 累加）。
        Phase I 结束后 \texttt{transition\_dual\_phase\_one\_to\_two} 要求原始界
        对偶缺陷计数为零，恢复 Phase II 界并重构；
        dual Phase I 失败（无法给出原界对偶可行基，或 Phase II 数值失败）时
        恢复 Phase I 前的基，改走原始人工 Phase I。
\end{enumerate}
对应源码为 \srcpath{src/engine/kernel/lp\_kernel/native\_dual/solver.cpp:solve\_impl}、
\srcpath{src/engine/kernel/lp\_kernel/native\_dual/state.cpp:decide\_cost\_shifted\_dual\_start}、
\srcpath{src/engine/kernel/lp\_kernel/native\_dual/state.cpp:initialize\_cost\_shifted\_dual\_start}、
\srcpath{src/engine/kernel/lp\_kernel/native\_dual/state.cpp:initialize\_dual\_phase\_one}、
\srcpath{src/engine/kernel/lp\_kernel/native\_dual/state.cpp:dual\_phase\_one\_objective}、
\srcpath{src/engine/kernel/lp\_kernel/native\_dual/state.cpp:transition\_dual\_phase\_one\_to\_two}、
\srcpath{src/engine/kernel/lp\_kernel/native\_dual/state.cpp:make\_dual\_phase\_one\_anchor}。

\paragraph{热启动.}
有基提示时：因子复用判定见 \ref{subsec:engine-lp-factor} 节；
\texttt{normalize\_nonbasic\_moves} 失败则尝试平移启动；
\texttt{allow\_phase\_one=false}（\texttt{solve\_phase2} 入口）直接以
\texttt{DualInfeasibleStart} 返回；否则视
\texttt{allow\_warm\_primal\_phase\_one} 走人工 Phase I 或报错。
对应源码为 \srcpath{src/engine/kernel/lp\_kernel/native\_dual/solver.cpp:solve\_impl}、
\srcpath{src/engine/kernel/lp\_kernel/native\_dual/solver.cpp:solve\_phase2}。

\paragraph{原始 Phase I 回退与 primal 内核.}
\texttt{run\_primal\_phase}（\srcpath{src/engine/kernel/lp\_kernel/native\_dual/primal.cpp:run\_primal\_phase}）
是同一个因子上的原始单纯形：入基取堆顶（度量
$\mathrm{gain}^2 / w$ 或纯 gain，\texttt{MIPSOLVERS\_PRIMAL\_DEVEX}
控制 Devex/PSE/shadow/off），比率测试 \texttt{choose\_leaving\_harris} 用
放宽步 $\max(0, \mathrm{dist} + \varepsilon_{\text{feas}})/|\mathrm{change}|$
加最小主元 $\sqrt{\varepsilon}\max(1, \|d\|_{\infty})$ 过滤，第二段取满足
放宽界的最大主元；有限 $q$ 范围小于步长时做界翻转不换基；
提交前以 \texttt{primal\_pivot\_identity\_acceptable} 校验行/列主元一致性
$|\alpha_{\text{col}} - \alpha_{\text{row}}| \le \sqrt{\varepsilon}
\max(|\alpha_{\text{col}}|, |\alpha_{\text{row}}|)$。
原始 INVERT 间隔硬编码 \texttt{kPrimalRebuildInterval}${}=256$；
Devex 坏权重比 $>3$ 累计 $>3$ 次重启框架；PSE 模式限 $n \le 50000$。
任何新鲜失败且在更新链中段（\texttt{updates\_since\_rebuild}$>0$）都先
\texttt{rebuild\_primal\_state} 再重试；\texttt{allow\_strong\_rebuild}
（仅原成本清理调用）允许把 Markowitz 主元阈值升到 $0.5$
（\texttt{strengthen\_pivot\_threshold}）后再重建一次。
对应源码为 \srcpath{src/engine/kernel/lp\_kernel/native\_dual/primal.cpp:run\_primal\_phase}、
\srcpath{src/engine/kernel/lp\_kernel/native\_dual/primal.cpp:choose\_leaving\_harris}（匿名命名空间）、
\srcpath{src/engine/kernel/lp\_kernel/native\_dual/primal.cpp:rebuild\_primal\_state}（匿名命名空间）、
\srcpath{src/engine/kernel/lp\_kernel/native\_dual/primal\_pricing.hpp:primal\_pivot\_identity\_acceptable}、
\srcpath{src/engine/kernel/lp\_kernel/native\_dual/primal\_pricing.hpp:prepare\_primal\_devex\_update}、
\srcpath{src/engine/kernel/lp\_kernel/native\_dual/primal\_pricing.hpp:prepare\_primal\_pse\_update}。

\paragraph{终止清理.}
Phase II 终端（新鲜重建后 CHUZR 无出基行）时，若工作成本被扰动/平移过
（\texttt{working\_cost\_is\_original} 为假），先
\texttt{restore\_original\_cost} 丢弃全部稳定化增量并重构；清理审计不过但
原始可行性仍在时，以 \texttt{run\_primal\_phase(allow\_strong\_rebuild=true)}
做无扰动原始清理；最终 \texttt{audit(require\_primal, require\_dual,
exact\_residual=true)} 通过才发布 \texttt{Optimal}。
对应源码为 \srcpath{src/engine/kernel/lp\_kernel/native\_dual/solver.cpp:run\_phase}、
\srcpath{src/engine/kernel/lp\_kernel/native\_dual/state.cpp:restore\_original\_cost}、
\srcpath{src/engine/kernel/lp\_kernel/native\_dual/state.cpp:working\_cost\_is\_original}。

% ─────────────────────────────────────────────────────────────────────────────
\subsection{portfolio 竞速与 abort 语义}
\label{subsec:engine-lp-portfolio}

\texttt{NativeAutoLPAdapter::solve\_with\_context} 把单次 LP 调用组织为
双工作者竞速：
\begin{itemize}
  \item \textbf{工作者}：线程 A 跑原生对偶单纯形
        （\texttt{NativeDualSimplexLPAdapter}：\texttt{FullExact} DSE 初始化 +
        HiGHS presolve 开启），线程 B 跑直连 IPM
        （\texttt{NativeIPMLPAdapter}，\texttt{presolve=false}）；
        两者共享同一个 \texttt{PortfolioSlot} 的 \texttt{cancel} 原子标志与
        调用级截止时刻；
  \item \textbf{胜出规则}（\texttt{publish}）：锁内 \texttt{finished} 加一；
        尚无胜者且（本结果 \texttt{success} 或已两人完成）时，本结果成为胜者，
        并置 \texttt{cancel}——即``首个成功即胜；双败时后到者胜''；
  \item \textbf{abort 语义}：取消是协作式的——\texttt{cancel} 经
        \texttt{SimplexOptions::cancel\_flag}/\texttt{IPMLPOptions::cancel\_flag}
        传入，内核在迭代与阶段边界轮询；进行中的库分解调用必须跑完。
        主线程在条件变量上等到胜者或截止；超时/取消时返回
        \texttt{Time limit}/\texttt{Cancelled}，置 \texttt{cancel} 后
        \textbf{总是 join 两个工作者再返回}（两者借用本调用的模型与截止期）；
  \item \textbf{降级}：\texttt{PortfolioMode::Throughput} 或线程预算为 1 时
        不竞速，直接跑 IPM；线程创建抛 \texttt{std::system\_error} 时置取消、
        join 已建线程、回退直连 IPM；截止期耗尽按 \texttt{Time limit} 短路。
\end{itemize}
对应源码为 \srcpath{src/engine/solver/native/native\_lp\_selector.cpp:NativeAutoLPAdapter::solve\_with\_context}、
\srcpath{src/engine/solver/native/native\_lp\_selector.cpp:publish}（匿名命名空间）、
\srcpath{src/engine/solver/native/native\_lp\_selector.cpp:NativeDualSimplexLPAdapter::solve\_lp}。

\begin{boundarynote}
\texttt{NativeDualSimplexLPAdapter::solve\_lp} 在成功时把公开目标改写为用户
方向的 $c^{\top}x$（内部统计里的 \texttt{objective} 是最小化方向
$\mathrm{obj}_{\text{const}} - \max$）。竞速两分支的 \texttt{time\_limit}
都是同一截止期的剩余量；环境变量 \texttt{MIPSOLVERS\_LP\_SELECTOR\_DEBUG}
只打印胜者名，不影响选择。
\end{boundarynote}

\paragraph{独立的逐 pivot 发布证据.}
\texttt{MIPSOLVERS\_LP\_PIVOT\_TRACE=1} 在已提交的 dual/primal
transaction 后输出 phase、row、entering/leaving 和 pivot/step 的
binary64 位串；primal bound flip 用 row=-1 区分，未定义的代数 pivot
与 dual step 字段为零。关闭时不格式化和写入，记录不回馈定价或比值检验。
对应源码为
\srcpath{src/engine/kernel/lp\_kernel/native\_dual/pivot\_trace.hpp:emit\_pivot\_trace}。
相同探针分别加入基准和候选二进制，逐字节比较固定语料的实际输出，
拒绝空 trace。该检查只证明所记录轨迹相同，不证明所有 LP 或内部全部
浮点状态相同；源码范围和迭代数不能替代这一证据。合同见测试手册 R3。


% source_equivalent_ipm_kkt.tex —— engine 模块实现转写章：IPM 主循环、KKT 装配与稀疏后端
% 本文件为实现转写章（非理论章），遵循 docs/modules/README.md 数学模型规范第 1/2/4 条。

\section{内点法、KKT 系统与稀疏线性代数：源码等价转写}
\label{sec:engine-source-ipm-kkt}

\subsection{转写范围与口径}
\label{subsec:engine-ipm-scope}

本章转写以下工作树文件在 2026-09-13 版本实际执行的计算：

\begin{itemize}
  \item NLP 内点主循环（filter 全局化）：\srcpath{src/engine/kernel/ipm/ipm\_solver.cpp}；
  \item LP 内点（Mehrotra/Gondzio）：\srcpath{src/engine/kernel/ipm/ipm\_lp\_solver.cpp}、
        \srcpath{src/engine/kernel/ipm/ipm\_lp\_solver\_cached.cpp}、
        \srcpath{src/engine/kernel/ipm/ipm\_lp\_solver\_internal.hpp}；
  \item 缩放、filter 数据结构、恢复阶段：\srcpath{src/engine/kernel/ipm/ipm\_scaling.cpp}、
        \srcpath{src/engine/kernel/ipm/ipm\_filter.cpp}、
        \srcpath{src/engine/kernel/ipm/ipm\_restoration.cpp}；
  \item KKT 装配与求解：\srcpath{src/engine/kernel/kkt/kkt\_system.cpp}；
  \item 稀疏后端：\srcpath{src/engine/kernel/linear\_algebra/linear\_solver.cpp}、
        \srcpath{src/engine/kernel/linear\_algebra/sparse\_lu\_factor.cpp}、
        \srcpath{src/engine/kernel/linear\_algebra/cholmod\_ldlt.cpp}、
        \srcpath{src/engine/kernel/linear\_algebra/hfactor\_backend.cpp}
        （vendored HFactor 位于
        \srcpath{src/engine/kernel/linear\_algebra/highs\_factor/}）。
\end{itemize}

素材文档 \texttt{docs/manual/06-numerical-methods.md} 的符号级描述已逐条回源码核验；
与代码不一致处在 \ref{subsec:engine-ipm-mismatch} 节集中声明。锚点规则同上一章。
编译期宏 \texttt{HACDCPF\_HAVE\_MUMPS} / \texttt{HACDCPF\_HAVE\_MKL\_PARDISO} /
\texttt{MIPSOLVERS\_HAVE\_CHOLMOD} / \texttt{MIPSOLVERS\_USE\_ACCELERATE} 控制的路径
按源码条件分支转写，并注明其编译开关。

\begin{boundarynote}
锥 IPM（\texttt{conic\_ipm\_solver.cpp}、\texttt{cones.cpp}）、
LCQP（\texttt{lcqp\_solver.cpp}）与 chordal 分解不在本章范围。
\texttt{MIPSOLVERS\_IPM\_VERBOSE} 等日志通道只读不写，不改变数值；
\texttt{lp\_timing} 计时插桩同理不转写。
\end{boundarynote}

% ─────────────────────────────────────────────────────────────────────────────
\subsection{NLP 内点主循环（filter 驱动）}
\label{subsec:engine-ipm-nlp}

\paragraph{问题与残差.}
\compmeta{NLPState}{src/engine/kernel/ipm/ipm\_solver.cpp}
求解器消费带回调的 \texttt{NLPModel}：$\min f(x)$ s.t. $g(x)=0$、$h(x)\le 0$、
变量界。不等式含非线性行与由变量界生成的行（\texttt{lb\_cols}/\texttt{ub\_cols}），
松弛 $s>0$ 满足 $h(x)+s=0$，乘子 $\mu>0$。KKT 残差的代码定义：
\begin{equation}
  r_d = \nabla f + J_g^{\top}\lambda + J_h^{\top}\mu,\qquad
  r_{eq} = g,\qquad
  r_{ineq} = h + s,\qquad
  r_{comp} = s \circ \mu - \bar\mu\,\mathbf{1},
  \label{eq:ipm:nlp-residuals}
\end{equation}
汇总为 $\theta_p = \max(\|r_{eq}\|_{\infty}, \|r_{ineq}\|_{\infty})$、
$\theta_d = \|r_d\|_{\infty}$、$\theta_c = \|s \circ \mu\|_{\infty}$；全局化用
缩放 merit $\max\{\theta_p/(1+\|x,s\|_{\infty}),\ \theta_d/(1+\|\lambda,\mu\|_{\infty},\
\theta_c/(1+\|\mu\|_{\infty})\}$，而终止判定用绝对量。
对应源码为 \srcpath{src/engine/kernel/ipm/ipm\_solver.cpp:summarize\_residuals}、
\srcpath{src/engine/kernel/ipm/ipm\_solver.cpp:evaluate\_nlp\_state}。

\paragraph{终止判据.}
\texttt{evaluate\_ipm\_termination} 计算乘子尺度
$s_d = \max(T, \bar m)/T$（$T$ = \texttt{multiplier\_scale\_threshold}，缺省
$\max(1, \|\nabla f\|_{\infty})$，$\bar m$ 为全部乘子绝对值均值）与
$s_c = \max(T, \bar\mu_{\mathrm{abs}})/T$；总体误差
$E = \max\{\theta_p,\ \theta_d / s_d,\ \theta_c / s_c\}$。
严格收敛（\texttt{strict}）要求
\begin{equation}
  E \le \texttt{tol\_overall}\ (\text{若启用}),\quad
  \theta_p \le \texttt{tol\_primal},\quad
  \theta_d^{*} \le \texttt{tol\_dual},\quad
  \theta_d \le \texttt{raw\_dual\_guard}\ (\text{若}>0),\quad
  \theta_c \le \texttt{tol\_complementarity},
  \label{eq:ipm:strict-term}
\end{equation}
其中 $\theta_d^{*} = \theta_d / s_d$ 当 \texttt{multiplier\_relative\_stationarity}
为真，否则 $\theta_d$。可接受收敛另用 \texttt{tol\_accept*} 门槛并要求
连续 \texttt{acceptable\_iter} 次满足。对应源码为
\srcpath{src/engine/kernel/ipm/ipm\_solver.cpp:evaluate\_ipm\_termination}、
\srcpath{src/engine/kernel/ipm/ipm\_solver.cpp:strict\_termination}、
\srcpath{src/engine/kernel/ipm/ipm\_solver.cpp:acceptable\_termination}。

\paragraph{双层结构.}
外层按障碍参数推进，内层在当前 $\bar\mu$ 上做 Newton/filter 迭代。
内层容差 $\tau_{\text{inner}} = \kappa_\varepsilon \bar\mu$
（\texttt{kappa\_epsilon} 缺省按 $1.0$ 处理）；处于最小障碍时再把
$\tau_{\text{inner}}$ 收紧到 $\min(\texttt{tol\_primal}, \texttt{tol\_dual},
\texttt{tol\_complementarity})$ 以留出最后一次 Newton 精化空间。内层收敛判据
\begin{equation}
  E_{\mu} = \max\bigl\{\|r_d\|_{\infty}, \|r_{eq}\|_{\infty},
  \|r_{ineq}\|_{\infty}, \|s \circ \mu - \bar\mu\|_{\infty}\bigr\}
  \le \tau_{\text{inner}},
  \label{eq:ipm:inner-conv}
\end{equation}
另有一条调度捷径：调用者坐标原始违反 $\le$ 调用者容差、$\theta_d^{*} \le$
\texttt{tol\_dual}、但 $\theta_c >$ \texttt{tol\_complementarity} 且
$\bar\mu >$ \texttt{tol\_complementarity} 时直接视为内层收敛、转入 $\mu$ 更新。
对应源码为 \srcpath{src/engine/kernel/ipm/ipm\_solver.cpp:solve\_nlp\_filter\_impl}。

\paragraph{$\mu$ 更新.}
显式策略（\texttt{kappa\_mu}<1 且 \texttt{theta\_mu}>1 均被设置时）：
\begin{equation}
  \bar\mu^{+} = \max\bigl(\mu_{\min},\;
  \min(\kappa_\mu \bar\mu,\; \bar\mu^{\theta_\mu})\bigr);
  \label{eq:ipm:mu-explicit}
\end{equation}
缺省策略取 $\max\bigl(\mu_{\min}, \min(\texttt{tol\_complementarity},
\min_i s_i \mu_i)\bigr)$；若该值不小于当前 $\bar\mu$（目标已到调用者门槛），
改为从当前中心残差推导最小可用降幅：
$\bar\mu^{+} = \max(\mu_{\min}, \bar\mu -
(\|s \circ \mu - \bar\mu\|_{\infty} + \sqrt{\varepsilon}\bar\mu))$；
无法再降时以``最小障碍未达 KKT 收敛''终止。$\mu$ 变化后 filter 清空
（旧 $(\theta, \varphi)$ 对属于另一障碍目标）。$\mu_{\min}$ 缺省
$\sqrt{\text{最小正规格化双精度}}$（\texttt{kMinPositive}）。
对应源码为 \srcpath{src/engine/kernel/ipm/ipm\_solver.cpp:solve\_nlp\_filter\_impl}、
\srcpath{src/engine/kernel/ipm/ipm\_solver.cpp:minimum\_safe\_positive}（匿名命名空间）。

\paragraph{初始点.}
\texttt{interiorize\_initial\_point} 把每个有界变量推入内部：
双边界取
$\ell' = \max(\mathrm{nextafter}(\ell,u), \ell + \min(0.01\max(1,|\ell|),
0.01(u-\ell)))$ 与对称的 $u'$（常量 \texttt{bound\_push=bound\_fraction=0.01}），
退化时取中点；单边界按 $0.01\max(1,|\cdot|)$ 推。
\texttt{primal\_feasible\_start}/\texttt{preserve\_initial\_point} 经
\texttt{solve\_nlp\_detail} 的原始违反审计（$\le$ \texttt{tol\_primal}）采纳后
跳过该内点化。松弛初值取线性化分辨率下界
（\texttt{initialize\_slacks\_from\_linearization\_resolution}），冷启动且行违反
超过调用者容差时把松弛抬到 $\max(s_i, \max(h_i, 0))$。
对应源码为 \srcpath{src/engine/kernel/ipm/ipm\_solver.cpp:interiorize\_initial\_point}、
\srcpath{src/engine/kernel/ipm/ipm\_solver.cpp:solve\_nlp\_filter\_impl}、
\srcpath{src/engine/kernel/ipm/ipm\_solver.cpp:NativeIPMAdapter::solve\_nlp\_detail}。

\paragraph{Newton 系统：凝聚形.}
凝聚路径组装 $W = H + J_h^{\top} \operatorname{diag}(\mu/s) J_h$
（$s$ 逐分量下截于 \texttt{kMinPositive}），右端与回代：
\begin{align}
  \mathrm{rhs}_x &= -r_d - J_h^{\top} v,\quad
  v_i = \frac{\mu_i r_{ineq,i} + \bar\mu - s_i \mu_i}{s_i},\qquad
  \mathrm{rhs}_{eq} = -r_{eq},
  \label{eq:ipm:condensed-rhs}\\
  d_s &= -r_{ineq} - J_h d_x,\qquad
  d\mu_i = \frac{\bar\mu - s_i \mu_i - \mu_i\, d s_i}{s_i}.
  \label{eq:ipm:condensed-back}
\end{align}
对应源码为 \srcpath{src/engine/kernel/ipm/ipm\_solver.cpp:solve\_nlp\_filter\_impl}
（lambda \texttt{build\_condensed\_w} 与同函数内联回代）。

\paragraph{Newton 系统：增广形.}
增广路径组装对称不定系统
\begin{equation}
  \begin{bmatrix}
    H + \delta_W I & J_g^{\top} & \tilde J_h^{\top}\\
    J_g & -\delta_C I & 0\\
    \tilde J_h & 0 & -I
  \end{bmatrix}
  \begin{bmatrix} d_x \\ d_\lambda \\ \tilde d_\mu \end{bmatrix}
  =
  \begin{bmatrix} -r_d \\ -r_{eq} \\
    t \circ \bigl(-r_{ineq} - (\bar\mu\mathbf{1} - s\circ\mu)/\mu\bigr)
  \end{bmatrix},
  \qquad t_i = \sqrt{\mu_i / s_i},
  \label{eq:ipm:augmented}
\end{equation}
其中 $\tilde J_h = \operatorname{diag}(t) J_h$：这是对未缩放增广形
$[H+\delta_W I, J_g^{\top}, J_h^{\top}; J_g, -\delta_C I, 0;
J_h, 0, -S M^{-1}]$ 施加合同变换 $T = \operatorname{diag}(I, I, t)$ 的结果；
解出后 $d_\mu = t \circ \tilde d_\mu$、$d_s = -r_{ineq} - J_h d_x$。
组装经 scatter map 直写缓存 CSC 值数组（结构逐问题不变），
$s_i, \mu_i$ 截断于 \texttt{kMinPositive}。
对应源码为 \srcpath{src/engine/kernel/ipm/ipm\_solver.cpp:assemble\_augmented\_newton}、
\srcpath{src/engine/kernel/ipm/ipm\_solver.cpp:solve\_augmented\_newton}、
\srcpath{src/engine/kernel/ipm/ipm\_solver.cpp:build\_augmented\_newton\_scatter}。

\paragraph{形态选择.}
首次组装时若 \texttt{use\_augmented\_newton} 或
\texttt{NewtonFormulation::Augmented} 强制则走增广；\texttt{Auto} 且有惯性后端
（MUMPS/PARDISO 宏）时以符号代价裁决：对两种形态各做一次 CHOLMOD 符号分析，
增广形当选当且仅当其符号 flops 与因子非零数\textbf{都}严格更小。
运行期回退：凝聚形证书失败（或凝聚全残差超过其块方程舍入限，见下）时
切到增广形重来；增广形失败且 \texttt{Auto} 时也允许切回凝聚。
对应源码为 \srcpath{src/engine/kernel/ipm/ipm\_solver.cpp:select\_augmented\_by\_symbolic\_cost}、
\srcpath{src/engine/kernel/ipm/ipm\_solver.cpp:solve\_nlp\_filter\_impl}
（lambda \texttt{solve\_augmented\_candidate}）。

\paragraph{凝聚残差审计.}
\texttt{Auto} 下凝聚解回收 $d_s, d\mu$ 后，重组四个块方程残差
（含 $\delta_W d_x$ 与 $-\delta_C d_\lambda$ 正则项）并取
$r_{\text{full}} = \max\{\|r_d + H d_x + J_g^{\top} d\lambda +
J_h^{\top} d\mu + \delta_W d_x\|_{\infty}, \ldots\}$；
限值为 $\sqrt{\varepsilon}\max(1, S)$，$S$ 为全部参与项的无穷范数最大值
（\texttt{comparison\_roundoff}）。超限即``凝聚精度耗尽''，转入增广形。
对应源码为 \srcpath{src/engine/kernel/ipm/ipm\_solver.cpp:solve\_nlp\_filter\_impl}、
\srcpath{src/engine/kernel/ipm/ipm\_solver.cpp:comparison\_roundoff}（匿名命名空间）。

% ─────────────────────────────────────────────────────────────────────────────
\subsection{正则化与惯性校正}
\label{subsec:engine-ipm-inertia}

\paragraph{盲正则化阶梯.}
\texttt{factor\_kkt\_with\_regularization}（无惯性后端时的默认路径）：
尺度 $S = \max(1, \|W\|_{\max}, \|J_g\|_{\max})$；先试 $\delta = 0$，然后
$\delta$ 从 $\sqrt{\varepsilon}\, S$ 起每次乘 $2$，上限 $S/\sqrt{\varepsilon}$；
全部失败则置 \texttt{factored=false} 并失败。
对应源码为 \srcpath{src/engine/kernel/ipm/ipm\_solver.cpp:factor\_kkt\_with\_regularization}。

\paragraph{惯性设置解析.}
\texttt{resolve\_inertia\_settings} 填充缺省值
\begin{equation}
  \delta_W^{0} = \sqrt{\varepsilon}\, S_W,\qquad
  \delta_W^{\max} = S_W / \sqrt{\varepsilon},\qquad
  \delta_C = \sqrt{\varepsilon}\, S_K,
  \label{eq:ipm:inertia-settings}
\end{equation}
$S_W = \max(1, \|W\|_{\max})$、$S_K = \max(S_W, \|J_g\|_{\max})$；
升级步进精确取 2 的幂：首次从 $0$ 踢到 $\delta_W^{0}$，其后
$\delta \leftarrow \max(2\delta, \delta_W^{0})$。
对应源码为 \srcpath{src/engine/kernel/kkt/kkt\_system.cpp:resolve\_inertia\_settings}、
\srcpath{src/engine/kernel/kkt/kkt\_system.cpp:increased\_primal\_regularization}（匿名命名空间）。

\paragraph{惯性校正分解.}
\texttt{factor\_kkt\_inertia\_corrected\_sparse} 的判定序列：
\begin{enumerate}
  \item 有惯性后端（MUMPS/PARDISO 宏）且未强制最小 $\delta_W$、未设
        \texttt{max\_tangent\_dimension} 时先试精确未正则化分解；后端报告
        负特征值数 $= m_{eq}$ 且亏秩 $=0$ 则直接接受
        （\texttt{factor\_unregularized\_with\_direct\_inertia}）；
  \item 否则构造简约空间证书：选切空间划分（自然/偏好/QR/结构四种策略），
        以切基求解构造 $Z$（$J_g Z = 0$ 按
        $\sqrt{\varepsilon}\max(1, \|J_g\|_{\infty}\|Z\|_{\infty})$ 审计），
        对 $(Z^{\top} W Z, Z^{\top} Z)$ 做对称化广义特征分解，得最小曲率
        $\lambda_{\min}$ 与裕量
        $M = \max(\texttt{reduced\_curvature\_tolerance},\;
        \gamma_{d}\, \lambda_{\max}^{|\cdot|})$，其中
        $\gamma_d = d\varepsilon/(1 - d\varepsilon)$（$d$ 为切维数）；
        所需平移 $\delta_W = \max(0, M - \lambda_{\min})$，需要平移时再加
        因子侧裕量 $\sqrt{\varepsilon}\max(1, |\lambda_{\min}|)$；
        然后以 $\delta_W$ 起步、2 幂升级到 $\delta_W^{\max}$，
        接受条件为``分解成功且简约惯性被认证''
        （$\lambda_{\min} + \delta_W \ge M$）；直接后端报告的枢轴计数此时
        只是诊断，不作否决；
  \item 证书构造失败（含超 \texttt{max\_tangent\_dimension} 上限）时退回
        直接惯性环：每个 $\delta_W$ 先试 $\delta_C = 0$ 再试
        $\delta_C = \delta_C^{\text{stripe}}$，接受条件为后端报告
        负特征值数 $= m_{eq}$ 且亏秩 $=0$；
  \item 无惯性后端时退回 \texttt{certify\_primal\_positive\_definite}
        （对 $W + \delta I$ 做 Eigen SimplicialLDLᵀ 并检查 \texttt{vectorD}
        全正），其后按块合同惯性 $(n, m_{eq}, 0)$ 分解，必要时加
        $\delta_C$ 条纹修复秩亏。
\end{enumerate}
对应源码为 \srcpath{src/engine/kernel/kkt/kkt\_system.cpp:factor\_kkt\_inertia\_corrected\_sparse}、
\srcpath{src/engine/kernel/kkt/kkt\_system.cpp:build\_reduced\_space\_certificate}（匿名命名空间）、
\srcpath{src/engine/kernel/kkt/kkt\_system.cpp:rounding\_gamma}（匿名命名空间）、
\srcpath{src/engine/kernel/kkt/kkt\_system.cpp:factor\_with\_direct\_inertia}（匿名命名空间）、
\srcpath{src/engine/kernel/kkt/kkt\_system.cpp:factor\_unregularized\_with\_direct\_inertia}（匿名命名空间）、
\srcpath{src/engine/kernel/kkt/kkt\_system.cpp:certify\_primal\_positive\_definite}（匿名命名空间）。

\paragraph{增广 Newton 的 $\delta_W$ 网格.}
\texttt{factor\_solve\_augmented\_newton} 对 \eqref{eq:ipm:augmented} 用
二分网格搜索最小可用指数：候选 $\delta_W = \mathrm{ldexp}(\delta_W^{0}, e)$，
成功指数与失败指数相邻时收敛；每个 $\delta_W$ 先试 $\delta_C=0$ 再试
$\delta_C$ 条纹（保持零对偶块非奇异时的精确等式 Newton 方程）。
接受要求分解成功、惯性 $(n, m_{eq} + m_{ineq}, 0)$（由
\texttt{augmented\_factor\_has\_correct\_inertia} 经后端
\texttt{negative\_eigenvalues}/\texttt{estimated\_deficiency} 读取）、
且求解通过共同的后向误差门。\texttt{primal\_feasible\_start}
（\texttt{require\_descent\_inertia=false}）时跳过惯性要求、先试完全
未正则化的系统。对应源码为
\srcpath{src/engine/kernel/ipm/ipm\_solver.cpp:factor\_solve\_augmented\_newton}、
\srcpath{src/engine/kernel/ipm/ipm\_solver.cpp:augmented\_factor\_has\_correct\_inertia}、
\srcpath{src/engine/kernel/ipm/ipm\_solver.cpp:exact\_condensed\_factor\_inertia\_ok}。

\paragraph{方向局部性强制.}
凝聚路径在原始残差已达容差（$\theta_p \le \texttt{tol\_primal} +
\sqrt{\varepsilon}\max(1,\texttt{tol\_primal})$）时，若归一化方向
$\max_j |d_{x,j}| / \max(1, |x_j|) > 1$，则以
$\delta_W \leftarrow \max(\delta_W^{\text{used}}, \sqrt{\varepsilon}\,
\|W\|_{\max})$ 起步、每次乘 $2$ 强制重分解，直到方向归一化无穷范数
$\le 1$。精确凝聚快速通道（\texttt{primal\_feasible\_start} 或原始残差
仍超差时先试的未正则化分解 + 惯性证书）同样受此门控，越线即弃用。
对应源码为 \srcpath{src/engine/kernel/ipm/ipm\_solver.cpp:solve\_nlp\_filter\_impl}。

% ─────────────────────────────────────────────────────────────────────────────
\subsection{filter 线搜索与步长规则}
\label{subsec:engine-ipm-filter}

\paragraph{量与斜率.}
\begin{equation}
  \theta(x, s) = \|g(x)\|_1 + \|h(x) + s\|_1,\qquad
  \varphi_{\bar\mu}(x, s) = f(x) - \bar\mu \sum_i \log s_i,
  \qquad \nabla\varphi \cdot d = \nabla f^{\top} d_x -
  \bar\mu \sum_i \frac{d s_i}{s_i},
  \label{eq:ipm:theta-phi}
\end{equation}
$\log$ 与除法中的 $s_i$ 下截于 \texttt{kMinPositive}。对应源码为
\srcpath{src/engine/kernel/ipm/ipm\_solver.cpp:compute\_theta}、
\srcpath{src/engine/kernel/ipm/ipm\_solver.cpp:compute\_barrier\_phi}、
\srcpath{src/engine/kernel/ipm/ipm\_solver.cpp:barrier\_descent\_slope}。

\paragraph{边界步长.}
分数边界规则 $\tau = \texttt{alpha\_max}$（合法时）否则
$1 - \sqrt{\varepsilon}$；最大正向步
\begin{equation}
  \alpha_{\max}^{p} = \min\Bigl(1, \min_{i:\, d s_i < 0}
  \frac{-\tau s_i}{d s_i}\Bigr),\qquad
  \alpha_{\max}^{d} = \min\Bigl(1, \min_{i:\, d\mu_i < 0}
  \frac{-\tau \mu_i}{d\mu_i}\Bigr),
  \label{eq:ipm:frac-boundary}
\end{equation}
$\alpha \leftarrow \min(1, \alpha_{\max}^{p})$。
\texttt{primal\_feasible\_start} 下原始步恒取 $1$ 且对偶步独立于 $\alpha$
（\texttt{trial\_alpha\_dual} $= \alpha_{\max}^{d}$ 而非
$\min(\alpha_{\max}^{d}, \alpha)$）。回溯每次乘 $0.5$，下限
\texttt{filter\_alpha\_min}，缺省为可表示性极限
$\varepsilon \cdot \mathrm{iterate\_scale} / \max(\mathrm{direction\_scale},
\texttt{kMinPositive})$（方向/迭代尺度取四组向量无穷范数的最大）。
对应源码为 \srcpath{src/engine/kernel/ipm/ipm\_solver.cpp:max\_positive\_step}、
\srcpath{src/engine/kernel/ipm/ipm\_solver.cpp:effective\_fraction\_to\_boundary}（匿名命名空间）、
\srcpath{src/engine/kernel/ipm/ipm\_solver.cpp:solve\_nlp\_filter\_impl}。

\paragraph{filter 数据结构.}
\texttt{Filter::is\_acceptable} 拒绝当（i）$\theta_t \ge$
\texttt{theta\_upper\_bound}，或（ii）存在已有条目 $(\theta_k, \varphi_k)$
使 $\theta_t \ge (1-\gamma_\theta)\theta_k$ 且
$\varphi_t \ge \varphi_k - \gamma_\varphi \theta_k$；
\texttt{add\_entry} 插入时剪去被新条目按同一裕量支配的旧条目；
\texttt{reset\_with\_theta\_upper\_bound} 设定 $\theta$ 上界。
$\gamma_\theta, \gamma_\varphi$ 缺省均取 $\varepsilon$（机器精度）。
主循环把上界置为
$\mathrm{nextafter}(\max(\theta_{\text{init}}, N_{\text{rows}} \cdot
\max(\|g\|_{\infty}, \|h+s\|_{\infty}, \texttt{tol\_primal})), +\infty)$。
对应源码为 \srcpath{src/engine/kernel/ipm/ipm\_filter.cpp:Filter::is\_acceptable}、
\srcpath{src/engine/kernel/ipm/ipm\_filter.cpp:Filter::add\_entry}、
\srcpath{src/engine/kernel/ipm/ipm\_filter.cpp:Filter::reset\_with\_theta\_upper\_bound}、
\srcpath{src/engine/kernel/ipm/ipm\_solver.cpp:solve\_nlp\_filter\_impl}。

\paragraph{切换条件与接受测试.}
切换条件（仅 \texttt{filter\_s\_theta}/\texttt{filter\_s\_phi}/
\texttt{filter\_delta}/\texttt{filter\_eta\_phi}/
\texttt{filter\_theta\_min\_scale} 全部为正时启用）：
\begin{equation}
  \mathrm{slope} < 0,\quad
  \alpha (-\mathrm{slope})^{s_\varphi} >
  \delta\, \theta_k^{s_\theta},\quad
  \theta_k \le \theta_{\min};
  \label{eq:ipm:switching}
\end{equation}
成立时按 Armijo 接受 $\varphi_t \le \varphi_k +
\eta_\varphi \alpha\, \mathrm{slope} + \sqrt{\varepsilon}\max(1,|\varphi_k|)$
（f 型步，不向 filter 加条目）；否则 h 型接受为
``$\theta_t \le (1-\gamma_\theta)\theta_k$ 或
$\varphi_t \le \varphi_k - \gamma_\varphi \theta_k$''（参数化形式），
缺省形式为``$\theta_t$ 或 $\varphi_t$ 以 $\varepsilon\max(1,|\cdot|)$ 裕量
严格下降''（\texttt{arithmetic\_roundoff}）。h 型接受后把
$(\theta_k, \varphi_k)$ 加入 filter。对应源码为
\srcpath{src/engine/kernel/ipm/ipm\_solver.cpp:solve\_nlp\_filter\_impl}。

\paragraph{中心性特例与 Phase-I 门.}
原始位移在 $\sqrt{\varepsilon}$ 相对尺度内时，filter 坐标可能不变；
此时仅当分量不恶化（原始/对偶残差各以 $\varepsilon\max(1,|\cdot|)$ 裕量
不增）、满足 $\theta$ 上界、merit 严格下降，且互补或对偶残差严格下降，
才接受该切向/中心性步。\texttt{primal\_feasible\_start} 的 Phase-I 门
另行要求调用者坐标原始违反严格小于 \texttt{phase1\_primal\_upper\_bound}、
对偶与互补不增（\texttt{comparison\_roundoff} 裕量）且二者之一严格下降
（\texttt{relative\_roundoff} 裕量）。对应源码为
\srcpath{src/engine/kernel/ipm/ipm\_solver.cpp:solve\_nlp\_filter\_impl}、
\srcpath{src/engine/kernel/ipm/ipm\_solver.cpp:arithmetic\_roundoff}（匿名命名空间）、
\srcpath{src/engine/kernel/ipm/ipm\_solver.cpp:relative\_roundoff}（匿名命名空间）。

\paragraph{二阶校正（SOC）.}
\texttt{use\_second\_order\_correction} 且（Phase-I 邻域被违反，或冷启动
全步使 $\theta$ 上升）时，以修正右端
\begin{equation}
  c_{\text{soc}} = g(x + \alpha d_x) - (1-\alpha)\, g(x)
  \label{eq:ipm:soc}
\end{equation}
替换等式块（增广布局取中段，凝聚布局取尾部；不再重复原始 Newton 强迫项），
复用当前分解求解；复合试验点 $x + \alpha d_x + d_x^{\text{soc}}$、
$s + \alpha d_s - J_h d_x^{\text{soc}}$，乘子叠加
$d\lambda^{\text{soc}}$、$d\mu^{\text{soc}}$（凝聚路径上
$d\mu^{\text{soc}} = -\mu \circ d s^{\text{soc}} / s$）。然后按同一套
filter/切换/Phase-I 门判定。对应源码为
\srcpath{src/engine/kernel/ipm/ipm\_solver.cpp:solve\_nlp\_filter\_impl}。

\paragraph{停滞与最优迭代点.}
接受的步若原始相对位移与对偶相对位移都不超过 $\sqrt{\varepsilon}$，
按停滞处理：原始停滞且原始与互补已达标的端点允许一次终端对偶投影
（\texttt{initialize\_primal\_feasible\_floor\_projection} +
\texttt{fit\_primal\_feasible\_duals\_active\_projection}），投影结果仍需通过
同一 \texttt{strict\_termination}；否则以停滞状态退出。线搜索崩溃
（$\alpha$ 跌破下限）或预算耗尽时回退到 merit 最优的历史迭代点。
严格收敛后可选 \texttt{try\_active\_kkt\_polish}：只有 polish 结果也满足
严格终止才替换当前解。对应源码为
\srcpath{src/engine/kernel/ipm/ipm\_solver.cpp:solve\_nlp\_filter\_impl}、
\srcpath{src/engine/kernel/ipm/ipm\_solver.cpp:try\_active\_set\_kkt\_polish}。

% ─────────────────────────────────────────────────────────────────────────────
\subsection{缩放与恢复阶段}
\label{subsec:engine-ipm-scaling-restoration}

\paragraph{目标/约束缩放.}
\texttt{compute\_scaling\_factors} 在内点化的 $x_0$ 处采样：
$s_f = g_{\max}/\|\nabla f\|_{\infty}$（范数非有限或 $\le g_{\max}$ 时取 $1$），
$s_g, s_h$ 逐行同理（$g_{\max}$ 缺省 $1.0$，来自
\texttt{scaling\_g\_max}）。\texttt{build\_scaled\_nlp\_model} 以 lambda 包装
全部回调：目标/梯度/Hessian 乘 $s_f$，约束值与 Jacobian 行乘 $s_{g,i}$/$s_{h,i}$；
Lagrangian Hessian 回调内先把乘子还原到原坐标
$\lambda^{\text{orig}}_i = (s_{g,i}/s_f)\lambda_i$、
$\mu^{\text{orig}}_i = (s_{h,i}/s_f)\mu_i$，结果再乘 $s_f$。
\texttt{solve\_nlp\_detail} 在 filter 求解前后分别用
\texttt{transform\_filter\_options\_to\_scaled\_coordinates} 与
\texttt{audit\_filter\_outcome\_in\_original\_coordinates} 做选项换算与
原坐标审计。对应源码为
\srcpath{src/engine/kernel/ipm/ipm\_scaling.cpp:compute\_scaling\_factors}、
\srcpath{src/engine/kernel/ipm/ipm\_scaling.cpp:build\_scaled\_nlp\_model}、
\srcpath{src/engine/kernel/ipm/ipm\_scaling.cpp:scale\_from\_norm}（匿名命名空间）、
\srcpath{src/engine/kernel/ipm/ipm\_solver.cpp:NativeIPMAdapter::solve\_nlp\_detail}。

\paragraph{恢复 NLP.}
filter 失败且状态属于可触发集合（线搜索步过小 / 接受步崩溃 / 迭代停滞 /
惯性校正超限 / 达到最大迭代）且模型有等式时，以当前（或输入）点 $x_R$
构造恢复问题
\begin{equation}
  \min_{x, p, n}\ \sum_i (p_i + n_i) + \frac{\zeta}{2}
  \|D_R (x - x_R)\|^{2}
  \quad \text{s.t.}\quad g(x) + p - n = 0,\ h(x) \le 0,\ p, n \ge 0,
  \label{eq:ipm:restoration}
\end{equation}
其中 $D_R = \operatorname{diag}(d_j)$、$d_j = \min(1, 1/\max(1, |x_{R,j}|))$；
$\zeta$ 缺省 $\sqrt{\varepsilon}$（\texttt{restoration\_zeta}）。
初始点 $p_i = \rho_i + \max(-g_i, 0)$、$n_i = \rho_i + \max(g_i, 0)$，
分辨率
\begin{equation}
  \rho_i = \max\Bigl(\sqrt{\nu_{\min}},\; \sqrt{\varepsilon}
  \max\bigl(|g_i(x_R)|,\; r_i\bigr)\Bigr),\qquad
  r_i = \Bigl(\sum_j (J_{g,ij} s^{\text{var}}_j)^2\Bigr)^{1/2},
  \label{eq:ipm:restoration-resolution}
\end{equation}
$\nu_{\min}$ 为最小正双精度，$s^{\text{var}}_j$ 取
$\max(1, |x_{R,j}|, \text{界宽}, |\ell_j|, |u_j|)$。
对应源码为 \srcpath{src/engine/kernel/ipm/ipm\_restoration.cpp:build\_restoration\_nlp}、
\srcpath{src/engine/kernel/ipm/ipm\_restoration.cpp:make\_dr}（匿名命名空间）、
\srcpath{src/engine/kernel/ipm/ipm\_restoration.cpp:restoration\_variable\_scale}（匿名命名空间）。

\paragraph{恢复接受准则.}
恢复子问题以求解同一 filter 驱动（禁缩放、禁 SOC、禁嵌套恢复，
预算 $\min(\texttt{max\_iter}, \max(1, \texttt{restoration\_max\_iter}))$）。
只有原模型法向残差（\texttt{raw\_primal\_violation}）严格下降——
$\mathrm{after} + \sqrt{\varepsilon}\max(1, \mathrm{before}) <
\mathrm{before}$——才接受恢复点并以其为 $x_0$ 重试原问题；
\texttt{recover\_restoration\_warm\_start} 在维度/正性检查齐全时把恢复
乘子与松弛映射回原模型作为暖启动（不等式对偶与松弛必须逐分量为正，
否则整套放弃）。重试仅一次（\texttt{use\_restoration\_phase=false}）。
对应源码为 \srcpath{src/engine/kernel/ipm/ipm\_solver.cpp:NativeIPMAdapter::solve\_nlp\_detail}、
\srcpath{src/engine/kernel/ipm/ipm\_restoration.cpp:recover\_restoration\_warm\_start}、
\srcpath{src/engine/kernel/ipm/ipm\_restoration.cpp:extract\_x\_from\_restoration}。

\paragraph{固定变量消元与暖启动审计.}
\texttt{solve\_nlp\_detail} 在进入 filter 前：
\texttt{build\_fixed\_nlp\_reduction} 把 $\ell_j = u_j$ 的变量消元，
求解后在 \texttt{restore\_fixed\_nlp\_certificate} 中映射回并做原模型证书；
\texttt{central\_warm\_start} 经 \texttt{audit\_central\_warm\_start} 全有或
全无审计（拒绝时把全部暖启动向量清空，观测上等同冷启动）；
\texttt{primal\_feasible\_start}/\texttt{preserve\_initial\_point} 仅当输入点的
原始违反 $\le$ \texttt{tol\_primal} 时采纳。对应源码为
\srcpath{src/engine/kernel/ipm/ipm\_solver.cpp:NativeIPMAdapter::solve\_nlp\_detail}、
\srcpath{src/engine/kernel/ipm/ipm\_solver.cpp:build\_fixed\_nlp\_reduction}、
\srcpath{src/engine/kernel/ipm/ipm\_solver.cpp:restore\_fixed\_nlp\_certificate}、
\srcpath{src/engine/kernel/ipm/ipm\_solver.cpp:audit\_central\_warm\_start}。

% ─────────────────────────────────────────────────────────────────────────────
\subsection{LP 内点（Mehrotra + Gondzio）}
\label{subsec:engine-ipm-lp}

\paragraph{有界变量标准形.}
\texttt{NativeIPMLPAdapter::solve\_lp\_impl} 求解
$\min \hat c^{\top} \hat x$ s.t. $\hat A \hat x = \hat b$、
$\ell \le \hat x \le u$ 的原始-对偶内点形式，其中 $\hat x = (x; s)$ 含
不等式松弛（$nn = n + m_i$）。G 行（$b = +\infty$ 而 lhs 有限）在输入边
翻号为 L 行；界哨兵 $\pm 10^{20}$（\texttt{kBig}）归一；
$|\ell| > 10^{8} \max_i |b_i|$ 级别的远界被停用（视为无界，
由最终原空间审计兜底）。$u_j - \ell_j < 10^{-9}$ 的固定变量移出障碍公式
（\texttt{flb=fub=0}，$\theta_j = 0$，增广对角钉 $10^{12}$，目标值锁定 $\ell_j$）。
对应源码为 \srcpath{src/engine/kernel/ipm/ipm\_lp\_solver.cpp:NativeIPMLPAdapter::solve\_lp\_impl}；
常量为 \srcpath{src/engine/kernel/ipm/ipm\_lp\_solver\_internal.hpp} 的
\texttt{kBig}/\texttt{kTau}/\texttt{kMinVal}（路径锚点）。

\paragraph{Ruiz 均衡.}
\texttt{ruiz\_equilibrate} 对 $[A; A_{eq}]$ 做固定轮数（\texttt{ruiz\_rounds}，
包装层缺省见 \texttt{IPMLPOptions}）的交替均衡：每轮
$D_r \mathrel{*}= \operatorname{diag}(1/\sqrt{\mathrm{rowmax}})$、
$D_c \mathrel{*}= \operatorname{diag}(1/\sqrt{\mathrm{colmax}})$（零行列跳过），
每行/列每轮只乘一次。还原关系 $x = D_c \hat x$、$y = D_r \hat y$、
$z = \hat z / D_c$。对应源码为
\srcpath{src/engine/kernel/ipm/ipm\_lp\_solver\_internal.hpp:ruiz\_equilibrate}。

\paragraph{初始点.}
冷启动以 $\Theta = I$ 的法矩阵 $N = A_e A_e^{\top} + 10^{-10} I$ 解
最小二乘点 $x = A_e^{\top} N^{-1} b$、$y = N^{-1} A_e c$（增广形态改用
等价的准定 KKT $[I, -A_e^{\top}; -A_e, -\delta I]$ 同一分解两次求解）；
分量非有限或模 $\ge 10^{10}$ 时退回 $x = 0.5\mathbf{1}$、$y = 0$。
然后把变量钳入界内（冷启动双边界留 $\min(0.01\,\mathrm{range}, 1)$、
单边界留 $1.0$；暖启动留 $10^{-6}$），重算松弛
$s_i = \mathrm{clamp}(b_i - a_i x, 10^{-4}, \max(10^{-4}, u^{s}_i - 10^{-4}))$，
间隙 $g_l = \max(x - \ell, 10^{-4})$、$g_u = \max(u - x, 10^{-4})$。
对偶初值为 Mehrotra 式：令 $r_c = c - A_e^{\top} y$，双边列
$z_l^{0} = \max(r_c, 0)$、$z_u^{0} = \max(-r_c, 0)$（单边列只取对应侧），
然后整体上移 $d_z = \max(-1.5 \min z^{0}, 0) + 0.5\, \overline{g z}/\bar g$
（$\overline{gz}/\bar g$ 为间隙加权均值比，分母 $\le 10^{-30}$ 时取 $1$），
最终逐分量下截于 $10^{-6}$。对应源码为
\srcpath{src/engine/kernel/ipm/ipm\_lp\_solver.cpp:NativeIPMLPAdapter::solve\_lp\_impl}。

\paragraph{每迭代残差与互补.}
记 $\mathrm{fl}_j, \mathrm{fu}_j \in \{0, 1\}$ 为界存在标志（分支自由算术）。
每迭代计算
\begin{equation}
  r_p = b - A_e \hat x,\qquad
  r_{d,j} = c_j - (A_e^{\top} y)_j - \mathrm{fl}_j z_{l,j}
  + \mathrm{fu}_j z_{u,j},
  \label{eq:ipmlp:residuals}
\end{equation}
互补均值 $\mu = \bigl(\sum_j \mathrm{fl}_j g_{l,j} z_{l,j} +
\mathrm{fu}_j g_{u,j} z_{u,j}\bigr) / n_{\text{compl}}$，
$\theta_j = 1 / (z_{l,j}/g_{l,j} + z_{u,j}/g_{u,j} + \mathrm{reg})$。
对应源码为 \srcpath{src/engine/kernel/ipm/ipm\_lp\_solver.cpp:NativeIPMLPAdapter::solve\_lp\_impl}
（lambda \texttt{compute\_residuals} 与 $\theta$ 装配段）。

\paragraph{正则化预算.}
$\mathrm{reg}$ 的下限 $\mathrm{reg}_{\text{floor}}$ 在 Auto-增广主动路径取
$\sqrt{\varepsilon}$，其余取 $\varepsilon$；每迭代
\begin{equation}
  \mathrm{reg} = \mathrm{clamp}\Bigl(
  \min\bigl(10^{-6} \mu,\;
  \tfrac{0.1 \max(\mathrm{pf}, \mathrm{df}, \mu)}{1 + \|(x,y)\|_{\infty}}
  \bigr),\; \mathrm{reg}_{\text{floor}},\; 10^{-2}\Bigr),
  \label{eq:ipmlp:reg}
\end{equation}
（PARDISO 枢轴探针未决或已证稳定时直接取下限）。
方向非有限时把 \texttt{reg\_direction\_boost} 乘 $100$（上限 $10^{8}$）、
PARDISO 路径同时把枢轴扰动指数提到 $10$（boost $\ge 10^{4}$ 时 $8$）并
重新分解；有限迭代令 boost 除以 $10$ 回充。
对应源码为 \srcpath{src/engine/kernel/ipm/ipm\_lp\_solver.cpp:NativeIPMLPAdapter::solve\_lp\_impl}。

\paragraph{分解路径与动态重试.}
Newton 系统按形态分派：带宽 $\le 128$（\texttt{kBandedThreshold}）走带状
Cholesky；散布表超过 $4\times 10^{6}$ 项（\texttt{kDenseScatterThreshold}，
另受 $256\,$MB \texttt{kDenseMaxBytes} 与 $10^{8}$ 项
\texttt{kMaxScatterEntries} 门控）走稠密 $LL^{\top}$（$N = A_e \Theta
A_e^{\top}$ 以 syrk 成形）；否则走稀疏法方程（CHOLMOD/PARDISO/Accelerate/
Eigen LDLᵀ，按编译宏与符号分析结果选择）；\texttt{ForceAugmented} 或
Auto 判定法方程过稠时走准定增广系统
\begin{equation}
  \begin{bmatrix} \operatorname{diag}(d) + \mathrm{reg}\, I & -A_e^{\top}\\
  -A_e & -\mathrm{reg}\, I \end{bmatrix}
  \begin{bmatrix} d_x \\ d_y \end{bmatrix} =
  \begin{bmatrix} \xi \\ -r_p \end{bmatrix},
  \qquad d_j = z_{l,j}/g_{l,j} + z_{u,j}/g_{u,j},
  \label{eq:ipmlp:augmented}
\end{equation}
由 \texttt{factor\_kkt\_sparse}/\texttt{solve\_kkt\_sparse} 或 CHOLMOD
simplicial LDLᵀ 分解。分解失败的重试阶梯统一为
$\mathrm{reg} \times 100^{k}$（$k = 1..4$），带状路径重试前先把工作缓冲
恢复为未分解的填充。对应源码为
\srcpath{src/engine/kernel/ipm/ipm\_lp\_solver.cpp:NativeIPMLPAdapter::solve\_lp\_impl}
（lambda \texttt{fill\_and\_factor\_banded} 与同函数分解段）、
\srcpath{src/engine/kernel/ipm/ipm\_lp\_solver\_internal.hpp:banded\_chol\_factor}、
\srcpath{src/engine/kernel/ipm/ipm\_lp\_solver\_internal.hpp}（常量的路径锚点）。

\paragraph{预测子（仿射）.}
仿射右端 $\xi^{\text{aff}}_j = -r_{d,j} - \mathrm{fl}_j z_{l,j} +
\mathrm{fu}_j z_{u,j}$，解出 $(d_x^{\text{aff}}, d_y^{\text{aff}})$ 后
\begin{equation}
  d z_l^{\text{aff}} = \mathrm{fl}\bigl(-z_l - z_l g_l^{-1} d x^{\text{aff}}\bigr),
  \qquad
  d z_u^{\text{aff}} = \mathrm{fu}\bigl(-z_u + z_u g_u^{-1} d x^{\text{aff}}\bigr),
  \label{eq:ipmlp:affine-dz}
\end{equation}
以最大正向步 $(\alpha_p^{\text{aff}}, \alpha_d^{\text{aff}})$
（分数取 $1-\varepsilon$，legacy 对照开关取 \texttt{kTau}${}=0.9995$；
结果下截于 \texttt{kMinVal}${}=10^{-14}$）得
\begin{equation}
  \mu_{\text{aff}} = \frac{1}{n_{\text{compl}}} \sum_j
  \Bigl[\mathrm{fl}_j (g_{l,j} + \alpha_p^{\text{aff}} d x^{\text{aff}}_j)
  (z_{l,j} + \alpha_d^{\text{aff}} d z^{\text{aff}}_{l,j}) +
  \mathrm{fu}_j (g_{u,j} - \alpha_p^{\text{aff}} d x^{\text{aff}}_j)
  (z_{u,j} + \alpha_d^{\text{aff}} d z^{\text{aff}}_{u,j})\Bigr],
  \label{eq:ipmlp:mu-aff}
\end{equation}
\begin{equation}
  \sigma = \mathrm{clamp}\Bigl(\bigl(\max(0, \mu_{\text{aff}}/\mu)\bigr)^{3},
  0, 1\Bigr)
  \quad(\text{legacy 开关下}\ \min(\cdot, 0.5)),
  \qquad \mu = 0 \Rightarrow \sigma = 0.
  \label{eq:ipmlp:sigma}
\end{equation}
对应源码为 \srcpath{src/engine/kernel/ipm/ipm\_lp\_solver.cpp:NativeIPMLPAdapter::solve\_lp\_impl}、
\srcpath{src/engine/kernel/ipm/ipm\_lp\_solver\_internal.hpp:ipm\_maximum\_step\_lengths}。

\paragraph{校正子.}
当 $\alpha_p^{\text{aff}} > 0.9$、$\alpha_d^{\text{aff}} > 0.9$ 且
$\sigma < 0.02$ 时跳过校正直接用仿射方向（硬编码阈值）；否则以
\begin{equation}
  \xi_j = -r_{d,j} + \mathrm{fl}_j\bigl(\sigma\mu\, g_{l,j}^{-1} - z_{l,j}
  - d z^{\text{aff}}_{l,j} d x^{\text{aff}}_j g_{l,j}^{-1}\bigr)
  - \mathrm{fu}_j\bigl(\sigma\mu\, g_{u,j}^{-1} - z_{u,j}
  + d z^{\text{aff}}_{u,j} d x^{\text{aff}}_j g_{u,j}^{-1}\bigr)
  \label{eq:ipmlp:corrector}
\end{equation}
再解同一分解，并按
$d z_l = \mathrm{fl}((\sigma\mu - d z^{\text{aff}} d x^{\text{aff}}) g_l^{-1}
- z_l - z_l g_l^{-1} d x)$、
$d z_u = \mathrm{fu}((\sigma\mu + d z^{\text{aff}} d x^{\text{aff}}) g_u^{-1}
- z_u + z_u g_u^{-1} d x)$ 回收乘子方向。对应源码为
\srcpath{src/engine/kernel/ipm/ipm\_lp\_solver.cpp:NativeIPMLPAdapter::solve\_lp\_impl}。

\paragraph{Gondzio 多中心校正.}
仅当 $\min(\alpha_p, \alpha_d) < 0.9$ 且 $\mu > 0$ 时启用，最多
\texttt{max\_correctors} 轮（缺省自动值 $3$；当所选因子的值数组达到
$4\times10^{6}$ 字节（\texttt{kAutoCorrectorFactorBytes}，不再缓存驻留）时
自动降为 $0$，\texttt{max\_correctors}${}\ge 0$ 的显式设置覆盖自动值）。
每轮试探步 $\alpha^{t} = \min(\alpha + 0.1, 1)$
（\texttt{kDeltaAlpha}${}=0.1$），把试探互补积
$v = (g_l + \alpha_p^t d x)(z_l + \alpha_d^t d z_l)$ 等对
$[0.1\mu, 10\mu]$（\texttt{kBetaMin}/\texttt{kBetaMax}）截断得残差
$r_{gl} = \mathrm{clamp}(v) - v$，以零原始右端解
$\xi^{gc}_j = \mathrm{fl}_j r_{gl,j} g_{l,j}^{-1} -
\mathrm{fu}_j r_{gu,j} g_{u,j}^{-1}$，
合成新方向后仅当总步长增益
$\alpha_p^{\text{new}} + \alpha_d^{\text{new}} > \alpha_p + \alpha_d + 0.01$
（\texttt{kGammaAccept}${}\times{}$\texttt{kDeltaAlpha}）才保留，否则回滚并
终止本轮循环。对应源码为
\srcpath{src/engine/kernel/ipm/ipm\_lp\_solver.cpp:NativeIPMLPAdapter::solve\_lp\_impl}。

\paragraph{步长规则.}
\texttt{compute\_step\_lengths} 调用 \texttt{ipm\_centrality\_step\_lengths}：
\texttt{centrality\_step\_control} 关闭时退化为 \texttt{kTau} 的
\texttt{ipm\_maximum\_step\_lengths}；开启时先取严格分数 $1-\varepsilon$ 的
最大步 $(\bar\alpha_p, \bar\alpha_d)$，计算满步互补均值 $\mu_{\text{full}}$，
在阻塞分量上为其配对互补积保留 $0.1 \mu_{\text{full}}$ 缓冲
（\texttt{kGamma}${}=0.9$），例如原始阻塞于下界侧时
\begin{equation}
  \alpha_p = \mathrm{clamp}\Bigl(
  \frac{g_{l,j^*} - (1 - 0.9)\mu_{\text{full}}/
  (z_{l,j^*} + \bar\alpha_d d z_{l,j^*})}{-d x_{j^*}},\;
  0.9\,\bar\alpha_p,\; 1\Bigr),
  \label{eq:ipmlp:centrality-step}
\end{equation}
最终两侧都封顶 $1 - 10^{-6}$。迭代更新
$x \mathrel{+}= \alpha_p d_x$、$y \mathrel{+}= \alpha_d d_y$、
$z \leftarrow \max(z + \alpha_d d_z, 10^{-14})$、
$g_l \leftarrow \max(x - \ell, 10^{-14})$、$g_u \leftarrow \max(u - x, 10^{-14})$；
固定变量随后重新锁定。对应源码为
\srcpath{src/engine/kernel/ipm/ipm\_lp\_solver\_internal.hpp:ipm\_centrality\_step\_lengths}、
\srcpath{src/engine/kernel/ipm/ipm\_lp\_solver.cpp:NativeIPMLPAdapter::solve\_lp\_impl}
（lambda \texttt{compute\_step\_lengths} 与更新段）。

\paragraph{Newton 方向的迭代改进.}
\texttt{solve\_step} 对未正则化 KKT 残差
$r_1 = \xi - \operatorname{diag}(d) d_x + A_e^{\top} d_y$、
$r_2 = -r_p + A_e d_x$ 取合并无穷范数，目标为 Dembo--Eisenstat--Steihaug
强制项 $\eta = 0.1$（\texttt{kNewtonRefinementEta}）乘右端尺度
$\max(1, \|r_p\|_{\infty}, \|\xi\|_{\infty})$；超目标时以同一分解回代校正，
最多 $5$ 轮，\textbf{仅当合并残差严格下降才提交校正}，否则回滚该校正并停止。
法方程路径另有至多 $2$ 轮的 $N$ 残差精化，门限
$10^{-12}\max(1, \|\mathrm{rhs}\|_{\infty})$。
方向相对残差 \texttt{iteration\_kkt\_relative} 非有限（求解链失败）时
按``Inaccurate Newton direction''/``Normal equations rejected''中断。
对应源码为 \srcpath{src/engine/kernel/ipm/ipm\_lp\_solver.cpp:NativeIPMLPAdapter::solve\_lp\_impl}
（lambda \texttt{solve\_step}、\texttt{raw\_kkt\_solve}、\texttt{solve\_normal}）。

\paragraph{终止与发布.}
候选门（廉价触发）：绝对候选 $\mathrm{pf} < \tau_p$、$\mathrm{df} < \tau_d$、
$\mu < \tau_g$，或相对候选
$\mathrm{pf}/(1+\|b\|_{\max}) < 100 \tau_p$ 等三条同时成立，其中
$\tau = \max(\texttt{tol}, \texttt{publication\_tol})$、
$\texttt{publication\_tol} = \max(10^{-10}, 10\,\texttt{tol}) \times
\texttt{publication\_tol\_scale}$；触发后以原模型审计
\texttt{audit\_ipm\_lp\_optimality} 裁决。该审计在原坐标重组原始违反、
对偶违反（含行乘子符号条件与固定列的最小范数乘子重建
$z_l = \max(0, s_j)$、$z_u = \max(0, -s_j)$）与原始/对偶目标，接受条件
\begin{equation}
  \frac{\mathrm{viol}_p}{\max(1, \mathrm{primal\_scale})} \le \tau_p,\quad
  \frac{\mathrm{viol}_d}{\max(1, \|c\|_{\infty})} \le \tau_d,\quad
  \frac{|\mathrm{obj}_p - \mathrm{obj}_d|}
  {1 + |\mathrm{obj}_p| + |\mathrm{obj}_d|} \le \tau_g .
  \label{eq:ipmlp:audit}
\end{equation}
发布是 fail-closed 的：循环内自称收敛但终审不过者改判
``Optimality audit rejected''；未收敛但终审通过者补判
``Optimal (final original KKT audit)''。法方程停滞守卫：
$\min(\alpha_p, \alpha_d) \le \sqrt{\varepsilon}$ 连续两次即以
``Normal equations stalled'' 中断（外层再决定是否以增广形重启）。
对应源码为 \srcpath{src/engine/kernel/ipm/ipm\_lp\_solver.cpp:audit\_ipm\_lp\_optimality}、
\srcpath{src/engine/kernel/ipm/ipm\_lp\_solver.cpp:IPMLPOptimalityAudit::acceptable}、
\srcpath{src/engine/kernel/ipm/ipm\_lp\_solver.cpp:NativeIPMLPAdapter::solve\_lp\_impl}。

\paragraph{外层回退阶梯.}
\texttt{solve\_lp\_with\_presolve\_snapshot} 的 \texttt{direct\_solve}
按序执行：配置轮数求解 $\to$（迭代 $0$ 即被拒且 Auto 时）无缩放法方程
$\to$ 增广形升级（停滞/分解失败时保留有限原始迭代点作恢复起点，
停滞用 \texttt{StructurePreserving} 后端策略，其余用
\texttt{PivotingPortfolio}）$\to$ 无缩放重试 $\to$
$\max(4\,\texttt{ruiz\_rounds}, 40)$ 轮的极端缩放重试；
\texttt{keep\_better} 以``成功优先，否则
$\max(\mathrm{pf}, \mathrm{df}, \mu)$ 小者''裁决。speculative
（presolve 减缩模型）求解在 $\le 8$ 迭代失败时直接放弃升级阶梯。
对应源码为 \srcpath{src/engine/kernel/ipm/ipm\_lp\_solver.cpp:NativeIPMLPAdapter::solve\_lp\_with\_presolve\_snapshot}
（lambda \texttt{direct\_solve}、\texttt{run\_variant}、\texttt{keep\_better}）。

LP 恢复线程限制是默认关闭的实验：只有环境变量
\texttt{MIPSOLVERS\_EXPERIMENTAL\_LP\_RECOVERY\_CAP2} 严格等于
\texttt{1} 时才在增广恢复期间安装至多两个 MKL 线程的局部作用域。
默认恢复继承调用方资源范围。该实验尚未通过 R3 发布合同。
\texttt{MIPSOLVERS\_LP\_FACTOR\_TIMING} 控制独立的只读诊断：关闭时
不遍历 seed；开启时输出版本为 2 的 factor/recovery 记录，包含 seed
范数、目标、位指纹及后端选择序列。超过记录容量会显式标记 overflow，
验收器拒绝不完整序列。诊断字段不参与数值裁决。对应源码为
\srcpath{src/engine/kernel/ipm/ipm\_lp\_solver.cpp:capture\_lp\_recovery\_transition}、
\srcpath{src/engine/kernel/ipm/ipm\_lp\_solver.cpp:emit\_lp\_recovery\_transition}。

\paragraph{缓存节点路径.}
\texttt{prepare\_for\_node\_solves} 缓存基 LP 副本、G 行翻号、Ruiz 缩放、
CSR/散布表与符号分解；\texttt{solve\_cached\_node\_lp} 每节点只重建界
（缩放坐标下除以 $D_c$）并复用全部结构，暖启动取上一节点的
\texttt{last\_dual}（否则 $y = 0$）与 \texttt{last\_zl}/\texttt{last\_zu}
（逐分量钳到 $[10^{-6}, 10^{4}]$，否则取 $1$），冷启动原始点取
界中点。其迭代循环与 \texttt{solve\_lp\_impl} 相同（无增广形态：
带状或稀疏法方程）。\texttt{solve\_structure\_aware\_node\_lp} 在
``非带状且 $m > 256$ 且带宽 $> m/4$''时改走
\texttt{ForceAugmented} + \texttt{StructurePreserving} 的完整
\texttt{solve\_lp\_impl}。\texttt{update\_cached\_cost} 只更新缩放成本
（松弛成本保持 $0$）。对应源码为
\srcpath{src/engine/kernel/ipm/ipm\_lp\_solver\_cached.cpp:NativeIPMLPAdapter::prepare\_for\_node\_solves}、
\srcpath{src/engine/kernel/ipm/ipm\_lp\_solver\_cached.cpp:NativeIPMLPAdapter::solve\_cached\_node\_lp}、
\srcpath{src/engine/kernel/ipm/ipm\_lp\_solver\_cached.cpp:NativeIPMLPAdapter::solve\_structure\_aware\_node\_lp}、
\srcpath{src/engine/kernel/ipm/ipm\_lp\_solver\_cached.cpp:NativeIPMLPAdapter::update\_cached\_cost}。

% ─────────────────────────────────────────────────────────────────────────────
\subsection{KKT 装配契约与求解闸门}
\label{subsec:engine-kkt-contract}

\paragraph{analyze-once / factorize-many.}
\texttt{assemble\_augmented\_kkt\_cached} 组装
$[W + \delta_W I, J_g^{\top}; J_g, -\delta_C I]$：结构变化（以
\texttt{memcmp} 比较外层/内层索引指纹判定）时三元组组装一次并建立
scatter map——每个 $W$/$J_g$ 源条目在 KKT CSC 值数组中的位置（含上/下
三角两个拷贝与两组对角槽）；之后每次装配把值数组清零再按 scatter 位置
累加，$O(\mathrm{nnz})$ 零分配。\texttt{factor\_current\_kkt} 同样以模式
指纹决定是否重做 \texttt{analyze\_pattern}；维度或非零数超 int32 时
响亮失败（后端为 int32 索引）。对应源码为
\srcpath{src/engine/kernel/kkt/kkt\_system.cpp:assemble\_augmented\_kkt\_cached}、
\srcpath{src/engine/kernel/kkt/kkt\_system.cpp:build\_augmented\_scatter}（匿名命名空间）、
\srcpath{src/engine/kernel/kkt/kkt\_system.cpp:factor\_current\_kkt}、
\srcpath{src/engine/kernel/kkt/kkt\_system.cpp:factor\_kkt\_sparse}。

\paragraph{求解闸门（迭代改进的接受准则）.}
\texttt{solve\_kkt\_sparse} 是全部 KKT 数值解的最终闸门。残差限
$L = \tau \max(1, \|b\|_{\infty})$，$\tau$ =
\texttt{refinement\_tolerance}（$>0$ 时）否则 $\sqrt{\varepsilon}$；
首次求解后循环：$\|b - K x\|_{\infty} \le L$ 即停；否则解校正
$K \delta = r$，候选 $x + \delta$ \textbf{仅当其残差严格小于当前残差才提交}，
否则停止精化；循环结束后再做一次认证残差检查，仍超 $L$ 则整体返回失败
（由上层正则化/恢复策略处理，绝不返回未认证方向）。
对应源码为 \srcpath{src/engine/kernel/kkt/kkt\_system.cpp:solve\_kkt\_sparse}；
Schur 简化求解（切空间证书等小规模用途）为
\srcpath{src/engine/kernel/kkt/kkt\_system.cpp:solve\_kkt\_reduced\_sparse}。

% ─────────────────────────────────────────────────────────────────────────────
\subsection{稀疏后端}
\label{subsec:engine-la-backends}

\paragraph{默认后端选择.}
\texttt{make\_default\_sparse\_solver} 的编译期优先级为
KLU $>$ UMFPACK $>$ PARDISO（非对称 LU 版）$>$ SuperLU $>$ Eigen SparseLU；
\texttt{MIPSOLVERS\_LINEAR\_BACKEND} 环境变量可强制单后端（调试用）。
MUMPS 被刻意排除在通用默认之外：它吸收零/小主元而不报奇异，会使
$\delta_W$/$\delta_C$ 升级机制失去触发信号；MUMPS/PARDISO LDLᵀ 只在
需要惯性的 KKT 路径上由 \texttt{ensure\_mumps\_augmented\_solver}/
\texttt{ensure\_pardiso\_augmented\_solver} 显式选用。
对应源码为 \srcpath{src/engine/kernel/linear\_algebra/linear\_solver.cpp:make\_default\_sparse\_solver}、
\srcpath{src/engine/kernel/kkt/kkt\_system.cpp:ensure\_mumps\_augmented\_solver}（匿名命名空间，宏内）、
\srcpath{src/engine/kernel/kkt/kkt\_system.cpp:ensure\_pardiso\_augmented\_solver}（匿名命名空间，宏内）。

\paragraph{CHOLMOD 包装.}
\texttt{CholmodLDLT} 实现 analyze-once：\texttt{analyze} 把 int32 模式
一次性加宽到 int64 并以按值包装结构体别名调用方的值缓冲；
\texttt{factorize} 只重挂值指针并做数值分解（\texttt{cholmod\_l\_factorize}），
\texttt{solve} 走 \texttt{cholmod\_l\_solve2} 并复用预分配的解与工作缓冲；
\texttt{set\_simplicial(true)} 强制 simplicial LDLᵀ（准定/不定 KKT 用，
接受负 $D$ 主元）。无 CHOLMOD 编译宏时全部操作返回 false（能力探针语义）。
对应源码为 \srcpath{src/engine/kernel/linear\_algebra/cholmod\_ldlt.cpp:CholmodLDLT::analyze}、
\srcpath{src/engine/kernel/linear\_algebra/cholmod\_ldlt.cpp:CholmodLDLT::factorize}、
\srcpath{src/engine/kernel/linear\_algebra/cholmod\_ldlt.cpp:CholmodLDLT::solve}、
\srcpath{src/engine/kernel/linear\_algebra/cholmod\_ldlt.cpp:CholmodLDLT::set\_simplicial}。

\paragraph{HFactor 后端（vendored Forrest--Tomlin）.}
\texttt{HFactorBackend} 包装
\srcpath{src/engine/kernel/linear\_algebra/highs\_factor/} 的 HFactor 移植：
\texttt{factorize\_impl} 以 \texttt{setupGeneral} + \texttt{build} 分解基矩阵，
Markowitz 主元阈值缺省 $0.1$（\texttt{pivot\_threshold\_}），
秩判定最小绝对主元 $10^{-10}$（\texttt{kDefaultPivotTolerance}，
\texttt{MIPSOLVERS\_HFACTOR\_PIVOT\_TOL} 可覆盖）；阈值上下界为
\srcpath{src/engine/kernel/linear\_algebra/highs\_factor/util/HFactorConst.h}
的 \texttt{kMinPivotThreshold}${}=8\times10^{-4}$ /
\texttt{kMaxPivotThreshold}${}=0.5$。
\texttt{strengthen\_pivot\_threshold} 把阈值一次性升到 $0.5$
（仅原始清理路径在新鲜分解验证失败后使用）。
\texttt{update\_captured} 要求 AQ/EP 捕获与因子代数一致，主元
$a_q[p]$ 为 $0$ 或倒数非有限即拒绝；更新后透传 HFactor 的
\texttt{hint} 到 \texttt{refactor\_hint\_}（\texttt{needs\_refactorise}）。
对应源码为 \srcpath{src/engine/kernel/linear\_algebra/hfactor\_backend.cpp:HFactorBackend::factorize\_impl}、
\srcpath{src/engine/kernel/linear\_algebra/hfactor\_backend.cpp:HFactorBackend::update\_captured}、
\srcpath{src/engine/kernel/linear\_algebra/hfactor\_backend.cpp:HFactorBackend::strengthen\_pivot\_threshold}、
\srcpath{include/mipsolvers/engine/kernel/linear\_algebra/hfactor\_backend.hpp:HFactorBackend}。

\paragraph{SparseLUFactor（UMFPACK 基上的 FT 更新）.}
\texttt{SparseLUFactor} 在 UMFPACK 分解 $B = P^{\top} L U Q^{\top}$
（带行尺度 $R_s$）上维护 Forrest--Tomlin 链：
\texttt{ftran} 依次执行``置换+尺度播种种子 $\to$ L 图 DFS 可达性与拓扑序
前代 $\to$ FT 项按序回放（先行交换 $S_i$ 再消元 $E_i$）$\to$ U 回代
$\to$ 逆置换写出''，消元丢弃 $|v| \le 10^{-20}$ 的项；
\texttt{btran} 为对偶序（$Q^{\top}$ 播种、$U^{\top}$ 前代、逆序回放、
$L^{\top}$ 回代）。\texttt{update\_ft} 执行列循环移位把 spike 移到末列
得上 Hessenberg 矩阵，再做相邻行消元：次对角 $|h_{\text{sub}}| > 10^{-20}$
且对角 $|h_{\text{diag}}| < 10^{-14}$ 时交换相邻行并记录交换标志；
全部新对角以 $10^{-14}$ 检查，任一不过则拒绝更新（返回 false 由调用方
重建）；同时跟踪主元增长比 \texttt{max\_growth}。
\texttt{compute\_spike} 对入基列执行同一播种/DFS/前代/回放序列。
对应源码为 \srcpath{src/engine/kernel/linear\_algebra/sparse\_lu\_factor.cpp:SparseLUFactor::ftran}、
\srcpath{src/engine/kernel/linear\_algebra/sparse\_lu\_factor.cpp:SparseLUFactor::btran}、
\srcpath{src/engine/kernel/linear\_algebra/sparse\_lu\_factor.cpp:SparseLUFactor::update\_ft}、
\srcpath{src/engine/kernel/linear\_algebra/sparse\_lu\_factor.cpp:SparseLUFactor::compute\_spike}。

% ─────────────────────────────────────────────────────────────────────────────
\subsection{与既有文档素材的核对结论}
\label{subsec:engine-ipm-mismatch}

素材 \texttt{docs/manual/06-numerical-methods.md} 与本章转写按``以代码为准''
核对，发现以下不一致（仅声明，不改代码）：

\begin{enumerate}
  \item 素材 §6.5.2 称 $\delta_W$``起步踢量 $10^{-8}\|L_{xx}\|$、每次重试
        $\times 8$、上限 $10^{-2}\|L_{xx}\|$''。代码实际为：缺省起步
        $\sqrt{\varepsilon}\, S_W \approx 1.49\times10^{-8} S_W$、
        精确 2 幂升级（$\times 2$）、上限 $S_W/\sqrt{\varepsilon}
        \approx 6.7\times10^{7} S_W$；
        盲阶梯 \texttt{factor\_kkt\_with\_regularization} 同为 $\times 2$。
        对应源码为 \srcpath{src/engine/kernel/kkt/kkt\_system.cpp:resolve\_inertia\_settings}、
        \srcpath{src/engine/kernel/kkt/kkt\_system.cpp:increased\_primal\_regularization}（匿名命名空间）、
        \srcpath{src/engine/kernel/ipm/ipm\_solver.cpp:factor\_kkt\_with\_regularization}。
  \item 素材 §6.3 称 \texttt{make\_default\_sparse\_solver}``保持
        UMFPACK $>$ KLU $>$ \ldots''。代码实际的编译期优先级是
        KLU $>$ UMFPACK $>$ PARDISO $>$ SuperLU $>$ Eigen。
        对应源码为 \srcpath{src/engine/kernel/linear\_algebra/linear\_solver.cpp:make\_default\_sparse\_solver}。
  \item 素材 §6.4 称残差门控``$\tau \approx 10^{-12}$''。该数值只适用于
        LP-IPM 法方程路径的至多 2 轮精化
        （$10^{-12}\max(1, \|\mathrm{rhs}\|_{\infty})$）；
        KKT 闸门 \texttt{solve\_kkt\_sparse} 的缺省 $\tau = \sqrt{\varepsilon}
        \approx 1.49\times10^{-8}$，LP 的 Newton 方向精化则以
        $\eta = 0.1$ 的相对强制项为目标。对应源码为
        \srcpath{src/engine/kernel/kkt/kkt\_system.cpp:solve\_kkt\_sparse}、
        \srcpath{src/engine/kernel/ipm/ipm\_lp\_solver.cpp:NativeIPMLPAdapter::solve\_lp\_impl}。
  \item 素材 §6.5.1 的增广形写法 $[H+\delta_W I, J_g^{\top}, J_h^{\top};
        J_g, -\delta_C I, 0; J_h, 0, -S M^{-1}]$ 是未缩放形式；
        代码实际组装的是经 $T = \operatorname{diag}(I, I, \sqrt{\mu/s})$
        合同变换后的 \eqref{eq:ipm:augmented}（(3,3) 块为 $-I$）。
        两者惯性相同，数值行为不同。对应源码为
        \srcpath{src/engine/kernel/ipm/ipm\_solver.cpp:assemble\_augmented\_newton}。
\end{enumerate}


% source_equivalent_milp_bc.tex —— engine 实现转写：原生 MILP Branch-and-Cut
% 本文件为实现转写章，遵循 docs/modules/README.md 的数学模型规范。
% 所有公式只转写当前源码实际执行的赋值、阈值与条件分支；锚点禁止行号。

\section{原生 MILP Branch-and-Cut：求解主流程与树搜索}\label{sec:impl-milpbc}

\subsection{转写范围与口径}\label{subsec:milpbc-scope}

本章转写 \srcpath{src/engine/solver/native/milp/bc/} 下原生 branch-and-cut（B\&C）
求解器的当前实现，核验日期 2026-09-13。覆盖文件：

\begin{itemize}
  \item 公共入口 \srcpath{src/engine/solver/native/milp/bc/api.cpp}（4 个
        \texttt{solve\_*milp/minlp*\_bc} 重载的转发层）；
  \item 主流程 \srcpath{src/engine/solver/native/milp/bc/legacy/branch\_and\_cut.cpp}
        及其文本包含的 14 段 \srcpath{src/engine/solver/native/milp/bc/legacy/bc\_run/}
        下的 \texttt{.inc} 片段（\texttt{01\_setup\_presolve\_root\_build.inc} 至
        \texttt{14\_parallel\_path\_and\_result.inc}）；
  \item 预求解 \srcpath{src/engine/solver/native/milp/bc/milp\_presolve.cpp}；
  \item legacy 支撑单元（均位于
        \srcpath{src/engine/solver/native/milp/bc/legacy/}）：
        分支评分 \texttt{bc\_branching.cpp}、割生成
        \texttt{bc\_cuts.cpp}/\texttt{bc\_cuts\_transformed.cpp}/\texttt{bc\_cuts\_common.hpp}、
        CGLP \texttt{bc\_cglp.cpp}/\texttt{bc\_cglp\_model.cpp}、clique 表
        \texttt{bc\_clique\_table.cpp}、域探测 \texttt{bc\_domain\_probe.cpp}、
        隐含界 \texttt{bc\_implied\_bounds.cpp}、并行树 \texttt{bc\_parallel.cpp}、
        对偶证明 \texttt{bc\_utils\_dual\_proof.cpp}/\texttt{bc\_utils\_dual\_proof\_conflict.cpp}、
        根审计 \texttt{bc\_root\_audit.cpp}、解认证 \texttt{bc\_validation.cpp}、
        HiGHS 风格数值规则 \texttt{bc\_highs\_style\_numerics.cpp}、通用工具
        \texttt{bc\_utils.cpp}/\texttt{bc\_legacy\_helpers.cpp} 等；
  \item 共享头文件 \srcpath{include/mipsolvers/engine/detail/} 下的
        \texttt{bc\_types.hpp}、\texttt{bc\_pools.hpp}、\texttt{bc\_utils\_core.hpp}、
        \texttt{bc\_restart.hpp}、\texttt{bc\_threading.hpp}、\texttt{bc\_status.hpp}，
        与公共选项/统计 \srcpath{include/mipsolvers/engine/bc/options.hpp}、\srcpath{include/mipsolvers/engine/bc/stats.hpp}。
\end{itemize}

\begin{boundarynote}
口径约定：（i）除明确标注外，本章所有目标值、界与 gap 均在 B\&C 的
\textbf{内部最小化坐标}下叙述；最大化模型在入口处把目标取负，结果装配时再
符号还原。（ii）\texttt{bc\_run/*.inc} 为 \texttt{BCSolveState::run()} 的文本
包含片段，不独立编译；其中规则以 ``段号 + lambda 名'' 锚定，宿主函数锚点统一为\srcpath{src/engine/solver/native/milp/bc/legacy/branch\_and\_cut.cpp:BCSolveState::run}。
（iii）本章不含通用 B\&C 理论；理论层见本手册理论章。（iv）数值证据只引用
\texttt{docs/archive/} 中已冻结的实验记录，正文不复述其未达标预测为已实现行为。
\end{boundarynote}

\subsection{入口与主流程骨架}\label{subsec:milpbc-entry}

公开 API 是纯转发层：\texttt{solve\_milp\_bc}/\texttt{solve\_minlp\_bc} 的两个
重载（无/有 warm-start 与回调）分别直接调用
\texttt{detail::solve\_milp\_bc\_legacy\_core} 与
\texttt{detail::solve\_minlp\_bc\_legacy\_core}，不附加任何逻辑。
对应源码为 \srcpath{src/engine/solver/native/milp/bc/api.cpp:solve\_milp\_bc}、\srcpath{src/engine/solver/native/milp/bc/api.cpp:solve\_minlp\_bc}。

带回调的重载在进入主流程前执行三步前置处理：（1）把
\texttt{BCWarmStart} 的原始点安装进模型副本；（2）若提供了
\texttt{hyperparam\_tuner} 回调，先计算实例特征并用其返回值整体替换
\texttt{BCOptions}；（3）校验回调与选项的相容性。求解后再调用
\texttt{post\_solve} 回调。
对应源码为\srcpath{src/engine/solver/native/milp/bc/legacy/branch\_and\_cut.cpp:detail::solve\_milp\_bc\_legacy\_core}、\srcpath{src/engine/solver/native/milp/bc/legacy/bc\_legacy\_entry\_helpers.hpp:install\_primal\_warm\_start}、\srcpath{src/engine/solver/native/milp/bc/legacy/bc\_instance\_features.cpp:compute\_instance\_features}。

核心状态对象是 \texttt{BCSolveState}：持有原模型引用、（可能被 auto-tune 修改的）
选项副本与回调指针，唯一的公开方法是 \texttt{run()}。该头文件同时固定了一组
HiGHS 对齐的硬编码默认常量：MIP 可行容差 $10^{-6}$、小矩阵元阈值 $10^{-9}$、
割池年龄上限 $30$、割池软上限 $10000$、LP 行年龄上限 $10$。
对应源码为\srcpath{src/engine/solver/native/milp/bc/legacy/bc\_solve\_state.hpp:BCSolveState}。

\texttt{run()} 的函数体由 14 个文本包含段按编号拼接而成，构成唯一主流程：

\begin{enumerate}
  \item \texttt{01\_setup\_presolve\_root\_build.inc}：参数/环境初始化、整数性
        规范化、strict-HiGHS 入口、presolve、根节点构建；
  \item \texttt{02}–\texttt{05}：根 LP 松弛与根割循环（HiGHS 侧状态、隐含界
        分离、tableau 源、crossover 与根收尾）；
  \item \texttt{06}–\texttt{08}：根启发式（feasibility pump、sub-MIP/舍入驱动、
        progressive rounding、LP diving、LNS、根 probing）与并行树并发启动决策；
  \item \texttt{09}：树基础设施（\texttt{NodeQueue}、全局下界、节点 LP 构造器）；
  \item \texttt{10}–\texttt{11}：节点域传播与冲突/蕴含学习、队列管理；
  \item \texttt{12}：单子节点处理（LP 求解、池分离、割生成）与对偶证明消费；
  \item \texttt{13}：顺序主循环；
  \item \texttt{14}：异构并行探索路径与最终结果装配/postsolve。
\end{enumerate}

对应源码为\srcpath{src/engine/solver/native/milp/bc/legacy/branch\_and\_cut.cpp:BCSolveState::run}
与 \srcpath{src/engine/solver/native/milp/bc/legacy/bc\_run/}。
MINLP 入口转发到同一遗产核心的 NLP 变体。
对应源码为\srcpath{src/engine/solver/native/milp/bc/legacy/branch\_and\_cut.cpp:detail::solve\_minlp\_bc\_legacy\_core}、\srcpath{src/engine/solver/native/milp/bc/legacy/bc\_minlp\_legacy.cpp}。

全程状态字符串集中在 \texttt{bc\_status} 命名空间，如
\texttt{kTimeLimitReached}（\texttt{"Time limit reached"}）、
\texttt{kNodeLimitReached}、\texttt{kSearchQueueExhausted}、
\texttt{kOptimalityGapReached}、\texttt{kOptimalTreeExhausted}、
\texttt{kFeasibleIncumbentFound}、\texttt{kNoFeasibleIntegerSolutionFound} 等；
结果装配只从这些常量取词，不另行拼接（审计失败等例外情形除外，见
\ref{subsec:milpbc-cert}）。
对应源码为\srcpath{include/mipsolvers/engine/detail/bc\_status.hpp}。

\subsection{全局时限合约与可选根预算}\label{subsec:milpbc-deadline}

段 01 在任何嵌套阶段之前启动唯一求解时钟 $t_0$，并定义三个贯穿全程的 lambda：
\begin{align}
  \mathrm{elapsed} &= \mathrm{now}()-t_0,\\
  \mathrm{remaining} &=
    \begin{cases}
      +\infty, & \texttt{time\_limit\_sec}\le 0\ \text{或非有限},\\
      \texttt{time\_limit\_sec}-\mathrm{elapsed}, & \text{否则},
    \end{cases}\\
  R_{\mathrm{finalize}} &=
    \begin{cases}
      \min\bigl(1.0,\ \max(0.05,\ 0.10\cdot\texttt{time\_limit\_sec})\bigr),
        & \texttt{time\_limit\_sec}>0\ \text{且有限},\\
      0, & \text{否则},
    \end{cases}
\end{align}
其中 $R_{\mathrm{finalize}}$ 是``根后收尾保留时间''。可选根工作（割轮次之间的
检查、树准备的各阶段）的截止判据为
\begin{equation}
  \texttt{bc\_optional\_root\_budget\_expired}:\quad
  \texttt{time\_limit\_sec}>0\ \land\ \mathrm{remaining}\le R_{\mathrm{finalize}} .
\end{equation}
$0.05$、$0.10$、$1.0$ 三个常数为硬编码。对应源码为\srcpath{src/engine/solver/native/milp/bc/legacy/bc\_run/01\_setup\_presolve\_root\_build.inc}
中的 lambda \texttt{bc\_remaining\_sec} 与 \texttt{bc\_optional\_root\_budget\_expired}
（宿主函数\srcpath{src/engine/solver/native/milp/bc/legacy/branch\_and\_cut.cpp:BCSolveState::run}）。
该``单一全局预算被所有嵌套阶段继承''的合约及其缺陷修复记录见
\srcpath{docs/archive/native\_milp\_cplex\_gap\_improvement\_2026-09-12.md}。

\subsection{Presolve：strict-HiGHS 入口与原生管线}\label{subsec:milpbc-presolve}

段 01 的 presolve 选择逻辑如下。当 \texttt{highs\_root\_native\_tree\_contract}
开启且剩余预算为正时，先以 \texttt{retain\_postsolve=true} 运行一次 HiGHS
presolve 预检并缓存其结果；同一原始模型只允许这一次权威 presolve，后续阶段
复用该侧状态而不再调用。strict HiGHS 合约（\texttt{strict\_highs\_mip\_contract}
或 LP 内核后端为 HiGHS 时的派生量）决定是否走直接 \texttt{evaluateRootNode}
式的上游根生命周期；非 strict 路径则在 PaPILO 与原生 \texttt{MILPPresolve}
之间按选项选择。
对应源码为\srcpath{src/engine/solver/native/milp/bc/legacy/bc\_run/01\_setup\_presolve\_root\_build.inc}、\srcpath{src/engine/solver/native/milp/bc/legacy/bc\_conformance\_trace.cpp:cached\_highs\_presolve\_side\_state}
（侧状态存储的实现在 \srcpath{src/engine/strategy/highs\_presolve\_side\_state.cpp}）。

原生 MILP presolve 的不动点循环为：每轮依次执行
固定变量消除 $\to$ 空行/空列消除 $\to$ 单行处理 $\to$ 单列处理 $\to$
界收紧 $\to$ 固定变量消除 $\to$ forcing row $\to$ 固定变量消除 $\to$
空行/空列消除 $\to$ doubleton 方程代入 $\to$ 被占优列 $\to$ 平行行
（仅前 3 轮）$\to$ 二进制 probing $\to$ 末尾再各做一次固定变量与空行/空列
消除；每步之后立刻检查不可行性，轮数上界为 \texttt{opts.max\_rounds}，
一轮零变更即终止。每个约简边界与每 256 个 probing 候选/行轮询一次全局
deadline。
对应源码为\srcpath{src/engine/solver/native/milp/bc/milp\_presolve.cpp:MILPPresolve::run}、\srcpath{src/engine/solver/native/milp/bc/milp\_presolve.cpp:MILPPresolve::deadline\_expired}、\srcpath{src/engine/solver/native/milp/bc/milp\_presolve.cpp:MILPPresolve::run\_probing}。

\begin{boundarynote}
失败语义（须在引用处保持可见）：\texttt{MILPPresolve::run} 在不可行或超时时
\textbf{不发布任何部分约简模型}——\texttt{rebuild\_model} 只在
既非不可行又未超时被跳过，调用方必须保留原坐标系。postsolve 与正向映射分别由
\srcpath{src/engine/solver/native/milp/bc/milp\_presolve.cpp:MILPPresolve::postsolve}
与\srcpath{src/engine/solver/native/milp/bc/milp\_presolve.cpp:MILPPresolve::forward\_map}
提供；另有原位变体
\srcpath{src/engine/solver/native/milp/bc/milp\_presolve.cpp:presolve\_inplace}。
\end{boundarynote}

HiGHS 保留 presolve 的坐标恒等式在代码中以两个工具函数固定：目标偏移
$f_{\mathrm{orig}} = f_{\mathrm{red}} + \Delta$ 的还原
\begin{equation}
  f_{\mathrm{red}} = f_{\mathrm{orig}} - \Delta_{\mathrm{presolve}},
\end{equation}
以及逐列仿射映射 $x_{\mathrm{orig}} = s\,x_{\mathrm{red}} + c$ 下变量界
$x_t \le a x_z + b$ 的约简空间转运
\begin{equation}
  y_t \le \frac{a\,s_z}{s_t}\,y_z + \frac{a\,c_z + b - c_t}{s_t},
  \qquad s_t<0\ \text{时上下界角色互换}.
\end{equation}
对应源码为\srcpath{src/engine/solver/native/milp/bc/legacy/bc\_highs\_style\_numerics.cpp:presolved\_objective\_value}、\srcpath{src/engine/solver/native/milp/bc/legacy/bc\_highs\_style\_numerics.cpp:highs\_style\_map\_variable\_bound\_to\_reduced}；
推导记录见
\srcpath{docs/archive/native\_milp\_root\_source\_eligibility\_2026-08-13.md} 与\srcpath{docs/archive/native\_milp\_presolve\_varbound\_coordinates\_2026-08-13.md}。

\subsection{根 LP 松弛与根割循环}\label{subsec:milpbc-root}

根 LP 松弛由 \texttt{solve\_lp\_relaxation} 系列完成（单纯形/IPM/vendored
HiGHS 三种后端），其结果类型 \texttt{LPRelaxationResult} 除原始解与对偶界外
还携带 \texttt{basis\_hint}、\texttt{row\_duals} 与标志
\texttt{requires\_ipm\_nodes}：当 IPM$\to$单纯形 crossover 未能产出可用基时，
后续节点 LP 必须路由到 IPM 而非从缺基状态冷启动单纯形——这是一条显式声明的
回退路径。
对应源码为\srcpath{src/engine/solver/native/milp/bc/legacy/bc\_relaxation.cpp:solve\_lp\_relaxation}、\srcpath{src/engine/solver/native/milp/bc/legacy/bc\_relaxation.cpp:solve\_lp\_relaxation\_with\_vendored\_highs}、\srcpath{include/mipsolvers/engine/detail/bc\_types.hpp:LPRelaxationResult}。

\paragraph{割预算与轮次上限。}
段 05 在割循环前把根剩余预算的有界份额分配给割：
\begin{align}
  B_{\mathrm{usable}} &= \max\bigl(0,\ \mathrm{remaining}-R_{\mathrm{finalize}}\bigr),
  \qquad
  B_{\mathrm{cut}} = 0.35\,B_{\mathrm{usable}},\\
  R_{\mathrm{eff}} &= \texttt{root\_cut\_rounds}\ \text{再取}\ 
  \begin{cases}
    \min(R_{\mathrm{eff}},3), & B_{\mathrm{cut}}<2\,\mathrm{s},\\
    \min\bigl(R_{\mathrm{eff}},\lfloor B_{\mathrm{cut}}/0.05\rfloor\bigr),
      & \text{否则},
  \end{cases}
\end{align}
且 $B_{\mathrm{cut}}\le 0.05$\,s 时 $R_{\mathrm{eff}}=0$。比例 $0.35$ 与每轮
$0.05$\,s 的下界估计均为硬编码。规模自适应进一步按 presolved 行数压减轮次与
每轮割数（阈值取 \texttt{cut\_budget\_*\_row\_threshold} 系列选项），SCUC 形态
（二进制占比高、行数适中的判定）则放宽 GMI 过滤器并把停滞上限压到 2。
对应源码为\srcpath{src/engine/solver/native/milp/bc/legacy/bc\_run/05\_root\_relaxation\_d.inc}
（宿主\srcpath{src/engine/solver/native/milp/bc/legacy/branch\_and\_cut.cpp:BCSolveState::run}）；
预算原则的推导见
\srcpath{docs/archive/native\_milp\_cplex\_gap\_improvement\_2026-09-12.md}。

\paragraph{根割源与停滞判据。}
根部分离按``变换源（tableau/path/mod-k）$\to$ 通用 Gomory/MIR 轮次 $\to$
根 probing''的顺序组织。变换源循环采用 HiGHS 风格的不动点：一轮所选行成功
追加并重优化后，即使目标零抬升也保留这些行；终止由平滑进展停滞规则拥有——
连续 \texttt{highs\_source\_stall}${}\ge 3$ 次无进展即停止。每轮之间检查
\texttt{bc\_optional\_root\_budget\_expired}。
对应源码为\srcpath{src/engine/solver/native/milp/bc/legacy/bc\_run/05\_root\_relaxation\_d.inc}、\srcpath{src/engine/solver/native/milp/bc/legacy/bc\_cuts\_transformed.cpp:add\_transformed\_tableau\_cuts}。
该退化不动点的缺陷史与``保留零抬升行''的实验结论（界增强被接受、端到端
运行时间合同失败而回滚生产策略）记录于
\srcpath{docs/archive/native\_milp\_degenerate\_cutpool\_fixed\_point\_2026-08-13.md}；
当前代码以选项与生命周期门槛体现其最终形态，正文不要把它读成``零抬升行
一律回滚''。

变换根的源合格性判据为
\begin{equation}
  \mathrm{eligible}\iff
  \mathrm{root\_bound}\ \text{有限}\ \land\
  \bigl(\mathrm{declared\_int}>0\ \lor\ \mathrm{implied\_int}>0\bigr),
\end{equation}
即声明整数列与 presolve 推断整数列的并集非空。
对应源码为\srcpath{src/engine/solver/native/milp/bc/legacy/bc\_highs\_style\_numerics.cpp:highs\_style\_root\_source\_eligible}；
资格修正的两轮实现保真记录见
\srcpath{docs/archive/native\_milp\_root\_source\_eligibility\_2026-08-13.md}。

\paragraph{根审计。}
三个诊断/审计函数在割提交路径上复核实现保真性：
（i）\texttt{audit\_root\_cut\_rows} 对每个新加行 $a^Tx\le b$ 计算 LP 点与
（可选）incumbent 点的违反量，缩放因子
$\mathrm{scale}=\max\{1,\|a\|,|b|\}$，LP 违反 $\le \texttt{tol}\cdot\mathrm{scale}$
记为未违反，incumbent 违反 $>100\,\texttt{tol}\cdot\mathrm{scale}$ 计数；
（ii）\texttt{audit\_root\_cut\_standard\_form\_rows} 把新加行重进标准形并逐系数
比对（含行符号翻转与下界平移后的右端项
$\mathrm{sign}\cdot(b-\sum_j a_j\,\mathrm{lbshift}_j)$）；
（iii）\texttt{audit\_root\_reduced\_costs} 统计非零 reduced cost 并枚举
reduced-cost fixing 候选：对处于下界的整数列 $j$，
\begin{equation}
  \mathrm{new\_ub}_j=\Bigl\lfloor \frac{\mathrm{gap}}{|\bar c_j|}
  + l_j + \texttt{int\_tol}\Bigr\rfloor,
  \qquad
  \mathrm{gap}=\mathrm{incumbent}-\mathrm{node\_bound},
\end{equation}
处于上界时对称地取
$\mathrm{new\_lb}_j=\lceil u_j-\mathrm{gap}/|\bar c_j|-\texttt{int\_tol}\rceil$；
仅当 $|\bar c_j|(u_j-l_j)>\mathrm{gap}+10^{-9}$ 且能收紧时计数。
对应源码为\srcpath{src/engine/solver/native/milp/bc/legacy/bc\_root\_audit.cpp:audit\_root\_cut\_rows}、\srcpath{src/engine/solver/native/milp/bc/legacy/bc\_root\_audit.cpp:audit\_root\_cut\_standard\_form\_rows}、\srcpath{src/engine/solver/native/milp/bc/legacy/bc\_root\_audit.cpp:audit\_root\_reduced\_costs}。

\paragraph{重启前置条件。}
根重启遵循 HiGHS 的``RENS 前置 + 非活跃整数率门槛''合约：低分数根上的 RENS
槽位（\texttt{enable\_root\_low\_fractionality\_rens}）调用受限 sub-MIP 助手，
其内部限制为至多 200 个子节点、500 个叶子、12 个无改进节点与
$\min(3\,\mathrm{s},\,5\%\,\text{剩余时间})$；strict HiGHS 根属主上的重 presolve
重启被显式关闭（\texttt{mip\_allow\_restart=false}），因为该合约下根坐标固定。
对应源码为\srcpath{src/engine/solver/native/milp/bc/legacy/bc\_run/05\_root\_relaxation\_d.inc}
（\texttt{highs\_root\_owner} 配置点，宿主
\srcpath{src/engine/solver/native/milp/bc/legacy/branch\_and\_cut.cpp:BCSolveState::run}）；
前提审计与多轮实测记录见
\srcpath{docs/archive/native\_milp\_root\_quality\_restart\_prerequisites\_2026-08-13.md}。
树级（非根）重启由状态化控制器掌管，见 \ref{subsec:milpbc-tree}。

\subsection{割生成与割池管理}\label{subsec:milpbc-cuts}

\paragraph{族调度。}
\texttt{add\_cuts} 是割生成的总分派器：按 \texttt{opt.cuts} 枚举
（\texttt{None}/\texttt{IntRounding}/\texttt{MIR}/\texttt{Gomory}/\texttt{Cover}
等）选择族序列，在共享预算 \texttt{max\_cuts} 内依次消费。每族的可选统计由
\texttt{CutFamilyTracker} 维护，admission 后按
\begin{equation}
  \mathrm{efficacy}(a^Tx\le b)=\frac{\max(0,\ a^Tx-b)}{\|a\|_2}
\end{equation}
记录该轮平均效力。
对应源码为\srcpath{src/engine/solver/native/milp/bc/legacy/bc\_cuts.cpp:add\_cuts}、\srcpath{include/mipsolvers/engine/detail/bc\_types.hpp:CutFamilyTracker}。

自适应门控：每族保留最近 $H=5$ 轮的环形效力历史；轮数达到 $H$ 后，若滑动平均
\begin{equation}
  \overline{\mathrm{eff}}=\frac{1}{H}\textstyle\sum_{i=0}^{H-1}\mathrm{history}_i
  \;<\; \varepsilon_{\min}=10^{-3},
\end{equation}
则该族被禁用（\texttt{disabled=true}），后续轮次跳过并计
\texttt{skipped\_rounds}；未满 $H$ 轮时滑动平均按乐观值 $1.0$ 处理。
$H=5$ 与 $\varepsilon_{\min}=10^{-3}$ 均为硬编码常量。
对应源码为\srcpath{include/mipsolvers/engine/detail/bc\_types.hpp:CutFamilyStats::add\_round}、\srcpath{include/mipsolvers/engine/detail/bc\_types.hpp:CutFamilyStats::should\_skip}。

\paragraph{各分离器。}
实现包含：基 MIR 与行 MIR
（\texttt{add\_basis\_mir\_cuts}/\texttt{add\_row\_mir\_cuts}）、基 Gomory
（\texttt{add\_basis\_gomory\_cuts}）、cover（\texttt{add\_cover\_cuts} 及投影
容量变体 \texttt{add\_projected\_capacity\_cover\_cuts}）、clique
（\texttt{add\_clique\_cuts}）、zero-half（\texttt{add\_zero\_half\_cuts}）、
整数舍入类 MIR（\texttt{add\_mir\_like\_cuts}）以及变换 tableau 割
（\texttt{add\_transformed\_tableau\_cuts}）。
对应源码为\srcpath{src/engine/solver/native/milp/bc/legacy/bc\_cuts.cpp:add\_basis\_mir\_cuts}、\srcpath{src/engine/solver/native/milp/bc/legacy/bc\_cuts.cpp:add\_row\_mir\_cuts}、\srcpath{src/engine/solver/native/milp/bc/legacy/bc\_cuts.cpp:add\_basis\_gomory\_cuts}、\srcpath{src/engine/solver/native/milp/bc/legacy/bc\_cuts.cpp:add\_cover\_cuts}、\srcpath{src/engine/solver/native/milp/bc/legacy/bc\_cuts.cpp:add\_clique\_cuts}、\srcpath{src/engine/solver/native/milp/bc/legacy/bc\_cuts.cpp:add\_zero\_half\_cuts}、\srcpath{src/engine/solver/native/milp/bc/legacy/bc\_cuts\_transformed.cpp:add\_transformed\_tableau\_cuts}。
GMI 准入阈值（\texttt{gmi\_min\_efficacy=1e-3}、\texttt{gmi\_max\_density=0.35}、
\texttt{gmi\_max\_parallelism=0.90} 等）为 \texttt{BCOptions} 中的硬编码默认，
对应源码为 \srcpath{include/mipsolvers/engine/bc/options.hpp}。

CGLP（cut-generating LP）析取割单独成链：候选选择
$\to$ 对每个候选列 $j$ 构造 master LP $\to$ 求解并抽取割 $\to$ 根处批量追加。
对应源码为\srcpath{src/engine/solver/native/milp/bc/legacy/bc\_cglp.cpp:select\_cglp\_candidates}、\srcpath{src/engine/solver/native/milp/bc/legacy/bc\_cglp\_model.cpp:build\_cglp\_lp}、\srcpath{src/engine/solver/native/milp/bc/legacy/bc\_cglp.cpp:generate\_cglp\_cut}、\srcpath{src/engine/solver/native/milp/bc/legacy/bc\_cglp.cpp:add\_cglp\_cuts\_root}；
采纳数经 \texttt{BCStats::cglp\_accepted} 透出到公共统计的
\texttt{cglp\_cuts\_added}。

\paragraph{割池（CutPool）。}
全局树割池 \texttt{CutPool} 默认容量 2000、年龄上限 50（由
\texttt{cut\_pool\_max\_size}/\texttt{cut\_pool\_max\_age} 选项给出）。插入规则：
\begin{itemize}
  \item 仅接受 \texttt{ValidityScope::GlobalCut}；rhs 必须有限，系数范数满足
        $\|a\|\ge 10^{-12}$ 且有限，维度必须匹配 \texttt{expected\_size}；
  \item 重复拒绝模仿 HiGHS cutpool 合约：同支撑（support hash 分桶后逐一比对）
        且平行度
        \begin{equation}
          \mathrm{parallelism}=\frac{a_1^{T}a_2}{\|a_1\|\,\|a_2\|}
          \;\ge\; 1-10^{-6}
        \end{equation}
        时视为重复；若新行 rhs 更紧（缩放后差超过
        $10^{-9}\max\{1,|\mathrm{rhs}|,|\mathrm{rhs}_{\mathrm{old}}|\}$）则替换旧行，
        或新行把 \texttt{domain\_only} 行升级为 LP 可见行时合并，否则拒绝；
  \item 池满时用新行替换年龄最大者并重建支撑桶。
\end{itemize}
分离-老化一体化接口 \texttt{select\_violated\_and\_age(x, min\_viol, max\_selected)}
以默认 $\texttt{min\_viol}=10^{-4}$ 选出违反行（重置其年龄为 0），其余行年龄
加 1，随后清除年龄 $>\texttt{max\_age}$ 的行。
对应源码为\srcpath{include/mipsolvers/engine/detail/bc\_pools.hpp:CutPool::add}、\srcpath{include/mipsolvers/engine/detail/bc\_pools.hpp:CutPool::select\_violated\_and\_age}、\srcpath{include/mipsolvers/engine/detail/bc\_pools.hpp:CutPool}。

\texttt{domain\_only=true} 的行只参与域传播（
\srcpath{include/mipsolvers/engine/detail/bc\_pools.hpp:CutPool::propagate\_with\_reasons}），
不注入节点 LP；这是代码中显式的近似/分工路径。

\paragraph{割的证明所有权。}
``两段式所有权''谓词固定为
\begin{equation}
  \mathrm{tree\_owned}(\mathrm{row})\iff
  \mathrm{alive}(\mathrm{row})\ \land\ \mathrm{in\_root\_lp}(\mathrm{row}),
\end{equation}
即一条生成行只有在通过割池选择、成功追加进根 LP 并重优化后才成为树全局的
证明构件；在 strict 根不动点或自动 HiGHS 根流水线
（\texttt{auto\_highs\_root\_pipeline}）任一激活时，候选行在根 LP 准入前保持
根私有：
\begin{equation}
  \mathrm{stays\_root\_owned}\iff
  \texttt{strict\_highs\_root\_fixed\_point}\ \lor\
  \texttt{auto\_highs\_root\_pipeline}.
\end{equation}
对应源码为\srcpath{src/engine/solver/native/milp/bc/legacy/bc\_highs\_style\_numerics.cpp:highs\_style\_cutpool\_row\_tree\_owned}、\srcpath{src/engine/solver/native/milp/bc/legacy/bc\_highs\_style\_numerics.cpp:highs\_style\_cutpool\_candidate\_stays\_root\_owned}；
所有权缺陷（最终播种才被当作所有权边界的实现失真）与修正记录见
\srcpath{docs/archive/native\_milp\_cutpool\_proof\_ownership\_2026-08-13.md}。

\subsection{域传播、clique 与冲突学习}\label{subsec:milpbc-domain}

节点域传播的组合件包括：clique 表 \texttt{CliqueTable}（边/字面量边插入
\texttt{add\_edges}/\texttt{add\_literal\_edges}/\texttt{add\_literal\_edges\_from\_cut}，
不动点传播 \texttt{propagate}）、域探测工作区
\texttt{BCDomainProbeWorkspace}（\texttt{probe(col, value\_one)} 试探、
\texttt{commit} 提交、\texttt{propagate} 传播）、隐含界与变量界源表
（\texttt{tighten\_integral\_row\_sides}、\texttt{append\_graph\_implied\_bound\_cuts}）。
对应源码为\srcpath{src/engine/solver/native/milp/bc/legacy/bc\_clique\_table.cpp:CliqueTable::propagate}、\srcpath{src/engine/solver/native/milp/bc/legacy/bc\_domain\_probe.cpp:BCDomainProbeWorkspace::probe}、\srcpath{src/engine/solver/native/milp/bc/legacy/bc\_implied\_bounds.cpp:tighten\_integral\_row\_sides}、\srcpath{src/engine/solver/native/milp/bc/legacy/bc\_proof\_artifacts.cpp:append\_graph\_implied\_bound\_cuts}。

冲突侧的数据结构在 \texttt{bc\_types.hpp}：域字面量
\texttt{BranchDomainLiteral}、冲突子句 \texttt{ConflictClause}、带子树作用域的
\texttt{ScopedConflictClause}、可全局消费的界提升证书
\texttt{BoundLiftingCertificate}，以及 HiGHS \texttt{domchgstack\_} 对应的节点
局部域轨迹 \texttt{LocalDomainTrailEntry}（每条局部界变更记录同侧前一活跃界）。
对应源码为 \srcpath{include/mipsolvers/engine/detail/bc\_types.hpp}。
子句池 \texttt{ConflictPool} 与重复拒绝工具（支撑/平行度判定）集中在
\srcpath{include/mipsolvers/engine/detail/bc\_pools.hpp:ConflictPool}。

对偶证明行把 LP 对偶转化为域侧证书：完整构造
\texttt{build\_full\_dual\_proof\_row}、reduced-cost 质量版
\texttt{build\_reduced\_cost\_mass\_dual\_proof\_row}、一致性校验
\texttt{dual\_proof\_identity\_ok}、域固定应用
\texttt{apply\_dual\_proof\_domain\_fixing}；冲突解析沿局部轨迹回解出
reason 侧前沿（HiGHS \texttt{ConflictSet} 的对应物），状态枚举为
\texttt{Success}/\texttt{InvalidInput}/\texttt{ActivityFailed}/\texttt{MissingReason}/%
\texttt{ScopeBlocked}/\texttt{LiteralLimit}。
对应源码为\srcpath{src/engine/solver/native/milp/bc/legacy/bc\_utils\_dual\_proof.cpp:build\_full\_dual\_proof\_row}、\srcpath{src/engine/solver/native/milp/bc/legacy/bc\_utils\_dual\_proof.cpp:apply\_dual\_proof\_domain\_fixing}、\srcpath{src/engine/solver/native/milp/bc/legacy/bc\_utils\_dual\_proof\_conflict.cpp:resolve\_dual\_proof\_conflict\_from\_local\_trail}、\srcpath{include/mipsolvers/engine/detail/bc\_utils\_core.hpp:DualProofResolutionStatus}。

\subsection{队列与节点选择}\label{subsec:milpbc-queue}

顺序树使用 \texttt{NodeQueue}：节点实体存于 id 映射，另维护 priority 视图、
DFS 栈视图与两个按界有序的多重索引。优先级键为``估计引导''量
\begin{equation}
  \mathrm{key}(\mathrm{node})=
  \begin{cases}
    +\infty, & \mathrm{bound}\ \text{非有限},\\
    \mathrm{bound}, & \mathrm{estimate}\ \text{非有限},\\
    \tfrac12\bigl(\mathrm{bound}+\max(\mathrm{bound},\mathrm{estimate})\bigr),
      & \text{否则},
  \end{cases}
\end{equation}
而精确下界索引始终单独维护（不混入估计）。
对应源码为\srcpath{include/mipsolvers/engine/detail/bc\_utils\_core.hpp:node\_selection\_priority}、\srcpath{include/mipsolvers/engine/detail/bc\_utils\_core.hpp:NodeQueue}。

三种 \texttt{NodeSelection} 模式的 pop 语义：
\begin{itemize}
  \item \texttt{DepthFirst}：恒走 DFS 栈；
  \item \texttt{BestBound}：从无 incumbent 起即按优先级/下界索引取最小；
  \item \texttt{Hybrid}（默认）：无 incumbent 时 DFS；出现 incumbent 后切换到
        优先级 pop，但每 $P=10$ 次出队（\texttt{kHybridBestBoundPeriod}，硬编码）
        或当前出队界不劣于切换前下界（容差
        $\max(10^{-9},10^{-12}\max(1,|\mathrm{lb}|))$）时强制一次 best-bound pop。
\end{itemize}
\texttt{pop\_bestbound} 在 priority 与 DFS 两个有序索引间取界更小者；
\texttt{lower\_bound()} 为两者最小值；\texttt{audit\_lower\_bound} 重算全部
存活节点以核对索引一致性（规模容差 $\max(10^{-9},\texttt{tol})\cdot
\max\{1,|\mathrm{indexed}|,|\mathrm{computed}|\}$）。
对应源码为\srcpath{include/mipsolvers/engine/detail/bc\_utils\_core.hpp:NodeQueue::pop}、\srcpath{include/mipsolvers/engine/detail/bc\_utils\_core.hpp:NodeQueue::pop\_bestbound}、\srcpath{include/mipsolvers/engine/detail/bc\_utils\_core.hpp:NodeQueue::lower\_bound}、\srcpath{include/mipsolvers/engine/detail/bc\_utils\_core.hpp:NodeQueue::audit\_lower\_bound\_index}。

入队节点的域存储采用稀疏增量格式：\texttt{compact\_node\_domain} 只保留与根域
不同的列界对并释放稠密向量，\texttt{materialize\_node\_domain} 在出队时还原；
\texttt{Node::branch\_child} 以 copy-on-write 共享父节点的域学习负载
（branch reasons、域轨迹、局部割等），子节点深度为父深度 $+1$。
对应源码为\srcpath{include/mipsolvers/engine/detail/bc\_types.hpp:compact\_node\_domain}、\srcpath{include/mipsolvers/engine/detail/bc\_types.hpp:materialize\_node\_domain}、\srcpath{include/mipsolvers/engine/detail/bc\_types.hpp:Node::branch\_child}。

\subsection{定界、剪枝与主循环}\label{subsec:milpbc-mainloop}

\paragraph{incumbent 剪枝。}
节点剪枝谓词（内部最小化坐标）为
\begin{equation}
  \mathrm{prune}\iff
  \mathrm{has\_incumbent}\ \land\
  \mathrm{node\_bound}\ \ge\ \mathrm{incumbent\_obj}-\texttt{prune\_tol},
  \qquad \texttt{prune\_tol}=\max(10^{-9},\ \texttt{lp\_tol}),
\end{equation}
队列级批量版把同一判据作用于全部排队节点（共享容差常量
\texttt{kIncumbentPruneTol}$=10^{-12}$ 用于部分队列维护路径——两处容差不同，
属实现现状）。
对应源码为\srcpath{src/engine/solver/native/milp/bc/legacy/bc\_utils.cpp:incumbent\_prunes\_node}、\srcpath{include/mipsolvers/engine/detail/bc\_utils\_core.hpp:NodeQueue::prune\_by\_lower\_bound\_limit}。

搜索所用的 incumbent 上限不是 incumbent 目标本身，而是 HiGHS
\texttt{computeNewUpperLimit} 的对应物：若目标在整数列上可按积分尺度
$s>0$ 缩放（全整数目标系数的最大可公度尺度，见
\texttt{highs\_style\_objective\_integral\_scale}），则
\begin{equation}
  U=\frac{\lceil s\cdot\mathrm{incumbent}-\mathrm{tol}\rceil-1}{s}+\mathrm{tol},
  \qquad \mathrm{tol}=\max(10^{-9},\texttt{feastol}),
\end{equation}
否则 $U=\min\{\mathrm{incumbent}-\mathrm{tol},\
\mathrm{nextafter}(\mathrm{incumbent},-\infty)\}$。
对应源码为\srcpath{src/engine/solver/native/milp/bc/legacy/bc\_highs\_style\_numerics.cpp:highs\_style\_incumbent\_upper\_limit}、\srcpath{src/engine/solver/native/milp/bc/legacy/bc\_highs\_style\_numerics.cpp:highs\_style\_integral\_scale}。

\paragraph{gap 剪枝。}
当 \texttt{require\_tree\_exhaustion\_certificate=false}（默认，即通常的
gap 容忍模式）且根 gap 证书就绪、incumbent 经审计有效时，最优性界限为
\begin{equation}
  L_{\mathrm{gap}}=\mathrm{incumbent\_obj}
  -\texttt{gap\_tol}\cdot\max\bigl(1,\ |\mathrm{incumbent\_obj}|\bigr),
\end{equation}
界 $\ge L_{\mathrm{gap}}-\tau$ 的节点被剪掉并计入
\texttt{gap\_suboptimal\_queue\_prunes}，其中
\begin{equation}
  \tau=\max\bigl\{10^{-7},\ 10\max(0,\texttt{lp\_tol}),\
  10^{-9}\max(1,|L_{\mathrm{gap}}|)\bigr\}.
\end{equation}
（代码注释明确禁止在此处把容差按 $|\mathrm{obj}|$ 放大到积分尺度以外。）
对应源码为\srcpath{src/engine/solver/native/milp/bc/legacy/bc\_run/11\_node\_domain\_conflict\_b.inc}
中的 lambda \texttt{gap\_optimality\_limit}、\texttt{gap\_optimality\_prune\_tol} 与
\texttt{prune\_queue\_by\_gap\_optimality\_limit}（宿主\srcpath{src/engine/solver/native/milp/bc/legacy/branch\_and\_cut.cpp:BCSolveState::run}）。

\paragraph{主循环。}
段 13 的顺序主循环每次迭代依次执行：
\begin{enumerate}
  \item 时限检查：$\mathrm{elapsed}>\texttt{time\_limit\_sec}$ 置
        \texttt{kTimeLimitReached} 退出；节点数达 \texttt{max\_nodes} 置
        \texttt{kNodeLimitReached} 退出；
  \item incumbent 驱动的树重启判定（见 \ref{subsec:milpbc-tree}）；
  \item IPM 节点模式下的人口控制：若剩余预算减收尾保留小于一对节点 LP 的估计
        代价
        \begin{equation}
          \hat T_{\mathrm{pair}}=2.2\max\bigl(0.05,\ T_{\mathrm{root\_lp}}\bigr),
        \end{equation}
        则刷新队列下界并按 \texttt{kTimeLimitReached} 退出（$2.2$ 与 $0.05$
        为硬编码）；
  \item proof-tail 模式激活时：触发域闭包、强制前沿回放的队列剪枝与 gap 剪枝，
        并以至多 8 轮的``前沿收缩''循环刷新近下界节点（每轮刷新带
        $\max(\mathrm{termination\_band},\,2\,\texttt{lp\_tol})$）；
  \item pop：proof-tail 模式恒 \texttt{pop\_bestbound}，否则按
        \ref{subsec:milpbc-queue} 的模式语义；队列空时若 incumbent 有限则按
        \begin{equation}
          \mathrm{gap}=\max\Bigl(0,\ \frac{\mathrm{incumbent}-g_{\mathrm{lb}}}
          {\max(1,|\mathrm{incumbent}|)}\Bigr)
        \end{equation}
        判定 \texttt{kOptimalityGapReached}（要求根证书就绪且
        $\mathrm{gap}\le\texttt{gap\_tol}$）或 \texttt{kSearchQueueExhausted}；
  \item 冲突池命中、incumbent 剪枝、gap 剪枝依次作用于弹出节点；
  \item 迟到学习的域回放：若节点落后于当前学习 epoch，则以
        \texttt{run\_node\_domain\_closure} 重放域传播，域变更后必须重解 LP
        （先单纯形、失败回退 IPM）方可提交更紧域——注释明确禁止用旧基/旧
        $x_{\mathrm{relax}}$ 搭配更紧域继续；
  \item 节点处理与分支（段 12 与段 13 后半）。
\end{enumerate}
对应源码为\srcpath{src/engine/solver/native/milp/bc/legacy/bc\_run/13\_main\_loop\_sequential.inc}
（宿主\srcpath{src/engine/solver/native/milp/bc/legacy/branch\_and\_cut.cpp:BCSolveState::run}）
与域闭包 lambda 所在段
\srcpath{src/engine/solver/native/milp/bc/legacy/bc\_run/11\_node\_domain\_conflict\_b.inc}。

\subsection{分支：伪费用与可靠性分支的实现公式}\label{subsec:milpbc-branch}

\paragraph{伪费用统计。}
每个可分支列维护方向性统计：观测和 $\Sigma$、平方和 $Q$、计数 $n$，以及
inference、cutoff、conflict 三路辅助计数。派生量：
\begin{align}
  \mathrm{avg} &= \Sigma/n\ (n=0\ \text{时取}\ 1.0),\\
  \mathrm{var} &= \max\Bigl(0,\ \frac{Q-n\bar x^{2}}{n-1}\Bigr)\ (n<2\ \text{时取}\ 0),\\
  \mathrm{lcb} &= \max\Bigl(10^{-6},\ \mathrm{avg}
    -1.96\,\sqrt{\mathrm{var}/n}\Bigr)\ (n<2\ \text{时取}\ \mathrm{avg}),\\
  \mathrm{hist} &= 1+0.05\ln(1+\overline{\mathrm{inf}})
    +0.25\,\mathrm{cutoff\_ratio}
    +0.05\ln\bigl(1+\mathrm{conflict\_score}\bigr),
\end{align}
其中 cutoff 比率为 $\mathrm{cutoff\_cnt}/\max(1,\mathrm{cutoff\_cnt}+n)$。
$1.96$（95\% 置信）与各权重 $0.05/0.25/0.05$ 均为硬编码。
对应源码为\srcpath{include/mipsolvers/engine/detail/bc\_types.hpp:PseudoCost}。

子节点目标增益归一化为单位距离费用：
\begin{equation}
  \widehat g=\frac{\max(0,g)}{\max(10^{-9},\ |\mathrm{dist}|)}
  \quad(g\ \text{非有限或}\ g\le 0\ \text{时记}\ 0).
\end{equation}
对应源码为\srcpath{include/mipsolvers/engine/detail/bc\_types.hpp:normalized\_pseudocost\_gain}。

\paragraph{方向评分与乘积评分。}
对分数变量 $j$（分数部分 $f=x_j-\lfloor x_j\rfloor$），方向评分按证据状态分派：
\begin{equation}
  s_{\mathrm{dir}}=
  \begin{cases}
    \max\bigl(\max(0,g_{\mathrm{obs}})\cdot\mathrm{hist},\ \varepsilon\bigr),
      & \text{Measured/ExactOptimal 证据},\\
    +\infty, & \text{ProvenCutoff 证据},\\
    \max\bigl(d\cdot\max(\mathrm{lcb}+\Delta,\ \varepsilon)\cdot
      \mathrm{hist},\ \varepsilon\bigr), & \text{否则（历史预测）},
  \end{cases}
\end{equation}
其中 $d$ 为该方向分支距离（下支 $f$、上支 $1-f$），$\varepsilon=10^{-6}$，
$\Delta$ 为预测偏移（调用方多传 0）。候选总分为两方向乘积
$S_j=s_{\mathrm{down}}\,s_{\mathrm{up}}$；选择时取
$\log(1+\max(0,S_j))$ 最大者，支持静态优先级覆盖与动态先验加性 overlay
（动态先验回调异常、空返回或维度不符时整体回退到无先验评分——显式回退路径）。
对应源码为\srcpath{src/engine/solver/native/milp/bc/legacy/bc\_branching.cpp:compute\_directional\_branch\_score}、\srcpath{src/engine/solver/native/milp/bc/legacy/bc\_branching.cpp:compute\_branch\_product\_score}、\srcpath{src/engine/solver/native/milp/bc/legacy/bc\_branching.cpp:choose\_branch\_var\_pseudocost\_with\_overlay}。

分支方向先走哪侧由 \texttt{choose\_up\_branch\_first} 决定：有 incumbent 且
恰好一侧被证 cutoff 时先走另一侧；无 incumbent 时比较两方向评分走更弱侧
（先探低增益分支）；否则遵循 advisory 方向或既有偏好。
对应源码为\srcpath{src/engine/solver/native/milp/bc/legacy/bc\_branching.cpp:choose\_up\_branch\_first}。

基础策略枚举：\texttt{FirstFractional}（首候选）、\texttt{MostInfeasible}
（$\tfrac12-|f-\tfrac12|$ 最大者，即最靠近 $0.5$）、\texttt{Pseudocost}
（默认，上述乘积评分）。
对应源码为\srcpath{src/engine/solver/native/milp/bc/legacy/bc\_branching.cpp:choose\_branch\_var}。

\paragraph{可靠性分支（reliability branching）。}
\texttt{choose\_branch\_var\_reliability} 的实现规则：
\begin{itemize}
  \item 深度 $>6$ 或 \texttt{max\_probes=0} 时退化为纯伪费用选择（硬编码深度
        上限）；
  \item 有效探测预算
        \begin{equation}
          P_{\mathrm{eff}}=
          \begin{cases}
            \texttt{max\_probes}, & \mathrm{depth}=0,\\
            \max\bigl(2,\ \lfloor\texttt{max\_probes}/2\rfloor\bigr), & \text{否则};
          \end{cases}
        \end{equation}
  \item 方向可靠门槛：若该列两方向均已有观测
        （$\mathrm{down\_cnt}\ge1\land\mathrm{up\_cnt}\ge1$），取
        $\max(2,\lfloor\texttt{reliability\_limit}/2\rfloor)$，否则取
        \texttt{reliability\_limit}；
  \item 探测 LP 在同一标准形的界事务（\texttt{StandardFormBoundTransaction}）
        上执行：应用分支界 $\to$ 持久单纯形重解 $\to$ 回滚；域已空
        （$\mathrm{lb}>\mathrm{ub}+10^{-12}$）直接记 ProvenCutoff；LP 结果按
        \texttt{classify\_lp\_result} 归为 Measured/ExactOptimal/ProvenCutoff
        （不可行且持 Farkas 证书或目标截断）/UnknownFailure 四态，再按
        \texttt{update\_pseudocost\_from\_branch\_evidence} 入账——失败态
        \textbf{不}向全局伪费用分布注入哨兵增益；
  \item 循环``按当前观测选最优候选 $\to$ 探测其不可靠方向''直至 $P_{\mathrm{eff}}$
        个候选探测完或候选全部可靠。
\end{itemize}
对应源码为\srcpath{src/engine/solver/native/milp/bc/legacy/bc\_branching.cpp:choose\_branch\_var\_reliability\_impl}、\srcpath{src/engine/solver/native/milp/bc/legacy/bc\_branching.cpp:update\_pseudocost\_from\_branch\_evidence}、\srcpath{src/engine/solver/native/milp/bc/legacy/bc\_branching.cpp:classify\_branch\_probe\_reuse}。
证明尾段曾绕过可靠性分支的缺陷及``局部强分支评分不蕴含固定节点数下
全局最优界推进''的失配结论见
\srcpath{docs/archive/native\_milp\_reliability\_branching\_2026-08-13.md}；
当前代码已无该绕过，但也不要把乘积评分读成全局界推进的保证。

\paragraph{节点估计。}
队列键所用的 \texttt{estimate} 由 \texttt{compute\_node\_estimate} 给出：对每个
分数可分支列取方向评分的较弱者 $\min(s_{\mathrm{down}},s_{\mathrm{up}})$ 作为
变量抬升，按 \texttt{NodeEstimateAggregation} 取和或取 max 加到节点界上：
\begin{equation}
  \mathrm{estimate}=\max\Bigl(\mathrm{bound},\ \mathrm{bound}
  +\!\!\sum_{j\,\mathrm{frac}}\!\!\min(s_{\mathrm{down}}^j,s_{\mathrm{up}}^j)\Bigr),
  \qquad
  \Delta=\max\Bigl(10^{-6},\ \frac{\texttt{int\_tol}\max(1,|\mathrm{bound}|)}
  {\max(1,n_{\mathrm{frac}})}\Bigr).
\end{equation}
对应源码为\srcpath{src/engine/solver/native/milp/bc/legacy/bc\_branching.cpp:compute\_node\_estimate}。

\subsection{并行树调度}\label{subsec:milpbc-parallel}

线程数解析：\texttt{num\_threads}$\ge 1$ 时照用；否则取硬件并发数并被
\texttt{auto\_parallel\_min\_threads}（默认 2）/%
\texttt{auto\_parallel\_max\_threads}（$0$ 表示不设上限）夹取；硬件并发不可用
时退回最小值。
对应源码为\srcpath{src/engine/solver/native/milp/bc/legacy/bc\_legacy\_helpers.cpp:resolve\_num\_threads}。

并行树的准入条件（段 05/07/13 共同决定）：
\begin{itemize}
  \item \texttt{par\_mode}$\iff$解析线程数 $>1$；
  \item \texttt{parallel\_serial\_on\_low\_fractionality\_root}（默认开）在
        ``大而低分数''的根上把树保持串行
        （\texttt{serial\_low\_fractionality\_root}）；但若已有审计 incumbent 且
        根 gap 未闭，\texttt{parallel\_proof\_on\_incumbent\_gap}（默认开）重新放行
        （\texttt{parallel\_proof\_low\_fractionality}）；
  \item \texttt{parallel\_delay\_until\_incumbent}（默认开）在尚无有限 incumbent
        时把主循环线程数压为 1（\texttt{serial\_until\_incumbent}）；
  \item 并发（早期）启动还要求：线程数 $\ge2$、剩余时间
        $\ge\max(0.25,\ 0.25\,\texttt{time\_limit\_sec})$ 秒、可选根预算未耗尽；
        否则推迟到根启发式之后的晚期启动（\texttt{late\_tree*}）。
\end{itemize}
$0.25$\,s 与 $25\%$ 为硬编码。调度原因的枚举值（\texttt{single\_thread}、
\texttt{enabled}、\texttt{concurrent\_tree}、\texttt{late\_tree} 等）写入
\texttt{BCStats::parallel\_schedule\_reason}。
对应源码为\srcpath{src/engine/solver/native/milp/bc/legacy/bc\_run/05\_root\_relaxation\_d.inc}、\srcpath{src/engine/solver/native/milp/bc/legacy/bc\_run/07\_root\_heuristics\_b.inc}、\srcpath{src/engine/solver/native/milp/bc/legacy/bc\_run/13\_main\_loop\_sequential.inc}
（宿主\srcpath{src/engine/solver/native/milp/bc/legacy/branch\_and\_cut.cpp:BCSolveState::run}）；
启动门槛的推导与 1/4 线程实测见
\srcpath{docs/archive/native\_milp\_cplex\_gap\_improvement\_2026-09-12.md}。

并行态由 \texttt{ConcurrentTreeState} 持有：共享队列
\texttt{ThreadSafeNodeQueue}、共享割池/冲突池/解池/共享 incumbent、原子统计与
停止标志。启动时把根割逐行播种进共享割池、把现有 incumbent 播种进共享
incumbent/解池，并把 gap 最优性界限
$\mathrm{incumbent}-\texttt{gap\_tol}\max(1,|\mathrm{incumbent}|)$ 写入
\texttt{optimality\_limit}（仅在根证书就绪且非树穷举证书模式时）。并行单纯形
选项被设为：迭代上限 $\max(10\,\texttt{max\_lp\_iter},2000)$、容差
$\max(10^{-10},0.1\,\texttt{lp\_tol})$、禁止冷启动、且\textbf{关闭}内核内
incumbent 截断（\texttt{incumbent\_bound=nullptr}）——注释说明并行节点 LP 是
证明载体，内核内截断只暴露中间目标而非已认证界。
对应源码为\srcpath{include/mipsolvers/engine/detail/bc\_concurrent\_tree\_state.hpp:ConcurrentTreeState}、\srcpath{include/mipsolvers/engine/detail/bc\_threading.hpp:ThreadSafeNodeQueue}、\srcpath{src/engine/solver/native/milp/bc/legacy/bc\_run/14\_parallel\_path\_and\_result.inc}。

工作者角色按线程号分派：启用混合 IPM（探索者 $\ge2$ 且根行数超过
\texttt{hybrid\_ipm\_threshold}）时 0 号线程为 \texttt{IPMDiver}；探索者
$\ge2$ 时末号线程为 \texttt{Prover}（best-bound 聚焦、可靠性分支）；其余为
\texttt{Diver}（DFS 聚焦、快找 incumbent）。当
\texttt{cuts}$\ne$\texttt{None} 且 \texttt{max\_cut\_depth}$>0$ 时预留一个线程跑
\texttt{cut\_worker\_thread}，探索者数 $=$ 线程数 $-1$。
对应源码为\srcpath{include/mipsolvers/engine/detail/bc\_threading.hpp:WorkerRole}、\srcpath{src/engine/solver/native/milp/bc/legacy/bc\_parallel.cpp:explorer\_thread}、\srcpath{src/engine/solver/native/milp/bc/legacy/bc\_parallel.cpp:cut\_worker\_thread}。

\begin{boundarynote}
并行正确性边界（代码注释明确要求保持可见）：
\texttt{share\_conflict\_learning\_across\_threads} 默认 \textbf{关闭}——跨线程
共享冲突子句/冲突割/二元蕴含会把任何作用域错误放大到整片森林，使并行运行
剪掉串行仍访问的区域；关闭时各探索者保留自己的局部冲突池与蕴含图，仅共享
根割（\texttt{shared\_cp}）、incumbent 与解池。
\texttt{deterministic\_parallel}（opt-in）用轮转令牌串行化队列 pop/子节点
push，使节点探索顺序只是线程数与队列键的函数；LP 求解仍并行。工作窃取池以
\texttt{random\_seed} 确定性播种。对应源码为\srcpath{include/mipsolvers/engine/bc/options.hpp} 与\srcpath{src/engine/solver/native/milp/bc/legacy/bc\_parallel.cpp:explorer\_thread}。
\end{boundarynote}

\subsection{树重启控制器}\label{subsec:milpbc-tree}

incumbent 驱动的顺序树重启由 \texttt{IncumbentTreeRestartController} 判定。
\texttt{note\_incumbent} 以相对容差
$10^{-10}\max\{1,|\mathrm{obj}|,|\mathrm{pending}|\}$ 去重后登记待决 incumbent；
\texttt{ready} 的全部前置条件（同时满足才允许重启）：
\begin{gather}
  \texttt{enabled}\ \land\ \mathrm{pending}\ \land\ \mathrm{proof\_valid}
  \ \land\ \mathrm{restart\_count}<\texttt{max\_restarts}\
  \land\ \mathrm{open\_nodes}\ge\texttt{min\_open\_nodes}\notag\\
  \land\ \mathrm{node}-\mathrm{last\_restart\_node}\ge
    \texttt{min\_nodes\_since\_restart}
  \ \land\ \mathrm{remaining}\ge\texttt{min\_remaining\_time\_sec},
\end{gather}
且当已有上次重启目标时还要求相对改进
\begin{equation}
  \frac{\mathrm{last\_obj}-\mathrm{pending\_obj}}
       {\max(1,|\mathrm{last\_obj}|)}+10^{-15}
  \ \ge\ \texttt{min\_relative\_incumbent\_improvement}.
\end{equation}
默认参数：\texttt{max\_restarts=2}、\texttt{min\_nodes\_since\_restart=256}、
\texttt{min\_open\_nodes=64}、\texttt{min\_relative\_incumbent\_improvement=0.01}、
\texttt{min\_remaining\_time\_sec=0.25}，默认 \texttt{enabled=false}（硬编码，
见 \texttt{BCTreeRestartOptions}）。控制器只判定时机；队列替换与证明状态归
调用点所有。
对应源码为\srcpath{include/mipsolvers/engine/detail/bc\_restart.hpp:IncumbentTreeRestartController::ready}、\srcpath{include/mipsolvers/engine/detail/bc\_restart.hpp:IncumbentTreeRestartController::note\_incumbent}、\srcpath{include/mipsolvers/engine/bc/options.hpp:BCTreeRestartOptions}。

\subsection{解的认证、gap 计算与结果装配}\label{subsec:milpbc-cert}

\paragraph{incumbent 审计。}
\texttt{validate\_incumbent\_solution} 在原模型（或指定的约简源模型）上复核
候选点：维度一致后逐项计算
\begin{align}
  \mathrm{viol}_{\mathrm{bnd}}&=\max_j\max(0,\ l_j-x_j,\ x_j-u_j),\\
  \mathrm{viol}_{\mathrm{int}}&=\max_{j\,\mathrm{int}}|x_j-\mathrm{round}(x_j)|,\\
  \mathrm{viol}_{\mathrm{ineq}}&=\max_i\max\bigl(0,\ (Ax)_i-b_i,\
  (\mathrm{lhs}_i\ \text{有限时})\ \mathrm{lhs}_i-(Ax)_i\bigr),\\
  \mathrm{viol}_{\mathrm{eq}}&=\max_i|(A_{\mathrm{eq}}x)_i-b_{\mathrm{eq},i}|,
\end{align}
四项均 $\le\texttt{tol}$ 才通过（结果装配处使用 $\texttt{tol}=10^{-6}$）。
对应源码为\srcpath{src/engine/solver/native/milp/bc/legacy/bc\_validation.cpp:validate\_incumbent\_solution}。

\paragraph{坐标还原与报告恒等式。}
段 14 的报告函数固定为
\begin{equation}
  \mathrm{report}(f)=\begin{cases}
    f, & f\ \text{非有限},\\
    s_{\mathrm{sense}}\,(f+\Delta), & \text{否则},
  \end{cases}
  \qquad s_{\mathrm{sense}}=
  \begin{cases}
    +1,&\text{Minimize},\\ -1,&\text{Maximize},
  \end{cases}
\end{equation}
其中 $\Delta$ 按激活的 presolve 路径取 strict-HiGHS 偏移、HiGHS MIP 偏移、
PaPILO 偏移或 0。即：内部界先加偏移回到原空间、再按目标方向还原符号。
对应源码为\srcpath{src/engine/solver/native/milp/bc/legacy/bc\_run/14\_parallel\_path\_and\_result.inc}
中的 lambda \texttt{report\_objective\_value}（宿主\srcpath{src/engine/solver/native/milp/bc/legacy/branch\_and\_cut.cpp:BCSolveState::run}）；
偏移坐标审计见
\srcpath{docs/archive/native\_milp\_root\_source\_eligibility\_2026-08-13.md}
（``Root-Gap Coordinate Audit''节）。

\paragraph{gap 与状态映射。}
有 incumbent 时若尚未登记 gap 且全局下界有限，则
\begin{equation}
  \texttt{bc\_stats.gap}=\max\Bigl(0,\ \frac{\mathrm{incumbent}-g_{\mathrm{lb}}}
  {\max(1,|\mathrm{incumbent}|)}\Bigr),\qquad
  \texttt{stats.mip\_gap}=\texttt{gap}\ (\text{有限时否则}\ 0);
\end{equation}
队列穷举（\texttt{kSearchQueueExhausted}）且无未安全回退剪枝
（\texttt{fallback\_unsafe\_prunes=0}）、gap 不超 \texttt{gap\_tol} 时，改报
\texttt{kOptimalTreeExhausted}、gap 置 0 并把 \texttt{best\_bound} 提升为
\texttt{best\_obj}（注释：避免已证最优仍显示陈旧树界造成的 7--21\%``假
gap''）；其余穷举情形保留实测 gap——队列耗尽但未闭合证明 gap 被视为
前沿失败而非最优性证明。无 incumbent 时 \texttt{mip\_gap=+$\infty$}、状态
\texttt{kNoFeasibleIntegerSolutionFound}、\texttt{success=false}。
对应源码为\srcpath{src/engine/solver/native/milp/bc/legacy/bc\_run/14\_parallel\_path\_and\_result.inc}
（宿主\srcpath{src/engine/solver/native/milp/bc/legacy/branch\_and\_cut.cpp:BCSolveState::run}）。

\paragraph{失败即拒发。}
发布前 incumbent 依次通过：约简（pre-cut 根模型）空间审计，失败则清空解、
置 \texttt{kInvalidReducedIncumbent} 系状态并把 \texttt{success} 翻为 false；
随后按 presolve 路径 postsolve 到原空间并再做
\texttt{validate\_incumbent\_solution} 原模型审计，失败同样拒发
（\texttt{kInvalidPapiloPostsolveIncumbent}/\texttt{kInvalidHighsPostsolveIncumbent}
等）。\texttt{success} 仅在审计通过时保持 true。
对应源码为\srcpath{src/engine/solver/native/milp/bc/legacy/bc\_run/14\_parallel\_path\_and\_result.inc}
与状态常量
\srcpath{include/mipsolvers/engine/detail/bc\_status.hpp}。

\begin{boundarynote}
边界与已知限制（均来自代码注释或冻结实验记录，非待办）：
（i）\texttt{use\_papilo\_presolve} 等选项控制的 PaPILO 路径在编译期依赖缺失
时整体缺席，运行时状态字符串会指明 presolve 来源（\texttt{kInfeasiblePapiloPresolve}
等）；（ii）HiGHS 根剖面 \texttt{auto\_highs\_root\_pipeline} 的行为范围不限于根
生命周期——基准剖面变更曾因下游树效应被否决，当前生产基准保持通用原生树策略
（\srcpath{docs/archive/native\_milp\_highs\_lp\_root\_profile\_2026-08-13.md}）；
（iii）浅节点池再分离的行数门槛在生产中被恢复为 500 行版本（
\srcpath{docs/archive/native\_milp\_degenerate\_cutpool\_fixed\_point\_2026-08-13.md}
末节的否决记录）；（iv）\texttt{presolve\_inplace} 不删列，仅把约简结果写回原
\texttt{LPModel}，其统计语义与 \texttt{MILPPresolve::run} 不同，引用时注意区分。
\end{boundarynote}


% source_equivalent_api_adapters.tex —— engine 实现转写：求解器门面与适配器调度
% 本文件为实现转写章，遵循 docs/modules/README.md 的数学模型规范。
% 所有规则只转写当前源码实际执行的逻辑；锚点禁止行号。

\section{求解器门面、适配器注册与分派}\label{sec:impl-api-adapters}

\subsection{转写范围与口径}\label{subsec:api-scope}

本章转写引擎对外 API 与适配器层的当前实现，核验日期 2026-09-13。覆盖文件：
\begin{itemize}
  \item 门面 \srcpath{src/engine/api/solver.cpp}（\texttt{SolverEngine} 全部成员）；
  \item 适配器基类与注册表
        \srcpath{src/engine/solver/solver\_adapter.cpp}、
        \srcpath{src/engine/solver/adapter\_registry.cpp} 及对应头文件
        \srcpath{include/mipsolvers/engine/solver/solver\_adapter.hpp}、
        \srcpath{include/mipsolvers/engine/api/result.hpp}、
        \srcpath{include/mipsolvers/engine/api/options.hpp}、
        \srcpath{include/mipsolvers/engine/solve\_context.hpp}；
  \item 分派器 \srcpath{src/engine/strategy/dispatcher.cpp}；
  \item 原生适配器
        \srcpath{src/engine/solver/native/native\_adapters.cpp}（LP 选择器
        \texttt{native\_lp\_selector.cpp}、PDLP \texttt{lp/pdlp\_solver.cpp}、IPM 系
        适配器仅述其注册名，实现转写见本手册 LP/IPM 章节）；
  \item 外部适配器 \srcpath{src/engine/solver/external/adapters.cpp}
        （HiGHS/Ipopt/SCIP/Gurobi/CPLEX）；
  \item 策略辅助件 \srcpath{src/engine/strategy/presolve\_manager.cpp}、
        \srcpath{src/engine/strategy/postsolve\_manager.cpp}、
        \srcpath{src/engine/strategy/warmstart.cpp}；
  \item 输入校验 \srcpath{src/engine/util/problem\_validation.cpp}。
\end{itemize}

\begin{boundarynote}
口径约定：（i）本章只描述``数据如何流动与何时回退''，各求解器内部的算法转写
属于各自章节（MILP B\&C 见 \ref{sec:impl-milpbc} 章）。（ii）与用户手册
\texttt{docs/manual/05-solvers-engines.md} \S 5.4 的两处行文差异在
\ref{subsec:api-docdiff} 集中声明；以代码为准。（iii）匿名命名空间内的自由函数
以 \texttt{文件名:函数名} 锚定（不出现 \texttt{::} 前缀），与仓库锚点约定一致。
\end{boundarynote}

\subsection{门面 \texttt{SolverEngine}：规范化、校验与结果映射}\label{subsec:api-facade}

公共求解入口的事务链固定为四步：
\begin{enumerate}
  \item \textbf{矩阵规范化}：对按值传入的 \texttt{ProblemVariant} 原地补齐零行
        矩阵的列数——\texttt{A}/\texttt{Aeq} 行数为 0 但列数不等于 $n$ 时
        resize 成 $0\times n$；QP 的空 Hessian resize 成 $n\times n$；MILP 作用于
        \texttt{linear\_part}。原因是 Eigen 稀疏矩阵默认构造为 $0\times0$，逐列
        迭代的适配器在 \texttt{cols()<n} 时会越界。
        对应源码为
        \srcpath{src/engine/api/solver.cpp:normalize\_lp\_matrices}、
        \srcpath{src/engine/api/solver.cpp:normalize\_problem}；
  \item \textbf{校验}：对规范化后的模型调用 \texttt{validate}（按 variant 分派到
        各模型重载），任一错误即以 \texttt{std::invalid\_argument} 抛出，
        消息为 \texttt{problem\_class\_name} 加分号连接的全部错误。
        对应源码为 \srcpath{src/engine/api/solver.cpp:throw\_if\_invalid}、
        \srcpath{src/engine/util/problem\_validation.cpp:validate}；
  \item \textbf{分派}：构造 \texttt{SolveContext} 并调用
        \texttt{StrategyDispatcher::solve}（见 \ref{subsec:api-dispatch}）；
  \item \textbf{结果映射}：\texttt{SolveResult}$\to$\texttt{api::Result} 为逐字段
        拷贝/移动，无任何重算；\texttt{x}、三组对偶向量与全部统计字段原样透传。
        对应源码为 \srcpath{src/engine/api/solver.cpp:to\_api\_result}、
        \srcpath{src/engine/api/solver.cpp:SolverEngine::solve\_normalized}。
\end{enumerate}
类型化便捷入口（\texttt{solve\_lp}/\texttt{solve\_milp}/\texttt{solve\_qp}/%
\texttt{solve\_conic} 等）只做对应的矩阵规范化后转入统一 \texttt{solve}。
对应源码为 \srcpath{src/engine/api/solver.cpp:SolverEngine::solve\_lp}、
\srcpath{src/engine/api/solver.cpp:SolverEngine::solve}。

输入校验的实现规则（\texttt{problem\_validation.cpp}）：通用变量元数据要求
界有限且 $l_j\le u_j+10^{-12}$（\texttt{kBoundEps}，硬编码）；binary 变量界超出
$[0,1]$ 只产生 warning。LP 要求 $c$ 非空、$\texttt{A.cols}=n$、
\texttt{row\_lhs} 为空或长度等于 \texttt{A.rows} 且每个有限两端行满足
$\mathrm{lhs}_i\le b_i+10^{-12}$；MILP/MINLP 额外要求整数/二进制指标集在界内、
无重复且两集合不交（排序后相邻判等）；Conic 要求
$\texttt{G.rows}=\texttt{h.size}=\texttt{dims.total()}$ 且锥维数非负。
对应源码为
\srcpath{src/engine/util/problem\_validation.cpp:validate\_variable\_meta}、
\srcpath{src/engine/util/problem\_validation.cpp:validate}。

\subsection{资源契约 \texttt{SolveContext}}\label{subsec:api-context}

\texttt{SolveContext} 在构造时固化一次调用级资源合约：时钟从构造点开始；
\texttt{time\_limit\_sec}$>0$ 时记录绝对 deadline。构造期即拒绝非法输入
（时限非有限或为负、线程数为负均抛 \texttt{std::invalid\_argument}）。
\begin{equation}
  \texttt{remaining\_time\_sec}=
  \max\bigl(0,\ \mathrm{deadline}-\mathrm{now}()\bigr),\qquad
  \texttt{stop\_requested}=
  \texttt{stop\_token.stop\_requested()}\ \lor\ \texttt{deadline\_expired()}.
\end{equation}
设计意图（头文件注释）：绝对 deadline 使 $N$ 次回退共享同一份时限，而不是各
拿一份相对时限。对应源码为
\srcpath{include/mipsolvers/engine/solve\_context.hpp:SolveContext}；
其 R1 条款出处为
\srcpath{docs/archive/general\_solver\_performance\_program\_2026-09-13.md}。

\subsection{适配器接口与注册表}\label{subsec:api-registry}

\texttt{SolverAdapter} 对每个问题类提供两个虚函数层：无 context 的
\texttt{solve\_*} 与带 \texttt{SolveContext} 的重载；基类默认实现全部返回
\texttt{unsupported\_result}（\texttt{success=false}、状态为
``Unsupported problem class for adapter <name>: <CLS>''），带 context 的默认实现
委托给无 context 版本（即不感知 deadline 的适配器仍能编译入注册表）。
对应源码为
\srcpath{src/engine/solver/solver\_adapter.cpp:SolverAdapter::unsupported\_result}、
\srcpath{src/engine/solver/solver\_adapter.cpp:SolverAdapter::solve\_lp}。

\texttt{AdapterRegistry} 是一个保序向量：\texttt{register\_adapter} 只追加
（空指针被静默忽略——显式的宽容路径）；\texttt{adapters\_for(cls)} 按注册序返回
全部 \texttt{supports(cls)} 为真的适配器；\texttt{find\_by\_name} 返回首个同名
适配器。对应源码为
\srcpath{src/engine/solver/adapter\_registry.cpp:AdapterRegistry::register\_adapter}、
\srcpath{src/engine/solver/adapter\_registry.cpp:AdapterRegistry::adapters\_for}、
\srcpath{src/engine/solver/adapter\_registry.cpp:AdapterRegistry::find\_by\_name}。

\paragraph{默认注册顺序。}
\texttt{register\_default\_adapters} 通过 \texttt{register\_if\_missing}（按名
幂等，重名跳过）依次注册：

\begin{center}
\small
\begin{tabular}{ll}
\toprule
顺序 & 适配器（注册名）\\
\midrule
1--12 & \texttt{NativeLinear}, \texttt{NativeNewton}, \texttt{NativeAutoLP},
        \texttt{NativeDualSimplex}, \texttt{NativeIPMLP}, \texttt{NativePDLP},
        \texttt{NativeLCQP}, \texttt{NativeIPM}, \texttt{NativeNLP},
        \texttt{NativeConicIPM}, \texttt{StrictHiGHS}, \texttt{NativeBranchAndCut}\\
13--17 & \texttt{Gurobi}, \texttt{CPLEX}, \texttt{HiGHS}, \texttt{Ipopt},
        \texttt{SCIP}\quad（仅当各自 \texttt{available()} 为真）\\
\bottomrule
\end{tabular}
\end{center}

原生 12 个无条件注册；外部 5 个先构造再探测 \texttt{available()}，不可用则不
入表。该顺序即注册表内序，影响 \texttt{adapters\_for} 的尾部枚举序（见
\ref{subsec:api-dispatch}）。注册名均不带 \texttt{Adapter} 后缀。
对应源码为
\srcpath{src/engine/api/solver.cpp:SolverEngine::register\_default\_adapters}、
\srcpath{src/engine/solver/native/native\_lp\_selector.cpp:NativeDualSimplexLPAdapter}。

\texttt{available()} 的探测语义分两类：编译期内嵌库宏
（\texttt{HACDCPF\_HAVE\_HIGHS\_LIB}/\texttt{HACDCPF\_HAVE\_IPOPT}/%
\texttt{HACDCPF\_HAVE\_SCIP\_LIB}）为真时恒可；否则 HiGHS/Ipopt/SCIP 退化为
``可执行文件路径非空''；Gurobi/CPLEX 无子进程路径，仅在对应编译宏且运行时
环境/许可初始化成功（\texttt{env\_ != nullptr}）时可用。注意
\texttt{GurobiAdapter::supports} 把 \texttt{available()} 并入判定，与其他适配器
``支持关系与可用性分离''的写法不同——属实现现状。
对应源码为
\srcpath{src/engine/solver/external/adapters.cpp:HighsAdapter::available}、
\srcpath{src/engine/solver/external/adapters.cpp:IpoptAdapter::available}、
\srcpath{src/engine/solver/external/adapters.cpp:ScipAdapter::available}、
\srcpath{src/engine/solver/external/adapters.cpp:GurobiAdapter::available}、
\srcpath{src/engine/solver/external/adapters.cpp:CplexAdapter::available}、
\srcpath{src/engine/solver/external/adapters.cpp:GurobiAdapter::supports}。

\subsection{候选排序与回退：\texttt{StrategyDispatcher}}\label{subsec:api-dispatch}

\paragraph{默认优先级表。}
每个问题类的内建候选名序列（\texttt{default\_priority\_for}）：
\begin{center}
\small
\begin{tabular}{ll}
\toprule
问题类 & 默认候选序（名前先）\\
\midrule
LE & \texttt{NativeLinear}\\
NLE & \texttt{NativeNewton}\\
LP & \texttt{NativeAutoLP}, \texttt{NativeIPMLP}, \texttt{NativePDLP},
     \texttt{NativeLCQP}, \texttt{HiGHS}\\
QP & \texttt{NativeLCQP}\\
NLP & \texttt{Ipopt}, \texttt{NativeIPM}, \texttt{NativeNLP}\\
MILP & \texttt{StrictHiGHS}, \texttt{HiGHS}, \texttt{NativeBranchAndCut}\\
MINLP & \texttt{Scip}, \texttt{NativeBranchAndCut}\\
CONIC & \texttt{NativeConicIPM}\\
\bottomrule
\end{tabular}
\end{center}
对应源码为
\srcpath{src/engine/strategy/dispatcher.cpp:default\_priority\_for}。

\paragraph{策略排序。}
\texttt{StrategyPolicy::ExternalFirst} 把名字属于
$\{\texttt{HiGHS},\texttt{Ipopt},\texttt{Scip},\texttt{Gurobi},\texttt{CPLEX}\}$
者排前（rank 0）、\texttt{StrictHiGHS} 或以 \texttt{Native} 前缀开头者次之
（rank 1），其余殿后（rank 2）；\texttt{NativeFirst} 对称交换；\texttt{Auto}
不动序。排序为 \texttt{stable\_sort}。按类覆盖
（\texttt{class\_strategy\_policy}）优先于全局策略。
对应源码为
\srcpath{src/engine/strategy/dispatcher.cpp:apply\_policy\_ordering}、
\srcpath{src/engine/strategy/dispatcher.cpp:effective\_policy}。

\paragraph{候选拼装。}
\texttt{candidate\_adapters} 依次并入（\texttt{seen} 集合去重）：
（1）\texttt{preferred\_solver}；（2）按类偏好
（\texttt{set\_solver\_preference} 登记）；（3）策略排序后的默认表；
（4）注册表中其余支持该类的适配器（注册序）。每一步都要求按名找到且
\texttt{supports(cls)} 为真。
对应源码为
\srcpath{src/engine/strategy/dispatcher.cpp:StrategyDispatcher::candidate\_adapters}、
\srcpath{src/engine/strategy/dispatcher.cpp:StrategyDispatcher::set\_solver\_preference}。

\paragraph{回退事务。}
\texttt{StrategyDispatcher::solve} 对候选依次：每次尝试\textbf{之前}检查绝对
deadline 与 stop token（命中则分别返回 \texttt{"Time limit reached before
fallback"}/\texttt{"Cancelled before fallback"}，\texttt{success=false}）；
随后经 \texttt{solve\_with\_adapter} 按 variant 取载荷并调用对应虚函数（载荷
类型不符返回 \texttt{"Problem payload mismatch ..."}）；结果
\texttt{stats.success=true} 即返回；否则记录为最近失败，且
\texttt{allow\_fallback=false} 时立即返回该失败。全部候选失败时返回
\textbf{最后一次}失败的结果（而非第一个）。候选为空返回
\texttt{"No adapter registered for class ..."}。
对应源码为
\srcpath{src/engine/strategy/dispatcher.cpp:StrategyDispatcher::solve}、
\srcpath{src/engine/strategy/dispatcher.cpp:solve\_with\_adapter}。

\begin{boundarynote}
回退的两个实现现状，引用时不得读成更强合约：
（i）回退只在 \texttt{success=false} 上触发；\texttt{success=true} 但
\texttt{status} 表明次优/截断的结果不会触发下一个候选。（ii）注册表尾部枚举
会把\textbf{所有}注册且支持该类的适配器追加为末位候选，包括 Gurobi/CPLEX
（只要 \texttt{available()} 使它们入了注册表）；它们只是不在
\texttt{default\_priority\_for} 的内建优先级表内。用户手册 \S 5.4.2
``Gurobi 与 CPLEX 不参与普通自动回退''一句与该尾部路径不符，见
\ref{subsec:api-docdiff}。
对应源码为
\srcpath{src/engine/strategy/dispatcher.cpp:StrategyDispatcher::solve}、
\srcpath{src/engine/strategy/dispatcher.cpp:StrategyDispatcher::candidate\_adapters}。
\end{boundarynote}

\subsection{原生适配器的适配层}\label{subsec:api-native}

\paragraph{LE/NLE/NLP 轻量适配器。}
\texttt{NativeLinear}（LE）直接稀疏 LU：分析-分解-求解后以原系统残差
$\|Ax-b\|_\infty$ 复核，状态 \texttt{"Optimal"}。
\texttt{NativeNewton}（NLE）是阻尼 Newton：正则化阶梯
$\lambda: \texttt{regularization0}\to 10^{-10}\to \times 10$（至多 4 次尝试，
硬编码）解 $(J+\lambda I)\,\Delta x=-f$，回溯线搜索要求
$\|f(x+\alpha\Delta x)\|_\infty\le\|f(x)\|_\infty$，步长回退因子
\texttt{step\_backoff}。
\texttt{NativeNLP} 是固定罚参数的二次罚 Newton（非 IPM）：罚梯度
$\nabla\phi=\nabla f+\rho J_g^{T}g+\rho J_h^{T}h^{+}$，罚 Hessian
$H=\nabla^2 f+\rho J_g^{T}J_g+\rho (J_h^{\mathrm{act}})^{T}J_h^{\mathrm{act}}$
再加对角正则 \texttt{regularization0}；收敛要求
$\|\nabla\phi\|_\infty$、$\|g\|_\infty$、$\|h^{+}\|_\infty$ 同时
$\le\texttt{tol\_grad}$，或 $\|\Delta x\|_\infty\le\texttt{tol\_step}$；
接受判据为罚函数不增。头文件级注释明确它不是 \texttt{NativeIPM}。
对应源码为
\srcpath{src/engine/solver/native/native\_adapters.cpp:NativeLinearAdapter::solve\_le}、
\srcpath{src/engine/solver/native/native\_adapters.cpp:NativeNewtonAdapter::solve\_nle}、
\srcpath{src/engine/solver/native/native\_adapters.cpp:NativeNLPAdapter::solve\_nlp}。

\paragraph{B\&C 适配器与 LP 快路径。}
\texttt{NativeBranchAndCut} 支持 MILP 与 MINLP。MILP 入口先走 LP 快路径：当
\texttt{binary\_idx} 与 \texttt{integer\_idx} 皆空时直接以单纯形求解
（迭代上限 $\max(\texttt{max\_lp\_iter},20000)$，容差
$\max(10^{-8},0.1\,\texttt{lp\_tol})$），并从最优基提取约束对偶：标准形空间
\begin{equation}
  y = B^{-T}c_B,\qquad
  d_i = s_{\mathrm{sense}}\cdot\mathrm{row\_sign}_i\cdot y_i,\quad
  s_{\mathrm{sense}}=\begin{cases}-1,&\text{Minimize}\\+1,&\text{Maximize}\end{cases}
\end{equation}
（标准形内部为最大化，故最小化要再翻一次符号）；优先用稠密基逆，缺省时用
稀疏 BTRAN，再缺省则对偶置零——显式回退链。布局统一为
$[\,\text{不等式行对偶}\;|\;\text{等式行对偶}\,]$。
带整数的模型转入 \texttt{solve\_milp\_bc}（见 \ref{sec:impl-milpbc} 章）。
context 版入口把绝对 deadline 的剩余时间与线程预算写进 \texttt{BCOptions}
副本（\texttt{time\_limit\_sec}/\texttt{num\_threads}）。
对应源码为
\srcpath{src/engine/solver/native/native\_adapters.cpp:NativeBranchAndCutAdapter::solve\_milp}、
\srcpath{src/engine/solver/native/milp/bc/api.cpp:solve\_milp\_bc}。

\paragraph{StrictHiGHS 选项改造。}
\texttt{StrictHiGHS} 是同名 B\&C 内核的选项变体（委托给
\texttt{NativeBranchAndCutAdapter} 执行）：\texttt{make\_strict\_highs\_production\_%
options} 固定 \texttt{lp\_kernel\_backend=HiGHS}、\texttt{strict\_highs\_mip\_%
contract=true}、\texttt{auto\_highs\_root\_pipeline=true}，开启
\texttt{highs\_mip\_run\_zi\_round}/\texttt{highs\_mip\_run\_shifting}，强制
presolve on 并把 \texttt{highs\_presolve\_substitution\_maxfillin} 提到 30，
\texttt{highs\_mip\_lp\_age\_limit=30}、\texttt{highs\_mip\_detect\_symmetry=true}。
\texttt{make\_strict\_highs\_problem\_options} 再按模型规模决定是否把根 LP
切成 IPM：当列数/行数进入
$[\texttt{min\_cols},\texttt{max\_cols})\times[\texttt{min\_rows},\texttt{max\_rows})$
窗口、全局开关
\texttt{highs\_strict\_auto\_ipm\_root\_for\_large\_models} 开启、且时限满足
最小预算
\begin{equation}
  T_{\min}=\max\Bigl(\texttt{min\_time\_sec},\
  \frac{n_{\mathrm{cols}}\cdot\texttt{secs\_per\_kcol}}{1000}\Bigr)
\end{equation}
（默认 \texttt{secs\_per\_kcol}$=2.0$ 时 6360 列对应 $30$\,s、26400 列对应
$52.8$\,s——代码注释中的两个参考点）时，置 \texttt{highs\_mip\_lp\_solver="ipm"}
并把 crossover 默认开。显式调用者设置不被该规模门覆盖。
对应源码为
\srcpath{src/engine/solver/native/native\_adapters.cpp:make\_strict\_highs\_production\_options}、
\srcpath{src/engine/solver/native/native\_adapters.cpp:make\_strict\_highs\_problem\_options}、
\srcpath{src/engine/solver/native/native\_adapters.cpp:StrictHighsBranchAndCutAdapter::solve\_milp}。

\subsection{外部适配器的模型映射与结果解析}\label{subsec:api-external}

五个外部适配器共享同一骨架：校验 $\to$ \texttt{available()} 门 $\to$
等价模型映射 $\to$ 内嵌库调用（编译期启用时优先）或子进程回退 $\to$
结果解析 $\to$ 统一 \texttt{SolveResult}。不可用/失败时返回
\texttt{unavailable\_result}（\texttt{success=false}、状态以
\texttt{"Unavailable:"} 等前缀说明原因）。
对应源码为
\srcpath{src/engine/solver/external/adapters.cpp:unavailable\_result}。

\paragraph{HiGHS。}
支持 LP/MILP。MILP 映射：把 \texttt{integer\_idx}/\texttt{binary\_idx} 提升为
\texttt{linear\_part} 上的变量类型，binary 界夹取到 $[0,1]$；
\texttt{initial\_solution} 尺寸匹配时作为 MIP start 传入。内嵌路径经
\texttt{solve\_lp\_with\_embedded\_highs}，调用前先
\texttt{Highs::resetGlobalScheduler(true)} 以消除先前 MILP 求解遗留的全局线程
调度器配置（注释说明原因）。子进程回退：写临时 MPS（列名 \texttt{X1..Xn}，
1-indexed）+ 选项文件（LP 写 \texttt{time\_limit}/\texttt{threads}/%
\texttt{random\_seed}；MILP 另写死 \texttt{mip\_rel\_gap = 1e-4}——硬编码），
调 \texttt{--model\_file/--solution\_file/--options\_file}，解析 .sol 的
\texttt{Model status} 与日志；\texttt{success}$\iff$状态为
\texttt{Optimal} 或 \texttt{Feasible}，解值按 \texttt{X<i+1>} 名回映射，缺失列
补 0。临时文件求解后删除。
对应源码为
\srcpath{src/engine/solver/external/adapters.cpp:HighsAdapter::solve\_lp\_impl}、
\srcpath{src/engine/solver/external/adapters.cpp:HighsAdapter::solve\_milp\_impl}。
另有确定性定价入口 \texttt{solve\_pricing\_lp}：只走内嵌库，参数非法直接抛
\texttt{std::invalid\_argument}。
对应源码为
\srcpath{src/engine/solver/external/adapters.cpp:HighsAdapter::solve\_pricing\_lp}。

\paragraph{SCIP。}
支持 MILP/MINLP。MILP 走 MPS 交换（同样的整数提升与 binary 夹取）；内嵌路径
\texttt{SCIPreadProb} 读入后求解，固定参数 \texttt{limits/gap=1e-3}、
\texttt{limits/time=300.0}（硬编码，注释说明与原生 UC 容差对齐）、
\texttt{display/verblevel=0}；有解即 \texttt{success=true}，
\texttt{status} 依 \texttt{SCIPgetStatus} 区分 \texttt{"Solved"}/%
\texttt{"Feasible (limit)"}/\texttt{"Infeasible"}，gap 取
\texttt{SCIPgetGap}；解值经 \texttt{SCIPgetOrigVars}（原始空间，而非 presolve
后的变换变量——注释说明后者名字对不上 MPS 导出）按名回映射。子进程回退用
\texttt{-c "set limits gap 0.001" -c read -c optimize -c "write solution"} 脚本。
MINLP：无符号目标时走 NLP 松弛 + 舍入回退（整数取整、binary 按 $0.5$ 阈值、
再夹回界内，$\|g\|_\infty$/$h^+$ 复核 $\le 10^{-4}$ 才算成功——硬编码阈值，
且这只是启发式，不是全局证明）；有符号目标时写 PIP 文件交给 SCIP。
对应源码为
\srcpath{src/engine/solver/external/adapters.cpp:ScipAdapter::solve\_milp}、
\srcpath{src/engine/solver/external/adapters.cpp:ScipAdapter::solve\_minlp}。

\paragraph{Ipopt。}
仅 NLP，且\textbf{只有内嵌库路径真正求解}：无内嵌时即使有可执行文件也返回
\texttt{"Unavailable: external Ipopt executable adapter is disabled by default"}
（显式禁用路径）。求解前：\texttt{x0} 尺寸不符时重构为界内中点/端点的默认
点；缺 \texttt{f}/\texttt{grad} 或有 \texttt{g} 无 \texttt{jac\_g}、有
\texttt{h} 无 \texttt{jac\_h} 直接 Unavailable；primal-dual warm start 需维度
匹配、全部有限且界对偶非负，否则报错。选项映射：
\texttt{hessian\_approximation} 按有无 Lagrangian Hessian 回调取
\texttt{exact}/\texttt{limited-memory}；\texttt{bound\_relax\_factor=0.0}
（注释：默认界松弛会在大乘子下破坏原模型互补性）；warm start 统一推送量
默认 $10^{-8}$。容差经 \texttt{positive\_finite\_or} 兜底（非有限或非正回退
默认），并强制 Ipopt 的不变量
$\texttt{acceptable\_tol}\ge\texttt{tol}$（对偶/约束/互补四组同法）。
结果由 \texttt{CallbackTNLP::build\_result} 从 \texttt{OptimizeTNLP} 状态装配。
对应源码为
\srcpath{src/engine/solver/external/adapters.cpp:IpoptAdapter::solve\_nlp}。

\paragraph{CPLEX。}
仅 MILP，Callable Library 内嵌（无子进程路径）。构造期打开环境并固定
\texttt{ScreenOutput=off}、\texttt{TimeLimit}、\texttt{MIPGap}、\texttt{Threads}、
\texttt{RandomSeed} 五个参数，任一失败即记 \texttt{initialization\_error\_} 且
环境置空。模型映射的两处精确转写：（i）双端行展开——每个不等式行按其有限端
拆成至多两行（上端 \texttt{'L'}、下端 \texttt{'G'}），等式行为 \texttt{'E'}，
列主序打包时同一列元素向所属各展开行各写一份；行数或非零元数超
\texttt{int} 上限时失败返回。（ii）列类型 \texttt{'I'}/\texttt{'B'}，binary
界夹取到 $[0,1]$，非有限界映射为 $\pm$\texttt{CPX\_INFBOUND}。结果解析：
\texttt{success} 以是否取到解为准（\texttt{CPXgetx} 失败会把 success 翻回
false），状态字符串原样取 \texttt{CPXgetstatstring}，gap 取
\texttt{CPXgetmiprelgap}；\texttt{cplex\_status\_proven}/\texttt{optimal}/%
\texttt{timed\_out} 三个谓词把 \texttt{CPXMIP\_*} 状态码分类进
\texttt{thread\_local} 诊断 \texttt{cplex\_solve\_info}。
对应源码为
\srcpath{src/engine/solver/external/adapters.cpp:CplexAdapter::solve\_milp}、
\srcpath{src/engine/solver/external/adapters.cpp:CplexAdapter::CplexAdapter}；
双端行展开的正确性论证见
\srcpath{docs/archive/cplex\_callable\_library\_integration\_2026-09-12.md}。

\paragraph{Gurobi。}
支持 LP/QP/MILP（\texttt{supports} 依赖 \texttt{available()}，见
\ref{subsec:api-registry}）。LP 映射：最大化经目标取负统一为
\texttt{GRB\_MINIMIZE}（$\texttt{obj\_sign}=-1$）；不等式行按下标
$|a_{ij}|>10^{-15}$ 过滤后逐行 \texttt{GRBaddconstr}（双端行拆成
\texttt{GRB\_LESS\_EQUAL} 与 \texttt{GRB\_GREATER\_EQUAL} 两条，$10^{-15}$
为硬编码）；行按\textbf{稠密提取}构造（代码保留了稀疏列主序迭代的空循环
与注释，说明为何不用它——已知的历史写法残留，不影响结果）。两个特化入口：
\texttt{solve\_pricing\_lp} 按规模固定方法（$n\ge10^{6}$ 用
\texttt{method=2}/\texttt{threads=8}/\texttt{crossover=0}，否则
\texttt{method=1}/单线程，硬编码阈值），并把 \texttt{OptimalityTol=1e-8} 写进
环境；\texttt{solve\_relaxation\_lp} 是不精确 barrier 入口
（\texttt{BarConvTol} 取参，合法域 $[10^{-8},10^{-2}]$，之外抛异常；注释明确
禁止用它认证价格或原始可行性）。
对应源码为
\srcpath{src/engine/solver/external/adapters.cpp:GurobiAdapter::solve\_lp}、
\srcpath{src/engine/solver/external/adapters.cpp:GurobiAdapter::solve\_pricing\_lp}、
\srcpath{src/engine/solver/external/adapters.cpp:GurobiAdapter::solve\_relaxation\_lp}。

\subsection{策略辅助件：presolve/postsolve/warmstart 管理器}\label{subsec:api-strategy-aux}

\begin{boundarynote}
接线现状（必须在引用处保持可见）：\texttt{PresolveManager}/%
\texttt{PostsolveManager}/\texttt{WarmStartManager} 是\textbf{独立工具类}，
当前未被 \texttt{StrategyDispatcher::solve} 调用——分派器直接调用适配器，主流程
的 presolve 发生在各求解器内部（MILP 的 presolve 见 \ref{sec:impl-milpbc} 章）。
\texttt{PresolveManager::should\_presolve} 恒返回 \texttt{false}，
\texttt{PresolveManager::presolve} 退化为恒等映射 \texttt{native\_presolve}
（只统计行列数，不做约简）；\texttt{PresolveManager::scale} 对
\texttt{"none"}/\texttt{"identity"} 之外的 scaling 直接抛
\texttt{std::logic\_error}（无通用缩放后端）。这些类的完整实现仅在测试中被
消费。对应源码为
\srcpath{src/engine/strategy/presolve\_manager.cpp:PresolveManager::should\_presolve}、
\srcpath{src/engine/strategy/presolve\_manager.cpp:PresolveManager::native\_presolve}、
\srcpath{src/engine/strategy/presolve\_manager.cpp:PresolveManager::scale}、
\srcpath{src/engine/strategy/dispatcher.cpp:StrategyDispatcher::solve}。
\end{boundarynote}

\texttt{PostsolveManager} 实现的三步还原（给定 \texttt{PresolveMapping}）：
\begin{enumerate}
  \item 固定变量还原：按 \texttt{col\_mapping} 把约简解散布回原列，再写入
        \texttt{fixed\_variables}；映射重复、固定列与活跃列重叠、或还原后存在
        非有限分量（未覆盖列）均抛 \texttt{std::invalid\_argument}；
  \item 平移还原：$x_{\mathrm{orig}} \mathrel{+}= \mathrm{shift}$（逐条
        \texttt{variable\_shifts}）；
  \item 反缩放：$x_j \mathrel{*}= s^{\mathrm{col}}_j$、
        $y_i \mathrel{/}= s^{\mathrm{row}}_i$，尺度必须有限且为正。
\end{enumerate}
一致性校验 \texttt{validate} 比较两结果的 \texttt{success}、解维数与有限性，
目标差容差为 $\texttt{feasibility\_tol}\cdot\max\{1,|f_1|,|f_2|\}$。
对应源码为
\srcpath{src/engine/strategy/postsolve\_manager.cpp:PostsolveManager::postsolve}、
\srcpath{src/engine/strategy/postsolve\_manager.cpp:PostsolveManager::restore\_fixed\_variables}、
\srcpath{src/engine/strategy/postsolve\_manager.cpp:PostsolveManager::unscale}、
\srcpath{src/engine/strategy/postsolve\_manager.cpp:PostsolveManager::validate}。

\texttt{WarmStartManager} 的规则：\texttt{is\_compatible} 逐槽检查 primal/约束
对偶/界对偶向量的维数；\texttt{extract\_from\_result} 只在尺寸匹配时搬运各槽；
\texttt{merge} 取各槽首个非空候选，\texttt{reuse\_root\_factorization} 取或；
\texttt{scale} 的坐标规则为 $x'=x/s_x$、$y'=y\odot s_c$、界对偶 $\odot\, s_x$；
\texttt{validate} 拒绝非有限 hint。
对应源码为
\srcpath{src/engine/strategy/warmstart.cpp:WarmStartManager::is\_compatible}、
\srcpath{src/engine/strategy/warmstart.cpp:WarmStartManager::extract\_from\_result}、
\srcpath{src/engine/strategy/warmstart.cpp:WarmStartManager::merge}、
\srcpath{src/engine/strategy/warmstart.cpp:WarmStartManager::scale}。

\subsection{结果契约}\label{subsec:api-result}

内部 \texttt{SolveResult}/\texttt{SolveStats} 与公共 \texttt{api::Result}/%
\texttt{api::Stats} 的字段一一对应（\texttt{to\_api\_result} 不透传的内部字段
仅有 \texttt{primal\_ray} 与稀疏代数/原生单纯形遥测等实现侧计数——未被公共
API 消费的字段在此声明）。契约要点：
\begin{itemize}
  \item \texttt{success} 只表示``该适配器成功终止''：B\&C 在有时限可行
        incumbent 时亦为 true（状态 \texttt{kFeasibleIncumbentFound} 系），
        与``已证最优''无关；验收必须同时看 \texttt{status}、可行度、
        \texttt{mip\_gap}；
  \item \texttt{status} 为人类可读字符串，无枚举约束；同义词集合以
        \texttt{bc\_status} 常量与各适配器字面量为准；
  \item LP 对偶布局跨适配器统一为 $[\,\text{不等式行}\;|\;\text{等式行}\,]$，
        界对偶 \texttt{box\_dual\_lb}/\texttt{box\_dual\_ub} 每列一个；
        MILP 等无证书路径置空；
  \item 不可行证书 \texttt{farkas\_ray}/\texttt{farkas\_ray\_eq}（配合
        \texttt{has\_farkas\_certificate}）：对
        $\min c^Tx\ \mathrm{s.t.}\ Ax\le b,\ A_{\mathrm{eq}}x=b_{\mathrm{eq}},\
        l\le x\le u$ 的不可行 LP，满足 $y\ge0$、$A^Ty+A_{\mathrm{eq}}^Ty_{eq}=0$、
        $b^Ty+b_{\mathrm{eq}}^Ty_{eq}<0$；
  \item \texttt{certified\_dual\_bound} 仅由原生对偶单纯形内核在中断时发布；
        中断结果的 \texttt{objective} 仍是未认证的工作值（头文件注释明确）。
\end{itemize}
对应源码为
\srcpath{include/mipsolvers/engine/solver/solver\_adapter.hpp}、
\srcpath{include/mipsolvers/engine/api/result.hpp:Stats}、
\srcpath{src/engine/api/solver.cpp:to\_api\_result}。

\subsection{与用户手册行文的两处差异}\label{subsec:api-docdiff}

以下差异按``以代码为准''声明，供手册维护侧对照：

\begin{enumerate}
  \item \texttt{docs/manual/05-solvers-engines.md} \S 5.4.2 把调度流程写成
        ``预处理 $\to$ 适配器调用 $\to$ 回退 $\to$ 后处理''六步；当前
        \texttt{StrategyDispatcher::solve} 不含门面级 presolve/postsolve 阶段
        （\texttt{PresolveManager} 未被接线，见
        \ref{subsec:api-strategy-aux}），实际为
        ``候选构造 $\to$ 逐个适配器调用（带绝对 deadline 门）$\to$ 成功即返回
        或回退''。模型规范化与校验在 \texttt{SolverEngine::solve} 一侧完成。
        对应源码为
        \srcpath{src/engine/strategy/dispatcher.cpp:StrategyDispatcher::solve}、
        \srcpath{src/engine/api/solver.cpp:SolverEngine::solve}；
  \item 同节称``Gurobi 与 CPLEX 不参与普通自动回退''；它们确实不在
        \texttt{default\_priority\_for} 内建表内，但一旦注册（可用性门通过），
        \texttt{candidate\_adapters} 的注册表尾部枚举会把它们追加为末位回退
        候选。对应源码为
        \srcpath{src/engine/strategy/dispatcher.cpp:StrategyDispatcher::candidate\_adapters}、
        \srcpath{src/engine/strategy/dispatcher.cpp:default\_priority\_for}。
\end{enumerate}


% numerical_validation.tex —— engine 手册：数值验证与基准证据
% 本文件为证据转写章：只复述仓库内真实存在的测试目标、基线文件、
% 报告目录与归档记录；每个数字注明出处（文档/目录、日期、平台、构建）
% 与验证层级。遵循 docs/modules/README.md「数值证据规范」与仓库
% AGENTS.md 的诚实口径：未执行的交叉验证不得写成已完成。

\section{数值验证与基准证据}

\begin{boundarynote}
本章不给出任何新的测量结果，也不对所列数字做超出原始记录的解读。每个
性能/精度结论都注明四项出处要素：证据文件（手册章节、归档推导记录或
\texttt{reports/} 原始数据）、测量日期、平台与构建配置、验证层级。
历史快照只描述其对应提交与机器的状态，不能直接代替当前提交的验证；
不同日期、不同机器、不同计时协议得到的绝对时间不可混算
（口径与 \texttt{docs/manual/09-testing-benchmarks.md} §9.5 一致）。
\end{boundarynote}

\subsection{验证层级与证据口径}

仓库的数值证据按验证层级由低到高分为四档：

\begin{enumerate}
  \item \textbf{CTest 单元/集成测试}（\texttt{unit}/\texttt{integration}
        标签）：正确性门禁，断言级验证；其 wall time 不构成性能证据。
  \item \textbf{冒烟基准}（\texttt{benchmark} 标签）：小样本、少重复的
        结果一致性检查，只证明路径可用，不支持性能排名。
  \item \textbf{数据集回归与协议化基准}：NETLIB 90 案例基准与
        MIPLIB 2017 协议，固定数据、固定时限、带原模型审计，
        结果以 \texttt{reports/} 下的机器可读 JSON/CSV 落盘。
  \item \textbf{独立内核微基准}：\texttt{benchmark/} 下的
        \texttt{native\_dual\_*} 等独立 harness，用于在改动进入生产路径
        之前先证明逐位一致性与字节/指令数收益。
\end{enumerate}

三条判读纪律贯穿全章（出处：\texttt{docs/manual/09-testing-benchmarks.md}
§9.5）：求解器的 \texttt{success}/\texttt{Optimal} 状态字符串不等于解正确，
正确性以原模型审计为准；性能结论只能来自 Release 构建、固定数据、固定线程
与固定重复次数，普通单元测试的 wall time 不能当作求解器性能基准；配对几何
均值、全集几何均值与全量总时间回答三个不同问题，不能互相替代。

\subsection{单元与集成测试矩阵}
\label{sec:engine-nv-testmatrix}

测试目标由根 \texttt{CMakeLists.txt} 中的宏
\texttt{mipsolvers\_add\_test} 统一注册：源文件存在于 \texttt{tests/} 时
生成同名可执行文件（输出目录为 \texttt{tests/}），链接 \texttt{mipsolvers}
与 Catch2，并以 \texttt{TIER} 参数写入 CTest 标签。SCUC 相关目标
（\texttt{test\_scuc\_module}、\texttt{test\_market\_simulation}）只在
\texttt{MIPSOLVERS\_BUILD\_SCUC=ON} 时注册，不属于本手册范围；
\texttt{regression\_nlp\_hs071} 为条件注册的集成测试（调用
\srcpath{benchmark/nlp\_open\_benchmark.cpp:main}，参数 \texttt{1 1e-7}，
环境固定 \texttt{OMP\_NUM\_THREADS=1;MKL\_NUM\_THREADS=1}）。
实际注册数以\texttt{ctest -N} 在对应构建上的列举为准；下表为当前
\texttt{CMakeLists.txt} 的注册快照（已逐一核对源文件存在性）。

\begin{footnotesize}
\begin{longtable}{@{}llp{6.6cm}@{}}
\toprule
CTest 目标 & 标签 & 覆盖面（源文件见 \texttt{tests/同名.cpp}）\\
\midrule
\endhead
\bottomrule
\endfoot
\texttt{test\_engine\_api} & unit &
  引擎公共 API、线性代数后端注册与行为；含 \texttt{EigenSparseLUSolver}
  接受空 $0\times0$ 系统不崩溃的 Windows 回归（59 个 \texttt{TEST\_CASE}）\\
\texttt{test\_adapter\_registry} & unit &
  外部求解器适配器注册表（HiGHS/SCIP/Ipopt 等）的查找与能力声明（7 个用例）\\
\texttt{test\_problem\_validation} & unit &
  模型校验入口：维数、系数、边界等结构化错误的拒绝路径（9 个用例）\\
\texttt{test\_presolve} & unit & 通用预求解变换与还原（19 个用例）\\
\texttt{test\_lp\_presolve} & unit & LP 专用预求解（35 个用例）\\
\texttt{test\_lp\_solver} & unit & LP 求解器门面与选项路由（6 个用例）\\
\texttt{test\_dual\_simplex} & unit &
  原生对偶单纯形内核（链接 \texttt{highs::highs} 以共用数据结构；
  81 个用例，为引擎侧最大单元套件）\\
\texttt{test\_ipm\_solver} & unit &
  LP/QP/NLP 原始-对偶内点法：收敛、终止门、正则化路径（54 个用例）\\
\texttt{test\_numerical\_stability} & unit &
  病态缩放与近奇异模型的数值稳定性（27 个用例）\\
\texttt{test\_conic\_ipm} & unit &
  SOCP/SDP 锥内点法；\texttt{MIPSOLVERS\_USE\_OPENMP=1} 下额外编译
  并行确定性区段（17 个用例）\\
\texttt{test\_milp\_solver} & unit &
  原生 MILP 求解器；含固定种子 \texttt{--rng-seed 1} 的可复现性回归
  （55 个用例）\\
\texttt{test\_branch\_and\_cut} & integration &
  分支定切全流程集成：根节点流水线、割平面、树搜索（35 个用例）\\
\texttt{test\_netlib\_regression} & integration &
  NETLIB 子集回归：\srcpath{tests/test\_netlib\_regression.cpp:read\_mps\_as\_lp}
  读取 \texttt{tests/data/netlib/} 下 afiro、adlittle、share2b、stocfor1、
  kb2 五个有公布最优目标的案例，经注册的 LP 求解器与原生内核求解并比对
  公布目标（3 个用例）\\
\texttt{test\_ipopt\_parameter\_stability} & integration &
  Ipopt 适配器参数稳定性（8 个用例）\\
\texttt{milp\_benchmark} & benchmark &
  \srcpath{benchmark/milp\_benchmark\_runner.cpp:main} 的
  \texttt{--smoke --no-highs-adapter} 确定性 6 母线冒烟，JSON 写入构建树\\
\texttt{milp\_benchmark\_regression} & benchmark &
  \srcpath{benchmark/check\_perf\_regression.py} 将上一目标输出的 JSON
  与已签入基线 \srcpath{benchmark/milp\_benchmark\_baseline.json} 按容差
  0.25 比对；\texttt{DEPENDS milp\_benchmark}\\
\texttt{native\_kernel\_comparison} & benchmark &
  \srcpath{benchmark/native\_kernel\_comparison.cpp:main} 的
  \texttt{--smoke --check --time-limit 60}：纯原生数值栈对 HiGHS 的
  目标一致性冒烟（详见 \ref{sec:engine-nv-microbench} 节）\\
\end{longtable}
\end{footnotesize}

\srcpath{tests/baselines/} 提供 JSON 基线加载工具
\srcpath{tests/baselines/load\_baselines.hpp:baselines::load} 与
\texttt{solver\_baselines.json}；当前由 CTest 实际消费的已签入回归基线是
\texttt{benchmark/milp\_benchmark\_baseline.json}（见上表）。

历史快照（均为 Windows Release 构建上的全量 CTest，出处：
\texttt{docs/manual/09-testing-benchmarks.md}）：

\begin{itemize}
  \item 2026-08-18（\texttt{build/windows-msvc-release}，Ninja
        Multi-Config/MSVC）：18/18 通过（unit 13 + integration 3 +
        benchmark 2）；历史对照 2026-08-06 为 16/18，当时的两个失败项
        （\texttt{test\_milp\_solver} 期望与数据不一致、
        \texttt{test\_ipopt\_parameter\_stability} 逐变量 margin 未满足）
        已不复现（§9.2.1）。
  \item 2026-09-11 离线复评（修复前提交 \texttt{75c6e192}，Windows 11 /
        i9-12900H / MSVC 19.44 Release，本地 sequential oneMKL，SCUC 与
        Python 未启用）：19/19 通过，14 unit + 3 integration + 2 benchmark
        （§9.6）。
  \item 2026-09-11 修复分支 \texttt{9652ddb9}（INTEL/4 线程层、串行运行、
        计时关闭）：20/20 通过（14 unit、4 integration、2 benchmark），
        305.41 秒（§9.7.1）；合并提交 \texttt{1d31f0eb} 的验收为 20/20
        通过、240.85 秒（§9.8）。这两个 wall time 与编译/其他负载并发，
        不能用于求解器速度排名（§9.9 同口径）。
\end{itemize}

各快照的目标数不同，原因是配置选项（SCUC、本地 Ipopt、Python 是否启用）
改变注册集合；引用时须连同日期与配置一起引用。

\subsection{NETLIB 90 案例基准}
\label{sec:engine-nv-netlib}

\subsubsection{数据来源与基准契约}

数据为 90 个解压后的标准 NETLIB 线性规划 MPS 算例，镜像自
\texttt{coin-or-tools/Data-Netlib}（提交 \texttt{f1cc4230}），由
\srcpath{tools/fetch\_netlib.py} 抓取；\texttt{tests/data/mps\_manifest.csv}
保存每个文件的 SHA-256 与参考目标（参考目标取 NETLIB 官方 readme，修订表中
有 CPLEX 双精度值的案例优先取 CPLEX 值）。出处与许可说明见
\srcpath{tests/data/NETLIB\_PROVENANCE.md}；数据目录默认被
\texttt{.gitignore} 排除，仅作本机测试数据。

基准程序为 \srcpath{benchmark/netlib\_solver\_benchmark.cpp:main}（构建产物
在 \texttt{tests/Release/} 下）。验收对象是\textbf{完全绕过 HiGHS
presolve 的 Native IPM direct 路径}：预处理、IPM 迭代、Newton 系统选择、
停止判定与解后审计全由原生路径完成；CHOLMOD/PARDISO 只是本地链接的数值
线性代数内核。计时只含求解调用，不含 MPS 解析。2026-08-06 全量复验的
契约（出处：\texttt{docs/manual/09-testing-benchmarks.md} §9.3.1）：
Windows 11 Pro 64-bit（10.0.26200）、Intel Core Ultra 9 285H（16 核 /
32 GiB）、MSVC 19.44、SuiteSparse/CHOLMOD + MKL PARDISO 后端、求解器时限
15 s、每案例独立子进程隔离（进程硬时限 20 s）、迭代上限 100{,}000。

正确性不依赖求解器状态字符串：基准在原始未 presolve 的 LP 上重算目标与
可行性（\srcpath{benchmark/netlib\_solver\_benchmark.cpp:audit}），同时满足
\texttt{success=true}、目标相对误差 $\le 10^{-5}$、归一化原始可行性违反
$\le 10^{-7}$ 才记为 \texttt{accurate}；阈值对原生与对照求解器完全相同
（§9.3.2）。

\subsubsection{主要结论与已知例外}

2026-08-06 全量复验（出处：§9.3.3）：Native IPM direct 为
\textbf{90/90 success、90/90 accurate}；对照 HiGHS IPM 为 89/90 success、
88/90 accurate。90 个原生解中最大目标相对误差 $1.250\times10^{-6}$、最大
归一化原始违反 $8.689\times10^{-8}$，均在发布阈值内（§9.3.2）。典型性能
以配对几何均值计：在 88 个双方都 accurate 的同名案例上 Native 11.7404 ms
对 HiGHS 15.2844 ms，配对加速 1.302x；全集几何均值 12.5773 ms 对
15.6772 ms。但\textbf{全量总时间}原生为 10{,}129.4127 ms，慢于 HiGHS 的
5{,}300.4629 ms——长尾案例 dfl001、greenbea、maros-r7、pilot87 决定总时间
差距（§9.3.4）。结论只能表述为「全部通过、配对典型性能更快且结果更准」，
不能宣称全量总时间或尾延迟领先。同批数据上 SCIP direct MPS reader 的单进程
对照为 90/90 success、89/90 accurate（唯一未过审计的是 forplan，其 MPS
reader 对 BOUNDS 段报警并返回目标 0）；该轮与隔离 HiGHS 报告计时协议不同
（单进程对独立子进程），两个表的绝对时间不能混合计算。

HiGHS IPM 未通过原模型审计的两个案例（§9.3.4，说明 \texttt{Optimal} 状态
不等于解正确）：ganges 目标相对误差 $1.801\times10^{-2}$；greenbea 返回
\texttt{Unknown} 且归一化原始违反 $1.432\times10^{-2}$。

2026-08-18 的 H1--H7 性能工作（推导记录：
\texttt{docs/archive/lp\_tail\_elimination\_2026-08-18.md}；结论转述自
§9.2.1）修复了 Windows 构建中 CHOLMOD 因缺 BLAS 降级为 simplicial 的缺陷
（H1，回落到静态 MKL，supernodal 恢复），并完成 INTEL 线程层升级（H4）：
dfl001 分解段 6.18 s 降至 1.64 s（16 线程 3.8x）；h7 轮最终口径为 Native
IPM direct 90/90 success + 90/90 accurate，同日 ctest 18/18。固定线程数下
迭代数逐位一致、目标值 run-to-run 相对差约 $4\times10^{-13}$；
\texttt{MKL\_CBWR=AUTO/AVX2} 会使 PARDISO 求解段劣化 4--5 倍，明确不得
启用。

2026-09-11 离线复评（出处：§9.6 与
\texttt{docs/archive/windows\_offline\_evaluation\_2026-09-11.md}；平台为
Windows 11 / i9-12900H / MSVC 19.44 Release / 本地 sequential oneMKL，
原始数据在 \texttt{reports/windows\_eval\_20260911/}）：NETLIB 90 案例
$\times$ 3 次重复下，原生 IPM direct 与 HiGHS simplex 均 270/270 准确；
关闭 crossover 的 HiGHS IPM 为 264/270（ganges、greenbea）。已知例外：
Native-Auto 批次为 264/270——\textbf{greenbea 超时、tuff 返回不准确的
Optimal}，且 tuff 在原生对偶单纯形加 HiGHS presolve 路径上复现。该记录
明确要求：不能仅凭 CTest 通过或 \texttt{success} 字段发布自动求解结果。

LP 分解计时与 INTEL 线程稳定性协议（出处：§9.7、
\texttt{docs/archive/windows\_lp\_stability\_2026-09-11.md}；脚本
\srcpath{tools/windows\_lp\_stability.py:main} 按 overhead/stability/gates
三阶段执行，固定 \texttt{OMP\_NUM\_THREADS=1}、\texttt{MKL\_DYNAMIC=FALSE}
并核对 \texttt{mkl\_max\_threads}）：对 dfl001、greenbea、maros-r7 各做
20 组独立进程测量（奇数组 2/4 线程、偶数组 4/2，预热不计入样本），报告
中位数、nearest-rank P95 与 IQR。修复分支 \texttt{9652ddb9} 上两档均
60/60 准确；4 线程相对 2 线程的区组三例合计中位数降 18.50\%、P95 降
17.72\%，通过预先固定的性能门（原预测 5--15\%，超出预测的调查已写回
推导记录 R6）。合并验收（提交 \texttt{1d31f0eb}，§9.8 与
\texttt{docs/archive/windows\_release\_merge\_2026-09-11.md}）中 NETLIB
两档全集 1080/1080 准确，4 比 2 线程区组合计中位数降 24.60\%、P95 降
22.41\%；但该批绝对耗时高于历史批次且波动明显，不能把 24.60\% 解释为
合并相对旧版本的加速。greenbea 迭代轨迹仍不稳定，长尾未被宣称消除。
P95 是本工作站的样本估计，不构成跨机器的尾延迟保证。

\subsubsection{机器可读证据}

\texttt{reports/}（默认被 Git 忽略）保存各轮原始结果：2026-08-06 基线轮
为 90 个逐案例隔离 JSON（\texttt{baseline\_<案例>.json}），另有
\texttt{h1\_*}--\texttt{h8\_*}、\texttt{p1\_*} 等后续轮次的同构序列；
\srcpath{reports/aggregate.py} 按前缀聚合隔离 JSON，计算每个求解器的
success/accurate 计数、总时间、中位数与两两配对几何均值。引用本章任何
NETLIB 数字时，应能回溯到对应前缀的逐案例文件。

\subsection{MIPLIB 2017 协议与 CPLEX 对照}
\label{sec:engine-nv-miplib}

\subsubsection{协议定义}

可复现比较协议固定于
\texttt{docs/archive/miplib2017\_benchmark\_protocol\_2026-08-25.md}：在
240 实例的 MIPLIB 2017 benchmark 集上做截尾单线程运行；一次运行记为
\texttt{solved} 当且仅当后端证明了参考状态且返回的 incumbent 通过原模型
可行性、整数性与目标审计（参考不可行被证明亦计 solved）。固定指标：

\begin{itemize}
  \item \textbf{PAR-10}：solved 运行取实测求解时间，其余取 $10T$（$T$ 为
        时限）；
  \item \textbf{shifted geomean}（时间）：$\exp(\bar{\ell})-1000$ ms，其中
        $\bar{\ell}$ 为 $\ln(t_i+1000)$ 的均值；
  \item \textbf{shifted geomean}（节点数）：在 solved 且有节点数的运行上，
        $\exp(\mathrm{mean}\ln(n_i+100))-100$；节点数只是描述量——各后端
        presolve、割分离与节点定义不同，节点比只在共同 solved 运行上计算，
        绝不替代 PAR-10；
  \item 两两不确定度：先在实例内聚合重复，再做 10{,}000 次实例级
        bootstrap 取 95\% 百分位区间。
\end{itemize}

协议同时记录了两处「测量与预测不符触发的再推导」（stability-gate 的
\texttt{HVector} 陈旧基维数缺陷、gap-certificate 归一化），其诊断顺序
（实现失真 → 机器/成本模型 → 假设违反 → 理论错误）即仓库 AGENTS.md 的
mismatch 协议实例。其实测验证在 Apple M4 Max / macOS 上完成：ASan 与
Release 的固定种子 100 次循环均 100/100，双调度 CTest 18/18；12 实例
10 秒 pilot 求解数 native/HiGHS/SCIP 为 0/2/2，与修正后预测完全一致。
该 pilot 累计仅 369.4 求解器秒、每后端 solved 0--2 例，协议文档明确将其
定性为冒烟/稳定性实验；\textbf{240 实例正式campaign 尚未执行}，仍待明确
的科学时限选择，不得把 pilot 数字当作求解器性能对比结论。

基准程序 \srcpath{benchmark/miplib2017\_benchmark.cpp:main} 递归扫描
\texttt{.mps[.gz]} 并用 \texttt{.solu} 参考文件审计；模型导入与优化分别
计为 \texttt{read\_ms}/\texttt{solve\_ms}，\texttt{--time-limit} 在导入
完成后才作用于优化阶段；审计驱动脚本为
\srcpath{tools/run\_miplib2017\_audit.py}。

\subsubsection{CPLEX 对照（2026-09-12，Windows/MSVC Release）}

出处：\texttt{docs/manual/09-testing-benchmarks.md} §9.4.3 与
\texttt{docs/archive/cplex\_callable\_library\_integration\_2026-09-12.md}；
机器可读结果为 \texttt{reports/miplib\_cplex\_pilot.\{json,csv\}} 与
\texttt{reports/miplib\_cplex\_12case.\{json,csv\}}。

两实例 pilot（\texttt{mas74}、\texttt{sct2}，seed 0、3 秒时限、单线程）：
CPLEX 最终 gap 10.3008\% 与 0.02690\%，HiGHS 17.4666\% 与 23.4491\%，
Native/HiGHS-LP 39.6664\% 与 3.1466\%；六个 incumbent 均通过原模型审计，
但\textbf{六次运行均未证明最优}，PAR-10 均为 30 秒；Native 在
\texttt{sct2} 上还出现 3 秒配置运行 8.739 秒的 deadline 失真。该轮只能
作为接口正确性与短时 gap pilot。

随后对本机已有全部 12 个样本按相同 seed/gap/3 秒设置复跑，并修正了
HiGHS 时限曾在 \texttt{readModel} 前生效的计时口径错误：CPLEX 为 1 个
proven、8 个 feasible、0 个审计失败，HiGHS 为 0、5、0；双方都有
incumbent 的 5 例中 CPLEX 最终 gap 较小 4 例。CPLEX 的 PAR-10 shifted
geomean 为 23460.1 ms，HiGHS 为 30000.0 ms（12 例均打满时限，两者都接近
上限 $10\times3000$ ms）；CPLEX 最大 import/(import+optimize) 为
4.6122\%，通过预注册的 5\% 门。该扩展仍不是严格等墙钟排名：Windows
runner 没有进程级 hard deadline（记录到 CPLEX 最大
\texttt{solve\_ms=3712}、HiGHS 最大 \texttt{solve\_ms=10033}），只有一个
seed，且本构建关闭 Gurobi——因此只能视为 CPLEX 的有利短预算信号，
\textbf{不能证明通用 MILP 性能已与 Gurobi 比肩}。

\subsection{微基准清单}
\label{sec:engine-nv-microbench}

\texttt{benchmark/} 下的基准目标均输出到 \texttt{tests/} 目录；是否注册为
CTest 见 \ref{sec:engine-nv-testmatrix} 节。下表已按 CMakeLists 注册段与
源文件头部注释逐一核对。

\begin{footnotesize}
\begin{longtable}{@{}p{4.6cm}p{9.6cm}@{}}
\toprule
目标（源文件） & 用途与运行方式 \\
\midrule
\endhead
\bottomrule
\endfoot
\texttt{native\_dual\_bfrt\_simd\_benchmark}\newline
\srcpath{benchmark/native\_dual\_bfrt\_simd\_benchmark.cpp:main} &
  对偶单纯形 Phase-II BFRT 预滤波的独立实验（自带
  \texttt{native\_dual\_bfrt\_simd\_kernel.cpp}）：先证明精确分类与更低的
  指令/访存计数，再谈墙钟；不链接进求解器，非 CTest，手动运行。\\
\texttt{native\_dual\_price\_resident\_benchmark}\newline
\srcpath{benchmark/native\_dual\_price\_resident\_benchmark.cpp:main} &
  S2 独立内核 harness：\texttt{row\_ep} 的 PRICE/dot-error-bound 消费者，
  「物化后读」对「直接读分解后端稠密驻留视图」；在生产分布 fixtures 上
  先证逐位一致，再报告解析字节模型与墙钟。非 CTest。\\
\texttt{native\_dual\_colaq\_pivot\_benchmark}\newline
  \srcpath{benchmark/native\_dual\_colaq\_pivot\_benchmark.cpp:main} &
  S2 独立内核 harness：\texttt{col\_aq} 主元元素读取，
  \texttt{IndexedVector::at()} 的 O(support) 建表对 O(1) 驻留稠密读；
  逐 pivot 固定开销，无法摊销。非 CTest。\\
\texttt{native\_dual\_price\_simd\_benchmark}\newline
  \srcpath{benchmark/native\_dual\_price\_simd\_benchmark.cpp:main} &
  S7 独立内核 harness：PRICE 步 $A^{\top} r_{EP}$ 的 SIMD 门检，在
  d2q06c 形态 fixture 上比较行向 SPA 散射与列向 gather-dot（标量对
  NEON），验收门为动态指令数与有效字节数同时下降，墙钟仅作次级信号。
  非 CTest。\\
\texttt{native\_dual\_workspace\_price\_benchmark}\newline
  \srcpath{benchmark/native\_dual\_workspace\_price\_benchmark.cpp:main} &
  P1+P2 独立内核 harness：整条 PRICE 数据通路两种形态对照（epoch 戳
  稀疏累加器三流多趟 对 驻留稠密单行单流单趟），报告门检指标加墙钟。
  非 CTest。\\
\texttt{kkt\_reuse\_benchmark}\newline
  \srcpath{benchmark/kkt\_reuse\_benchmark.cpp:main} &
  固定模式 KKT 生命周期基准：一次符号分析、重复数值分解、每次分解多
  右端项。非 CTest，手动运行。\\
\texttt{klu\_refactor\_benchmark}\newline
  \srcpath{benchmark/klu\_refactor\_benchmark.cpp:main} &
  KLU 固定模式数值重分解基准；CMake 注释明确其墙钟只作发布证据记录，
  不作机器相关的正确性门，故刻意不注册为 CTest。\\
\texttt{netlib\_solver\_benchmark}\newline
  \srcpath{benchmark/netlib\_solver\_benchmark.cpp:main} &
  NETLIB 跨算法端到端基准（见 \ref{sec:engine-nv-netlib} 节）；完整矩阵
  含一阶与通用 NLP 方法、耗时远超集成回归子集，刻意只做显式基准目标。
  常用复现命令见 \texttt{docs/manual/09-testing-benchmarks.md} §9.4.2。\\
\texttt{miplib2017\_benchmark}\newline
  \srcpath{benchmark/miplib2017\_benchmark.cpp:main} &
  MIPLIB 2017 标准 MILP 比较（见 \ref{sec:engine-nv-miplib} 节）；保留
  MPS 整数性并审计每个 incumbent。非 CTest。\\
\texttt{milp\_benchmark\_runner}\newline
  \srcpath{benchmark/milp\_benchmark\_runner.cpp:main} &
  StrictHiGHS 改进项（warm-start 注入、目标截断、stall 早停、根节点
  IPM 等）在合成 SCUC 实例上的效果测量；CTest 只注册其
  \texttt{--smoke} 形态（\texttt{milp\_benchmark}）。\\
\texttt{native\_kernel\_comparison}\newline
  \srcpath{benchmark/native\_kernel\_comparison.cpp:main} &
  纯原生数值栈（原生 B\&C + 原生对偶单纯形/IPM 作 LP 内核）对裸 API
  HiGHS 的 MILP/LP 双层对照，含对抗性缩放 LP 探针；CTest 注册
  \texttt{--smoke --check --time-limit 60}，\texttt{--check} 把必过行
  转为退出码。其单次样本太小、未重复，只作正确性烟雾测试，不能作性能
  排名依据（§9.4.4）。\\
\texttt{solver\_comparison}\newline
  \srcpath{benchmark/solver\_comparison.cpp:main} &
  同一合成 SCUC MIP 在 Native B\&C、StrictHiGHS、Gurobi（如有许可）、
  HiGHS、SCIP（经 MINLP 重表达）上的并排比较；交互式诊断工具，
  非 CTest。\\
\texttt{nlp\_open\_benchmark}\newline
  \srcpath{benchmark/nlp\_open\_benchmark.cpp:main} &
  原生 Filter-IPM 对 Ipopt 的 HS071 对照（同回调、导数、起点与容差）；
  条件注册为集成测试 \texttt{regression\_nlp\_hs071}。\\
\texttt{conic\_benchmark}\newline
  \srcpath{benchmark/conic\_benchmark.cpp:main} &
  合成 SOCP/SDP 套件（植入自对偶解）的锥 IPM 驱动：逐案例计时与线程
  伸缩表，可写 JSON；兼作冒烟（任一运行非 \texttt{optimal} 即非零
  退出）。非 CTest 注册目标。\\
\end{longtable}
\end{footnotesize}

另有 \texttt{opf\_scale\_probe}（\srcpath{benchmark/opf\_scale\_probe.cpp:main}），
CMake 注释标注为临时 OPF 可伸缩性探针，用后即删，不构成长期证据链。

\subsection{独立性声明与未覆盖边界}

\begin{boundarynote}
下列边界与 \texttt{docs/manual/09-testing-benchmarks.md} §9.5--§9.9 及
各归档记录口径一致；超出边界的结论一概不成立。
\end{boundarynote}

\textbf{已验证的平台/后端组合。}本章全部 Windows 数字来自两台本机
Windows 11 / MSVC 19.44 / Release 构建：Core Ultra 9 285H（2026-08 批次）
与 i9-12900H（2026-09 批次）。稀疏线性代数的验证基线为本地
\textbf{sequential oneMKL} 静态链接（§9.6 基线）；INTEL 线程层
（2/4 线程）只在 LP direct 的 dfl001/greenbea/maros-r7 长尾与 NETLIB 全集
上做过 20 组稳定性验收（§9.7--§9.8），Auto/混合负载没有对应验收，保守
建议仍为 2 线程。MIPLIB 协议的稳定性与归一化验证在 Apple M4 Max / macOS
上完成（见 \ref{sec:engine-nv-miplib} 节）。LP 分解计时诊断不覆盖 macOS
Accelerate 内部分解细分，其验收范围仅为 Windows 后端（§9.7）。

\textbf{未验证的组合。}以下情形没有任何已执行的交叉验证，不得外推：
Linux 平台与 GCC/Clang 构建；macOS 上的 NETLIB 90 全量基准与 INTEL 线程
协议；Gurobi 后端（Windows 各验证批次均在 configure 时关闭）；PaPILO
并发预求解；\texttt{MKL\_CBWR} 开启下的任何性能或复现性结论（实测劣化
4--5 倍，明确禁用）；SCUC 与 Python 目标在 2026-09-11 各批次中未启用，
其最新状态须以相应配置的实际 CTest 列举为准。

\textbf{结论强度的上限。}
(1) 单元/集成测试的 wall time 与编译、其他负载并发，不构成求解器性能
证据；(2) \texttt{success}/\texttt{Optimal} 状态字符串不构成正确性证据，
发布以原模型审计为准，greenbea/tuff 个案（§9.6）是现成的反例；
(3) 配对几何均值、全集几何均值、全量总时间回答三个不同问题，Native 的
NETLIB 结论只在「配对典型性能」口径下领先，全量总时间与尾延迟不领先；
(4) 隔离子进程与单进程计时协议的绝对时间不可混算；(5) MIPLIB 侧的
CPLEX 对照是单 seed、短预算、无进程级 hard deadline 的 pilot 信号，
且 240 实例正式协议campaign 尚未执行，不支持任何总体排名结论；
(6) \texttt{reports/} 默认被 Git 忽略，跨机复现须按
\texttt{docs/manual/09-testing-benchmarks.md} §9.4 的命令重新生成原始
证据，而不是引用他机落盘的数字。


\section{接口与参数}

本章列出调用方最常接触的五组选项结构。所有字段均已对照 2026-09-13 工作树的
头文件逐一核实；未列出即不代表不存在——\texttt{BCOptions} 等大型选项束仅收录关键
字段，完整清单以头文件为准。``默认值''为 C++ 结构体的成员初始化值。

\subsection{\texttt{SolveOptions}（调用级，\srcpath{include/mipsolvers/engine/api/options.hpp}）}

\begin{paramtable}
\fld{preferred\_solver} & — & — & 已注册适配器名 & 空串 & 首选求解器；空表示按策略自动\\
\fld{allow\_fallback} & — & — & true/false & true & 首选失败时允许回退下一候选\\
\fld{strategy\_policy} & — & — & Auto/NativeFirst/ ExternalFirst & Auto & 全局原生/外部优先策略\\
\fld{class\_strategy\_policy} & — & — & 按问题类覆盖 & 空 & 逐 \texttt{ProblemClass} 的策略覆盖表\\
\fld{time\_limit\_sec} & — & s & $\ge 0$ & 0.0 & 整次调用的协作式墙钟预算；0 = 不限，回退只继承剩余时间\\
\fld{threads} & — & 个 & $\ge 0$ & 0 & 求解器工作线程上限；0 = 交由后端自选\\
\fld{random\_seed} & — & — & $\ge 0$ & 0 & 透传给暴露种子接口的后端\\
\fld{memory\_limit\_bytes} & — & 字节 & $\ge 0$ & 0 & 建议性内存预算；0 = 未指定，不能强制执行的后端必须如实报告\\
\fld{portfolio\_mode} & — & — & Latency/ Throughput & Latency & 原生 LP 选择器优化单次时延或聚合吞吐\\
\fld{stop\_token} & — & — & — & 空 & 调用方协作取消令牌，与时限一同检查\\
\end{paramtable}

\subsection{\texttt{IPMLPOptions}（LP 内点，\srcpath{include/mipsolvers/engine/kernel/ipm/ipm\_lp\_solver.hpp}）}

\begin{paramtable}
\fld{max\_iter} & — & 次 & 50--500 & 200 & IPM 最大迭代数\\
\fld{time\_limit\_sec} & — & s & $\ge 0$ & 0.0 & 本次 LP 求解墙钟上限；0 = 不限\\
\fld{tol\_primal} & $\varepsilon_p$ & — & $10^{-10}$--$10^{-6}$ & $10^{-8}$ & 原始可行容差\\
\fld{tol\_dual} & $\varepsilon_d$ & — & $10^{-10}$--$10^{-6}$ & $10^{-8}$ & 对偶可行容差\\
\fld{tol\_gap} & $\varepsilon_g$ & — & $10^{-10}$--$10^{-6}$ & $10^{-8}$ & 互补间隙容差\\
\fld{ruiz\_rounds} & — & 轮 & 0--20 & 10 & Ruiz 均衡轮数；0 = 关闭\\
\fld{max\_correctors} & — & 个 & $-1$--若干 & $-1$ & Gondzio 校正上限；$-1$=自动，0=关闭\\
\fld{centrality\_step\_control} & — & — & — & true & 自适应向中心步长规则；false 回退固定 0.9995 边界比例\\
\fld{presolve} & — & — & — & true & 内点自带预求解（定变量、单例行等）\\
\fld{crossover} & — & — & — & false & 纯化解到顶点（供单纯形热启动）\\
\fld{verbose} & — & — & — & false & 逐迭代日志\\
\fld{use\_highs\_presolve} & — & — & — & false & 显式开启 HiGHS 预求解前置（仅冷启动，失败回退直接求解）\\
\fld{newton\_formulation} & — & — & Auto/ForceNormal/ ForceAugmented & Auto & 牛顿系统形态：法方程或增广系统\\
\fld{cancel\_flag} & — & — & — & nullptr & 协作中止标志（portfolio 竞速取消败者）；空 = 不改变行为\\
\end{paramtable}

\subsection{\texttt{IPMOptions}（NLP 内点关键字段，\srcpath{include/mipsolvers/engine/kernel/ipm/ipm\_solver.hpp}）}

\begin{paramtable}
\fld{max\_iter} & — & 次 & 100--1000 & 400 & 最大迭代数\\
\fld{tol\_primal} & $\varepsilon_p$ & — & $10^{-8}$--$10^{-4}$ & $10^{-6}$ & 原始可行容差\\
\fld{tol\_dual} & $\varepsilon_d$ & — & $10^{-8}$--$10^{-4}$ & $10^{-6}$ & 平稳性容差（配合乘子缩放口径，见头文件注释）\\
\fld{tol\_complementarity} & $\varepsilon_c$ & — & $10^{-8}$--$10^{-4}$ & $10^{-6}$ & 互补守卫（未缩放 $\max_i|s_i\mu_i|$）\\
\fld{tol\_accept} & — & — & $\ge 0$ & 0.0 & 可接受水平的缩放总误差；$\le 0$ 关闭可接受判据\\
\fld{acceptable\_iter} & — & 次 & $\ge 0$ & 0 & 可接受水平需连续满足的迭代数\\
\fld{alpha\_max} & $\bar\alpha$ & — & $(0,1]$ & 0.0 & 步长上限；$\le 0$ 取机器精度可辨的最大严格内点比例\\
\fld{verbose} & — & — & — & false & 逐迭代日志\\
\fld{allow\_external\_fallback} & — & — & — & false & 原生求解失败后可回退 Ipopt（结果中以后端名标识）\\
\end{paramtable}

\subsection{\texttt{ConicIPMOptions}（锥内点，\srcpath{include/mipsolvers/engine/kernel/ipm/conic\_ipm\_solver.hpp}）}

\begin{paramtable}
\fld{max\_iterations} & — & 次 & 20--200 & 100 & 牛顿迭代上限\\
\fld{abstol} / \fld{reltol} & $\varepsilon_a,\varepsilon_r$ & — & $10^{-9}$--$10^{-4}$ & $10^{-7}$ / $10^{-6}$ & 对偶间隙绝对/相对容差\\
\fld{feastol} & $\varepsilon_f$ & — & $10^{-9}$--$10^{-5}$ & $10^{-7}$ & 可行（原始/对偶残差）容差\\
\fld{refinement} & — & 步 & 0--3 & 1 & KKT 精化步数上限（带同坐标 backward-error 闸门）\\
\fld{num\_threads} & — & 个 & $\ge 0$ & 0 & OpenMP 线程数；0 = 库默认\\
\fld{chordal\_decomposition} & — & — & — & true & 对有益稀疏 SDP 块做 chordal 分解\\
\fld{chordal\_min\_order} & — & 阶 & $\ge 2$ & 32 & 参与分解的最小 SDP 阶\\
\fld{chordal\_max\_clique\_ratio} & — & — & $(0,1)$ & 0.75 & 要求的最大团/阶压缩比\\
\fld{chordal\_max\_expansion\_ratio} & — & — & $\ge 1$ & 2.0 & 团/原条目数膨胀比上限\\
\fld{verbose} & — & — & — & false & 逐迭代进度表\\
\end{paramtable}

\subsection{\texttt{SimplexOptions}（对偶单纯形关键字段，\srcpath{include/mipsolvers/engine/kernel/lp\_kernel/dual\_simplex.hpp}）}

\begin{paramtable}
\fld{max\_iter} & — & 次 & $10^3$--$10^6$ & 2000 & 单纯形迭代上限\\
\fld{time\_limit\_sec} & — & s & $\ge 0$ & 0.0 & 墙钟上限；0 = 不限\\
\fld{feasibility\_tol} & $\varepsilon_f$ & — & $10^{-9}$--$10^{-6}$ & $10^{-8}$ & 可行容差\\
\fld{optimality\_tol} & $\varepsilon_o$ & — & $10^{-9}$--$10^{-6}$ & $10^{-8}$ & 最优（对偶可行）容差\\
\fld{verbose} & — & — & — & false & 日志开关\\
\fld{allow\_warm\_primal\_phase\_one} & — & — & — & false & 显式允许对非对偶可行的热启动基先跑原始 Phase I\\
\fld{incumbent\_bound} / \fld{incumbent\_check\_interval} & — & — / 次 & — & nullptr / 50 & 并行 B\&C 在位值早停：原子指针与检查间隔\\
\fld{exact\_dse\_initialization} & — & — & — & false & 显式开启冷启动全精确 DSE 权重初始化（一次检查性 BTRAN/行）\\
\fld{dual\_edge\_weight\_initiali\-zation} & — & — & Production/Devex/ StructuralExact/ FullExact/ CertifiedExact & Production & 对偶边权初始化模式\\
\fld{lp\_kernel\_threads} & — & 个 & $\ge 0$ & 0 & PRICE/BFRT 并行预算；0 = 有界硬件默认，1 = 串行\\
\fld{dual\_shift\_start\_max\_count} / \fld{dual\_shift\_start\_max\_fraction} & — & 个 / — & — & 256 / 0.125 & 对偶可行化移项预算：超过即改用真正的对偶 Phase I\\
\fld{allow\_incremental\_row\_append\_factor} & — & — & — & false & 根割追加行时复用父基因子（块三角结构）\\
\fld{retain\_standard\_form} & — & — & — & false & 结果中保留标准形副本（供割生成）；瞬态求解应关闭\\
\fld{intermediate\_audit\_interval} & — & 次 & $\ge 0$ & 32 & 非终止 INVERT 的完全残差扫描间隔；$\le 0$ 关闭周期扫描\\
\fld{lp\_kernel\_backend} & — & — & HiGHS/ ExperimentalNative & HiGHS & 连续 LP 数值内核；HiGHS 被拒即终止，不回退\\
\fld{factor\_backend} & — & — & BackendA\_Umfpack\-Native 等 & BackendA & 基因子后端（A 为当前生产路径，B 系为实验）\\
\fld{escalation\_residual\_tol} & — & — & $10^{-8}$--$10^{-4}$ & $10^{-6}$ & 冷启动结果原空间可行审计的接受阈值\\
\fld{use\_native\_presolve} / \fld{use\_highs\_presolve} & — & — & — & false / false & 显式开启原生 P1 / HiGHS 预求解（仅冷启动，失败回退直解）\\
\end{paramtable}

遗留 ABI 字段（\fld{allow\_cold\_start}、\fld{use\_partial\_pricing}、
\fld{perturb\_degenerate\_primal}、\fld{fallback\_basis}、
\fld{enable\_degenerate\_frontier\_remap}、\fld{suppress\_degenerate\_frontier\_remap}、
\fld{escalation\_*} 其余字段）仅为源码兼容保留，重写后的原生内核不消费；以头文件
注释为准。

\subsection{\texttt{BCOptions}（原生 B\&C 关键字段，\srcpath{include/mipsolvers/engine/bc/options.hpp}）}

\texttt{BCOptions} 为大型平铺选项束（核心/割/启发式/并行/回退/自动调参等分组），
下表只收录关键字段；完整清单与环境变量覆盖见头文件注释。

\begin{paramtable}
\fld{max\_nodes} & — & 个 & $10^2$--$10^6$ & 50000 & B\&B 节点上限\\
\fld{max\_lp\_iter} & — & 次 & — & 500 & 每个节点 LP/NLP 松弛的迭代上限\\
\fld{time\_limit\_sec} & — & s & $\ge 0$ & 120.0 & 墙钟时限\\
\fld{int\_tol} & $\varepsilon_{\mathrm{int}}$ & — & $10^{-6}$--$10^{-4}$ & $10^{-5}$ & 整性容差\\
\fld{gap\_tol} & $\varepsilon_{\mathrm{gap}}$ & — & $10^{-6}$--$10^{-2}$ & $10^{-4}$ & 相对原始-对偶间隙容差\\
\fld{lp\_tol} & — & — & — & $10^{-6}$ & LP 松弛收敛容差\\
\fld{root\_cut\_rounds} / \fld{cuts\_per\_round} & — & 轮 / 行 & 0--50 & 10 / 20 & 根节点割轮数 / 每轮每节点新增割上限\\
\fld{max\_cut\_depth} & — & 层 & $\ge 0$ & 0 & 树节点割生成的最大深度（0 = 仅根）\\
\fld{gmi\_min\_efficacy} / \fld{gmi\_max\_density} / \fld{gmi\_max\_parallelism} & — & — & — & $10^{-3}$ / 0.35 / 0.90 & GMI 割准入：最小规范违反度、最大非零率、近平行拒绝\\
\fld{cuts} & — & — & None/IntRounding/ MIR/Gomory/ Cover/All & All & 割平面族选择\\
\fld{branching} & — & — & MostInfeasible/ Pseudocost/ FirstFractional & Pseudocost & 分支变量选择策略\\
\fld{node\_sel} & — & — & BestFirst/DepthFirst/ Hybrid & Hybrid & 节点选择策略\\
\fld{use\_simplex\_lp\_nodes} & — & — & — & true & 节点 LP 用对偶单纯形\\
\fld{use\_ipm\_root} / \fld{use\_ipm\_nodes} & — & — & — & false / false & 根/节点松弛改用 IPM\\
\fld{ipm\_root\_probe} & — & — & — & true & 根节点 IPM 试探开关\\
\fld{lp\_kernel\_backend} & — & — & HiGHS/ ExperimentalNative & HiGHS & 节点 LP 数值内核\\
\fld{use\_feasibility\_pump} / \fld{use\_progressive\_rounding} & — & — & — & true / true & 可行性泵 / 渐进舍入启发式\\
\fld{use\_papilo\_presolve} / \fld{native\_presolve\_probing} & — & — & — & true / true & PaPILO 预求解 / 原生探测\\
\fld{feasibility\_pump\_iters} & — & 次 & — & 5 & 可行性泵迭代数\\
\fld{max\_dive\_lps} / \fld{max\_probe\_vars} & — & 次 / 个 & — & 60 / 50 & 潜水 LP 数 / 探测变量数\\
\fld{bound\_propagation\_rounds} & — & 轮 & — & 3 & 界传播轮数\\
\fld{root\_cut\_stall\_tol} / \fld{root\_cut\_max\_stalls} & — & — / 次 & — & $10^{-4}$ / 2 & 根割停滞判据与最大停滞轮数\\
\fld{cut\_pool\_max\_size} / \fld{cut\_pool\_max\_age} & — & 行 / 轮 & — & 2000 / 50 & 割池容量与割龄上限\\
\fld{solution\_pool\_size} & — & 个 & — & 10 & 解池容量\\
\fld{num\_threads} & — & 个 & $-1$--若干 & $-1$ & 树并行线程数；$-1$ = 自动\\
\fld{deterministic\_parallel} & — & — & — & false & 确定性并行模式\\
\fld{enable\_lp\_fallback} & — & — & — & true & 节点 LP 失败的分级回退\\
\fld{verbose} & — & — & — & false & 日志开关\\
\end{paramtable}

\section{验证与边界局限}

\subsection{验证状态}

模块的数值验证证据集中在第 11 章：CTest 单元/集成测试矩阵、NETLIB 90 案例基准、
MIPLIB 2017 协议与 CPLEX 对照、以及 \texttt{benchmark/} 下的独立内核微基准，均注明
证据文件、日期、平台、构建配置与验证层级。历史快照只描述其对应提交与机器的状态，
不直接代替当前提交的验证。

\subsection{边界与局限}

\begin{boundarynote}
以下口径与 \texttt{docs/manual/05-solvers-engines.md} \S 5.8 一致：
\begin{itemize}
  \item \textbf{默认 LP 内核}：部分平台默认内核为 HiGHS；原生对偶单纯形的冷启动、
        极端退化与大型 NETLIB 覆盖需逐平台验证（见第 11 章证据口径）；
  \item \textbf{PDLP}：不提供 basis，不能替代依赖基热启动与精确约化成本的 B\&C
        节点内核；
  \item \textbf{NativeIPMLP}：无 HSD（齐次自对偶）状态机，不能从普通迭代失败推出
        不可行/无界结论；
  \item \textbf{NativeNLP}：固定罚参数不构成一般 NLP 的强收敛保证；``步长很小''
        状态必须结合可行度审计判读；
  \item \textbf{原生 B\&C}：主要实现依赖 HiGHS 库编译；非凸 MINLP 经空间分支定界
        求解，无全局最优性保证；
  \item \textbf{性能结论}：强依赖编译器、BLAS、稀疏后端、线程数与数据集；可复现
        数据见第 11 章，源码用例数不等于通过数。
\end{itemize}
\end{boundarynote}

\section{跨模块工业级技术评价基线}

本章规定本手册所述模块进入工程算例、批量研究或生产服务前必须共同满足的评价基线。
具体模型公式、算法步骤、选项和终止判据以正文相应章节为准；本章不扩大源码已经实现的范围。

\subsection{输入、输出、边界条件与数据结构审计}
输入审计应逐项检查问题维度一致性、系数有限性、上下界合法性（含 $\ell\le u$）、
变量类型标记（连续/整数/二进制）、锥成员归属与选项枚举值；非法输入必须在进入求解内核前由
\srcpath{src/engine/util/problem\_validation.cpp} 一类入口拒绝并返回结构化错误。
输出审计必须记录求解状态（最优/不可行/无界/时限/数值异常且相互区分）、目标值、
原始与对偶残差、间隙、后端名称、容差、迭代/节点计数和告警。
空问题、零行零列、退化约束、固定变量、重复索引与不可用可选后端均属于边界测试集。

\subsection{工程假设与数值稳定性分析}
模型使用的凸性、光滑性、约束规范、中心路径存在性、基非奇异性等假设必须与所求解的问题类匹配。
数值评价至少包含数据缩放（scaling）、原始/对偶不可行度、互补间隙、KKT 残差、
基或因子分解的条件数/枢轴质量、步长/阻尼、迭代停滞与容差复现性。
若采用正则化、扰动、惯性校正、回退后端或近似（如容差放宽），应同时报告触发条件、参数值、
对主指标的敏感性以及是否改变结果口径；不得用``已返回结果''替代最优性或收敛性证明。

\subsection{复杂度与资源分析}
令 $m$、$n$、$n_z$ 分别表示约束数、变量数与矩阵非零元数。
报告应分别给出建模、预求解、求解、后处理的时间与内存。
单纯形单次迭代的稠密代价为 $O(m^2)$ 量级，实际取决于基因子更新与定价策略；
内点法每步代价由 KKT 稀疏因子分解主导，应报告结构非零数、填充率和因子分解峰值内存；
MILP 不给出虚假的多项式最坏界，而报告节点数、间隙、时限和实测扩展曲线。
复杂度结论必须基于固定硬件、固定后端、固定构建配置和可复现算例族。

\subsection{异常工况处理}
不可行、无界、数值失败、时限、内存不足和后端缺失必须相互区分并沿 API 如实上报；
不得把数值失败静默改写为不可行，不得把次优解标为全局最优，不得用零值或 NaN 伪装未知量。
异常路径必须保证问题对象可恢复、可重试；回退（fallback）到替代后端时必须在结果或日志中
声明实际使用的求解器。

\subsection{与其他模块的接口关系}
接口链路遵循``AML/文件输入 $\to$ 问题校验 $\to$ 预求解 $\to$ 求解内核或外部后端
$\to$ 后处理/解回写''。跨模块传递时保留变量/约束的稳定身份、索引空间、符号约定与单位；
消费 SCUC 等上层应用结果前，先检查其求解状态与容差口径。Python 绑定层只做语义保持适配，
不重新解释求解结果。

\subsection{测试与验证建议}
最低验证矩阵包括：解析可解小例、标准算例集（如 NETLIB、MIPLIB 2017）、边界/异常输入、
算法或后端交叉校核、序列化往返、重复运行确定性和规模扩展测试。
数值算法还应加入残差独立重算；近似路径应与高保真路径比较并给出误差分布，而非只报告平均值。
回归报告必须列明提交哈希、构建配置、算例版本、容差、通过/失败/跳过数及未验证范围。

\subsection{工业级评估指标}
\begin{itemize}
  \item \textbf{正确性}：KKT 残差、约束最大违反、互补间隙、解析/基准目标误差、不可行证书正确性；
  \item \textbf{鲁棒性}：标准算例集收敛率、不可行/无界识别率、数值失败率、容差敏感性；
  \item \textbf{性能}：端到端时延、迭代/节点计数、峰值内存、并行效率与规模增长斜率；
  \item \textbf{可追溯性}：选项快照、后端与版本记录、容差口径、告警可定位率；
  \item \textbf{工程有效性}：与参考求解器或历史基线的目标值/时间偏差，及其对业务指标的影响。
\end{itemize}
指标应预先给出验收阈值和失败处置；未达到阈值时只能标记为实验性或受限适用。

\subsection{适用范围、局限性与后续改进}
本手册只评价当前源码覆盖的问题类、后端和选项，不对未实现的问题类、未安装求解器或未执行的
外部验证作保证。后续改进应按``缺口 $\to$ 可量化指标 $\to$ 源码/测试变更 $\to$ 独立验证''闭环，
优先消除会造成错误结论的语义缺口，其次提升数值鲁棒性与性能，最后扩展问题类范围。


\end{document}
